% Source presentations for register row N-c-012.

\subsection{Temporal expressiveness, complexity, safety, and liveness}
\label{sec:p09-presentation-n-c-012}

\subsubsection{Kamp's 1968 thesis presentation}

\paragraph{Source and rights boundary.}
Kamp's University of Stuttgart institutional bibliography supplies the
complete UCLA dissertation: Johan Anthony Willem Kamp, \emph{Tense Logic and
the Theory of Linear Order}, 1968.  The copy's reproduction notice does not
grant permission for further republication; the mathematics below is
reconstructed in independent prose, and citations use the printed thesis pages
\cite{Kamp1968}.

\paragraph{The four source classes of orders.}
Kamp writes \(\mathcal J=\langle T,<\rangle\), with nonempty \(T\), and takes
a linear order to be a strict asymmetric, transitive relation which compares
every two distinct points.  His class names are
\[
\begin{array}{c@{\quad}l}
 \mathsf{Lin} & \text{all strict linear orders},\\
 \mathsf{Den} & \text{those in which }i<j\text{ implies some }i<k<j,\\
 \mathsf{Dis} & \text{those with an immediate successor whenever a later
                      point exists, and an immediate predecessor whenever an
                      earlier point exists},\\
 \mathsf{Com} & \text{those in which every nonempty lower-bounded subset has
                      a greatest lower bound and every nonempty
                      upper-bounded subset has a least upper bound}.
\end{array}
\]
These names stand for classes of order structures.  The last is
order-theoretic completeness, not deductive completeness; it neither entails
density nor excludes endpoints or discreteness.  This nonredundancy is used by
the main theorem: finite orders, \(\langle\mathbb Z,<\rangle\), and
\(\langle\mathbb R,<\rangle\) belong to \(\mathsf{Com}\), whereas
\(\langle\mathbb Q,<\rangle\) does not.  Kamp's conditional predecessor and
successor clauses also show that none of the four class names silently imposes
endpointlessness
\cite[ch.~I, pp.~12--14]{Kamp1968}.

\paragraph{First-order standing and the current-point assignment.}
The displayed first-order language \(L\) has variables
\(t_0,t_1,\ldots\), the ordinary Boolean connectives and quantifiers, one
binary predicate \(<\), and unary predicates \(P_0,P_1,\ldots\).  No equality
symbol occurs in the displayed vocabulary.  An \(n\)-place
\emph{tense-defining formula} is an \(L\)-formula whose only free variable is
\(t_0\) and whose only unary predicate symbols are
\(P_0,\ldots,P_{n-1}\).  This is a metalinguistic class of formulas, not a new
kind of point in \(T\).  It is needed to distinguish the source's
first-order-definable operations from arbitrary functions on
\(\mathcal P(T)\), and it is the quantified class in the main theorem.

For such a formula \(\varphi\), the operation it defines on
\(\mathcal J\) is
\[
\begin{aligned}
 f_{\varphi,\mathcal J}:
 \mathcal P(T)^n&\longrightarrow\mathcal P(T),\\
 f_{\varphi,\mathcal J}(Q_0,\ldots,Q_{n-1})
 &=
 \{t\in T:
   \langle T,<,Q_0,\ldots,Q_{n-1},\ldots\rangle
   \models\varphi[t_0\mapsto t]\}.
\end{aligned}
\]
Thus \(P_i\) receives \(Q_i\) and the distinguished variable \(t_0\) receives
the current point.  There is no additional temporal assignment convention.
For a class of orders, one formula supplies this operation uniformly on every
member of the class; the uniformity, rather than a separate formula for each
order, is part of Kamp's expressibility claim
\cite[ch.~II, defs.~2--3, pp.~21--22]{Kamp1968}.

\paragraph{The temporal language and its exact argument order.}
Kamp's language \(\mathsf{TL}_1\) is generated from proposition letters by
negation, conjunction, and the binary constructors \(S,U\).  Their semantics
are those of Definition~\ref{def:p08-tense-until-since}:
\[
\begin{aligned}
 \mathcal J,Q,t\models S(A,B)
 &\Longleftrightarrow
 \exists s<t\,
 \bigl(\mathcal J,Q,s\models A\land
       \forall u\,(s<u<t\Rightarrow\mathcal J,Q,u\models B)\bigr),\\
 \mathcal J,Q,t\models U(A,B)
 &\Longleftrightarrow
 \exists s>t\,
 \bigl(\mathcal J,Q,s\models A\land
       \forall u\,(t<u<s\Rightarrow\mathcal J,Q,u\models B)\bigr).
\end{aligned}
\]
The witness is strict, the first argument holds at that witness, and the
second holds only throughout the open interval.  On a discrete order that
interval may be empty.  These source constructors differ from the later
non-strict LTL convention, and their reversed-looking argument order may not
be normalized silently.  Ordinary strict existential past and future are
definable from them
\cite[ch.~II, pp.~25--28; ch.~III, pp.~57--58]{Kamp1968}.

\paragraph{The two translation directions.}
The temporal-to-first-order map \(D\) is a metalinguistic map from
\(\mathsf{TL}_1\)-formulas to first-order formulas having the sole free
variable \(t_0\).  In the two temporal clauses below, choose distinct
variables \(x,y\) which do not occur in either translated argument, rename
bound variables first when necessary, and let brackets denote
capture-avoiding replacement of the displayed free occurrences of \(t_0\):
\[
\begin{aligned}
 D(q_i)&=P_i(t_0),\\
 D(\neg A)&=\neg D(A),\\
 D(A\land B)&=D(A)\land D(B),\\
 D(S(A,B))
 &=
 \exists x\Bigl(
   x<t_0\land D(A)[x/t_0]\land
   \forall y\bigl((x<y\land y<t_0)\rightarrow D(B)[y/t_0]\bigr)
 \Bigr),\\
 D(U(A,B))
 &=
 \exists x\Bigl(
   t_0<x\land D(A)[x/t_0]\land
   \forall y\bigl((t_0<y\land y<x)\rightarrow D(B)[y/t_0]\bigr)
 \Bigr).
\end{aligned}
\]
Thus the clauses preserve the strict witness, the open interval, and the
source argument order without introducing two further connectives.  Result
(R2) proves by structural induction that a temporal formula and its
\(D\)-translation agree at every point of every binary structure.  The
forward direction therefore needs neither linearity nor order completeness
\cite[ch.~III, pp.~39--42]{Kamp1968}.

The converse begins with the \(\mathsf{Lin}\) normal-form construction of
Theorem~III.1, given next with its complete quantifier-position proof and
sentence corollary.  This preliminary transformation retains first-order
regional conditions.

% Kamp III.1: source-specific normal form over Lin.

\paragraph{The first-order normal form over \(\mathsf{Lin}\).}
The following construction gives the preliminary result of Kamp's
Theorem~III.1 and its sentence corollary.  The order remains an arbitrary
nonempty strict linear order, with arbitrary interpretations of the unary
predicates.  All intervals below are open; they can be empty.  Neither
density, absence of endpoints, nor order completeness is assumed
\cite[ch.~III, defs.~4--9 and thm.~1, pp.~43--56]{Kamp1968}.

\begin{hbdefinition}{def:p09-kamp-iii1-prefix}{Kamp's temporal prefixes and their representatives}
Work in the equality-free language \(L\) already specified.  In this
construction the source's displayed equality is represented by the
abbreviation
\[
 x\doteq y:=\neg(x<y)\land\neg(y<x).
\]
It expresses coincidence of the assigned points over \(\mathsf{Lin}\).
It is not an additional predicate symbol.  Also put
\[
 \top_x:=\neg(x<x),\qquad \bot_x:=x<x.
\]
These are respectively true and false at every point of every order in
\(\mathsf{Lin}\), and both retain the displayed free variable.

A \emph{temporal prefix} is a finite conjunction of comparisons
\(x<y\) or \(x\doteq y\).  It is \emph{complete} on its variables if it
decides every such comparison over \(\mathsf{Lin}\), and
\emph{consistent} if some linear-order assignment satisfies it.
Order variable names by their indices in \(t_0,t_1,\ldots\).
For a finite nonempty set \(X\) of variables, a complete consistent prefix
specifies an ordered partition
\[
 (B_1,\ldots,B_s)\quad\text{of }X.
\]
Variables in one block denote the same point, and every point of \(B_k\)
precedes every point of \(B_{k+1}\).  Let \(y_k\) be the first variable
name in \(B_k\).  These are Kamp's \emph{essential variables}; their
\emph{order in the prefix} is \(y_1,\ldots,y_s\), the order of their
assigned points, which need not agree with their variable-name order.

The \emph{standard prefix} for this partition contains the chain
\(y_1<\cdots<y_s\), followed by
\(y_k\doteq v\) for each \(v\in B_k\setminus\{y_k\}\).
The chain is omitted when \(s=1\).
Write the equalities block by block, and within each block in variable-name
order.  When \(X=\{y_1\}\), use \(y_1\doteq y_1\).
This writes only adjacent strict comparisons, puts strict comparisons before
equalities, and uses representatives on the left of equalities; a redundant
variable occurs only on the right of its one equality.  It is the source's
standard-prefix convention with an explicit ordering of same-block
equalities and of parentheses.  The number \(s\) is the source's
\(l(\pi)\); \(y_1\) and \(y_s\) are its first and last essential
variables.  The prefix does not assert the existence or absence of other
points between them.

Before a construction, rename bound variables so that no variable is both
free and bound or bound by two quantifiers.  Kamp additionally orders free
variables before bound variables and outer bound variables before inner
ones when speaking of precedence in a formula.  In a temporal prefix all
its variables are free, so representative selection uses the variable-name
order just specified.  Every substitution below concerns free occurrences
and is capture-avoiding.
\end{hbdefinition}

\begin{hbdefinition}{def:p09-kamp-iii1-components}{The \(\operatorname{At}\), \(\operatorname{Bef}\), \(\operatorname{Aft}\), and \(\operatorname{Bet}\) components}
Kamp's literal class \(\operatorname{At}(x)\) consists of finite nonempty
conjunctions of unary literals \(P_i(x)\) and \(\neg P_i(x)\), together
with the particular formula \(P_0(x)\lor\neg P_0(x)\).
A literal is an atomic unary formula or its negation.  Write
\[
 T^K_x:=P_0(x)\lor\neg P_0(x),\qquad
 F^K_x:=P_0(x)\land\neg P_0(x).
\]
Both belong to the source's \(\operatorname{At}(x)\).  The order formulas
\(\top_x,\bot_x\) of Definition~\ref{def:p09-kamp-iii1-prefix} do not
belong to that unextended syntactic class.

Initially
\[
 \operatorname{Bef}_0(x)=\operatorname{Aft}_0(x)
 =\operatorname{Bet}_0(x,y)=\varnothing.
\]
For each row of the following display, retain all components of stage \(n\)
and add the two displayed forms, where \(m\geq1\), each \(\theta_i\)
belongs to the indicated set, and \(z\) is a fresh variable distinct from
the endpoints:
\[
 \exists z\bigl(\beta\land\bigwedge_{i=1}^{m}\theta_i\bigr),
 \qquad
 \forall z\bigl(\beta\longrightarrow\bigvee_{i=1}^{m}\theta_i\bigr).
\]
\[
\begin{array}{c|c|l}
 \text{new class}&\beta&\text{permitted components }\theta_i\\ \hline
 \operatorname{Bef}_{n+1}(x)&z<x&
 \operatorname{At}(z)\cup\operatorname{Bef}_n(z)
                    \cup\operatorname{Bet}_n(z,x)\\
 \operatorname{Aft}_{n+1}(x)&x<z&
 \operatorname{At}(z)\cup\operatorname{Aft}_n(z)
                    \cup\operatorname{Bet}_n(x,z)\\
 \operatorname{Bet}_{n+1}(x,y)&x<z\land z<y&
 \operatorname{At}(z)\cup\operatorname{Bet}_n(x,z)
                    \cup\operatorname{Bet}_n(z,y)
\end{array}
\]
The unindexed classes are the unions over finite stages \(n\).
In the literal source grammar, an empty conjunction inside an existential
constructor is written as \(T^K_z\), and an empty disjunction inside a
universal constructor as \(F^K_z\).  These are formulas in the source's
fixed language \(L\), which contains \(P_0\).
The \(\operatorname{Bef}\) and \(\operatorname{Aft}\) components have only
their one endpoint free, and a \(\operatorname{Bet}\) component has only
its two displayed endpoints free.  The constructors describe nested
first-order conditions on regions; they are not yet temporal operators.

For a separate predicate-support-preserving rendition, define
\(\operatorname{At}^{*}(x)\) to contain the unary-literal conjunctions
and the two distinguished constants \(\top_x,\bot_x\), and repeat the
same recursion with every component class starred.  In that recursion use
\(\top_z\) and \(\bot_z\) for the two empty-list cases.  These
starred classes are an explicitly different syntactic presentation, not a
claim that an order atom is a source \(\operatorname{At}\) component.
The syntactic map \(\varepsilon\) replaces each distinguished starred
truth/falsehood slot by \(T^K_x/F^K_x\) at the same variable and commutes
with the regional constructors.  It maps the starred classes into the
unstarred source classes and preserves truth over \(\mathsf{Lin}\).
\end{hbdefinition}

\begin{hbdefinition}{def:p09-kamp-iii1-standard}{Standard formulas and standard sentences}
For a standard prefix \(\pi\) on \(X\), with essential variables
\(y_1,\ldots,y_s\), let
\[
\begin{split}
 \mathcal C_\pi={}&
 \bigcup_{k=1}^{s}\operatorname{At}(y_k)
 \ \cup\ \bigcup_{k=1}^{s-1}\operatorname{Bet}(y_k,y_{k+1})\\
 &{}\cup\operatorname{Bef}(y_1)\cup\operatorname{Aft}(y_s).
\end{split}
\]
A \emph{standard clause} on \(X\) is
\(\pi\land\bigwedge_j\theta_j\), with every
\(\theta_j\in\mathcal C_\pi\).  If it has no local condition, insert
\(T^K_{y_1}\).  A \emph{standard formula} on \(X\) is a finite
disjunction of standard clauses.  An always-false result can be written by
putting \(F^K_{y_1}\) in a clause for each standard prefix on \(X\).
Thus its named free variables are preserved as well.

A \emph{standard sentence} is \(\exists x\,N(x)\) or
\(\forall x\,N(x)\), where \(N\) is standard on \(\{x\}\).
These are Kamp's Definitions~8--9
\cite[ch.~III, pp.~48, 56]{Kamp1968}.  A \emph{starred standard formula}
or \emph{starred standard sentence} uses the starred component classes and
\(\top_x,\bot_x\) in the same positions.  Its image under
\(\varepsilon\) is a standard formula or sentence in the literal source
grammar.
\end{hbdefinition}

\begin{hbtheorem}{thm:p09-kamp-iii1}{proved=annex:p9.kamp.iii1}{Kamp's linear-order normal form and sentence corollary}
Let \(X=\{x_1,\ldots,x_r\}\), with \(r>0\), be the free variables of
an \(L\)-formula \(\varphi\).  There is an effectively constructed
standard formula \(N_\varphi\) on \(X\) such that
\[
 \mathsf{Lin}\models\varphi\longleftrightarrow N_\varphi.
\]
Equivalence holds for every assignment of \(X\) and every interpretation
of the unary predicates of \(L\).  The source grammar permits its fixed
\(P_0\) in constant components even if the input did not contain it.
Every sentence of \(L\) is \(\mathsf{Lin}\)-equivalent to a Boolean
combination of standard sentences.

\emph{Additional handbook property.}  There is also an effectively
constructed starred standard formula \(N^*_{\varphi}\), with exactly the
same free variables and no additional unary predicate, equivalent to
\(\varphi\) over \(\mathsf{Lin}\).  Taking
\(N_{\varphi}=\varepsilon(N^*_{\varphi})\) gives the literal
source-standard result above.  Predicate-support preservation is a property
of the starred rendition; it is not an additional restriction silently
imposed on Kamp's syntactic classes.

The construction first puts the input into prenex form.  It normalizes the
quantifier-free matrix by complete order types, then removes the original
quantifiers from the outside of the already normalized inner formula, using
the induction in Annex~\ref{annex:p9.kamp.iii1}.  For an existential
quantifier it retains one of the equal-point, left-ray, right-ray, or
intervening-gap possibilities.  For a universal quantifier it groups clauses
by the surviving prefix and by all possible positions of the quantified
point, including positions for which the matrix has no disjunct.  All
required cases, including empty regions, are proved in the annex
\cite[ch.~III, thm.~1 and cor.~1, pp.~48--56]{Kamp1968}.
\end{hbtheorem}

\paragraph{Witness existence at an endpoint or in a gap.}
The starred standard form of \(\exists x\,(x<y\land P(y))\) is
\[
 (y\doteq y)\land P(y)\land
 \exists x\,(x<y\land\top_x).
\]
The last component belongs to \(\operatorname{Bef}^{*}_1(y)\); its
\(\varepsilon\)-image belongs to \(\operatorname{Bef}_1(y)\).
On a one-point order with \(P(y)\) true the formula is false; on
\(\{0,1\}\) with \(y=1\) and \(P(1)\) true it is true.
The mirror \(\exists x\,(y<x\land P(y))\) uses
\(\operatorname{Aft}^{*}_1(y)\) and is false when \(y\) is the last point.
For a gap the component
\[
 (y<z)\land\exists x\,(y<x\land x<z\land\top_x)
\]
is false at consecutive integers \(y=0,z=1\), and true for the same
endpoints in the rational order.  A strict comparison between named
endpoints does not itself supply an intervening witness.

The universal counterpart has the form
\[
 \forall x\,(x<y\longrightarrow P(y))
 \quad\longleftrightarrow\quad
 P(y)\lor\forall x\,(x<y\longrightarrow\bot_x).
\]
The second disjunct records the empty left ray.  The analogous formulas with
\(y<x\) and \(y<x<z\) record a last point and an empty intervening gap.
These examples explain why the universal construction retains an emptiness
condition when its remaining disjuncts are independent of the quantified
variable.

\paragraph{Scope of this normal form.}
The output still contains nested first-order quantification in its
\(\operatorname{Bef}\), \(\operatorname{Aft}\), and
\(\operatorname{Bet}\) components.  Theorem~\ref{thm:p09-kamp-iii1}
establishes the preliminary \(\mathsf{Lin}\) normal form and the sentence
corollary.  Expressing the one-variable components in the image of \(D\)
is the further task of Theorem~III.2 and its completeness-dependent lemmas.


For one free current-point variable,
Theorem~III.2 and Lemmas~1--10 perform the further reduction to formulas built
from current-point, strict-past, strict-future, and intervening-interval
conditions.  The cut and interval arguments use greatest lower bounds and
least upper bounds, so this second stage is proved only over
\(\mathsf{Com}\).  Its output lies in the image of \(D\), whose unique inverse
image is the required temporal formula
\cite[ch.~III, pp.~48--94]{Kamp1968}.

% Kamp III.2, first supporting lemma: basic past and future patterns.

\paragraph{The first basic-pattern lemma for Theorem III.2.}
Kamp's Theorem~III.2 begins with special finite interval patterns.  Their
existential past and future closures can already be expressed by the
strict operators \(S,U\) over \(\mathsf{Lin}\).  This is Lemma~1 of the
source's proof of III.2; its statement and proof are on printed
pp.~69--70, following the basic-pattern definitions on pp.~63--66
\cite[ch.~III, pp.~60--70]{Kamp1968}.

\begin{hbdefinition}{def:p09-kamp-basic-expressibility}{Expression at a named current point}
Keep the language \(L\), temporal language \(\mathsf{TL}_1\), and map
\(D\) fixed above.  If \(R\) is the unary predicate letter \(P_i\),
write \(\widehat R=q_i\) for its corresponding temporal proposition
letter.  For an \(L\)-formula \(\psi\) with exactly the free variable
\(w\), a temporal formula \(\eta\) \emph{expresses \(\psi\) over
\(\mathsf{Lin}\)} when
\[
 \mathsf{Lin}\models
 \psi\longleftrightarrow D(\eta)[w/t_0].
\]
Rename bound variables before replacing the free current-point variable
\(t_0\) by \(w\).  Equivalence includes every interpretation of the
unary predicates and every assignment to \(w\).  This is the
named-variable form of the source's expressibility convention on p.~59.
It uses Kamp's strict argument order: \(S(A,B)\) or \(U(A,B)\) places
\(A\) at its witness and \(B\) throughout the intervening open interval.
\end{hbdefinition}

\begin{hbdefinition}{def:p09-kamp-basic-classes}{Basic past and future interval patterns}
Let \(x,y\) be distinct variable names.  All bound variables below are
fresh, and all substitutions of endpoint names concern free occurrences.
The letters \(A,B,C\) range over unary predicate letters of \(L\);
they need not be different.

The first binary stages are identical:
\[
\begin{split}
 \operatorname{Bp}_1(x,y)=\operatorname{Bf}_1(x,y)
 =\bigl\{\,&A(x)\land
       \forall v\bigl((x<v\land v<y)\longrightarrow B(v)\bigr)
       \land C(y):\\[-2pt]
       &A,B,C\text{ are unary predicate letters}\,\bigr\}.
\end{split}
\]
For \(n\geq1\), \(\operatorname{Bp}_{n+1}(x,y)\) consists of the
formulas
\[
\begin{gathered}
 A(x)\land\exists z\Bigl(
 x<z\land z<y\land
 \forall v\bigl((x<v\land v<z)\longrightarrow B(v)\bigr)
 \land\phi\Bigr),
 \\
 \phi\in\operatorname{Bp}_n(z,y),
\end{gathered}
\]
and \(\operatorname{Bf}_{n+1}(x,y)\) consists of the formulas
\[
\begin{gathered}
 \exists z\Bigl(
 x<z\land z<y\land\phi\land
 \forall v\bigl((z<v\land v<y)\longrightarrow B(v)\bigr)
 \Bigr)\land A(y),
 \\
 \phi\in\operatorname{Bf}_n(x,z).
\end{gathered}
\]
Each binary formula has exactly its two displayed endpoints free.
The subscripts count the open interval pieces of the pattern.  These are
exact stages, without a clause retaining all formulas of earlier stages.
Put
\[
 \operatorname{Bp}(x,y)=\bigcup_{n\geq1}\operatorname{Bp}_n(x,y),
 \qquad
 \operatorname{Bf}(x,y)=\bigcup_{n\geq1}\operatorname{Bf}_n(x,y).
\]

The corresponding one-variable classes are
\[
\begin{split}
 \operatorname{Bp}_n(y)
   &=\{\exists x\,(x<y\land\phi):
              x\ne y,\ \phi\in\operatorname{Bp}_n(x,y)\},\\
 \operatorname{Bf}_n(x)
   &=\{\exists y\,(x<y\land\phi):
              y\ne x,\ \phi\in\operatorname{Bf}_n(x,y)\},
\end{split}
\]
where \(x\ne y\) here says that the variable names are distinct.
Let \(\operatorname{Bp}(y)\) and \(\operatorname{Bf}(x)\) be their
unions over \(n\geq1\).  The notations
\((\exists x<y)\phi\) and \((\exists y>x)\phi\) abbreviate the two
displayed guarded quantifiers.

The binary base formula does not itself assert \(x<y\).  That guard
is supplied in the one-variable closures.  In particular, an empty open
interval does not remove the requirement for two strictly ordered
endpoint witnesses.  These \(\operatorname{Bp},\operatorname{Bf}\)
classes are distinct from the earlier
\(\operatorname{Bef},\operatorname{Aft},\operatorname{Bet}\) classes
\cite[ch.~III, pp.~63--66]{Kamp1968}.
\end{hbdefinition}

\begin{hbdefinition}{def:p09-kamp-basic-word}{The predicate word and its realizations}
For \(\phi\in\operatorname{Bp}_n(x,y)\cup
\operatorname{Bf}_n(x,y)\), list the unary predicate occurrences from
left to right, including repeated occurrences.  Denote the resulting word by
\[
 \sigma(\phi)=(A_0,B_1,A_1,\ldots,B_n,A_n).
\]
It has \(2n+1\) positions.  Kamp denotes this occurrence sequence by
\(S(\phi)\); \(\sigma\) here distinguishes it from the temporal
operator.  The binary base has word \((A,B,C)\).  A past-stage
constructor prepends \((A,B)\) to its constituent's word; a future-stage
constructor appends \((B,A)\).

In a linear structure \(\mathcal M=\langle T,<,(P_i^{\mathcal M})_i\rangle\),
a \emph{realization} of this word is a sequence of points
\(a_0<\cdots<a_n\) such that
\[
\begin{aligned}
 &a_i\in A_i^{\mathcal M} &&(0\leq i\leq n),\\
 &(a_{i-1},a_i)\subseteq B_i^{\mathcal M} &&(1\leq i\leq n).
\end{aligned}
\]
Here \((a,b)=\{t\in T:a<t<b\}\), and \(R^{\mathcal M}\) is the
interpretation of a unary letter \(R\).  The interval requirement is
vacuous when that interval is empty.  This is the source's sequence
satisfaction relation on pp.~65--66, indexed by \(n\) intervals and
\(n+1\) points.
\end{hbdefinition}

\begin{hblemma}{thm:p09-kamp-iii2-lemma1}{proved=annex:p9.kamp.iii2.lemma1}{Kamp's basic past and future lemma}
Every member of \(\operatorname{Bp}(w)\cup\operatorname{Bf}(w)\)
is expressible over \(\mathsf{Lin}\).

More explicitly, let
\(\sigma=(A_0,B_1,A_1,\ldots,B_n,A_n)\), where \(n\geq1\), be the
word of the binary pattern inside the guarded existential.  Form the
temporal formulas
\[
\begin{aligned}
 L_0&=\widehat A_0,&
 L_i&=\widehat A_i\land S(L_{i-1},\widehat B_i)
            &&(1\leq i\leq n),\\
 R_n&=\widehat A_n,&
 R_i&=\widehat A_i\land U(R_{i+1},\widehat B_{i+1})
            &&(0\leq i<n).
\end{aligned}
\]
Construct the \(L_i\) in increasing index order and the \(R_i\) in
decreasing index order.  If \(\phi\in\operatorname{Bp}_n(x,y)\)
and \(\psi\in\operatorname{Bf}_n(x,y)\) have word \(\sigma\), then
\[
\begin{aligned}
 \mathsf{Lin}\models
  \exists x\,(x<y\land\phi)
       &\longleftrightarrow D(L_n)[y/t_0],\\
 \mathsf{Lin}\models
  \exists y\,(x<y\land\psi)
       &\longleftrightarrow D(R_0)[x/t_0].
\end{aligned}
\]
These are finite constructions using only the input predicate letters and
their corresponding proposition letters.  The seeds \(L_0,R_n\) are
single proposition letters; they do not introduce a stage zero into the
source's basic-pattern classes.  No density, endpointlessness or order
completeness is required \cite[ch.~III, lemma~1, pp.~69--70]{Kamp1968}.
\end{hblemma}

\paragraph{Two interval pieces.}
For a pattern with word \((A_0,B_1,A_1,B_2,A_2)\), its past closure is
expressed by
\[
 \widehat A_2\land
 S\bigl(\widehat A_1\land S(\widehat A_0,\widehat B_1),
         \widehat B_2\bigr),
\]
and its future closure by
\[
 \widehat A_0\land
 U\bigl(\widehat A_1\land U(\widehat A_2,\widehat B_2),
         \widehat B_1\bigr).
\]
Each requires three strictly increasing points.  On a two-point order both
are false, even if every displayed predicate holds everywhere.  For a
one-piece pattern the past formula is false at a first point and the future
formula is false at a last point.  At the other endpoint of a two-point
order, the one-piece formula can be true: its sole open interval is empty,
so its interval predicate imposes no condition.

\paragraph{Role in the remaining construction.}
Theorem~\ref{thm:p09-kamp-iii1} reduces a one-variable formula to Boolean
combinations of current-point and regional components.  The lemma above
supplies the temporal expressions for the basic finite-chain patterns used
in Kamp's subsequent analysis of those components.  The source next treats
the limiting past and future classes in Lemma~2, then develops the remaining
interval lemmas.  Expressibility of all
\(\operatorname{Bef},\operatorname{Aft}\) components over
\(\mathsf{Com}\), and hence Theorem~III.2 itself, is a further result
\cite[ch.~III, pp.~60--62, 66--70]{Kamp1968}.

% Kamp III.2, Lemma 2: limiting finite patterns over all linear orders.

\subsubsection{Limiting past and future patterns}
\label{sec:p09-kamp-iii2-lemma2}

The next lemma expresses finite patterns that can be realized inside every
past or future interval adjoining the current point.  Kamp's Lemma~2 states
expressibility of the limiting classes over \(\mathsf{Lin}\); its proof
occupies printed pp.~70--73 \cite[ch.~III, lemma~2, pp.~66--73]{Kamp1968}.
The construction below supplies both directions and every possible position
at which a provisional finite chain reaches the current point.

\begin{hbdefinition}{def:p09-kamp-limit-operators}{The two strict limit operations}
For a temporal formula \(\eta\), use the following abbreviations in
\(\mathsf{TL}_1\):
\[
\begin{aligned}
 H'(\eta)&\defeq S(\neg(\eta\land\neg\eta),\eta),&
 G'(\eta)&\defeq U(\neg(\eta\land\neg\eta),\eta),\\
 P'(\eta)&\defeq\neg H'(\neg\eta),&
 F'(\eta)&\defeq\neg G'(\neg\eta).
\end{aligned}
\]
These are Kamp's derived operations, with the strict witness and open-interval
argument order already fixed for \(S,U\)
\cite[ch.~III, defs.~10--11, pp.~57--59]{Kamp1968}.
For a first-order formula \(\theta\) whose sole free variable is \(w\),
write
\[
\begin{aligned}
 \partial^-\theta(w)&\defeq
   \forall x\bigl(x<w\to\exists z\,
        (x<z\land z<w\land\theta[z/w])\bigr),\\
 \partial^+\theta(w)&\defeq
   \forall x\bigl(w<x\to\exists z\,
        (w<z\land z<x\land\theta[z/w])\bigr).
\end{aligned}
\]
The bound variables \(x,z\) are fresh.  The symbols \(\partial^-,\partial^+\)
name these formula operations; they do not denote analytic derivatives.
The annex proves that if \(\eta\) expresses \(\theta(w)\), then
\(P'(\eta)\) and \(F'(\eta)\) express these two formulas.

The universal conditions include empty past and future domains.  In
particular, \(\partial^-\theta\) holds at a first point and
\(\partial^+\theta\) holds at a last point, whatever \(\theta\) says.
They impose no condition on \(\theta\) at the current point itself.
\end{hbdefinition}

\begin{hbdefinition}{def:p09-kamp-limiting-classes}{Limiting basic-pattern classes}
Keep the basic classes and occurrence words of
\cref{def:p09-kamp-basic-classes,def:p09-kamp-basic-word}.
For each unary predicate letter \(A\), the base formulas are
\(\partial^-A(w)\) and \(\partial^+A(w)\).  Set
\[
 \operatorname{Lp}_0(w)\defeq\{\partial^-A(w):A\text{ unary}\},
 \qquad
 \operatorname{Lf}_0(w)\defeq\{\partial^+A(w):A\text{ unary}\}.
\]
For every \(n\geq1\), define
\[
\begin{split}
 \operatorname{Lp}_n(w)\defeq\bigl\{
  &\forall x\bigl(x<w\to\exists u\,
     (x<u\land u<w\land\exists v\,
       (u<v\land v<w\land\phi))\bigr):\\
  &\hspace{35mm}\phi\in\operatorname{Bp}_n(u,v)\bigr\},\\
 \operatorname{Lf}_n(w)\defeq\bigl\{
  &\forall x\bigl(w<x\to\exists u\,
     (w<u\land u<x\land\exists v\,
       (u<v\land v<x\land\phi))\bigr):\\
  &\hspace{35mm}\phi\in\operatorname{Bp}_n(u,v)\bigr\}.
\end{split}
\]
All displayed endpoint names are distinct and all other bound variables are
fresh.  Thus each resulting formula has exactly \(w\) free.  The positive
index counts the \(n\) interval pieces of its constituent basic pattern;
the base index zero concerns one predicate letter.  There is no retention
of earlier stages inside a later stage.  Put
\[
 \operatorname{Lp}(w)\defeq\bigcup_{n\geq0}\operatorname{Lp}_n(w),
 \qquad
 \operatorname{Lf}(w)\defeq\bigcup_{n\geq0}\operatorname{Lf}_n(w),
\]
and let \(\operatorname{Lp},\operatorname{Lf}\) be their unions over the
choice of free variable.  Both source base definitions test their predicate
at the quantified witness point.  Moving that test to the current point
changes the condition, as the handbook comparison in
\cref{def:p09-kamp-current-point-comparison} demonstrates.
The higher-stage indexing is stated by interval count throughout
\cite[ch.~III, pp.~63, 66--67, 70--72]{Kamp1968}.
\end{hbdefinition}

\begin{hbdefinition}{def:p09-kamp-limit-translation}{The finite temporal construction}
Let
\(\sigma=(A_0,B_1,A_1,\ldots,B_n,A_n)\) be an occurrence word.
For \(n\geq1\), let \(E^-_\sigma\) be the formula \(L_n\) and
\(E^+_\sigma\) the formula \(R_0\) constructed in
\cref{thm:p09-kamp-iii2-lemma1}.  They say respectively that a realization
of the whole word ends or begins at the current point.  For \(n=0\),
set \(E^-_{(A_0)}=E^+_{(A_0)}=\widehat A_0\).

For \(n\geq1\), the suffix and prefix words are
\[
 \sigma^-=(A_1,B_2,A_2,\ldots,B_n,A_n),\qquad
 \sigma^+=(A_0,B_1,A_1,\ldots,B_{n-1},A_{n-1}).
\]
They each have one fewer interval; at \(n=1\) they are the one-letter
words \((A_1)\) and \((A_0)\).  Define recursively
\[
\begin{aligned}
 \Lambda^-_{(A)}&\defeq P'(\widehat A),&
 \Lambda^+_{(A)}&\defeq F'(\widehat A),\\
 \Lambda^-_\sigma&\defeq
            P'(E^+_\sigma)\land\Lambda^-_{\sigma^-},&
 \Lambda^+_\sigma&\defeq
            F'(E^-_\sigma)\land\Lambda^+_{\sigma^+}
                    \qquad(n\geq1).
\end{aligned}
\]
This is a finite construction: each recursive call removes exactly two
predicate occurrences and one interval.  All its leaves are proposition
letters already occurring in the input word.  The abbreviations \(H',G',P',F'\)
expand into negation, conjunction and the original \(S,U\).
\end{hbdefinition}

\begin{hblemma}{thm:p09-kamp-iii2-lemma2}{proved=annex:p9.kamp.iii2.lemma2}{Kamp's limiting past and future lemma}
Every member of \(\operatorname{Lp}\cup\operatorname{Lf}\) is expressible
over \(\mathsf{Lin}\).  More explicitly, a member of
\(\operatorname{Lp}_n(w)\), with constituent occurrence word \(\sigma\),
is expressed by \(\Lambda^-_\sigma\); a member of
\(\operatorname{Lf}_n(w)\) is expressed by \(\Lambda^+_\sigma\).
At stage zero, the constituent word is its single predicate letter.

These assertions use the named-variable \(D\)-translation of
\cref{def:p09-kamp-basic-expressibility}, with capture-avoiding replacement
of \(t_0\) by \(w\).  They hold on arbitrary nonempty strict linear
orders, with arbitrary interpretations of the unary predicates.  Density,
endpointlessness and order completeness are not hypotheses
\cite[ch.~III, lemma~2, pp.~70--73]{Kamp1968}.
\end{hblemma}

\begin{hbexample}{def:p09-kamp-limit-one-piece}{One interval and the required second conjunct}
For the word \((A,B,C)\), the two formulas are
\[
\begin{aligned}
 \Lambda^-_{(A,B,C)}
   &=P'(\widehat A\land U(\widehat C,\widehat B))
                                      \land P'(\widehat C),\\
 \Lambda^+_{(A,B,C)}
   &=F'(\widehat C\land S(\widehat A,\widehat B))
                                      \land F'(\widehat A).
\end{aligned}
\]
The first conjunct alone permits a provisional chain whose other endpoint
lies beyond the current point.  For example, on \(\mathbb R\) at zero,
take \(A=(-\infty,0)\), \(B=\mathbb R\), and \(C=\{1\}\).
Then \(P'(\widehat A\land U(\widehat C,\widehat B))\) is true, while
the limiting past pattern and \(P'(\widehat C)\) are false.
The future counterpart takes \(A=\{-1\}\), \(B=\mathbb R\), and
\(C=(0,\infty)\): its first future conjunct is true and its second
conjunct and limiting future pattern are false.
\end{hbexample}

\begin{hbexample}{def:p09-kamp-limit-endpoints}{Endpoints, adjacency and an extremal crossing}
At a first point every limiting past formula above is true; at a last point
every limiting future formula is true.  If a point has an immediate
predecessor, every limiting past formula is false.  If it has an immediate
successor, every limiting future formula is false.  In particular both
families hold at the point of a singleton order, independently of the
predicate interpretation.

The last possible crossing must be included in the past splice.  With all
predicates true on \(\mathbb R\), the two-interval chain
\((-3,-1,0)\) begins above the cut \(-4\) but reaches the current point
zero.  A suffix realization \((-1/2,-1/4)\) lies above its last strictly
past point \(-1\).  Splicing gives \((-3,-1,-1/4)\), wholly below zero.
The crossing is between the last two original points and includes equality
of the old last point with the current point.  The proof covers this case
and its future counterpart explicitly.
\end{hbexample}

\begin{hbexample}{def:p09-kamp-current-point-comparison}{A current-point condition and a strict limit}
Replacing the witness-point predicate \(A(z)\) by \(A(w)\) produces a
different formula.  For this additional handbook comparison, write
\[
 \operatorname{CurF}_A(w)\defeq
  \forall x\bigl(w<x\to\exists z\,
                  (w<z\land z<x\land A(w))\bigr).
\]
This also is expressible over \(\mathsf{Lin}\).  Put
\(\tau_A\defeq\neg(\widehat A\land\neg\widehat A)\); its expression is
\[
 \neg U(\tau_A,\tau_A)\ \lor\
                  (\widehat A\land F'(\tau_A)),
\]
where disjunction abbreviates \(\neg(\neg\alpha\land\neg\beta)\).
On \(\mathbb R\) at zero with \(A=\{0\}\), this current-point condition is
true and the source's witness-point base \(\partial^+A\) is false.
Thus a condition on \(A\) at the current point cannot replace the
witness-point clause in the limiting definition.  The annex proves this
comparison as well as the source-based limiting construction.
\end{hbexample}

\paragraph{The next dependency.}
The next source result is Lemma~3 on printed p.~73.  It concerns the closure
of current-point, basic and limiting classes under the Boolean operations
and proper predicate substitution, called \(E\) in the source.  Its
substitution step and the subsequent interval lemmas remain separate work.
Lemma~3 still uses \(\mathsf{Lin}\); only after it does Kamp change the
default convention to \(\mathsf{Com}\).  Lemma~2 alone does not establish
the full expressibility result of Theorem~III.2.

% Kamp III.2, Lemma 3: the Boolean and proper-substitution closure E.

\subsubsection{Kamp's Boolean and substitution closure}
\label{sec:p09-kamp-iii2-lemma3}

Kamp next closes the current-point, basic and limiting classes under Boolean
operations and proper substitution for unary predicate letters. His
Lemma~3 states that every formula in the resulting class is expressible.
The proof invokes Lemmas~1--2 and the substitution observation R4
\cite[ch.~III, pp.~59--60, 67, 73]{Kamp1968}.
The construction below supplies that observation with its variable
conditions and gives an expression at every finite construction step.

\begin{hbdefinition}{def:p09-kamp-proper-substitution}{Proper substitution for a unary predicate}
Keep the equality-free first-order language \(L\) fixed above. Let
\(\phi\) be an \(L\)-formula, let \(P_i\) be a unary predicate letter,
and let \(\psi\) be an \(L\)-formula with exactly one free variable,
called \(v\). An \emph{alphabetic variant} changes bound variable names
and their bound occurrences without changing the binding or the free
variables. Write
\[
 \operatorname{psub}_{P_i}(\phi,\psi)
\]
for the following operation. In the fixed enumeration of \(L\)-formulas
used by Kamp, choose the first alphabetic variant \(\psi'\) of
\(\psi\) whose bound variables do not occur in \(\phi\). Replace each
original atomic occurrence \(P_i(y)\) in \(\phi\) by the result of
replacing the variable \(v\) throughout \(\psi'\) by the argument
variable \(y\), in the source's variable-replacement convention.
Leave all other predicate atoms and all order
atoms unchanged. This is Kamp's proper predicate substitution, written in
the source as \([\phi]^{\psi}_{P_i}\)
\cite[ch.~III, def.~1, p.~40]{Kamp1968}.

The enumeration here selects a formula spelling; it is a different object
from the formula class \(\operatorname{E}\) defined below. Fresh bound
names can also be chosen directly, separating any bound occurrences of
the parameter name \(v\) first. Then substitute \(y\) only for its
free occurrences, written \(\psi'[y/v]\). Any such choice gives an
alphabetic variant of the source-selected result, as proved in the annex.
Thus the construction need not identify the first entry in an unspecified
enumeration, and does not claim literal identity of selected bound names.

Substitution visits the original atoms of \(\phi\) once. It does not
then replace occurrences of \(P_i\) inside an inserted copy of
\(\psi'\). The replacement may itself contain \(P_i\), and it may
contain other unary letters already used by \(\phi\). No disjointness
or nonoccurrence condition is imposed on predicate letters. If \(P_i\)
does not occur in \(\phi\), the result is \(\phi\).

The one-free-variable condition is exact: a sentence is not an admissible
replacement here, nor is a formula with additional free parameters. The
binary relation symbol \(<\) is not a substitution target. The name
\(v\) need not be the free variable of \(\phi\); it is the argument
place of the unary replacement. When an original occurrence \(P_i(y)\)
is under a quantifier binding \(y\), that same quantifier binds the
inserted free argument occurrences of \(y\). Freshness protects the
replacement's own bound variables; it does not remove this intended binding.
\end{hbdefinition}

\begin{hbdefinition}{def:p09-kamp-temporal-substitution}{Substitution in the temporal language}
For \(\eta,\zeta\in\mathsf{TL}_1\), define
\(\eta[\zeta/q_i]\) by recursion on \(\eta\):
\[
\begin{aligned}
 q_j[\zeta/q_i]&\defeq
   \begin{cases}\zeta,&j=i,\\q_j,&j\ne i,\end{cases}\\
 (\neg\alpha)[\zeta/q_i]&\defeq\neg(\alpha[\zeta/q_i]),\\
 (\alpha\land\beta)[\zeta/q_i]
   &\defeq\alpha[\zeta/q_i]\land\beta[\zeta/q_i],\\
 S(\alpha,\beta)[\zeta/q_i]
   &\defeq S(\alpha[\zeta/q_i],\beta[\zeta/q_i]),\\
 U(\alpha,\beta)[\zeta/q_i]
   &\defeq U(\alpha[\zeta/q_i],\beta[\zeta/q_i]).
\end{aligned}
\]
The temporal language has no variable-binding constructor. At the atomic
case the output is the supplied \(\zeta\), without another substitution
inside it. The two operands of \(S,U\) retain their strict witness and
open-interval roles. Derived temporal operations are expanded into these
primitive constructors before this recursion is applied.

Disjunction is the abbreviation
\(\alpha\lor\beta\defeq\neg(\neg\alpha\land\neg\beta)\).
A finite sequence of substitutions means repeated applications of the
single operation in the stated order; the empty sequence is the identity.
It is not a simultaneous replacement or a recursive system of predicate
equations.
\end{hbdefinition}

\begin{hblemma}{thm:p09-kamp-proper-substitution}{proved=annex:p9.kamp.proper.substitution}{Proper substitution preserves expressibility}
Let \(\phi\) have exactly the free variable \(w\), and let
\(\psi\) have exactly the free variable \(v\). Suppose that
\(\eta,\zeta\in\mathsf{TL}_1\) express them over
\(\mathsf{Lin}\), respectively, with the named-variable convention of
\cref{def:p09-kamp-basic-expressibility}. Then
\(\operatorname{psub}_{P_i}(\phi,\psi)\) has exactly \(w\) free and
is expressed by \(\eta[\zeta/q_i]\). In particular,
\[
 \mathsf{Lin}\models
 \operatorname{psub}_{P_i}(\phi,\psi)
 \longleftrightarrow D(\eta[\zeta/q_i])[w/t_0].
\]
This includes \(v=w\), a predicate absent from \(\phi\), and a
replacement containing that same predicate. Each substitution is finite.
The annex also proves compatibility directly when the two first-order
inputs are the named-variable \(D\)-translations themselves. This is the
full substitution step needed from Kamp's R4
\cite[ch.~III, pp.~59--60]{Kamp1968}.
\end{hblemma}

\begin{hbdefinition}{def:p09-kamp-elementary-class}{The finite closure \(\operatorname{E}\)}
For a variable \(w\), let
\[
 \mathcal G(w)\defeq
  \operatorname{At}(w)\cup\operatorname{Bp}(w)
  \cup\operatorname{Bf}(w)\cup\operatorname{Lp}(w)
  \cup\operatorname{Lf}(w).
\]
Here \(\operatorname{At}\) is the unstarred source class from
\cref{def:p09-kamp-iii1-components}; it consists of finite nonempty
conjunctions of unary literals and the particular tautology
\(P_0(w)\lor\neg P_0(w)\). The basic and limiting classes are those
of \cref{def:p09-kamp-basic-classes,def:p09-kamp-limiting-classes}.

Kamp defines \(\operatorname{E}(w)\) as the closure of this union under
negation, conjunction, disjunction and proper predicate substitution
\cite[ch.~III, p.~67]{Kamp1968}. Explicit finite stages for that least
closure are as follows. Set \(\operatorname{E}_0(w)\defeq\mathcal G(w)\).
For each \(m\in\mathbb N\), at stage \(m+1\) retain
\(\operatorname{E}_m(w)\) and add:
\[
\begin{gathered}
 \neg\phi,\qquad \phi\land\psi,\qquad \phi\lor\psi,
 \qquad \operatorname{psub}_{P_i}(\phi,\psi),\\
 \phi,\psi\in\operatorname{E}_m(w),\qquad i\in\mathbb N.
\end{gathered}
\]
Put
\[
 \operatorname{E}(w)\defeq\bigcup_{m\geq0}\operatorname{E}_m(w),
 \qquad
 \operatorname{E}\defeq\bigcup_{j\geq0}\operatorname{E}(t_j).
\]
The substitution operands in this fixed-variable closure both belong to
\(\operatorname{E}_m(w)\); the preceding substitution lemma permits the
more general parameter name \(v\) as well. Each constructed formula has
exactly its named variable free. The stages describe finite constructions,
not a closure under semantic equivalence, arbitrary quantification,
infinitary Boolean operations, or substitution of arbitrary formulas from
outside the class. A formula can have several construction histories.

There are no primitive zero-operand Boolean constructors in this grammar.
When an empty finite conjunction or disjunction is requested as an
abbreviation, use respectively
\(T^K_w=P_0(w)\lor\neg P_0(w)\) or
\(F^K_w=P_0(w)\land\neg P_0(w)\), which already belong to
\(\operatorname{At}(w)\). This uses the declared source letter \(P_0\).
\end{hbdefinition}

\begin{hbmethod}{def:p09-kamp-elementary-translation}{A temporal expression from a finite construction}
The input is a finite construction of \(\phi\in\operatorname{E}(w)\),
including the constituent basic or limiting pattern at a seed. Read that
construction from its leaves upward and assign a temporal formula to each
node:
\begin{enumerate}
\item For an \(\operatorname{At}(w)\) seed, replace \(P_i(w)\) by
      \(q_i\) and preserve its literal negations and conjunctions.
      For its particular tautology use
      \(\neg(q_0\land\neg q_0)\).
\item For a basic past or future seed, use respectively the \(L_n\) or
      \(R_0\) construction in \cref{thm:p09-kamp-iii2-lemma1}.
      For a limiting past or future seed, use respectively
      \(\Lambda^-_\sigma\) or \(\Lambda^+_\sigma\) from
      \cref{thm:p09-kamp-iii2-lemma2}, including their one-letter seeds.
\item For a negation node with child expression \(\eta\), return
      \(\neg\eta\). For conjunction and disjunction nodes with child
      expressions \(\eta,\zeta\), return respectively
      \(\eta\land\zeta\) and
      \(\neg(\neg\eta\land\neg\zeta)\).
\item For a node \(\operatorname{psub}_{P_i}(\phi_1,\phi_2)\),
      first obtain its two child expressions \(\eta,\zeta\), then
      return \(\eta[\zeta/q_i]\) by the finite structural recursion
      just defined.
\end{enumerate}
Stage retention introduces no new operation. The output is one finite
\(\mathsf{TL}_1\)-formula. A construction history is an explicit witness
of membership; this method does not assert a decision procedure for
membership of an arbitrary first-order formula in \(\operatorname{E}\).
\end{hbmethod}

\begin{hblemma}{thm:p09-kamp-iii2-lemma3}{proved=annex:p9.kamp.iii2.lemma3}{Kamp's expressibility lemma for \(\operatorname{E}\)}
Every member of \(\operatorname{E}\) is expressible over
\(\mathsf{Lin}\). More explicitly, if a finite construction of
\(\phi\in\operatorname{E}(w)\) returns \(\eta\) by
\cref{def:p09-kamp-elementary-translation}, then
\[
             \mathsf{Lin}\models
                  \phi\longleftrightarrow D(\eta)[w/t_0].
\]
The conclusion holds for every interpretation of the unary predicates
on every nonempty strict linear order. The output uses only proposition
letters corresponding to unary predicate letters of the constructed
formula, including the source's \(P_0\) when a source truth/falsehood
seed contains it. No order-completeness, density, or endpointlessness
assumption is used \cite[ch.~III, lemma~3, p.~73]{Kamp1968}.
\end{hblemma}

\begin{hbexample}{def:p09-kamp-substitution-example}{A limiting predicate inside a basic future pattern}
Let \(A,B,C\) be distinct unary predicate letters. Take the basic future
formula with word \((A,B,C)\), expressed by
\(\widehat A\land U(\widehat C,\widehat B)\), and substitute
\(\partial^-C(w)\) properly for \(C\). Both operands belong to
\(\operatorname{E}(w)\). The resulting temporal formula is
\[
                  \widehat A\land U(P'(\widehat C),\widehat B).
\]
The inserted limit is tested at the future endpoint chosen by \(U\),
not at the original current point. Its internal occurrence of \(C\)
is retained: a single substitution step does not call itself again inside
the inserted formula. The annex writes the first-order result with fresh
variables and verifies this example.
\end{hbexample}

\paragraph{The following order-completeness boundary.}
Immediately after Lemma~3, Kamp changes the default interpretation of
``valid'', ``equivalent'' and ``expressible'' from \(\mathsf{Lin}\) to
\(\mathsf{Com}\). The next result is Lemma~4, concerning conjunction
for the interval property \(I(x,y)\); its proof begins on printed p.~73.
Its elementary interval classes, interval property and remaining cases
are the next separate dependency. The present closure lemma supplies no
proof that every one-variable first-order formula belongs to
\(\operatorname{E}\), and does not complete Theorem~III.2
\cite[ch.~III, pp.~67--68, 73]{Kamp1968}.

% Kamp III.2, Lemma 4: conjunction by finite fusion of interval patterns.

\subsubsection{Conjunction of elementary interval conditions}
\label{sec:p09-kamp-iii2-lemma4}

Kamp's next step combines two finite interval descriptions with the same
ordered endpoints. Their internal witnesses can interleave or coincide.
Lemma~4 accounts for every such arrangement by a finite disjunction of
elementary interval formulas. The source introduces the classes and the
property on printed pp.~67--68; the complete proof runs from p.~73 to the
explicit end of Lemma~4 on p.~78 \cite[ch.~III, pp.~67--68, 73--78]{Kamp1968}.

\begin{hbdefinition}{def:p09-kamp-elementary-intervals}{Elementary past and future interval classes}
Use the basic classes of \cref{def:p09-kamp-basic-classes}, the class
\(\operatorname{E}\) of \cref{def:p09-kamp-elementary-class}, and proper
unary-predicate substitution from \cref{def:p09-kamp-proper-substitution}.
For distinct variable names \(x,y\), a member of
\(\operatorname{Ep}(x,y)\) is obtained from a member of
\(\operatorname{Bp}(x,y)\) by a finite sequence of proper substitutions
whose replacements belong to \(\operatorname{E}\). The corresponding
construction from \(\operatorname{Bf}(x,y)\) defines
\(\operatorname{Ef}(x,y)\).

More literally, the source writes the successive targets as
\(P_0,P_1,\ldots,P_N\), for some finite \(N\), with a replacement in
\(\operatorname{E}\) for each target. Identity replacements
\(P_i(w)\) allow unused positions in such a list. The equivalence of this
initial-segment description and arbitrary finite substitution histories is
proved in the annex. Replacements are applied successively in the stated
order; a replacement inserted at one step can therefore be affected by a
later step. Within a single step, its newly inserted atoms are not visited
again. Repeated and shared predicate letters are allowed.

As in the preceding substitution construction, bound variables are chosen
fresh before inserting a replacement, and formulas obtained by different
such choices are compared up to alphabetic variation. Every elementary
interval formula has exactly \(x,y\) free. Neither binary class inserts
an outer assertion \(x<y\).
These are Kamp's elementary past and elementary future classes, distinct
from the one-variable class \(\operatorname{E}\)
\cite[ch.~III, p.~67]{Kamp1968}.
\end{hbdefinition}

\begin{hbdefinition}{def:p09-kamp-labelled-intervals}{Elementary labels on a finite interval word}
Fix a variable \(w\). An \emph{elementary labelled word} is a finite list
\[
 W=(\alpha_0,\beta_1,\alpha_1,\ldots,\beta_m,\alpha_m),
 \qquad m\geq1,
\]
whose entries belong to \(\operatorname{E}(w)\). Repetitions are
permitted and retain their separate positions. For a label \(\theta\)
and a variable \(v\), write \(\theta(v)\) for the capture-avoiding
replacement \(\theta[v/w]\), after making its bound variables fresh.
This is a formula abbreviation, not application of a new predicate symbol.

Let \(\mathcal B_p(W;x,y)\) be the basic past pattern with these
one-variable formulas in its successive point and interval positions. Its
explicit recursion is
\[
\begin{split}
 \mathcal B_p((\alpha_0,\beta_1,\alpha_1);x,y)
 &\defeq \alpha_0(x)\land
   \forall v\,(x<v<y\to\beta_1(v))\land\alpha_1(y),\\
 \mathcal B_p(W;x,y)
 &\defeq\alpha_0(x)\land\exists z\bigl(x<z<y\land
   \forall v\,(x<v<z\to\beta_1(v))\\
 &\hspace{43mm}{}\land\mathcal B_p(W^-;z,y)\bigr)\quad(m>1),
\end{split}
\]
where \(W^-=(\alpha_1,\beta_2,\alpha_2,\ldots,\beta_m,\alpha_m)\).
The notation \(a<b<c\) abbreviates \(a<b\land b<c\).
Every displayed bound variable is fresh from the endpoints and from all
labels at that occurrence. The future version has the same one-interval
base and the recursion
\[
\begin{split}
 \mathcal B_f(W;x,y)\defeq
 \exists z\bigl(x<z<y\land\mathcal B_f(W^+;x,z)
 &\land\forall v\,(z<v<y\to\beta_m(v))\bigr)\\
 &{}\land\alpha_m(y)\quad(m>1),
\end{split}
\]
where \(W^+=(\alpha_0,\beta_1,\ldots,\beta_{m-1},\alpha_{m-1})\).

In a linear structure \(\mathcal M\), a realization of \(W\) is a
strictly increasing sequence \(a_0<\cdots<a_m\) such that
\(\alpha_i\) holds with \(w=a_i\), and \(\beta_i\) holds with
\(w=t\) for every \(a_{i-1}<t<a_i\). A label is evaluated in the
whole structure with its stated free argument. Its internal quantifiers
are not restricted to the displayed interval. Only the outer interval
quantifier in the pattern has that restriction.
\end{hbdefinition}

\begin{hblemma}{thm:p09-kamp-elementary-intervals}{proved=annex:p9.kamp.elementary.intervals}{Elementary words, substitution and the interval remarks}
Every elementary interval formula has an elementary labelled word obtained
by following its finite substitution history at each original predicate
occurrence. Conversely, every \(\mathcal B_p(W;x,y)\) has a
\(\operatorname{Ep}(x,y)\) presentation, and every
\(\mathcal B_f(W;x,y)\) has an \(\operatorname{Ef}(x,y)\)
presentation, with only fresh bound names possibly differing.
Whenever \(a<b\), either presentation holds at \((a,b)\) exactly when
\(W\) has a realization beginning at \(a\) and ending at \(b\).
These assertions hold over \(\mathsf{Lin}\).

In particular the source's remarks R5--R6 hold: if
\(\phi\in\operatorname{Ep}(x,y)\), its guarded past closure
\(\exists x\,(x<y\land\phi)\) has an \(\operatorname{E}(y)\)
presentation; if \(\phi\in\operatorname{Ef}(x,y)\), its guarded
future closure \(\exists y\,(x<y\land\phi)\) has an
\(\operatorname{E}(x)\) presentation. Proper substitution of another
member of \(\operatorname{E}\) preserves either elementary interval
class. These are statements about the displayed constructions with allowed
bound renaming, and do not add semantic-equivalence closure to the grammar
of \(\operatorname{E}\) \cite[ch.~III, pp.~67--68]{Kamp1968}.
\end{hblemma}

\begin{hbdefinition}{def:p09-kamp-interval-property}{Kamp's interval property \(I(x,y)\)}
Let \(x,y\) be distinct variable names. An \(L\)-formula \(\phi\)
with no free variables other than possibly
\(x,y\) has \emph{property \(I(x,y)\)} if there is a finite
disjunction of formulas \(\theta_i\in\operatorname{Ep}(x,y)\) such that
\[
 \mathsf{Com}\models
 x<y\longrightarrow
             \left(\phi\longleftrightarrow\bigvee_i\theta_i\right).
\]
The equivalence concerns every unary-predicate interpretation and every
assignment on every complete linear order. It makes no assertion about
assignments with \(x\not<y\). No condition that both variables actually
occur freely in \(\phi\) is imposed.

A disjunction may be taken nonempty: a false interval formula is represented
by \(\mathcal B_p((F^K_w,T^K_w,T^K_w);x,y)\), with the source
truth and falsehood formulas of \cref{def:p09-kamp-elementary-class}.
Thus an empty disjunction, if used as a finite Boolean abbreviation, can
be replaced by this one disjunct. This convention introduces no new
primitive connective \cite[ch.~III, p.~68, display~(13)]{Kamp1968}.
\end{hbdefinition}

\begin{hblemma}{thm:p09-kamp-interval-substitution}{proved=annex:p9.kamp.interval.substitution}{Substitution and one-point instances of the interval property}
If \(\phi\) has property \(I(x,y)\) and
\(\psi\in\operatorname{E}\), then
\(\operatorname{psub}_{P_i}(\phi,\psi)\) has property \(I(x,y)\).
If \(\theta\in\operatorname{E}(x)\) and \(y\) is a different
variable name, then \(\theta\) has both property \(I(x,y)\) and
property \(I(y,x)\). These are the source's remarks R7--R8; the annex
proves their substitution semantics and both endpoint placements
\cite[ch.~III, p.~68]{Kamp1968}.
\end{hblemma}

\begin{hbdefinition}{def:p09-kamp-fusions}{Fusions of two endpoint sequences}
Let the two input words be
\[
 W=(\alpha_0,\beta_1,\ldots,\beta_m,\alpha_m),\qquad
 V=(\gamma_0,\delta_1,\ldots,\delta_n,\gamma_n),
 \qquad m,n\geq1.
\]
Give their point positions distinct tags
\((e_0,\ldots,e_m)\) and \((f_0,\ldots,f_n)\). Tags record
occurrences, even when the formulas attached to two positions coincide.

A \emph{total preorder} on these tags is a reflexive, transitive binary
relation \(F\) comparing every two tags. Put
\(a\sim_F b\) when both \(aFb\) and \(bFa\) hold. A
\emph{fusion} is such a relation whose restriction to each input sequence
is exactly that sequence's reflexive linear order, and which satisfies
\[
                      e_0\sim_F f_0,
                \qquad e_m\sim_F f_n.
\]
Thus different positions of the same input cannot coincide, while a
position of one input may coincide with a position of the other. The
endpoint identifications are mandatory. Write \(\mathcal F_{m,n}\)
for the finite set of these fusions.

For a fusion \(F\), enumerate its equivalence classes in their strict
order as \(C_1,\ldots,C_{r_F}\). Define \(h_F(a)=j\) when
\(a\in C_j\). In particular,
\[
 h_F(e_0)=h_F(f_0)=1,\qquad
 h_F(e_m)=h_F(f_n)=r_F,\qquad r_F\geq2.
\]
Kamp calls a total preorder a \emph{quasi-ordering}, a finite sequence
without repetition a \emph{pure sequence}, and two such sequences with
disjoint members a \emph{pure pair}. His general fusion preserves the
two input orders; the endpoint identifications select the fusions used
in Lemma~4. His \emph{independent} members of the second
sequence are those not equivalent to any member of the first. Taking all
members of the first sequence and those independent members of the second
selects exactly one representative of each class. Hence his fusion length
is the number \(r_F\) of classes used here
\cite[ch.~III, pp.~74--75]{Kamp1968}.
\end{hbdefinition}

\begin{hbmethod}{def:p09-kamp-fusion-construction}{The formula associated with a fusion}
In addition to the point tags, give the interval positions tags
\(b_1,\ldots,b_m,d_1,\ldots,d_n\), all distinct. Let \(\mathcal T\)
be the set of all point and interval tags. Attach the labels
\[
 \theta_{e_i}=\alpha_i,\quad \theta_{b_i}=\beta_i,\quad
 \theta_{f_j}=\gamma_j,\quad \theta_{d_j}=\delta_j.
\]
For a fixed fusion \(F\), there are \(2r_F-1\) alternating slots:
slot \(2k-1\) is its \(k\)-th point, and slot \(2k\) is the open
gap between its \(k\)-th and \((k+1)\)-st points. Define the set of
slots required by each input tag:
\[
\begin{aligned}
 f_F(e_i)&\defeq\{2h_F(e_i)-1\},&
 f_F(f_j)&\defeq\{2h_F(f_j)-1\},\\
 f_F(b_i)&\defeq
   \{s:2h_F(e_{i-1})-1<s<2h_F(e_i)-1\},\\
 f_F(d_j)&\defeq
   \{s:2h_F(f_{j-1})-1<s<2h_F(f_j)-1\}.
\end{aligned}
\]
The variable \(s\) ranges over integers in \(\{1,\ldots,2r_F-1\}\).
The notation \(f_F\) is a set-valued map on tags; the tag \(f_j\)
is a point position of the second input. These have different declared
roles. In particular an interval tag can require both even gap slots and
odd point slots strictly between its endpoints.

For each slot put
\[
 \rho_{F,s}(w)\defeq
       \bigwedge_{\substack{\tau\in\mathcal T\\s\in f_F(\tau)}}\theta_\tau(w).
\]
Use the tag order
\(e_0,b_1,e_1,\ldots,b_m,e_m,f_0,d_1,f_1,\ldots,d_n,f_n\)
to fix the order of this finite conjunction.
Every slot has at least one required tag, as proved in the annex, so its
label belongs to \(\operatorname{E}(w)\). All listed conjuncts are
required; no negative requirement is imposed on any unlisted label.
Finally set
\[
 \chi_F(x,y)\defeq
  \mathcal B_p((\rho_{F,1},\ldots,\rho_{F,2r_F-1});x,y).
\]
It has an \(\operatorname{Ep}(x,y)\) presentation by
\cref{thm:p09-kamp-elementary-intervals}. Enumerate all fusions and
return \(\bigvee_{F\in\mathcal F_{m,n}}\chi_F(x,y)\).
This is the slot construction of source pp.~75--76, with point positions
and alternating slots indexed separately.
\end{hbmethod}

\begin{hblemma}{thm:p09-kamp-iii2-lemma4}{proved=annex:p9.kamp.iii2.lemma4}{Kamp's conjunction lemma for \(I(x,y)\)}
If \(\phi,\psi\) have property \(I(x,y)\), then \(\phi\land\psi\)
has property \(I(x,y)\).

The elementary construction used in the proof has the stronger, explicit
linear-order identity
\[
\begin{split}
 \mathsf{Lin}\models x<y\longrightarrow
 \Bigl(&\mathcal B_p(W;x,y)\land\mathcal B_p(V;x,y)\\
       &\longleftrightarrow
                   \bigvee_{F\in\mathcal F_{m,n}}\chi_F(x,y)\Bigr).
\end{split}
\]
For given finite \(I(x,y)\) representations, distribute their conjunction
over every pair of disjuncts, construct their labelled words, and apply this
fusion method to each pair. Restricting the displayed identity to
\(\mathsf{Com}\) gives the required representation of
\(\phi\land\psi\). All constructions are finite and use only
predicate letters in the input elementary labels; temporary fresh letters
used to exhibit membership are eliminated
\cite[ch.~III, lemma~4, pp.~73--78]{Kamp1968}.
\end{hblemma}

\begin{hbexample}{def:p09-kamp-fusion-three-cases}{Two interior points: either order or coincidence}
Take the two words
\(W=(A_0,B_1,A_1,B_2,A_2)\) and
\(V=(C_0,D_1,C_1,D_2,C_2)\), with all entries in
\(\operatorname{E}(w)\). Write \(a\) and \(c\) for their respective
interior witnesses. The three fusions have the following labelled words;
conjunction here combines labels at the same parameter \(w\):
\[
\begin{array}{c|l}
 a<c & (A_0\land C_0,\ B_1\land D_1,\ A_1\land D_1,\
           B_2\land D_1,\\[-2pt]
     & \hspace{12mm} C_1\land B_2,\ B_2\land D_2,\ A_2\land C_2)\\[3pt]
 a=c & (A_0\land C_0,\ B_1\land D_1,\ A_1\land C_1,\
           B_2\land D_2,\ A_2\land C_2)\\[3pt]
 c<a & (A_0\land C_0,\ B_1\land D_1,\ C_1\land B_1,\
           B_1\land D_2,\\[-2pt]
     & \hspace{12mm} A_1\land D_2,\ B_2\land D_2,\ A_2\land C_2).
\end{array}
\]
The coincident case cannot be discarded: on the three-point order
\(0<1<2\), take both middle labels true only at 1 and every other
label true everywhere. The two input patterns hold at \((0,2)\), and
only the coincident fusion can be realized.

An interval label must also hold at an inserted interior point. On
\(0<1<2<3\), take \(A_1\) true only at 1, \(C_1\) true only at 2,
\(B_2\) false everywhere, and all other labels true everywhere. The
first input fails because its interval \((1,3)\) contains 2. In the
\(a<c\) word the conjunct \(B_2\) at the \(C_1\) position detects
this failure; the two neighboring open gaps are empty and cannot detect it.
The annex verifies these cases and the one-interval and endpoint cases.
\end{hbexample}

\paragraph{Order scope and the next source result.}
Property \(I(x,y)\) is defined with \(\mathsf{Com}\) on p.~68, and
the convention immediately before Lemma~4 assigns that same scope to
unqualified equivalence. The opening proof on p.~73 nevertheless prints
\(\mathsf{Lin}\) on the representation premises. The argument here
uses those premises with their defining \(\mathsf{Com}\) scope and
proves the elementary fusion identity separately over \(\mathsf{Lin}\).
Its restriction to complete orders is sufficient; no converse transfer
from a complete-order equivalence to all linear orders is assumed.

The next result is Lemma~5, beginning on printed p.~78. It concerns the
guarded existential composition through an interior point \(z\), with
\(\operatorname{E}(z)\) conditions at that point and interval properties
\(I(x,z)\), \(I(z,y)\) on either side. That construction and the later
lemmas remain separate from the conjunction result proved here
\cite[ch.~III, pp.~68, 73, 78--79]{Kamp1968}.

% Kamp III.2, Lemma 5: existential composition through a shared midpoint.

\subsubsection{Composition through an interior point}
\label{sec:p09-kamp-iii2-lemma5}

Kamp's Lemma~5 combines elementary conditions on two intervals meeting at
one quantified point. The point must lie strictly between the outer
endpoints and satisfy its own elementary conditions. The source uses
Lemma~4 for the conjunction on each side, distributes over their
elementary disjuncts, and concatenates each pair of finite patterns
\cite[ch.~III, lemma~5, pp.~78--79]{Kamp1968}.

\begin{hbdefinition}{def:p09-kamp-midpoint-word}{Concatenating two words at an elementary midpoint}
Use the elementary labelled words and formula recursion
\(\mathcal B_p\) of \cref{def:p09-kamp-labelled-intervals}.
Fix a common free parameter \(w\), a label
\(\eta\in\operatorname{E}(w)\), and words
\[
\begin{aligned}
 W&=(\alpha_0,\beta_1,\alpha_1,\ldots,\beta_a,\alpha_a),
       &&a\geq1,\\
 V&=(\gamma_0,\delta_1,\gamma_1,\ldots,\delta_b,\gamma_b),
       &&b\geq1.
\end{aligned}
\]
All entries belong to \(\operatorname{E}(w)\). Define
\[
\begin{split}
 W\mathbin{\star_\eta}V\defeq
  (&\alpha_0,\beta_1,\alpha_1,\ldots,\beta_a,\,
                  \eta\land\alpha_a\land\gamma_0,\\
   &\delta_1,\gamma_1,\ldots,\delta_b,\gamma_b).
\end{split}
\]
The symbol \(\star_\eta\) is a finite word operation: it retains the
left word up to its last point, combines the three midpoint labels at
that point, and then appends the right word after its first point.
It is not an object-language connective.
The conjunction is associated as \((\eta\land\alpha_a)\land\gamma_0\).

The resulting word has \(a+b\) intervals, \(a+b+1\) point positions
and \(2(a+b)+1\) alternating point and interval positions. Its midpoint
is point position \(a\), counting from zero. Every interval label from
both inputs is retained, and both labels formerly attached to the shared
point occur in its new conjunction. In particular \(a,b\geq1\) make
the midpoint an interior position. Labels remain formulas on the whole
structure at their sole free argument; their internal quantifiers are
not restricted to either constituent interval.
\end{hbdefinition}

\begin{hblemma}{thm:p09-kamp-elementary-midpoint}{proved=annex:p9.kamp.elementary.midpoint}{The elementary midpoint concatenation}
Let \(x,y,z\) be distinct variable names, let
\(\theta\in\operatorname{E}(z)\), and let
\(\rho\in\operatorname{Ep}(x,z)\),
\(\sigma\in\operatorname{Ep}(z,y)\).
Obtain their elementary words \(W,V\) with a common label parameter
\(w\) distinct from \(x,y,z\), using
\cref{thm:p09-kamp-elementary-intervals}, and put
\(\eta\defeq\theta[w/z]\), renaming bound variables before replacing
free occurrences. Then
\[
                 \mathcal B_p(W\mathbin{\star_\eta}V;x,y)
\]
has an \(\operatorname{Ep}(x,y)\) presentation and
\[
\begin{split}
 \mathsf{Lin}\models x<y\longrightarrow
 \Bigl(&\exists z\,(x<z<y\land\theta\land\rho\land\sigma)\\
       &\longleftrightarrow
              \mathcal B_p(W\mathbin{\star_\eta}V;x,y)\Bigr).
\end{split}
\]
Every bound variable in the constituent formulas is first made fresh
from the other constituents and from \(x,y,z\). The outer
\(\exists z\) binds the original free occurrences of \(z\) in
\(\theta,\rho,\sigma\). The output has exactly \(x,y\) free; the
parameter \(w\) and the intermediate endpoint \(z\) do not remain free.
This is the elementary concatenation underlying the source's
display~(21), with both directions supplied in the annex
\cite[ch.~III, p.~79]{Kamp1968}.
Its \(\mathsf{Lin}\) scope is an additional conclusion of that
finite-realization proof.
\end{hblemma}

\begin{hbconvention}{def:p09-kamp-midpoint-families}{Finite families and their empty cases}
Let \(x,y,z\) be distinct variable names. For finite family sizes
\(m,n,p\geq0\), take
\[
\begin{gathered}
 \chi_i\in\operatorname{E}(z)\quad(1\leq i\leq m),\\
 \phi_j\text{ with property }I(x,z)\quad(1\leq j\leq n),\qquad
 \psi_k\text{ with property }I(z,y)\quad(1\leq k\leq p).
\end{gathered}
\]
The property \(I\) is exactly the guarded \(\mathsf{Com}\) property
of \cref{def:p09-kamp-interval-property}. In particular a \(\phi_j\)
has no free variables other than possibly \(x,z\), and a \(\psi_k\)
has no free variables other than possibly \(z,y\).

Write the three aggregate formulas as
\[
\begin{aligned}
 \theta(z)&\defeq
  \begin{cases}\bigwedge_{i=1}^{m}\chi_i,&m>0,\\T^K_z,&m=0,\end{cases}\\
 \Phi(x,z)&\defeq
  \begin{cases}\bigwedge_{j=1}^{n}\phi_j,&n>0,\\T^K_x,&n=0,\end{cases}\\
 \Psi(z,y)&\defeq
  \begin{cases}\bigwedge_{k=1}^{p}\psi_k,&p>0,\\T^K_y,&p=0,\end{cases}
\end{aligned}
\]
where finite nonempty conjunctions are associated from left to right.
The arguments of \(\Phi,\Psi\) display their permitted free variables;
they do not assert that every argument actually occurs freely.
The formulas \(T^K_v=P_0(v)\lor\neg P_0(v)\) are the source
tautologies already fixed in \cref{def:p09-kamp-elementary-class}.
Define the binary formula
\[
 \Gamma(x,y)\defeq
 \exists z\bigl(x<z<y\land\theta(z)\land\Phi(x,z)\land\Psi(z,y)\bigr).
\]
Bound variables inside all inputs are renamed fresh before this
construction. The outer order guard supplies both free endpoints,
including when every family is empty.

The source displays the indexed families and their three conjunctions
in (20), and reduces the midpoint family to one formula. It gives no
separate zero-family analysis there. The definitions for zero sizes
above make the handbook's empty-family reading explicit; they are
proved below together with the displayed positive-size case. A
zero-size family adds no semantic condition, but it does not delete
the existential midpoint or either strict inequality
\cite[ch.~III, p.~78]{Kamp1968}.
\end{hbconvention}

\begin{hbmethod}{def:p09-kamp-midpoint-construction}{Constructing the interval disjunction}
The input consists of the families above, their finite
\(\operatorname{E}\) constructions, and finite elementary disjunctions
witnessing each stated property \(I\). The construction uses those
witnesses; it does not decide whether a bare formula has property \(I\).

Choose \(w\) distinct from \(x,y,z\). Form \(\theta(z)\) and its renamed label
\(\eta(w)=\theta[w/z]\). Construct finite lists \(\mathcal L\) and
\(\mathcal R\) of elementary interval formulas as follows.

For \(n>0\), initialize \(\mathcal L\) with the given representation
of \(\phi_1\). For \(j=2,\ldots,n\), apply the finite fusion construction
of \cref{thm:p09-kamp-iii2-lemma4} to every pair consisting of a
current list member and a disjunct in the representation of \(\phi_j\).
Replace the list by all resulting elementary fusion disjuncts.
For \(n=0\), take the one-member list whose member is
\(\mathcal B_p((T^K_w,T^K_w,T^K_w);x,z)\).
Construct \(\mathcal R\) in the same way from
\(\psi_1,\ldots,\psi_p\) at endpoints \(z,y\); for \(p=0\), use
\(\mathcal B_p((T^K_w,T^K_w,T^K_w);z,y)\).

A supplied empty disjunction means false on the ordered endpoints.
It can first be replaced by the one false elementary pattern in
\cref{def:p09-kamp-interval-property}. Thus all lists in this method
may be taken nonempty. A singleton family needs no conjunction step.
Repeated disjuncts may be retained.

For each \(\rho\in\mathcal L\), obtain its elementary word \(W_\rho\);
for each \(\sigma\in\mathcal R\), obtain its word \(V_\sigma\).
Return the finite disjunction
\[
 \bigvee_{\rho\in\mathcal L}\ \bigvee_{\sigma\in\mathcal R}
       \mathcal B_p(W_\rho\mathbin{\star_\eta}V_\sigma;x,y).
\]
Every original choice of disjuncts on the left and on the right,
and every intervening fusion branch, is included. The final step takes
every pair from the two resulting lists. The shared midpoint fixes
the order of the two constituent pieces.
\end{hbmethod}

\begin{hblemma}{thm:p09-kamp-iii2-lemma5}{proved=annex:p9.kamp.iii2.lemma5}{Kamp's guarded interval composition lemma}
For positive \(m,n,p\), the source formula
\[
 \exists z\left(x<z<y\land
       \bigwedge_{i=1}^{m}\chi_i\land
       \bigwedge_{j=1}^{n}\phi_j\land
       \bigwedge_{k=1}^{p}\psi_k\right)
\]
has property \(I(x,y)\), under the hypotheses of
\cref{def:p09-kamp-midpoint-families}.

With the explicit empty-family convention there, the conclusion also
holds for every \(m,n,p\geq0\). More precisely, the construction in
\cref{def:p09-kamp-midpoint-construction} gives
\[
\begin{split}
 \mathsf{Com}\models x<y\longrightarrow
 \Bigl(&\Gamma(x,y)\\
 &\longleftrightarrow
   \bigvee_{\rho\in\mathcal L}\ \bigvee_{\sigma\in\mathcal R}
      \mathcal B_p(W_\rho\mathbin{\star_\eta}V_\sigma;x,y)\Bigr).
\end{split}
\]
Every disjunct has an \(\operatorname{Ep}(x,y)\) presentation.
The construction is finite. It uses the predicate letters of the
supplied elementary labels, together with the declared \(P_0\)
tautology or falsehood when an empty-family or empty-disjunction
convention requires it. Fresh letters used to exhibit the source
substitution construction are eliminated.

The given \(I\)-representations and their conjunctions retain
\(\mathsf{Com}\) scope. The elementary midpoint identity is applied
to their actual elementary disjuncts and then restricted to complete
orders. This proves the source's conclusion without promoting an
arbitrary \(I\)-representation to \(\mathsf{Lin}\)
\cite[ch.~III, lemma~5, pp.~78--79]{Kamp1968}.
\end{hblemma}

\begin{hbexample}{def:p09-kamp-midpoint-one-piece}{Two one-interval patterns and their shared point}
Let \(A,B,C,D,E,F\) be unary predicate letters and let
\(H\in\operatorname{E}(w)\). View \(A(w),\ldots,F(w)\) as elementary
labels. The two words \((A,B,C)\) and \((D,E,F)\), joined with
midpoint label \(H\), give
\[
                         (A,B,H\land C\land D,E,F).
\]
With fresh bound variables this pattern is
\[
\begin{split}
 A(x)\land\exists z\bigl(&x<z<y\land
          \forall u\,(x<u<z\to B(u))\\
      &{}\land H(z)\land C(z)\land D(z)\\
      &{}\land\forall v\,(z<v<y\to E(v))\land F(y)\bigr).
\end{split}
\]
All three middle labels concern the same chosen \(z\).
The letter \(E\) in this example is a unary predicate; the class
\(\operatorname{E}\) continues to denote the class of formulas.

On \(0<1<2\), take \(A,B,C,D,E,F\) true everywhere. The formula at
\((0,2)\) holds exactly when \(H\) holds at 1; its two open intervals
are empty. Requiring \(H\) only at either outer endpoint would change
the condition. If all three input families are empty, the conclusion
instead represents precisely \(\exists z\,(x<z<y)\). It is false
at adjacent endpoints and true at \((0,2)\) in this three-point order.
An empty family therefore leaves a substantive strict-witness condition.
\end{hbexample}

\paragraph{The next source boundary.}
Lemma~6 begins on printed p.~79. It fixes one basic pattern and four
strictly increasing points, assumes realizations from the first point
to the second and fourth, and characterizes realization to the third
using a guarded past closure. That separate finite-pattern argument
is the next dependency. Lemma~5 supplies the midpoint composition
above; it does not complete the remaining interval lemmas or
Theorem~III.2 \cite[ch.~III, p.~79]{Kamp1968}.

% Kamp III.2, Lemma 6: recovering a fixed initial endpoint between two realizations.

\subsubsection{A fixed initial endpoint between two realizations}
\label{sec:p09-kamp-iii2-lemma6}

Kamp's Lemma~6 fixes one basic pattern and two realizations beginning at
the same point. Between their final points, a realization with that
fixed beginning exists exactly when some past realization of the pattern
ends there. The proof joins portions of two finite witness sequences
\cite[ch.~III, lemma~6, pp.~79--80]{Kamp1968}.

\begin{hbconvention}{def:p09-kamp-fixed-initial-standing}{One fixed pattern and its guarded past closure}
Fix distinct variable names \(x,y\), an integer \(n\geq1\), and
\(\phi\in\operatorname{Bp}_n(x,y)\). Its occurrence word is
\[
                 \sigma(\phi)=(A_0,B_1,A_1,\ldots,B_n,A_n)
\]
as in \cref{def:p09-kamp-basic-word}. All point and interval letters
have one fixed interpretation in the structure under consideration,
and repeated letters remain allowed. The integer \(n\) counts
intervals; each realizing sequence has exactly \(n+1\) points.

For a linear structure
\(\mathcal M=\langle T,<,(P_i^{\mathcal M})_i\rangle\), write
\(\phi(r,s)\) as a semantic abbreviation for
\(\mathcal M\models\phi[x\mapsto r,y\mapsto s]\), where \(r,s\in T\).
This does not introduce an additional relation symbol.
Define the one-variable formula
\[
 \phi^{\exists<}(y)\defeq
             \exists v\,(v<y\land\phi[v/x]),
\]
where \(v\) is fresh and the replacement concerns only free occurrences
of \(x\). Bound variables are renamed fresh beforehand if needed.
Write \(\phi^{\exists<}(s)\) for its truth with \(y=s\).
The superscript names this formula operation. It is the source's
guarded closure \((\exists x<y)\phi(x,y)\), with the bound initial
endpoint renamed to distinguish it from a separately fixed point.

By the basic-pattern realization equivalence already proved in
\cref{annex:p9.kamp.iii2.lemma1}, for \(r<s\) the assertion
\(\phi(r,s)\) means that the word has a realization from \(r\)
to \(s\). The assertion \(\phi^{\exists<}(s)\) means that some
realization ends at \(s\), with any strictly earlier initial point.
In particular this quantifier is not restricted by a fixed initial
point occurring in another premise
\cite[ch.~III, pp.~63--66, 79]{Kamp1968}.
\end{hbconvention}

\begin{hblemma}{thm:p09-kamp-iii2-lemma6}{proved=annex:p9.kamp.iii2.lemma6}{Kamp's fixed-initial-point lemma}
Let \(\mathcal M\) be any structure over a nonempty strict linear
order, and keep the fixed pattern above. Let
\[
                         a<b<c<d
\]
be four points of its domain. Suppose
\[
                         \phi(a,b),\qquad \phi(a,d).
\]
Then
\[
                     \phi(a,c)
                         \quad\Longleftrightarrow\quad
                                  \phi^{\exists<}(c).
\]
Both premises concern the same formula \(\phi\) with the same
predicate interpretation. The source's four underlined point symbols
\(\underline{t_1},\underline{t_2},\underline{t_3},\underline{t_4}\)
correspond respectively to \(a,b,c,d\); they are distinct from the
object-language variable names.

The statement explicitly assumes a linear order in the source, and
the proof applies throughout \(\mathsf{Lin}\). Completeness, density
and absence of endpoints are not hypotheses. The result therefore also
holds on \(\mathsf{Com}\), without transferring any earlier
complete-order equivalence to incomplete orders
\cite[ch.~III, lemma~6, pp.~79--80]{Kamp1968}.
\end{hblemma}

\begin{hbmethod}{def:p09-kamp-fixed-initial-splice}{Constructing the realization with the fixed beginning}
Suppose the right side of the lemma holds. The input to this finite
construction consists of a realizing sequence
\[
                  \mathbf v=(v_0,\ldots,v_n),\qquad v_n=c,
\]
and realizations \(\mathbf w^-\) from \(a\) to \(b\) and
\(\mathbf w^+\) from \(a\) to \(d\), supplied by the two premises.
Every sequence realizes the same fixed word \(\sigma(\phi)\).

If \(a\leq v_0\), take \(\mathbf w=\mathbf w^+\) and choose the
least index \(k\in\{0,\ldots,n-1\}\) satisfying
\[
                         w_k\leq v_k<w_{k+1}.
\]
If \(v_0<a\), take \(\mathbf w=\mathbf w^-\) and choose the
least index in that same set satisfying
\[
                         v_k\leq w_k<v_{k+1}.
\]
The annex proves existence of an eligible index in each case, and
proves the second splice also valid under its weak comparison
\(v_0\leq a\). Thus equality of the two initial points is covered
by either source argument.

Return the sequence
\[
                    (w_0,\ldots,w_k,v_{k+1},\ldots,v_n).
\]
It has exactly \(n+1\) points and realizes the same word from \(a\)
to \(c\). The index range includes both \(k=0\) and \(k=n-1\).
In the source's numbering from 1, its switch index is \(k+1\),
so this is exactly the printed range \(1,\ldots,n\).
The choice of the least eligible index merely makes the finite
construction definite; the proof works for every index meeting the
indicated source condition.

The method is a construction from supplied witnesses. It does not
claim to find arbitrary witnesses in a structure from a formula alone.
It uses only points in the two selected realizations, and it never
changes the number of pattern intervals
\cite[ch.~III, p.~80]{Kamp1968}.
\end{hbmethod}

\begin{hbexample}{def:p09-kamp-splice-overlap}{An old endpoint inside the new joining interval}
Take \(T=\{0,1,\ldots,8\}\) in its usual order and a two-interval
pattern with word \((A_0,B_1,A_1,B_2,A_2)\), using distinct
predicate letters interpreted by
\[
 \begin{gathered}
 A_0=\{0,1\},\qquad A_1=\{2,4\},\qquad A_2=\{3,6,8\},\\
 B_1=\{1,2,3\},\qquad B_2=T.
 \end{gathered}
\]
Here the displayed sets are the predicate interpretations.
Choose \((a,b,c,d)=(0,3,6,8)\).
The realizations \((0,2,3)\) and \((0,2,8)\) supply the two
premises, and \(\mathbf v=(1,4,6)\) supplies the guarded past
closure at 6.

The first case uses \(\mathbf w=(0,2,8)\). Its first eligible
switch is \(k=0\), since \(0\leq1<2\), and the output is
\((0,4,6)\). The new first gap \((0,4)\) contains both 1 and 2.
The old gap \((0,2)\) supplies \(B_1\) at 1, and the old gap
\((1,4)\) supplies it at 2 and 3. Thus the old point \(w_1=2\),
which is now inside the new gap, receives its interval predicate
from the other realization. Its old point predicate \(A_1\)
would not by itself supply that requirement. The annex checks the
whole example and the general joining interval.
\end{hbexample}

\begin{hbexample}{def:p09-kamp-splice-final-label}{The guarded past condition retains the final point label}
For a one-interval word \((A_0,B_1,A_1)\) with distinct predicate letters,
take \(T=\{0,1,2,3\}\) in its usual order. Interpret \(A_0,B_1\)
by \(T\), and \(A_1\) by \(\{1,3\}\).
The four points \(a=0,b=1,c=2,d=3\) satisfy both premises of
the lemma. Nevertheless \(\phi(0,2)\) and
\(\phi^{\exists<}(2)\) are both false, because their final
point must satisfy \(A_1\). The two premises therefore do not
assert that the pattern holds from \(a\) to every intermediate
point without the stated guarded past condition.
\end{hbexample}

\paragraph{The next source boundary.}
Lemma~7 begins on printed p.~80, immediately after this proof.
Its two parts transform guarded formulas with a universal interior
condition into finite alternatives of basic and elementary interval
conditions, with specified restrictions on the lengths of the resulting
disjuncts. Its statement continues on p.~81, and its further proof
is the next separate construction. The present result establishes
only the fixed-pattern assertion of Lemma~6
\cite[ch.~III, pp.~80--81]{Kamp1968}.

% Kamp III.2, Lemma 7: finite reductions across a fixed interior cut.

\subsubsection{Reducing a universal interior condition at a fixed cut}
\label{sec:p09-kamp-iii2-lemma7}

The next construction replaces a basic pattern reaching across an
interior point by finitely many pairs of patterns meeting at that point.
When one endpoint varies under a universal quantifier, a finite choice
of the fixed-side conditions determines the available disjuncts on the
varying side. Kamp's Lemma~7 uses this reduction with an additional
restriction: the varying-side disjunction contains no longer pattern
and at most one pattern as long as the original
\cite[ch.~III, lemma~7, pp.~80--83]{Kamp1968}.

\begin{hbconvention}{def:p09-kamp-cut-standing}{Patterns, lengths and finite aggregates}
Use the basic grammar and occurrence semantics of
\cref{def:p09-kamp-basic-classes,def:p09-kamp-basic-word}.
Fix a basic past pattern with occurrence word
\[
                W=(A_0,B_1,A_1,\ldots,B_n,A_n),\qquad n\geq1.
\]
Predicate letters may repeat. Each occurrence retains its point or
open-interval position. For any word with \(m\) intervals, put
\(\ell(W)=2m+1\); the length \(\ell(\phi)\) of a basic pattern
is the length of its occurrence word. Thus this length counts predicate
occurrences, rather than interval pieces or printed formula symbols.
The source makes this length convention explicit at the start of
Lemma~8, on p.~83.

For distinct endpoint names \(r,s\), write
\(\mathcal B_p(W;r,s)\) for the past-pattern recursion of
\cref{def:p09-kamp-labelled-intervals}, with the letters in \(W\)
as its unary labels. It is a member of \(\operatorname{Bp}_m(r,s)\).
All new bound names are fresh. In particular a pattern is obtained from
its occurrence word without changing its interval count.

For a finite family, an empty conjunction means the previously declared
\(T^K_w=P_0(w)\lor\neg P_0(w)\), and an empty disjunction means
\(F^K_w=P_0(w)\land\neg P_0(w)\), at the indicated variable.
These are formula abbreviations. An empty aggregate is not an additional
stage-zero basic pattern. Nonempty aggregates are associated from left
to right in their specified index order.
\end{hbconvention}

\begin{hbdefinition}{def:p09-kamp-cut-words}{All point cuts and gap cuts of one word}
Let \(\mathcal C_n\) be the finite index set
\[
 \mathcal C_n\defeq
 \{(\mathrm g,k):1\leq k\leq n\}\ \cup\
 \{(\mathrm p,k):1\leq k<n\}.
\]
Order it by increasing \(k\), with the gap tag before the point tag
at the same \(k\). At \(n=1\) it contains only \((\mathrm g,1)\).
There are \(2n-1\) entries. The tag \(\mathrm p\) means an
interior point position and \(\mathrm g\) an open-gap position;
these tags are not predicate letters or object-language constants.
For each entry \(c\), define two words \(W_c^L,W_c^R\).

For \(c=(\mathrm p,k)\), where \(1\leq k\leq n-1\), set
\[
\begin{aligned}
 W_c^L&=(A_0,B_1,A_1,\ldots,B_k,A_k),\\
 W_c^R&=(A_k,B_{k+1},A_{k+1},\ldots,B_n,A_n).
\end{aligned}
\]
The two words share the point label \(A_k\).
For \(c=(\mathrm g,k)\), where \(1\leq k\leq n\), set
\[
\begin{aligned}
 W_c^L&=(A_0,B_1,A_1,\ldots,B_{k-1},A_{k-1},B_k,B_k),\\
 W_c^R&=(B_k,B_k,A_k,B_{k+1},A_{k+1},\ldots,B_n,A_n).
\end{aligned}
\]
At \(k=1\), the first word is \((A_0,B_1,B_1)\); at \(k=n\),
the second is \((B_n,B_n,A_n)\). Each duplicate \(B_k\) has a
specific role: one occurrence labels the new endpoint and the other
labels its adjoining gap. This requires the old interval predicate at
the cut point itself.

Finally define the two basic formulas
\[
 L_c(r,s)\defeq\mathcal B_p(W_c^L;r,s),\qquad
 R_c(r,s)\defeq\mathcal B_p(W_c^R;r,s).
\]
Their displayed arguments identify their free endpoints; they are formula
abbreviations, not new relation symbols. Every word used here has at
least one interval. A point cut gives respectively \(k,n-k\)
intervals, and a gap cut gives \(k,n-k+1\).

These are the source's two word operations \(R_i^*,S_i^*\), written
with explicit occurrence positions. A split before an interior point
duplicates its point letter between the two words. A split before an
interval occurrence duplicates its interval letter inside each adjoining
word. Occurrence positions determine the case even when the same letter
also appears in the other kind of position
\cite[ch.~III, pp.~64, 81--82]{Kamp1968}.
\end{hbdefinition}

\begin{hblemma}{thm:p09-kamp-cut-factorization}{proved=annex:p9.kamp.cut.factorization}{Factorization at any strict interior point}
Let \(\phi\in\operatorname{Bp}_n(r,t)\) have word \(W\), and
let \(s\) be a distinct variable. Then
\[
 \mathsf{Lin}\models r<s<t\longrightarrow
 \left(\phi\longleftrightarrow
       \bigvee_{c\in\mathcal C_n}(L_c(r,s)\land R_c(s,t))\right).
\]
For each side, all listed basic patterns have length at most
\(2n+1\). Exactly one candidate on the left has that length,
namely \(L_{(\mathrm g,n)}\), and exactly one on the right has
that length, namely \(R_{(\mathrm g,1)}\). Every other candidate
on its respective side is shorter.

The equivalence holds on every nonempty strict linear order. There is
no density, completeness or endpointlessness assumption. In particular
the cut is an already supplied domain point, not a point created inside
an empty interval.
\end{hblemma}

\begin{hbmethod}{def:p09-kamp-cut-aggregates}{The finite alternatives for the two orientations}
Let \(x,y,u,v\) be distinct variable names, and let \(\psi_1,\psi_2\)
be arbitrary formulas of \(L\). Additional free parameters are allowed
and are left unchanged. The displayed outer \(\exists u\) and
inner \(\forall v\) have exactly their displayed scopes, so the
inner quantifier does not bind any free occurrence of \(v\) in
\(\psi_1\). Bound variables internal to newly constructed patterns
are chosen fresh from both formulas and from all four named variables.

List every subset \(J\subseteq\mathcal C_n\), including the empty
subset. Equivalently list all functions
\(\alpha:\mathcal C_n\to\{0,1\}\), with
\(J=\{c:\alpha(c)=1\}\). There are \(2^{2n-1}\) indexed
alternatives before any repetitions are removed.

For a pattern \(\phi(v,y)\) spanning the left interior interval,
form
\[
 \Sigma_J^L(u,y)\defeq\bigwedge_{c\in J}R_c(u,y),\qquad
 \rho_J^L(v,u)\defeq\bigvee_{c\in J}L_c(v,u).
\]
The empty cases are \(\Sigma_\varnothing^L=T^K_u\) and
\(\rho_\varnothing^L=F^K_v\).
For a pattern \(\phi(x,v)\) spanning the right interior interval,
form
\[
 \Sigma_J^R(x,u)\defeq\bigwedge_{c\in J}L_c(x,u),\qquad
 \rho_J^R(u,v)\defeq\bigvee_{c\in J}R_c(u,v).
\]
Here the empty cases are \(\Sigma_\varnothing^R=T^K_u\) and
\(\rho_\varnothing^R=F^K_v\).

Each conjunction concerns the fixed side of the cut. Each disjunction
concerns the endpoint that will vary under \(\forall v\).
The construction uses only the word \(W\), its two kinds of cuts,
finite conjunction and finite disjunction. In particular it does not
require a decision about the truth of \(\psi_1\) or \(\psi_2\).
Truth assignments to the fixed-side patterns are used to prove that the
complete list of alternatives works in every structure
\cite[ch.~III, pp.~81--83]{Kamp1968}.
\end{hbmethod}

\begin{hblemma}{thm:p09-kamp-iii2-lemma7-left}{proved=annex:p9.kamp.iii2.lemma7.left}{Kamp's left-interior reduction}
Suppose \(\phi\in\operatorname{Bp}_n(v,y)\), with \(n\geq1\),
and use its word to construct the left aggregates above. The formula
\[
 \exists u\left(x<u<y\land\psi_1\land
       \forall v\,(x<v<u\to\psi_2\lor\phi(v,y))\right)
\]
is equivalent over \(\mathsf{Lin}\) to
\[
 \bigvee_{J\subseteq\mathcal C_n}
 \exists u\left(x<u<y\land\psi_1\land\Sigma_J^L(u,y)\land
       \forall v\,(x<v<u\to\psi_2\lor\rho_J^L(v,u))\right).
\]
Every \(\Sigma_J^L\) is a finite conjunction of members of
\(\operatorname{Bp}(u,y)\), and every nonempty \(\rho_J^L\)
is a disjunction of members of \(\operatorname{Bp}(v,u)\).
Every such disjunct is also an elementary interval formula.
The empty aggregate conventions are as stated above.

Each \(\rho_J^L\) has at most one basic disjunct of length
\(\ell(\phi)\), and none of greater length. These bounds apply
to the actual disjunct list, including if different cuts happen to
produce equivalent formulas. No larger replacement bound is used
\cite[ch.~III, lemma~7(a), pp.~80--83]{Kamp1968}.
\end{hblemma}

\begin{hblemma}{thm:p09-kamp-iii2-lemma7-right}{proved=annex:p9.kamp.iii2.lemma7.right}{The right-interior reduction}
Suppose \(\phi\in\operatorname{Bp}_n(x,v)\), with \(n\geq1\),
and use its word to construct the right aggregates above. The formula
\[
 \exists u\left(x<u<y\land\psi_1\land
       \forall v\,(u<v<y\to\psi_2\lor\phi(x,v))\right)
\]
is equivalent over \(\mathsf{Lin}\) to
\[
 \bigvee_{J\subseteq\mathcal C_n}
 \exists u\left(x<u<y\land\psi_1\land\Sigma_J^R(x,u)\land
       \forall v\,(u<v<y\to\psi_2\lor\rho_J^R(u,v))\right).
\]
Each fixed-side condition \(\Sigma_J^R\) is a finite conjunction
of members of \(\operatorname{Bp}(x,u)\), and each nonempty
\(\rho_J^R\) is a finite disjunction of members of
\(\operatorname{Bp}(u,v)\), with the declared empty cases.
Each residual disjunction has at most one disjunct of length
\(\ell(\phi)\), and no disjunct of greater length.

This is the handbook's corrected aggregate formulation of the orientation
in Kamp's Lemma~7(b). The source refers its proof to the analogous left case;
the annex supplies both directions for this orientation in full
\cite[ch.~III, lemma~7(b), p.~81]{Kamp1968}.
\end{hblemma}

\begin{hbexample}{def:p09-kamp-cut-two-intervals}{All three cuts of a two-interval word}
For \(W=(A_0,B_1,A_1,B_2,A_2)\), the complete cut list is
\[
\begin{array}{c|c|c}
 c&W_c^L&W_c^R\\ \hline
 (\mathrm g,1)&(A_0,B_1,B_1)&(B_1,B_1,A_1,B_2,A_2)\\
 (\mathrm p,1)&(A_0,B_1,A_1)&(A_1,B_2,A_2)\\
 (\mathrm g,2)&(A_0,B_1,A_1,B_2,B_2)&(B_2,B_2,A_2).
\end{array}
\]
There are eight subset alternatives. Only the last left word has
length five, and only the first right word has length five.

To see why a single cut need not work for all universally quantified
points, take \(T=\{0,1,2,3,4,5,6\}\) in its usual order and set
\[
 A_0=\{1,2\},\quad B_1=\{3\},\quad A_1=\{2,4\},\quad
 B_2=\{3,4,5\},\quad A_2=\{6\}.
\]
These are the interpretations of five distinct predicate letters.
Let \(x=0,u=3,y=6\), take \(\psi_1=T^K_u\) and
\(\psi_2=F^K_v\). At \(v=1\), the realization \((1,2,6)\)
crosses the second gap; at \(v=2\), the realization \((2,4,6)\)
crosses the first. The point cut is unavailable because \(A_1(3)\)
is false. Neither gap cut alone works for both values of \(v\).
The alternative \(J=\{(\mathrm g,1),(\mathrm g,2)\}\) includes
both required residual disjuncts, and its two fixed-side conjuncts
both hold at \((3,6)\). The annex verifies every predicate test.
\end{hbexample}

\begin{hbexample}{def:p09-kamp-cut-empty-choice}{The empty alternative and the smallest word}
For \(W=(A,B,C)\), there is only the gap cut
\[
                    W^L=(A,B,B),\qquad W^R=(B,B,C).
\]
There are two alternatives, corresponding to the empty subset and
the singleton subset. Both words have the original length three.

If \(\psi_2\) holds throughout the universally tested interval,
the empty alternative applies regardless of the pattern letters.
For example, take \(0<1<2\), let every unary predicate be false,
and take \(\psi_1=T^K_u,\psi_2=T^K_v\). For \(x=0,y=2\)
the witness \(u=1\) makes each of the two original quantified
formulas true. The corresponding empty alternative has true fixed-side
condition and a false residual; the tautological \(\psi_2\)
supplies its universal condition. Every nonempty conjunction of basic
patterns is false in this structure because its initial point predicate
is false. Thus the empty alternative is substantive and cannot be
replaced by requiring a basic pattern on the fixed side.

An empty tested open interval likewise leaves the universal condition
vacuously true. The outer strict witness \(x<u<y\) still remains:
if \(x,y\) are adjacent, both displayed existential formulas and
every constructed alternative are false.
\end{hbexample}

\paragraph{Use in the following induction.}
The cut theorem and both universal reductions have the source's explicit
\(\mathsf{Lin}\) scope. Their restrictions to \(\mathsf{Com}\)
are therefore available for the following interval-property argument.
Their at-most-one bound controls the longest varying-side disjuncts
while retaining every shorter case. Lemma~8 begins on printed p.~83 and applies a
well-founded comparison of finite families to establish property
\(I(x,y)\) for a universal interior disjunction. That induction is
the next separate result; no conclusion of Lemma~8 is used here
\cite[ch.~III, pp.~83--84]{Kamp1968}.

% Kamp III.2, Lemma 8: completed induction foundations and boundary decomposition.

\subsubsection{Finite families and the eleven interval boundary configurations}
\label{sec:p09-kamp-iii2-lemma8-foundations}

Kamp's Lemma~8 treats a universal disjunction whose members describe
the current point or a basic pattern connecting it to either endpoint.
Its proof uses a comparison of finite families, separates eleven boundary
configurations of a distinguished pattern, and then applies the preceding
cut construction to obtain elementary interval representations. The results
below supply the family comparison and a complete boundary decomposition
with its bracketing conditions retained
\cite[ch.~III, lemma~8, pp.~83--92]{Kamp1968}.

\begin{hbdefinition}{def:p09-kamp-universal-family}{The exact finite-family problem}
Let \(x,y,z\) be distinct variable names, and let
\[
 F\subseteq\operatorname{At}(z)\cup
                   \operatorname{Bp}(x,z)\cup\operatorname{Bp}(z,y)
\]
be finite. The unstarred \(\operatorname{At}\) is the literal class
of \cref{def:p09-kamp-iii1-components}; it contains no order atoms.
Define
\[
 D_F(z)\defeq\bigvee_{\psi\in F}\psi,\qquad
 \mathcal U_F(x,y)\defeq\forall z\,(x<z<y\to D_F(z)).
\]
An enumeration fixes the association of a nonempty disjunction; for an
empty family use \(D_\varnothing(z)=F^K_z\). Every occurrence
of \(x,y\) inside a member retains its original free-endpoint role.

For \(\psi\in\operatorname{At}\), put \(\ell(\psi)=1\).
For a basic pattern, \(\ell(\psi)\) is its occurrence-word length
from \cref{def:p09-kamp-cut-standing}. In particular a basic pattern
has odd length at least three. For a nonempty family set
\[
 M(F)\defeq\max_{\psi\in F}\ell(\psi),\qquad
 c(F)\defeq\bigl|\{\psi\in F:\ell(\psi)=M(F)\}\bigr|.
\]
For the empty family set \(M(F)=c(F)=0\).
The bars denote finite cardinality.

The source comparison is
\[
 G\prec F\quad\Longleftrightarrow\quad
 \exists m\in\mathbb N\left(
   (\forall\theta\in G\;\ell(\theta)\leq m)\ \land\
   |\{\psi\in F:\ell(\psi)=m\}|>
   |\{\theta\in G:\ell(\theta)=m\}|
                        \right).
\]
It compares finite formula sets by their lengths. It is not semantic
implication and does not count multiple occurrences of an identical
formula as distinct set members
\cite[ch.~III, p.~83]{Kamp1968}.
\end{hbdefinition}

\begin{hblemma}{thm:p09-kamp-family-rank}{proved=annex:p9.kamp.family.rank}{The source comparison and strict decrease}
For finite families in the declared classes,
\[
 G\prec F\quad\Longleftrightarrow\quad
 M(G)<M(F)\ \lor\
       \bigl(M(G)=M(F)\land c(G)<c(F)\bigr).
\]
Thus \(\prec\) is well founded. If \(\phi\) has maximal
length in a nonempty \(F\), removing \(\phi\) strictly
decreases this comparison.

More generally, remove such a \(\phi\), and replace each remaining
basic member \(\rho\) by any finite set of patterns of length at
most \(\ell(\rho)\), at most one of which has length
\(\ell(\rho)\). Retain the other length-one members, and
optionally add finitely many length-one members. If \(M(F)>1\),
every resulting family \(G\) satisfies \(G\prec F\).
This includes empty replacement sets and coincident output formulas.
The numerical condition is exactly the residual restriction proved for
both orientations of Lemma~7.
\end{hblemma}

\begin{hbdefinition}{def:p09-kamp-family-cut-data}{All choices for a finite family crossing one cut}
Let \(\rho_1,\ldots,\rho_k\), with \(k\geq0\), be
a finite list of basic patterns, with \(\rho_j\) having
\(n_j\geq1\) intervals. For each word separately, use the
cut index set \(\mathcal C_{n_j}\) and cut formulas
\(L_{j,c},R_{j,c}\) of \cref{def:p09-kamp-cut-words}.
Define the finite set of choices
\[
 \mathcal J\defeq\prod_{j=1}^k\mathcal P(\mathcal C_{n_j}).
\]
An element \(\mathbf J=(J_1,\ldots,J_k)\) chooses one
subset for each original word. For \(k=0\), \(\mathcal J\)
has one element, the empty tuple. For positive \(k\), its
cardinality is \(2^{\sum_{j=1}^k(2n_j-1)}\).

For the left-interior orientation \(\rho_j(v,y)\), define
\[
\begin{aligned}
 S_{\mathbf J}^L(u,y)&\defeq
                  \bigwedge_{j=1}^k\ \bigwedge_{c\in J_j}R_{j,c}(u,y),\\
 R_{\mathbf J}^L(v,u)&\defeq
                  \bigvee_{j=1}^k\ \bigvee_{c\in J_j}L_{j,c}(v,u).
\end{aligned}
\]
For the right-interior orientation \(\rho_j(x,v)\), define
\[
\begin{aligned}
 S_{\mathbf J}^R(x,u)&\defeq
                  \bigwedge_{j=1}^k\ \bigwedge_{c\in J_j}L_{j,c}(x,u),\\
 R_{\mathbf J}^R(u,v)&\defeq
                  \bigvee_{j=1}^k\ \bigvee_{c\in J_j}R_{j,c}(u,v).
\end{aligned}
\]
An empty total conjunction is \(T^K_u\) and an empty total
disjunction is \(F^K_v\). The two levels of indices retain
the originating pattern for every residual. \(S_{\mathbf J}\)
and \(R_{\mathbf J}\) are aggregate formula abbreviations;
they are distinct from the individual cut formulas and from the
truth set \(R_{a,b}\) introduced below.
\end{hbdefinition}

\begin{hblemma}{thm:p09-kamp-family-cut}{proved=annex:p9.kamp.family.cut}{The complete finite-family cut reduction}
Let \(x,y,u,v\) be distinct variable names and let
\(\psi_0,\psi_1\) be arbitrary formulas, with all new
internal bound names made fresh. For the left-interior patterns,
the following two formulas are equivalent over \(\mathsf{Lin}\):
\[
 \exists u\left(x<u<y\land\psi_1\land
   \forall v\left(x<v<u\to\psi_0\lor\bigvee_{j=1}^k\rho_j(v,y)\right)\right)
\]
and
\[
 \bigvee_{\mathbf J\in\mathcal J}
 \exists u\left(x<u<y\land\psi_1\land S_{\mathbf J}^L(u,y)\land
       \forall v\,(x<v<u\to\psi_0\lor R_{\mathbf J}^L(v,u))\right).
\]
For the right-interior patterns the corresponding equivalence is
between
\[
 \exists u\left(x<u<y\land\psi_1\land
   \forall v\left(u<v<y\to\psi_0\lor\bigvee_{j=1}^k\rho_j(x,v)\right)\right)
\]
and
\[
 \bigvee_{\mathbf J\in\mathcal J}
 \exists u\left(x<u<y\land\psi_1\land S_{\mathbf J}^R(x,u)\land
       \forall v\,(u<v<y\to\psi_0\lor R_{\mathbf J}^R(u,v))\right).
\]
Both statements include \(k=0\), empty individual choices,
and a varying cut at different universal witnesses. Each original
\(\rho_j\) contributes no longer residual and at most one
residual of its original length. Thus the comparison decrease of
\cref{thm:p09-kamp-family-rank} applies when the distinguished
maximal member has already been removed and the other retained
members satisfy its stated length conditions. For that application the
list enumerates each replaced family member once, and the actual
members represented by \(\psi_0\) must satisfy the rank lemma's
retention conditions. The formula equivalences themselves permit
arbitrary \(\psi_0\). Every fixed-side
conjunct remains in the output; the rank statement concerns the
residual family, not those additional conjuncts.

The method enumerates the displayed finite product and constructs
every aggregate. It performs exactly one word-cut construction per
original crossing pattern, providing the finite-family step used
when source pp.~90--91 repeat Lemma~7.
\end{hblemma}

\begin{hblemma}{thm:p09-kamp-family-normalization}{proved=annex:p9.kamp.family.normalization}{Atomic base, fresh occurrences and a removable disjunct}
If \(F\subseteq\operatorname{At}(z)\), including \(F=\varnothing\),
then \(\mathcal U_F\) has property \(I(x,y)\). An actual
elementary representative on ordered endpoints is the one-interval word
\[
                         (T^K_w,D_F(w),T^K_w).
\]

For an arbitrary finite \(F\), one can construct a family \(G\)
of the same size with the same length counts, in which every predicate
letter occurs at most once across the entire family, together with a
finite proper substitution sequence taking \(\mathcal U_G\)
back to \(\mathcal U_F\). Each \(\operatorname{At}\) member
is first represented by one fresh unary atom; each occurrence in a
basic pattern receives its own fresh letter. The substitutions restore
the complete original \(\operatorname{At}\) formulas and original
basic-pattern letters. Consequently an \(I(x,y)\) representation
for \(\mathcal U_G\) yields one for \(\mathcal U_F\).

If \(M(G)>1\), choose a further fresh letter \(Q\) and put
\(G^Q=G\cup\{Q(z)\}\). Then \(M(G^Q)=M(G)\), \(c(G^Q)=c(G)\),
and the two families have exactly the same \(\prec\)-predecessors.
Properly substituting \(F^K_w\) for \(Q\) in
\(\mathcal U_{G^Q}\) recovers a formula equivalent to
\(\mathcal U_G\). Hence an \(I(x,y)\) representation for
\(\mathcal U_{G^Q}\) also gives one for \(\mathcal U_G\).
This is the corrected falsehood substitution for the fresh-disjunct
step on source p.~84.
\end{hblemma}

\begin{hbdefinition}{def:p09-kamp-fixed-endpoint-limits}{Accumulation next to the fixed endpoint}
For a word \(W=(A_0,B_1,A_1,\ldots,B_n,A_n)\), define
\[
\begin{aligned}
 C_W^+(x)&\defeq\forall v\bigl(x<v\to
       \exists z\,(x<z<v\land\mathcal B_p(W;x,z))\bigr),\\
 C_W^-(y)&\defeq\forall v\bigl(v<y\to
       \exists z\,(v<z<y\land\mathcal B_p(W;z,y))\bigr).
\end{aligned}
\]
The first says that realizations beginning at \(x\) end
arbitrarily close to its right; the second says that realizations
ending at \(y\) begin arbitrarily close to its left.
They retain the fixed endpoint in the basic pattern.

Define labels \(\lambda_W^+(w),\lambda_W^-(w)\) as follows.
If \(n=1\), use respectively
\(\partial^+A_1(w)\) and \(\partial^-A_0(w)\).
If \(n>1\), use the actual \(\operatorname{Lf}_{n-1}(w)\)
formula with suffix word \((A_1,B_2,A_2,\ldots,B_n,A_n)\)
and the actual \(\operatorname{Lp}_{n-1}(w)\) formula with
prefix word \((A_0,B_1,A_1,\ldots,B_{n-1},A_{n-1})\),
as defined in \cref{def:p09-kamp-limiting-classes}.
Both labels belong to \(\operatorname{E}(w)\). Their
positive-stage versions require the whole indicated suffix or prefix
strictly inside every corresponding interval adjoining \(w\).
\end{hbdefinition}

\begin{hblemma}{thm:p09-kamp-fixed-endpoint-limits}{proved=annex:p9.kamp.fixed.endpoint.limits}{Elementary representations for accumulation at the fixed endpoint}
For every word \(W\) with \(n\geq1\), both \(C_W^+(x)\)
and \(C_W^-(y)\) have property \(I(x,y)\).
Explicit elementary representatives, with all labels written at their
common free argument \(w\), are
\[
\begin{aligned}
 V_W^+&=(A_0\land\lambda_W^+,\ B_1,\ T^K_w,\ T^K_w,\ T^K_w),\\
 V_W^-&=(T^K_w,\ T^K_w,\ T^K_w,\ B_n,\ A_n\land\lambda_W^-).
\end{aligned}
\]
More precisely,
\[
\begin{aligned}
 \mathsf{Lin}\models x<y\to
    (C_W^+(x)&\longleftrightarrow\mathcal B_p(V_W^+;x,y)),\\
 \mathsf{Lin}\models x<y\to
    (C_W^-(y)&\longleftrightarrow\mathcal B_p(V_W^-;x,y)).
\end{aligned}
\]
Each representative has exactly two interval pieces. The distinction
between a one-letter limiting label at \(n=1\) and a positive-stage
limiting word at \(n>1\) is retained. The first construction
supplies the source's \(\alpha_1\) step on pp.~89--90;
the second is its explicit opposite-endpoint version. Neither is a
claim about accumulation at the other endpoint while keeping a remote
initial point fixed.
\end{hblemma}

\begin{hbdefinition}{def:p09-kamp-boundary-set}{The distinguished truth set and its two orientations}
Choose a maximal-length basic member \(\phi\) of \(F\), and
put \(D(z)=D_{F\setminus\{\phi\}}(z)\). There are two
orientations. With ordered domain endpoints \(a<b\), define
\[
\begin{array}{c|c|c}
 \text{member}&H(t)&P(t)\\ \hline
 \phi(x,z)&\phi(a,t)&\exists r<t\;\phi(r,t)\\
 \phi(z,y)&\phi(t,b)&\exists r>t\;\phi(t,r).
\end{array}
\]
The table uses semantic endpoint assignment; the existential forms are
the guarded past or future closures of the same fixed word. Their
object-language versions use fresh bound endpoint names. They do not
constrain the new endpoint by \(a\) or \(b\) beyond the displayed
strict guard. Let
\[
                   R_{a,b}\defeq\{t:a<t<b\text{ and }H(t)\}.
\]
The symbol \(R_{a,b}\) is a set of domain points; it is distinct
from the cut formulas \(R_c\) of Lemma~7 and the class
\(\operatorname{E}\).
Point-argument displays below are satisfaction schemata. Their formula
versions use the original free endpoint variables in place of \(a,b\)
and fresh variables for the displayed quantified points; they add no
constants or relation symbols to \(L\). A validity display with
\(a,b\) means that its schema holds for every endpoint assignment.

For the following formulas, \(D(t)\) means evaluation at the
current-point variable \(z=t\), retaining both original endpoints
and all their occurrences inside \(D\). Write
\[
 \mathcal V_D(r,s)\defeq\forall t\,(r<t<s\to D(t)),\qquad
 \mathcal V_{D\lor P}(r,s)\defeq
                         \forall t\,(r<t<s\to D(t)\lor P(t)).
\]
These are formula abbreviations with fresh quantified variables, or
their semantic evaluations. Changing \(r,s\) changes only the
outer interval guard: it does not replace the original free \(x,y\)
inside \(D\) or the fixed endpoint inside \(H\).
\end{hbdefinition}

\begin{hbdefinition}{def:p09-kamp-boundary-types}{Exact lower and upper boundary conditions}
Assume \(R_{a,b}\ne\varnothing\). A lower boundary \(l\)
has one of the following three types:
\[
\begin{array}{c|l}
 0&l=a\text{ and }a=\inf R_{a,b},\\
 1&a<l<b\text{ and }l=\min R_{a,b},\\
 2&a<l<b,\quad l=\inf R_{a,b},\quad l\notin R_{a,b}.
\end{array}
\]
An upper boundary \(h\) has the corresponding types
\[
\begin{array}{c|l}
 0&h=b\text{ and }b=\sup R_{a,b},\\
 1&a<h<b\text{ and }h=\max R_{a,b},\\
 2&a<h<b,\quad h=\sup R_{a,b},\quad h\notin R_{a,b}.
\end{array}
\]
Infimum and supremum here are order-theoretic extrema, including
whether the boundary belongs to the set.

These are first-order conditions in the present setting. For example,
type 2 at the lower boundary is expressed by
\[
\begin{split}
 a<l<b\ \land\ \neg H(l)
 &\land\forall t\,(a<t<l\to\neg H(t))\\
 &\land\forall s\bigl(l<s\to
        \exists t\,(l<t<b\land t<s\land H(t))\bigr).
\end{split}
\]
The final quantifier ranges over all \(s>l\), so it includes
\(s=b\) and supplies an actual member. Type 2 at the upper
boundary is expressed by
\[
\begin{split}
 a<h<b\ \land\ \neg H(h)
 &\land\forall t\,(h<t<b\to\neg H(t))\\
 &\land\forall s\bigl(s<h\to
        \exists t\,(a<t<h\land s<t\land H(t))\bigr).
\end{split}
\]
Type 0 uses the corresponding coinitial or cofinal clause at \(a\)
or \(b\), with no point-membership test there. Type 1 uses
\(H(l)\) and absence of earlier members, or \(H(h)\) and
absence of later members. The annex proves every equivalence and makes
the type-0 and type-1 formulas explicit. The original fixed endpoint
in \(H\) remains fixed throughout.
\end{hbdefinition}

\begin{hblemma}{thm:p09-kamp-eleven-boundaries}{proved=annex:p9.kamp.eleven.boundaries}{The eleven complete-order configurations}
On every order in \(\mathsf{Com}\), for every \(a<b\), exactly
one of the following configurations holds:
\[
\begin{array}{c|c|c|l}
 \text{source case}&\text{lower type}&\text{upper type}&\text{additional condition}\\ \hline
 1&\multicolumn{2}{c|}{\text{empty}}&R_{a,b}=\varnothing\\
 2&0&0&l=a<h=b\\
 3&0&1&l=a<h<b\\
 4&0&2&l=a<h<b\\
 5&\multicolumn{2}{c|}{\text{singleton}}&R_{a,b}=\{c\},\ a<c<b\\
 6&1&0&a<l<h=b\\
 7&1&1&a<l<h<b\\
 8&1&2&a<l<h<b\\
 9&2&0&a<l<h=b\\
 10&2&1&a<l<h<b\\
 11&2&2&a<l<h<b.
\end{array}
\]
The bounds in each nonempty case are unique. The nine two-boundary
cases have \(l<h\); the coincident minimum and maximum belong to
the separate singleton case. These are the eleven configurations
organized by the source's \(\gamma_1,\ldots,\gamma_{11}\),
with all bounds interpreted using the same original \(H\).
The exhaustiveness assertion requires \(\mathsf{Com}\).
\end{hblemma}

\begin{hblemma}{thm:p09-kamp-hull-interpolation}{proved=annex:p9.kamp.hull.interpolation}{The fixed endpoint inside its truth-set bounds}
In either orientation, suppose the nonempty set \(R_{a,b}\)
has infimum \(l\), supremum \(h\), and \(l<h\).
Then for every \(l<t<h\),
\[
                            H(t)\longleftrightarrow P(t).
\]
There is a member of \(R_{a,b}\) strictly before \(t\) and
one strictly after \(t\); these are precisely the bracketing
premises needed for the fixed-endpoint splice. The assertion holds
on any linear order where the stated bounds exist. It does not
infer existence of those bounds on all of \(\mathsf{Lin}\).
\end{hblemma}

\begin{hbdefinition}{def:p09-kamp-corrected-boundary-branches}{The complete guarded branch formulas}
Let \(\Gamma_1\) assert emptiness of \(R_{a,b}\), and let
\(\Gamma_5(c)\) assert \(R_{a,b}=\{c\}\), with \(a<c<b\).
For each of the other nine rows let
\(\Gamma_i(l,h)\) assert exactly its lower and upper types and
its displayed strict inequalities. These symbols abbreviate the
first-order conditions just defined; they introduce no new predicates.

Define
\[
\begin{aligned}
 Z_1&\defeq\Gamma_1\land\mathcal V_D(a,b),\\
 Z_5&\defeq\exists c\bigl(\Gamma_5(c)\land
                     \mathcal V_D(a,c)\land\mathcal V_D(c,b)\bigr).
\end{aligned}
\]
For the other rows, with lower type \(\lambda\) and upper type
\(\mu\), define the endpoint tests, for \(j\in\{0,1,2\}\), by
\[
 E_j(s)\defeq
 \begin{cases}
 T^K_s,&j=0\text{ or }j=1,\\
 D(s),&j=2.
\end{cases}
\]
Next define the formula abbreviation
\[
\begin{split}
 K_i(l,h)\defeq{}&\Gamma_i(l,h)\land\mathcal V_D(a,l)
       \land E_\lambda(l)\\
 &{}\land\mathcal V_{D\lor P}(l,h)
       \land E_\mu(h)\land\mathcal V_D(h,b),
\end{split}
\]
and set
\[
 Z_i\defeq
 \begin{cases}
 K_i(a,b),&\lambda=0,\ \mu=0,\\
 \exists h\,K_i(a,h),&\lambda=0,\ \mu\ne0,\\
 \exists l\,K_i(l,b),&\lambda\ne0,\ \mu=0,\\
 \exists l\exists h\,K_i(l,h),&\lambda\ne0,\ \mu\ne0.
 \end{cases}
\]
The arguments of \(K_i\) replace only the displayed boundary
parameters \(l,h\); the original endpoints inside its constituents
stay fixed. Thus a type-0 boundary is its outer endpoint and is not
quantified. All new bound names are fresh from the original formulas.

Thus an attained interior boundary is already covered by \(H\),
an unattained interior boundary must satisfy \(D\), and an outer
endpoint is outside the original universal test. Both side intervals
require \(D\). All bracketing conditions remain in \(\Gamma_i\).
\end{hbdefinition}

\begin{hblemma}{thm:p09-kamp-corrected-boundary-decomposition}{proved=annex:p9.kamp.corrected.boundary.decomposition}{The full eleven-branch universal decomposition}
For either orientation of the distinguished maximal basic member,
\[
 \mathsf{Com}\models a<b\longrightarrow
 \left(\forall t\,(a<t<b\to D(t)\lor H(t))
                 \longleftrightarrow\bigvee_{i=1}^{11}Z_i\right).
\]
Each branch retains its actual boundary descriptor and all required
endpoint tests. Conditional on a descriptor, the corresponding branch
equivalence holds over \(\mathsf{Lin}\); complete orders are used
to guarantee that some descriptor exists. This is a corrected
boundary decomposition for the case analysis of source pp.~85--89.
\end{hblemma}

\begin{hbexample}{def:p09-kamp-bracketing-required}{A later past realization does not keep the original beginning}
Take \(T=\{0,1,2,3,4,5\}\) in its usual order, a one-interval
word \((A,B,C)\), and a remaining atom \(Q(z)\), with
\[
 A=\{0,3\},\qquad B=T\setminus\{2\},\qquad
 C=\{1,4\},\qquad Q=T\setminus\{4\}.
\]
At \(a=0,b=5\), the fixed-start truth set is \(R_{0,5}=\{1\}\).
The target universal disjunction is false at \(t=4\):
\(Q(4)\) is false and \(\phi(0,4)\) crosses the point 2
where its interval predicate fails. Nevertheless the unrestricted past
closure \(P(4)\) is true via the realization \((3,4)\).

The condition consisting only of a first witness at 1, \(Q\) below
1, and \(Q\lor P\) above 1 is therefore true. It cannot replace
the original formula. The corrected singleton branch instead requires
\(Q\) throughout \((1,5)\) and is false at 4. A branch using
\(Q\lor P\) on a two-boundary interval requires members of the
fixed-start truth set on both sides of each tested point.
\end{hbexample}

\begin{proofboundary}{def:p09-kamp-lemma8-remaining}{From boundary conditions to elementary representations}
The results above supply the comparison, normalization, guarded boundary
decomposition and fixed-endpoint accumulation formulas.
Section~\ref{sec:p09-kamp-iii2-lemma8-completion} supplies the remote
and interior approach formulas, the positive elementary supports, and the
complete smaller-family induction. Its
Theorem~\ref{thm:p09-kamp-iii2-lemma8} proves the full \(I(x,y)\)
conclusion with an explicit finite elementary disjunction.
That construction retains every fixed-side cut condition and the
bracketing needed by Lemma~6. A first-order boundary descriptor does not
acquire property \(I\) merely by being first order; the required
elementary presentations are constructed in that proof.
\end{proofboundary}

% Completion of the corrected elementary-representation induction for Kamp III.2, Lemma 8.

\subsubsection{Elementary boundary formulas and the universal-family induction}
\label{sec:p09-kamp-iii2-lemma8-completion}

The boundary decomposition becomes an elementary interval construction once
its bracketing conditions have actual elementary presentations. We first
construct those presentations, then localize each remaining universal family
at its interval and combine the resulting interval conditions. The final
induction proves the original conclusion of Kamp's Lemma~8 over
\(\mathsf{Com}\), with the corrected bracketing conditions
\cite[ch.~III, pp.~85--92]{Kamp1968}.

\begin{hbdefinition}{def:p09-kamp-order-dual}{Order reversal of a formula or word}
For a formula \(\theta\), let \(\theta^{\mathrm d}\) be obtained
by replacing every order atom \(r<s\) by \(s<r\), retaining
unary atoms and every Boolean operation and quantifier. This is an
operation on formulas, not an additional connective. For a structure
\(\mathcal M\), write \(\mathcal M^{\mathrm{op}}\) for the
same domain and unary interpretations with the order reversed.
For a labelled word \(W=(\alpha_0,\beta_1,\ldots,\beta_n,\alpha_n)\),
put
\[
 W^{\mathrm d}\defeq
 (\alpha_n^{\mathrm d},\beta_n^{\mathrm d},\alpha_{n-1}^{\mathrm d},
                    \ldots,\beta_1^{\mathrm d},\alpha_0^{\mathrm d}).
\]
Each label is reversed as a formula on its free parameter. Bound variables
are made fresh in every subsequent construction.
\end{hbdefinition}

\begin{hblemma}{thm:p09-kamp-elementary-duality}{proved=annex:p9.kamp.elementary.duality}{Order reversal preserves the elementary constructions}
Every \(\theta\in\operatorname{E}(w)\) has an explicitly
constructed \(\operatorname{E}(w)\) presentation of
\(\theta^{\mathrm d}\). Given an elementary word presentation
of \(\rho(x,y)\), reversing its word and its labels gives an
elementary presentation, using the constructed elementary representatives
for its dual labels, satisfying
\[
 \mathcal M\models\mathcal B_p(W^{\mathrm d};b,a)
 \quad\Longleftrightarrow\quad
 \mathcal M^{\mathrm{op}}\models\mathcal B_p(W;a,b)
 \qquad(b<a\text{ in }\mathcal M).
\]
Consequently a supplied property \(I(x,y)\) representation can be
reversed into an \(I(y,x)\) representation of its order dual.
Both \(\mathsf{Lin}\) and \(\mathsf{Com}\) are preserved
by order reversal. Basic words retain their occurrence lengths.
\end{hblemma}

\begin{hbdefinition}{def:p09-kamp-neighbor-labels}{One-sided neighborhoods as elementary labels}
For a unary letter or elementary label \(B\), define
\[
\begin{aligned}
 N_B^+(w)&\defeq\exists u\,(w<u\land\forall t\,(w<t<u\to B(t))),\\
 N_B^-(w)&\defeq\exists u\,(u<w\land\forall t\,(u<t<w\to B(t))).
\end{aligned}
\]
Use \(T^K\) for \(B\) to obtain the strict existence of a
later or earlier point. Each formula has an actual
\(\operatorname{E}(w)\) presentation: take the guarded future
or past closure of the one-interval elementary word \((T^K,B,T^K)\).

For a basic word \(W=(A_0,B_1,\ldots,B_n,A_n)\), keep the
limiting labels \(\lambda_W^+,\lambda_W^-\) of
\cref{def:p09-kamp-fixed-endpoint-limits}, and define
\[
\begin{aligned}
 e_W^+&\defeq\neg N_{T^K}^+\lor
                  (A_0\land N_{B_1}^+\land\lambda_W^+),\\
 e_W^-&\defeq\neg N_{T^K}^-\lor
                  (A_n\land N_{B_n}^-\land\lambda_W^-).
\end{aligned}
\]
All labels have the same free parameter \(w\). Their displayed
formulas stand for the corresponding explicit elementary constructions.
\end{hbdefinition}

\begin{hblemma}{thm:p09-kamp-unary-accumulation}{proved=annex:p9.kamp.unary.accumulation}{Global elementary labels for fixed-endpoint accumulation}
The labels \(e_W^+,e_W^-\) belong to \(\operatorname{E}\)
and, over \(\mathsf{Lin}\), express respectively
\(C_W^+\) and \(C_W^-\) from
\cref{def:p09-kamp-fixed-endpoint-limits}. The negative-neighborhood
disjunct handles a last or first point, where the defining universal
accumulation condition is vacuously true. This includes every
\(n\geq1\) and the separate one-letter limiting case.
\end{hblemma}

\begin{hbdefinition}{def:p09-kamp-remote-limits}{Realizations approaching the remote endpoint}
For a basic word \(W\), put
\[
\begin{aligned}
 B_W^+(x,y)&\defeq\forall v\bigl(v<y\to
  \exists t\,(x<t<y\land v<t\land\mathcal B_p(W;x,t))\bigr),\\
 B_W^-(x,y)&\defeq\forall v\bigl(x<v\to
  \exists t\,(x<t<y\land t<v\land\mathcal B_p(W;t,y))\bigr).
\end{aligned}
\]
Under \(x<y\), the instances \(v=x\) and \(v=y\)
respectively force an actual realization. These definitions are therefore
not vacuous when the outer endpoints are adjacent.

For \(1\leq k\leq n\), let \(p_{W,k}^-(w)\) be the
limiting-past formula for the suffix beginning at \(A_k\).
It is \(\partial^-A_n(w)\) when \(k=n\); otherwise it is
the \(\operatorname{Lp}_{n-k}(w)\) formula with word
\((A_k,B_{k+1},\ldots,B_n,A_n)\).
Let \(p_{W,k}^+(w)\) be the limiting-future formula for the
prefix ending at \(A_{k-1}\): it is \(\partial^+A_0(w)\)
when \(k=1\), and otherwise the corresponding
\(\operatorname{Lf}_{k-1}(w)\) formula. Define elementary words
\[
\begin{aligned}
 V_{W,k}^+&\defeq(A_0,B_1,\ldots,B_{k-1},A_{k-1},B_k,p_{W,k}^-),\\
 V_{W,k}^-&\defeq(p_{W,k}^+,B_k,A_k,B_{k+1},\ldots,B_n,A_n).
\end{aligned}
\]
The first word has \(k\) intervals; the second has \(n-k+1\).
The limiting labels are evaluated at the endpoint where they occur,
in the whole structure.
\end{hbdefinition}

\begin{hblemma}{thm:p09-kamp-remote-limits}{proved=annex:p9.kamp.remote.limits}{The complete remote-boundary representations}
Over \(\mathsf{Lin}\), under the guard \(x<y\),
\[
\begin{aligned}
 B_W^+(x,y)&\longleftrightarrow
             \bigvee_{k=1}^n\mathcal B_p(V_{W,k}^+;x,y),\\
 B_W^-(x,y)&\longleftrightarrow
             \bigvee_{k=1}^n\mathcal B_p(V_{W,k}^-;x,y).
\end{aligned}
\]
In particular both conditions have property \(I(x,y)\).
For a fixed initial endpoint, the first condition supplies the
remote cofinal boundary used as \(\beta_1\) in the source.
Its explicit \(x<t\) guard agrees with that source condition
under \(x<y\), as proved in the annex. There is no assumption
that different realizing sequences share their intermediate points
\cite[ch.~III, pp.~85, 89--90]{Kamp1968}.
\end{hblemma}

\begin{hbdefinition}{def:p09-kamp-interior-approach}{Approaching an interior point from beyond it}
For \(x<u\), define
\[
 A_W^+(x,u)\defeq\forall v\bigl(u<v\to
       \exists t\,(u<t<v\land\mathcal B_p(W;x,t))\bigr).
\]
When \(u\) has a later point, use every cut \(c\in\mathcal C_n\)
of \(W\) to form
\[
                    L_c(x,u)\land e_{W_c^R}^+(u).
\]
The second factor has an \(\operatorname{E}(u)\) presentation.
The finite disjunction of these factors has an \(I(x,u)\)
presentation by endpoint insertion and conjunction.
\end{hbdefinition}

\begin{hblemma}{thm:p09-kamp-interior-approach}{proved=annex:p9.kamp.interior.approach}{The unattained lower-boundary approach formula}
For every \(x<u\) in a linear structure,
\[
 A_W^+(x,u)\longleftrightarrow
 \left(\neg N_{T^K}^+(u)\ \lor\
   \bigvee_{c\in\mathcal C_n}
          (L_c(x,u)\land e_{W_c^R}^+(u))\right).
\]
The displayed formula has property \(I(x,u)\), including when
\(u\) is a last point. In an interior-boundary application
\(x<u<y\), the no-future disjunct is false.
Order reversal gives the corresponding fixed-final-endpoint formula
approaching an interior point from its left, with its explicit
elementary presentation. The finitely many cut alternatives are all
retained; one does not require the same cut at every initial witness.
\end{hblemma}

\begin{hblemma}{thm:p09-kamp-segmented-assembly}{proved=annex:p9.kamp.segmented.assembly}{Assembling interval conditions across finitely many ordered points}
Fix pairwise distinct variable names \(x_0,\ldots,x_m\), with \(m\geq1\),
and the strict order \(x_0<\cdots<x_m\).
Suppose finitely many conditions are given, each either an
\(\operatorname{E}(x_i)\) formula at a point, or a formula
with a supplied property \(I(x_i,x_j)\) representation for
some \(i<j\). Then the conjunction of those conditions,
existentially quantifying \(x_1,\ldots,x_{m-1}\) with the
strict order guard, has property \(I(x_0,x_m)\).

The finite construction expands every supplied interval disjunction,
cuts each elementary word at every intervening named point, distributes
over all cuts, fuses the resulting conditions on each adjacent interval,
and concatenates at the elementary point labels. Empty conjunctions use
the declared truth labels. Every imposed interval edge and every point
condition is retained. Input \(I\) equivalences retain their
\(\mathsf{Com}\) scope throughout this construction.
\end{hblemma}

\begin{hbdefinition}{def:p09-kamp-localization-packages}{Localizing a smaller family at one of the boundary intervals}
Let \(F\) have maximal basic member \(\phi\), write
\(D=D_{F\setminus\{\phi\}}\), and fix ordered endpoints
\(a\leq r<s\leq b\), where equality means that the endpoint
name is literally reused. Let \(\eta\in\operatorname{E}(t)\)
be a chosen additional current-point label, possibly \(F^K_t\).
The current-point variable \(t\) is distinct from all endpoint names.
Localize
\[
             \forall t\,(r<t<s\to D(t)\lor\eta(t))
\]
as follows. For an original \(\rho(a,t)\), retain it when
\(r=a\); otherwise apply all cut choices at \(r\), recording
fixed prefixes on \((a,r)\) and residuals on \((r,t)\).
For an original \(\rho(t,b)\), retain it when \(s=b\);
otherwise record all fixed suffix choices on \((s,b)\) and
residuals on \((t,s)\). Take every tuple of choices for every
crossing input. Retain the original \(\operatorname{At}(t)\)
members.

For each tuple, let \(G\) be the resulting local residual family
together with one fresh atom \(Q(t)\). It belongs to
\(\operatorname{At}(t)\cup\operatorname{Bp}(r,t)\cup
\operatorname{Bp}(t,s)\). The original member \(\phi\) is
absent. A \emph{localization package} consists of all selected
fixed prefix/suffix conjuncts and the local universal formula for
\(G\), properly substituting \(\eta\) for \(Q\).
The fresh letter occurs nowhere in the input or in \(\eta\).
\end{hbdefinition}

\begin{hblemma}{thm:p09-kamp-localization}{proved=annex:p9.kamp.localization}{Localization with an actual smaller-family induction hypothesis}
Assume the Lemma~8 conclusion for every finite family \(G\prec F\)
and every choice of its three endpoint/current-point variables.
Then the preceding localization construction gives a finite disjunction
of packages equivalent over \(\mathsf{Com}\) to the indicated
universal formula on \((r,s)\). Each package contains finitely
many basic fixed-side interval conditions and an actual
\(\operatorname{Ep}(r,s)\) disjunction obtained from the
induction hypothesis and proper substitution.

Every induction call has \(G\prec F\): deletion removes one
maximal member, each other original basic member contributes at most
one residual of its length and none longer, and the fresh atom has
length one. Equal outputs can only decrease the set count. The
elementary label \(\eta\) is inserted after that call; it is
not reclassified as an \(\operatorname{At}\) formula.
\end{hblemma}

\begin{hbdefinition}{def:p09-kamp-positive-boundaries}{Elementary support conditions for the eleven cases}
First take the left orientation \(H(t)=\mathcal B_p(W;a,t)\),
and let \(P(t)\) be its guarded past closure. Use the regional
matrices \(K_i\) and endpoint tests of
\cref{def:p09-kamp-corrected-boundary-branches}, but replace the
full descriptor \(\Gamma_i\) by the following positive supports.
For a two-boundary case with lower type \(\lambda\) and upper
type \(\mu\), the lower supports are
\[
 \begin{array}{c|l}
 0&l=a,\quad e_W^+(a),\\
 1&a<l,\quad\mathcal B_p(W;a,l),\\
 2&a<l,\quad A_W^+(a,l),
 \end{array}
\]
and the upper supports are
\[
 \begin{array}{c|l}
 0&h=b,\quad B_W^+(a,b),\\
 1&h<b,\quad\mathcal B_p(W;a,h),\\
 2&h<b,\quad B_W^+(a,h).
 \end{array}
\]
Retain \(a\leq l<h\leq b\), all outer-region \(D\) tests,
all middle-region \(D\lor P\) tests, and the \(D\) test at
each type-2 interior boundary. Quantify the same interior bounds
as in the eleven-case construction. Call the resulting formulas
\(Y_i\), for \(i=2,3,4,6,7,8,9,10,11\).
Set
\[
\begin{aligned}
 Y_1&\defeq\mathcal V_D(a,b),\\
 Y_5&\defeq\exists c\,(a<c<b\land\mathcal B_p(W;a,c)
                    \land\mathcal V_D(a,c)\land\mathcal V_D(c,b)).
\end{aligned}
\]
These supports retain actual fixed-start witnesses on both sides of
every middle point. They omit absence assertions only where the
unchanged \(D\) conditions already suffice for the target.
Their relation to the full descriptors is proved below, rather than
assumed as an identity of printed source formulas.
\end{hbdefinition}

\begin{hblemma}{thm:p09-kamp-positive-boundaries}{proved=annex:p9.kamp.positive.boundaries}{Correctness of the elementary-supported branch replacement}
For the left orientation, over \(\mathsf{Com}\) and under
\(a<b\),
\[
 \forall t\,(a<t<b\to D(t)\lor H(t))
                  \longleftrightarrow\bigvee_{i=1}^{11}Y_i.
\]
Each positive support has an actual \(\operatorname{E}\) or
\(I\) presentation as supplied above. This assertion does not
discard a bracketing hypothesis: the lower support supplies a
fixed-start witness strictly below every middle point, and the
upper support one strictly above it. Each \(Y_i\) individually
implies the target on every linear order. Complete orders supply
the full boundary descriptor needed to choose a branch in the
forward direction.
\end{hblemma}

\begin{hbmethod}{def:p09-kamp-lemma8-algorithm}{The terminating elementary-representation construction}
Given the exact finite family \(F\) of
\cref{def:p09-kamp-universal-family}, use its pair
\((M(F),c(F))\) for the recursion. For the point-only or empty
family, return the elementary word supplied by
\cref{thm:p09-kamp-family-normalization}. Otherwise choose a maximal
basic member. If it is right-oriented, first reverse the order
presentation and the basic words, perform the following left construction
with that distinguished reversed word, and reverse every resulting
elementary word back.

For the left construction form all eleven \(Y_i\). Replace the
guarded closure \(P\) by its actual \(\operatorname{E}\)
presentation from Lemma~1. For every regional universal condition
use the localization construction, with additional label \(P\)
in the middle regions and \(F^K\) in the outer regions. Recursively
obtain each smaller local family's representation, then perform the
specified proper substitution.

Expand every point condition \(D(l)\) or \(D(h)\) into its
original finite disjunction. Each alternative is either an elementary
point label or one original basic interval condition at its endpoints.
Expand every boundary support and every localization package into its
supplied interval disjunctions, distributing every finite choice.
Each resulting conjunction has only elementary point conditions and
interval edges between its at most four ordered named points.
Apply \cref{thm:p09-kamp-segmented-assembly} to each conjunction
and concatenate all returned disjunctions. No truth decision about the
structure is part of this syntax construction.
\end{hbmethod}

\begin{hblemma}{thm:p09-kamp-iii2-lemma8}{proved=annex:p9.kamp.iii2.lemma8}{Kamp's universal interior-disjunction lemma}
Let \(x,y,z\) be distinct variables and let
\[
 F\subseteq\operatorname{At}(z)\cup\operatorname{Bp}(x,z)
                                      \cup\operatorname{Bp}(z,y)
\]
be finite, including the empty family with the stated false disjunction.
Then
\[
                   \forall z\,(x<z<y\to\bigvee_{\psi\in F}\psi)
\]
has property \(I(x,y)\). The algorithm above returns a finite
disjunction of actual \(\operatorname{Ep}(x,y)\) formulas
equivalent to it under \(x<y\) over \(\mathsf{Com}\).
All original finite-family, variable and order hypotheses are retained.
The full proof includes both orientations and every one of the eleven
boundary cases, with strict decrease at every recursive call
\cite[ch.~III, lemma~8, pp.~83--92]{Kamp1968}.
\end{hblemma}

\begin{hbexample}{def:p09-kamp-remote-single-interval}{The terminal label is a strict limit}
For \(W=(A,B,C)\), the remote cofinal representative is
\[
                      (A,B,\partial^-C).
\]
Thus it requires \(A(a)\), \(B\) throughout \((a,b)\),
and \(C\)-points arbitrarily close below \(b\); it does not
require \(C(b)\). On \(\mathbb R\), take \(a=0,b=1\),
\(A=\{0\}\), \(B=\mathbb R\), and \(C=(0,1)\).
Realizations from 0 finish arbitrarily close below 1 although
\(C(1)\) is false. On the two-point order \(0<1\), the
remote condition is false even if every predicate is true, because
its instance at the initial endpoint requires a point strictly between
0 and 1. The elementary terminal limit is false there as well.
\end{hbexample}

\paragraph{The next source result.}
Lemma~9 begins on printed p.~92. It uses the completed universal-family
lemma to treat elementary point and interval conditions under a universal
interior quantifier. Its substitution and distribution construction is
the next separate result; Lemma~8 above no longer depends on it.

% Kamp III.2, Lemma 9: universal composition of elementary interval conditions.

\subsubsection{Universal composition of elementary interval conditions}
\label{sec:p09-kamp-iii2-lemma9}

Lemma~9 extends the basic universal-family construction to elementary
point conditions and supplied interval representations. The source first
combines each of its three finite disjunctions, then uses Lemma~8 and
proper predicate substitution
\cite[ch.~III, lemma~9, pp.~92--93]{Kamp1968}.

\begin{hbdefinition}{def:p09-kamp-universal-elementary-data}{The three universal families}
Let \(x,y,z\) be pairwise distinct variable names. For finite sizes
\(m,n,p\geq0\), take
\[
\begin{aligned}
 \chi_i&\in\operatorname{E}(z) &&(1\leq i\leq m),\\
 \phi_j&\text{ has property }I(x,z) &&(1\leq j\leq n),\\
 \psi_k&\text{ has property }I(z,y) &&(1\leq k\leq p).
\end{aligned}
\]
Each interval condition has no free variables outside its two indicated
endpoints. A supplied \(I\) representation is part of the construction
data. Let \(\chi,\Phi,\Psi\) be the respective disjunctions,
and define
\[
 U(x,y)\defeq\forall z\,(x<z<y\to\chi(z)\lor\Phi(x,z)\lor\Psi(z,y)).
\]
An empty disjunction is the declared falsehood at an indicated endpoint;
for the unary family it is \(F^K_z\). No empty family deletes the
universal quantifier or either strict inequality. The source's indexed
lists \(0,\ldots,m\), \(0,\ldots,n\), \(0,\ldots,p\)
are its nonempty-family cases; the zero-size convention here is explicit.
\end{hbdefinition}

\begin{hbmethod}{def:p09-kamp-universal-elementary-construction}{Fresh basic words and complete restoration}
Combine all supplied left representations by concatenating their disjunct
lists, and do the same on the right. Obtain the elementary labelled word
of every resulting interval disjunct. The unary disjunction \(\chi\)
has an \(\operatorname{E}(z)\) construction, including its false
empty case.

Choose a fresh unary letter \(Q\) for \(\chi\). At each
occurrence position of every elementary word choose another fresh unary
letter; all chosen letters are distinct and absent from every input
label. Replace each labelled word by the basic past word on those fresh
letters. Apply Lemma~8 to the finite family consisting of \(Q(z)\),
all left basic words at \((x,z)\), and all right basic words at
\((z,y)\).

In its resulting \(I(x,y)\) representation, properly substitute
\(\chi\) for \(Q\), and each original elementary label for its
fresh occurrence letter. Perform these substitutions in a fixed finite
order, with bound variables made fresh before insertion. The replacement
labels contain none of the chosen targets, so later steps do not alter
earlier replacements. Return the resulting finite elementary disjunction.
Every input interval disjunct and every occurrence label is restored.
\end{hbmethod}

\begin{hblemma}{thm:p09-kamp-iii2-lemma9}{proved=annex:p9.kamp.iii2.lemma9}{Kamp's universal elementary-composition lemma}
Under the hypotheses above, \(U(x,y)\) has property \(I(x,y)\).
The stated construction returns actual
\(\operatorname{Ep}(x,y)\) disjuncts whose disjunction is equivalent
to \(U\) under \(x<y\) over \(\mathsf{Com}\).
The claim includes every nonempty family of the source and the explicit
empty-family cases. It uses the given \(I\) equivalences only at
their complete-order scope; no equivalence on incomplete orders is inferred
\cite[ch.~III, lemma~9, pp.~92--93]{Kamp1968}.
\end{hblemma}

\begin{hbexample}{def:p09-kamp-universal-different-disjuncts}{Different disjuncts at different interior points}
Take \(0<1<2<3\), with outer endpoints 0 and 3. Let the unary
condition \(Q(z)\) be false everywhere. Let the left word be
\((A,B,C)\), with \(A=\{0\}\), \(B=T\), \(C=\{1\}\),
and the right word be \((D,E,F)\), with
\(D=\{2\}\), \(E=T\), \(F=\{3\}\), where \(T\)
is this four-point domain. The left disjunct holds at the tested point
1 and fails at 2; the right holds at 2 and fails at 1. Thus the
universal disjunction holds, while neither individual interval disjunct
holds at every interior point. The construction preserves this varying
choice; it does not distribute a universal quantifier over disjunction.

If all three families are empty, the result instead asserts that
\((x,y)\) is empty. On adjacent endpoints it is true; at \((0,3)\)
in this example it is false. The annex verifies both cases.
\end{hbexample}

% Kamp III.2, Lemma 10 and the finite Boolean normalization used in its final application.

\subsubsection{All regional interval components}
\label{sec:p09-kamp-iii2-lemma10}

The recursive \(\operatorname{Bet}\) classes use an existential
conjunction or a universal disjunction over a strict interior point.
Lemmas~5 and 9 handle those two constructors. Lemma~10 therefore
provides the interval property for every finite stage, and conjunction
adds the endpoint literals needed in the concluding proof
\cite[ch.~III, lemma~10, p.~93]{Kamp1968}.

\begin{hbmethod}{def:p09-kamp-bet-construction}{The stage construction for regional interval formulas}
Use the exact unstarred recursion of
\cref{def:p09-kamp-iii1-components}. Stage zero is empty.
At a retained component reuse its earlier representation. At an
existential constructor, separate the finite matrix into its At
conditions at the quantified point and its two Bet families on either
side. Apply Lemma~5 to their supplied representations. At a universal
constructor apply Lemma~9 to those three disjunction families.

The At conditions have elementary histories; the constituent Bet
conditions come from earlier stages. Keep the strict guard
\(x<z<y\), including when a matrix family is empty. For a
conjunction of finished Bet components and At conditions at \(x\)
or \(y\), insert each endpoint formula using the interval endpoint
lemma and apply the complete finite fusion construction to all factors.
An empty outer conjunction is represented by the one-interval truth word.
\end{hbmethod}

\begin{hblemma}{thm:p09-kamp-iii2-lemma10}{proved=annex:p9.kamp.iii2.lemma10}{Kamp's interval-component lemma}
Every member of \(\operatorname{Bet}(x,y)\) has property
\(I(x,y)\). Every finite conjunction of members of
\[
       \operatorname{Bet}(x,y)\cup\operatorname{At}(x)
                                      \cup\operatorname{At}(y)
\]
also has that property. For a supplied finite construction history,
the stated algorithm returns a finite disjunction of actual
\(\operatorname{Ep}(x,y)\) formulas, equivalent under \(x<y\)
over \(\mathsf{Com}\). This includes the declared empty-family
conventions and retains both source quantifier constructors
\cite[ch.~III, lemma~10, p.~93]{Kamp1968}.
\end{hblemma}

\begin{hbdefinition}{def:p09-kamp-interval-boolean}{Positive Boolean expressions in interval components}
An interval Boolean expression on \((x,y)\) is a finite Boolean
combination of members of
\(\operatorname{Bet}(x,y)\cup\operatorname{At}(x)\cup
\operatorname{At}(y)\). It is \emph{positive} when the external
Boolean constructors are only conjunction and disjunction. A negative
unary literal inside an At component remains permitted.
Use \(T^K_x\) and \(F^K_x\) for empty Boolean truth and
falsehood. No claim is made that these Boolean expressions themselves
are single members of \(\operatorname{Bet}\).
\end{hbdefinition}

\begin{hblemma}{thm:p09-kamp-interval-boolean-normalization}{proved=annex:p9.kamp.interval.boolean.normalization}{Finite interval negation normalization}
Every interval Boolean expression has an effectively constructed positive
interval Boolean expression equivalent to it over \(\mathsf{Lin}\).
In particular the negation of a Bet component can be replaced by such
a finite expression without adding a Bef or Aft component or an extra
free variable.

Consequently every interval Boolean expression has property \(I(x,y)\):
apply Lemma~10 and finite fusion to the positive expression. The
normalization supplies explicitly the negation and Boolean distribution
step used by Kamp before his final existential reduction
\cite[ch.~III, pp.~61--62, 93]{Kamp1968}.
\end{hblemma}

\begin{hbexample}{def:p09-kamp-bet-empty-interior}{Existence and universal vacuity are retained}
The stage-one formulas
\[
 \exists z\,(x<z<y\land T^K_z),\qquad
 \forall z\,(x<z<y\to F^K_z)
\]
assert respectively that the open interval has a point and that it
has no point. They are complements on every linear order. The first
has a two-interval truth word as an elementary representative; the
second has the one-interval word \((T^K,F^K,T^K)\).
At adjacent endpoints only the second holds. The normalization keeps
these strict witness and vacuity conditions when moving a negation
across a regional quantifier.
\end{hbexample}

% Completion of Kamp's Theorem III.2 from the full ten-lemma interval construction.

\subsubsection{Completion of the first-order-to-temporal construction}
\label{sec:p09-kamp-iii2-completion}

The final step converts the regional components in the one-variable
standard form into elementary one-variable formulas, then follows their
elementary construction histories to temporal formulas. The initial
reduction on source pp.~60--62 and the conclusion on pp.~93--94 are
both needed: a general Bef or Aft component can contain nested regional
components and can begin with either quantifier
\cite[ch.~III, pp.~60--62, 93--94]{Kamp1968}.

\begin{hblemma}{thm:p09-kamp-one-sided-closure}{proved=annex:p9.kamp.one.sided.closure}{Closing a supplied interval representation on one side}
Let \(x,w\) be distinct variable names.
If \(\rho(x,w)\) has a supplied \(I(x,w)\) representation,
then \(\exists x\,(x<w\land\rho(x,w))\) has an actual
\(\operatorname{E}(w)\) presentation equivalent over
\(\mathsf{Com}\). If \(\rho(w,x)\) has a supplied
\(I(w,x)\) representation, its guarded future existential
\(\exists x\,(w<x\land\rho(w,x))\) has such a presentation.

For the past case, distribute the guarded existential over the given
elementary disjuncts and use the elementary past-closure construction.
For the future case first convert each elementary past word into its
same-word elementary future presentation, then close it on the future
side. Retain every disjunct. The empty disjunction gives \(F^K_w\).
\end{hblemma}

\begin{hbdefinition}{def:p09-kamp-primitive-regions}{The source's primitive regional class}
Let \(\mathcal D(w)\) consist of the members of
\(\operatorname{Bef}(w)\cup\operatorname{Aft}(w)\) that have
no proper subformula occurrence belonging to
\(\operatorname{Bef}\cup\operatorname{Aft}\). This is the
source's class \(\mathcal D\), with the sole free variable
displayed. Its members begin with an existential or universal regional
quantifier. Their matrices contain only At conditions at the quantified
point and Bet conditions between it and the fixed endpoint.

An \emph{existential envelope} also permits finitely many At conditions
at the fixed endpoint. Its two forms are
\[
 \exists x\,(x<w\land\Theta(x,w)),\qquad
 \exists x\,(w<x\land\Theta(w,x)),
\]
where \(\Theta\) is a finite conjunction of Bet components and
At components at either endpoint. This supplementary envelope notation
records the source's final display~(28); it does not enlarge the
syntactic definition of \(\mathcal D\).
\end{hbdefinition}

\begin{hblemma}{thm:p09-kamp-primitive-regions}{proved=annex:p9.kamp.primitive.regions}{Elementary presentations of all primitive regional cases}
Every existential envelope and every member of \(\mathcal D(w)\)
has an effectively constructed \(\operatorname{E}(w)\)
presentation equivalent over \(\mathsf{Com}\).

For an existential envelope, apply Lemma~10 to its matrix, then use
the appropriate one-sided closure. For a universal primitive member,
negate the universal, normalize its complemented interval matrix by
\cref{thm:p09-kamp-interval-boolean-normalization}, keep all resulting
finite disjunctive alternatives, and close their existential envelopes.
Negate the resulting elementary formula to recover the original universal.
The construction does not treat a finite disjunction of primitive
formulas as one primitive formula by definition
\cite[ch.~III, pp.~61--62, 93--94]{Kamp1968}.
\end{hblemma}

\begin{hbmethod}{def:p09-kamp-regional-restoration}{Restoring maximal proper regional occurrences}
For a Bef or Aft formula \(\varphi(w)\), select all maximal proper
subformula occurrences \(\varphi_1(y_1),\ldots,\varphi_k(y_k)\)
that belong to \(\operatorname{Bef}\cup\operatorname{Aft}\).
Maximal means not contained in any larger proper occurrence of that kind.
These are occurrences in the finite syntax tree, including distinct
occurrences with identical text. They are disjoint.

Recursively obtain an elementary presentation for each selected occurrence.
Choose distinct fresh unary letters \(Q_1,\ldots,Q_k\), absent
from the original formula and every obtained presentation, and replace
the selected occurrence \(\varphi_i(y_i)\) by \(Q_i(y_i)\).
The resulting skeleton belongs to \(\mathcal D(w)\).
Construct its elementary presentation using the preceding theorem.
Properly substitute the obtained elementary formula for each \(Q_i\),
renaming its free argument and bound variables as required.
All target letters are absent from all replacements. If \(k=0\),
the formula already belongs to \(\mathcal D(w)\).
\end{hbmethod}

\begin{hblemma}{thm:p09-kamp-regional-elementary}{proved=annex:p9.kamp.regional.elementary}{Elementary presentations of all past and future regional components}
Every member of \(\operatorname{Bef}(w)\cup\operatorname{Aft}(w)\)
has an effectively constructed \(\operatorname{E}(w)\)
presentation equivalent over \(\mathsf{Com}\). The construction
just given terminates by recursion on proper original subformula occurrences.
It includes existential and universal roots, repeated occurrences, nested
quantified arguments, and both temporal directions. All substitutions
preserve the intended binding of the sole free argument of each replaced
regional component \cite[ch.~III, pp.~60--62, 93--94]{Kamp1968}.
\end{hblemma}

\begin{hbdefinition}{def:p09-kamp-translation-range}{The image of the temporal translation}
Keep the source temporal language \(\mathsf{TL}_1\) and its
first-order translation \(D\), with distinguished current-point
variable \(t_0\). Define
\[
             \mathcal C\defeq\{D(\eta):\eta\in\mathsf{TL}_1\}.
\]
This is the source's image class, distinct from
\(\operatorname{E}\). An intermediate elementary formula need
not literally belong to \(\mathcal C\); its constructed temporal
expression \(\eta\) supplies the actual image member \(D(\eta)\).
For another free variable \(w\), use capture-avoiding replacement
of the free \(t_0\) by \(w\), after freshening bound names.
\end{hbdefinition}

\begin{hbmethod}{def:p09-kamp-final-translation}{The complete first-order-to-temporal procedure}
Input an \(L\)-formula \(\varphi\) with exactly one free variable
\(w\). Construct its one-variable standard form by Theorem~III.1.
Its sole-variable prefix is true over linear orders; replace that
prefix by the source truth label. Its other components belong to
\(\operatorname{At}(w)\), \(\operatorname{Bef}(w)\) or
\(\operatorname{Aft}(w)\).

Keep each At component as an elementary seed. Replace each Bef and Aft
component by the elementary presentation just constructed, and keep
the standard form's finite conjunctions and disjunctions. This gives
an actual \(\operatorname{E}(w)\) construction history.
Apply the elementary-to-temporal procedure of Lemma~3 to that history.
Return its \(\mathsf{TL}_1\) formula \(\eta\) and the literal
image member \(D(\eta)\). Every proper-substitution placeholder
has been eliminated before this final output.
\end{hbmethod}

\begin{hbtheorem}{thm:p09-kamp-iii2}{proved=annex:p9.kamp.iii2}{Kamp's expressive-completeness theorem over complete linear orders}
For every \(L\)-formula \(\varphi(w)\) with exactly one free
variable, the stated procedure constructs a temporal formula
\(\eta\in\mathsf{TL}_1\) such that
\[
          \mathsf{Com}\models
                \varphi(w)\longleftrightarrow D(\eta)[w/t_0].
\]
In particular every tense-defining formula at \(t_0\) is
\(\mathsf{Com}\)-equivalent to a member of \(\mathcal C\),
the source's Theorem~III.2. One formula works uniformly for every
complete linear order, every unary interpretation and every current
point. Density and absence of endpoints are not assumed.

The temporal constructors are exactly negation, conjunction and the
source's strict \(S,U\), with their witness operand first and
their open-interval operand second. Thus they express every uniformly
first-order-definable tense over \(\mathsf{Com}\).
For a positive finite arity \(n\), if the allowed unary signature
is \(P_0,\ldots,P_{n-1}\), the output can use only the corresponding
\(q_0,\ldots,q_{n-1}\). The literal source constant convention
may use \(P_0\) even when it was absent from the input formula.
The proof establishes an effective terminating construction; no
output-size or running-time bound is asserted here
\cite[ch.~III, thm.~2, p.~60 and pp.~93--94]{Kamp1968}.
\end{hbtheorem}

\begin{hbexample}{def:p09-kamp-final-universal-example}{A universal past condition and its finite disjunction}
Let \(Q,R\) be unary letters, with temporal counterparts \(q,r\).
For
\[
 \varphi(w)=\forall z\,(z<w\to\neg Q(z)\land\neg R(z)),
\]
the complemented universal gives
\[
 \neg\varphi(w)\longleftrightarrow
       \exists z\,(z<w\land Q(z))\ \lor\
       \exists z\,(z<w\land R(z)).
\]
With \(\tau=\neg(q\land\neg q)\), a temporal expression
for \(\varphi\) is
\[
                        \neg S(q,\tau)\land\neg S(r,\tau).
\]
At a first point both strict past existentials are false, so this
formula is true, matching universal vacuity. If the sole predecessor
in a two-point order satisfies either \(Q\) or \(R\), it is
false. The disjunctive alternatives produced during normalization
must therefore be retained.
\end{hbexample}

\paragraph{End of the source construction.}
The concluding argument ends on printed p.~94. Chapter~IV begins on
p.~95 with other expressibility results. The ten interval lemmas and
the preliminary reduction now supply the complete constructive argument
for the stated Theorem~III.2.

% The explicit primitive-recursive strengthening of Kamp's constructive translation.

\subsubsection{A primitive-recursive translation with explicit bounds}
\label{sec:p09-kamp-primitive-recursive-translation}

Chapter~II Theorem~2 asserts a primitive-recursive function
\(\operatorname{Exp}\) from tense-defining formulas to temporal
schemata. The preceding proof supplied the translation's logical
correctness. We now specify finite codes, deterministic choices and
primitive-recursive bounds for its construction. The source states the
strengthening on printed p.~28; its formula and schema conventions
are on pp.~21--23 \cite[ch.~II, pp.~21--23, 27--29]{Kamp1968}.

\begin{hbdefinition}{def:p09-kamp-pr-codes}{Natural-number codes for finite syntax}
Encode a finite list of natural-number tokens by
\[
 [\ ]^\#=0,\qquad
 [a_0,\ldots,a_{r-1}]^\#
       =2^{a_0}\bigl(2[a_1,\ldots,a_{r-1}]^\#+1\bigr).
\]
The decoded token list is written in prefix order. First-order tags are
\[
\begin{array}{c|l}
 0,i,j&P_i(t_j)\\
 1,i,j&t_i<t_j\\
 2&\neg\text{ followed by its formula}\\
 3,4,7,8&\land,\lor,\to,\leftrightarrow\text{ followed by two formulas}\\
 5,i\ ;\ 6,i&\exists t_i\ ;\ \forall t_i\text{ followed by a formula}.
\end{array}
\]
Temporal tags are \(0,i\) for \(q_i\), tag 1 for negation,
and tags \(2,3,4\) for conjunction, \(S,U\), followed by
their two operands. Tag numbers are syntactic data, not new operators.
Implication and biconditional are expanded before normal-form processing.
The source language remains equality-free.

For a decoded list \(a\), define its weight
\[
                  w(a)=1+\operatorname{length}(a)+\sum_j a_j.
\]
It bounds its token count and every name index. Intermediate finite
records use tags 9 through 31 and length-delimited inline fields in
the same flat-list coding. Formulas and construction histories are
stored inline, not as unexplained oracle objects. The weight of data
means the weight of this flattened serialization.
\end{hbdefinition}

\begin{hblemma}{thm:p09-kamp-pr-codes}{proved=annex:p9.kamp.pr.codes}{Primitive-recursive finite syntax operations}
Encoding, decoding, parsing, free-variable tests, occurrence lists,
lexicographic comparison, capture-avoiding renaming, finite substitution,
and the finite enumerations used below are primitive recursive on these
codes. A token list of weight \(s\) has code less than \(2^s\).
Malformed codes or inputs with a free-variable set other than
\(\{t_0\}\) receive the fixed default output \(q_0\).
This makes the eventual numerical function total without extending
its semantic assertion to invalid inputs.
\end{hblemma}

\begin{hbdefinition}{def:p09-kamp-pr-deterministic}{The deterministic implementation}
Use the proved translation algorithms in the preceding sections with
the following fixed choices. Scan input syntax left to right; list
subformula occurrences in postorder. Associate new Boolean lists to
the left. Sort finite families by their token lists and remove exact
duplicates. Choose the first maximal-length member. List cuts in the
already defined gap/point order, subsets by their binary masks, and
tuples lexicographically. Enumerate ordered partitions and fusions by
their finite relation matrices in binary order and retain exactly the
matrices passing their finite defining checks.

Fresh indices are the first indices above every index forbidden in
the current operation. In a pure Lemma~8 child family the one fresh
predicate is one above the maximum predicate index of its parent
family. This choice is local to that child computation; its placeholder
is eliminated before outputs from separate children are combined.
Bound variables are freshened before insertion. Pure basic family
members are stored as their orientation and predicate word, with
canonical endpoint slots; At members retain their literal trees.
Expansion to an actual basic formula uses its declared recursion.

At a Lemma~8 node, compute each of its at most 33 regional-location
types and each tuple of cut choices once, before distributing the
products of their returned packages. Store the returned lists and copy
them when forming those products; distribution does not launch new
recursive calls. The other structural procedures visit only their
original input nodes, never newly generated output as a new recursive
input. Operation tags 9 through 20 mean, in order: input, standard
form, interval negation, interval representation, Lemma~8, localization,
assembly, primitive region, regional restoration, elementary translation,
\(D\)-image and padding. Tags 21 and 22 mean return and halt.
A frame is the flat record
\[
 (\text{tag},\text{phase},\text{input length},\text{input tokens},
   \text{child count},\text{length-delimited children},
   \text{next child},\text{length-delimited results}).
\]
Phase 0 constructs the ordered child list by that operation's stated
recipe; phase 1 evaluates the next child and appends its result;
phase 2 performs the operation's local combination and returns.
An operation with no children goes directly from phase 0 to phase 2.
The input frame schedules standard form, regional restoration,
elementary translation, padding and the \(D\)-image in that order.
Return stores both the temporal output and its image; \(\operatorname{Exp}\)
projects the former. A halted state remains fixed. These rules fix
the evaluation order of the previously specified constructions.
\end{hbdefinition}

\begin{hbdefinition}{def:p09-kamp-pr-local-envelope}{A uniform finite-operation envelope}
For a unary function \(f\), let
\(f^{[0]}(s)=s\) and \(f^{[k+1]}(s)=f(f^{[k]}(s))\).
Set
\[
 R(s)=2^{\,2^{\,2^{(s+16)^4}}},\qquad
 A(s)=R^{[32(s+1)^2]}(s+16).
\]
These are explicit primitive-recursive functions. A construction step
below is a finite-list operation or bounded loop instruction; arithmetic
operations on its natural-number codes are themselves primitive recursive.
The bounds concern this specified construction, not a practical estimate.
\end{hbdefinition}

\begin{hblemma}{thm:p09-kamp-pr-local-envelope}{proved=annex:p9.kamp.pr.local.envelope}{Bounds for all nonrecursive syntax phases}
Starting with total data weight at most \(s\), each nonrecursive
phase of the deterministic construction has output and temporary data
weight at most \(A(s)\) and uses at most \(A(s)\)
construction steps. The phases include Boolean distribution, finite
relation enumeration, all word cuts, elementary-word assembly, proper
substitution, freshening and construction of basic or limiting histories.
They exclude recursive Lemma~8 calls and visits to original nodes of
an outer structural recursion, which are bounded separately below.
\end{hblemma}

\begin{hbdefinition}{def:p09-kamp-pr-family-mass}{Family mass and the explicit Lemma 8 bounds}
For a pure family record let
\[
 \mu(F)=1+\sum_{\theta\in F\cap\operatorname{At}}
                      (1+\operatorname{nodes}(\theta))
             +\sum_{\rho\in F\setminus\operatorname{At}}(1+\ell(\rho)).
\]
At-node count treats a unary literal as a leaf and retains its Boolean
tree; basic length \(\ell\) is the original predicate-word length.
Predicate indices do not contribute to this structural mass and are
bounded separately. Let \(M\geq1\), \(s\geq1\) and \(p\geq0\)
bound the initial maximum basic length, mass and predicate index.
Define
\[
\begin{gathered}
 b_0=s+1,\qquad b_{i+1}=(M+3)^{2(b_i+1)},\qquad B=b_{M+1},\\
 h=(M+1)(B+1)+1,\qquad
 c=34\,2^{M(B+1)}+1,\qquad z=(c+1)^{h+1},\\
 s_*=64(B+h+p+M+16)^2,\qquad u_0=A(s_*),\\
 u_{i+1}=A\bigl(A(s_*)+c(u_i+1)+16\bigr),\\
 L_8(M,s,p)=(z+1)(u_h+s_*+h+c+16)^2.
\end{gathered}
\]
Every recurrence is ordinary primitive recursion with its displayed
parameters fixed. Finally put \(L(t)=L_8(t+1,t+1,t+1)\).
\end{hbdefinition}

\begin{hblemma}{thm:p09-kamp-pr-lemma8}{proved=annex:p9.kamp.pr.lemma8}{A primitive-recursive bound for the entire Lemma 8 construction}
Every child satisfies
\[
                    \mu(G)\leq(M+3)(\mu(F)+1).
\]
A branch finishes a current maximum-length level after at most its
entry mass many edges, and never returns to a higher level. Consequently
\(B\) bounds all branch masses, \(h\) bounds depth, \(c\)
bounds immediate recursive branching, and \(z\) bounds the whole
recursive tree. Predicate indices along a branch are at most \(p+h\).
The function \(L_8\) bounds the output weight and the construction
steps of the full deterministic Lemma~8 evaluator. In particular that
evaluator is primitive recursive. This conclusion uses the growth bound
and level analysis; lexicographic well-foundedness alone would not suffice.
\end{hblemma}

\begin{hbdefinition}{def:p09-kamp-pr-outer-bounds}{Bounds for the outer construction}
Define fixed unary primitive-recursive functions
\[
\begin{aligned}
 J(t)&=A\bigl(t+A(t)+L(A(t))+16\bigr),\\
 S(t)&=J^{[4(t+1)^2]}(t+16),\\
 C(t)&=S^{[6]}(t+16),\\
 V(t)&=A\bigl(t+C(A(t))+16\bigr),\\
 W(t)&=V^{[t+1]}(t+16).
\end{aligned}
\]
For an input of weight \(n\), set
\[
\begin{gathered}
 n_0=n,\quad n_1=A(n_0),\quad n_2=S(n_1),\quad n_3=W(n_2),\\
 n_4=S(n_3),\quad n_5=A(n_4+n+16),\quad n_6=S(n_5),\\
                       T(n)=2\sum_{i=0}^{6}(n_i+16).
\end{gathered}
\]
The stages are respectively initial expansion, standard form, elementary
regional construction, temporal construction, schema padding and the
literal \(D\)-image. These functions are deliberately coarse.
\end{hbdefinition}

\begin{hbmethod}{def:p09-kamp-pr-exp}{The total numerical translation}
Parse the input code and check its sole free variable. For a valid input,
run the deterministic first-order-to-temporal construction, using the
stated stack frames and local rules. An atomic input \(P_i(t_0)\)
uses the fixed initial branch \(\eta=q_i\) and then proceeds to padding.
For a nonempty unary signature set \(n=1+\max i\), where the maximum
is over its predicate indices. On a source tense-defining formula this
is its arity and all \(P_0,\ldots,P_{n-1}\) occur; the rule also
defines the numerical extension on a gapped signature. After obtaining \(\eta\),
append, in increasing index order, the conjunct
\(\neg(q_i\land\neg q_i)\) for each permitted \(q_i\)
that does not occur. This is true at every point and ensures that
the result contains every required propositional letter and no others,
as the source's schema definition requires. Return its temporal code.
When applied to a formula rather than its numerical code,
\(\operatorname{Exp}(\varphi)\) denotes the decoded temporal output.

Malformed inputs receive the default code of \(q_0\). A one-free-variable
input with no unary letters uses padding arity 1 and the same construction
with the inessential letter \(q_0\); its output is a unary schema independent
of that letter. It is not a variable-free temporal formula: the source
\(\mathsf{TL}_1\) grammar has no such formula.
\end{hbmethod}

\begin{hbtheorem}{thm:p09-kamp-primitive-recursive-translation}{proved=annex:p9.kamp.primitive.recursive.translation}{Kamp's primitive-recursive translation and its exact schema scope}
The numerical function just defined is primitive recursive. On each
positive-arity tense-defining formula \(\varphi\), its output
\(\operatorname{Exp}(\varphi)\) is an \(n\)-place schema of
the original \(\mathsf{TL}_1\) and satisfies
\[
 \mathsf{Com}\models
       \varphi\longleftrightarrow D(\operatorname{Exp}(\varphi)).
\]
The strict \(S,U\) convention and the literal \(D\)-image are
unchanged. The stated list program has at most \(T(w(\varphi))\)
construction steps, and its output code and \(D\)-image code are
less than \(2^{T(w(\varphi))}\). Thus the assertion is uniform
over all input formulas and all complete linear orders.

The source's definitions require an \(n\)-place schema to contain
all \(n\) designated letters. If their arity parameter is read as
including zero, the literal zero-place assertion has a syntactic
exception: no formula of \(\mathsf{TL}_1\) contains zero proposition
letters. The explicit unary constant extension above supplies equivalent
truth conditions for zero-predicate inputs without mislabelling its arity
\cite[ch.~II, defs.~2,4 and thm.~2, pp.~21--23,28]{Kamp1968}.
\end{hbtheorem}

\begin{hbexample}{def:p09-kamp-pr-example}{A concrete code and a length-level bound}
The first-order token list for \(P_0(t_0)\) is \([0,0,0]\),
with code 7 and weight 4. Its atomic translation is the temporal
list \([0,0]\), with code 3 and weight 3. No padding is
needed, and \(D(q_0)=P_0(t_0)\).

For the Lemma~8 family with At atom \(P_0\) and basic word
\((P_1,P_2,P_3)\), \(M=3\) and \(\mu(F)=1+2+4=7\).
In its whole-interval localization after the maximal word is removed,
the actual child is the two-atom family \(\{P_0,P_4\}\), of mass 5;
\(P_4\) is its fresh placeholder. Every child in the general bound
has mass at most \(6(7+1)=48\). The bound on a
whole level beginning with mass 7 is \(6^{16}\); the actual
maximal-word count here is only one, so that example leaves its
maximal level after one edge. The general bound remains valid when
many maximal words and all their shorter residuals are present.
\end{hbexample}


Consequently the positive-arity conclusion of Chapter~II Theorem~1
says that every uniformly first-order-definable tense of that arity
over \(\mathsf{Com}\) is expressible from
negation, conjunction, \(S\), and \(U\).  Theorem~2 strengthens this with a
primitive-recursive first-order-to-temporal translation.
\Cref{thm:p09-kamp-primitive-recursive-translation} supplies explicit
primitive-recursive bounds for the construction, its output and its
evaluation, together with the exact positive-arity schema convention
and the nullary syntax boundary. These bounds are supplied by the
handbook; the thesis states primitive recursiveness without a numerical
bound. The result concerns expressive completeness. Kamp's remarks
about axiomatization and decidability cite other authors and require
their separate arguments
\cite[ch.~I, p.~14; ch.~II, pp.~27--29]{Kamp1968}.

% Kamp 1968, Chapter IV, Theorem 1 and its three lemmas, printed pp. 95--99.

\subsection{Why unary real-order tenses cannot express SINCE or UNTIL}
\label{sec:p09-kamp-iv1-unary-separation}

The expressive-completeness theorem permits binary temporal operations.
Kamp's next result shows that their arity cannot be reduced to one, even
when every first-order-definable unary tense on the real order is available.
The proof has three steps: one discrete proposition generates only four
truth sets; finitely many disjoint such propositions generate a finite
Boolean algebra; every formula using unary tenses preserves that algebra.
An even/odd-integer interpretation then separates this language from each
of the strict binary operations \(S,U\)
\cite[ch.~IV, thm.~1 and lemmas~1--3, pp.~95--99]{Kamp1968}.

\begin{hbdefinition}{def:p09-kamp-real-unary-language}{The quantified class of unary real-order tenses}
Throughout this section the time structure is exactly
\(\mathcal R=\langle\mathbb R,<\rangle\). A proposition is a subset
of \(\mathbb R\). For a subset \(A\), write
\(A^c=\mathbb R\setminus A\); complements are always relative to
this fixed time domain.

Let \(\mathfrak M_{\mathcal R}\) be the set of all operations
\(O:\mathcal P(\mathbb R)\to\mathcal P(\mathbb R)\) defined by
a source \(L\)-formula \(\varphi(t_0)\) whose only free variable
is \(t_0\), which contains \(P_0\), and which contains no
other unary predicate symbol:
\[
 O(A)=\{t\in\mathbb R:
       \langle\mathbb R,<,A\rangle\models\varphi[t_0\mapsto t]\}.
\]
The signature contains no real constants or additional predicate
parameters. The required occurrence of \(P_0\) may be
inessential: conjoining \(P_0(t_0)\lor\neg P_0(t_0)\)
inserts one without changing the defined operation. The operation
need only have the displayed definition on \(\mathcal R\).

Write \(\mathsf{TM}_{\mathcal R}\) for the language of finite
formulas generated by
\[
 \eta ::= q_i\mid\neg\eta\mid(\eta\land\eta)
                 \mid O(\eta),\qquad O\in\mathfrak M_{\mathcal R}.
\]
For a valuation \(V\) assigning propositions to letters, let
\(\llbracket\eta\rrbracket_V\) denote the truth set of \(\eta\).
In particular,
\(\llbracket O(\eta)\rrbracket_V
 =O(\llbracket\eta\rrbracket_V)\).
The alphabet contains all the permitted unary operations at once;
each individual formula is finite. Arbitrary finite truth-functional
sentential connectives may also be used, since negation and conjunction
express them. The constants are expressible using
\(q_0\land\neg q_0\) and its negation.

For propositions \(A,B\), the binary target operations have the
source's witness-first, open-interval semantics:
\[
\begin{aligned}
 S(A,B)&=\{t:\exists s<t\ (s\in A\land(s,t)\subseteq B)\},\\
 U(A,B)&=\{t:\exists s>t\ (s\in A\land(t,s)\subseteq B)\}.
\end{aligned}
\]
These are the same strict operations used in
\cref{thm:p09-kamp-iii2}. The quantification over unary tenses below
is over \(\mathfrak M_{\mathcal R}\), with the stated first-order
definitions and without additional parameters
\cite[ch.~II, defs.~2--3, pp.~21--22; ch.~IV, p.~95]{Kamp1968}.
\end{hbdefinition}

\begin{hbdefinition}{def:p09-kamp-du}{Discrete and unbounded real propositions}
A proposition \(A\subseteq\mathbb R\) is \emph{discrete} in
Kamp's sense when
\[
 \forall t\in\mathbb R\ \exists z,y\in\mathbb R\quad
 z<t<y,\qquad (z,t)\cap A=(t,y)\cap A=\varnothing.
\]
The condition is imposed at every real point, including points outside
\(A\). It excludes finite accumulation points.

The source calls \(A\) \emph{unbounded} when either
\(A=\varnothing\), or
\[
 \forall t\in\mathbb R\ \exists z,y\in A\quad z<t<y.
\]
Thus every nonempty unbounded proposition is unbounded in both
directions, and the empty proposition is expressly included.
Write \(\operatorname{DU}(A)\) when either \(A\) is
discrete and unbounded or its complement \(A^c\) is discrete
and unbounded, and put
\[
                  \mathcal F(A)=\{\varnothing,\mathbb R,A,A^c\}.
\]
The family \(\mathcal F(A)\) can have two members when
\(A=\varnothing\) or \(A=\mathbb R\)
\cite[ch.~IV, def.~1, p.~96]{Kamp1968}.
\end{hbdefinition}

\begin{hblemma}{thm:p09-kamp-iv1-lemma1}{proved=annex:p9.kamp.iv1.lemma1}{The four-set lemma for a single proposition}
Let \(\operatorname{DU}(A)\). If \(\eta\) is a
\(\mathsf{TL}_1\)-formula and every proposition letter occurring
in \(\eta\) is interpreted by \(A\), then
\[
                    \llbracket\eta\rrbracket\in\mathcal F(A).
\]
More generally the same conclusion holds when every occurring letter
is interpreted by any member of \(\mathcal F(A)\). The latter
form records the closure used in the induction; the former is Kamp's
Lemma~1. Both \(S\) and \(U\) send any two members of
\(\mathcal F(A)\) to either \(\varnothing\) or
\(\mathbb R\)
\cite[ch.~IV, lemma~1, pp.~96--97]{Kamp1968}.
\end{hblemma}

\begin{hbcorollary}{thm:p09-kamp-iv1-unary-corollary}{proved=annex:p9.kamp.iv1.unary.corollary}{Every definable unary tense preserves the four choices}
For every \(O\in\mathfrak M_{\mathcal R}\) and every
\(A\subseteq\mathbb R\),
\[
      \operatorname{DU}(A)\quad\Longrightarrow\quad
      O(A)\in\{\varnothing,\mathbb R,A,A^c\}.
\]
The choice among the four sets may depend on \(O\) and on \(A\).
This is a statement about every permitted unary tense, irrespective
of the number or alternation of quantifiers in its defining formula
\cite[ch.~IV, corollary to lemma~1, p.~97]{Kamp1968}.
\end{hbcorollary}

\begin{hbdefinition}{def:p09-kamp-discrete-generated-field}{The field generated by finitely many propositions}
For \(n\geq0\) and \(A_1,\ldots,A_n\subseteq\mathbb R\),
let \(\mathcal B(A_1,\ldots,A_n)\) be the smallest family of
subsets of \(\mathbb R\) containing these sets,
\(\varnothing\), and \(\mathbb R\), and closed under
complement and intersection. It is a Boolean algebra of sets, called
a field of subsets in the source. For
\(I\subseteq\{1,\ldots,n\}\), put
\[
                         A_I=\bigcup_{i\in I}A_i,
                  \qquad A_{\varnothing}=\varnothing.
\]
When \(n=0\), the generated field is
\(\{\varnothing,\mathbb R\}\).
\end{hbdefinition}

\begin{hblemma}{thm:p09-kamp-iv1-lemma2}{proved=annex:p9.kamp.iv1.lemma2}{Normal form in the finite generated field}
Let \(n\geq0\), and let \(A_1,\ldots,A_n\) be pairwise
disjoint, discrete, unbounded real propositions. Then
\[
\begin{aligned}
 \mathcal B(A_1,\ldots,A_n)
 &=\{A_I:I\subseteq\{1,\ldots,n\}\}\\
 &\quad\cup\{A_I^c:I\subseteq\{1,\ldots,n\}\}.
\end{aligned}
\]
In particular every member of the generated field satisfies
\(\operatorname{DU}\). Empty generators and an empty family are
allowed. The normal form is not asserted to be unique
\cite[ch.~IV, lemma~2, pp.~97--98]{Kamp1968}.
\end{hblemma}

\begin{hblemma}{thm:p09-kamp-iv1-lemma3}{proved=annex:p9.kamp.iv1.lemma3}{Closure under every unary-tense formula}
Under the hypotheses of \cref{thm:p09-kamp-iv1-lemma2}, let
\(\eta\in\mathsf{TM}_{\mathcal R}\) use only the letters
\(q_1,\ldots,q_n\), interpreted respectively by
\(A_1,\ldots,A_n\). Then
\[
              \llbracket\eta\rrbracket_V
                    \in\mathcal B(A_1,\ldots,A_n).
\]
The same conclusion holds if further letters occur and each of their
interpretations already belongs to this generated field. Consequently,
if two real points have the same membership in every generator, every
such formula has the same truth value at those two points.
The condition concerns every letter actually occurring in the formula
\cite[ch.~IV, lemma~3, p.~98]{Kamp1968}.
\end{hblemma}

\begin{hbtheorem}{thm:p09-kamp-iv1}{proved=annex:p9.kamp.iv1}{Kamp's real-order unary separation theorem}
Neither SINCE nor UNTIL is expressible on
\(\langle\mathbb R,<\rangle\) using sentential connectives and
arbitrary unary real-order tenses. Precisely, for each
\(T\in\{S,U\}\), there is no finite formula
\(\eta\in\mathsf{TM}_{\mathcal R}\), with no proposition
letters other than \(q_0,q_1\), such that
\[
 \forall A,B\subseteq\mathbb R\quad
       \llbracket\eta\rrbracket_{q_0\mapsto A,\,q_1\mapsto B}
                           =T(A,B).
\]
The impossibility holds with the full alphabet
\(\mathfrak M_{\mathcal R}\) available, and therefore with any
chosen family of its operations, finite or infinite.

For the fixed interpretation
\[
 E=\{2k:k\in\mathbb Z\},\qquad
 H=\{2k+1:k\in\mathbb Z\},\qquad
 q_0\mapsto E,\quad q_1\mapsto H,
\]
all \(\mathsf{TM}_{\mathcal R}\)-formulas in these two letters
agree at every pair of noninteger points. In contrast,
\[
\begin{array}{c|cc}
 & -\tfrac12 & +\tfrac12\\ \hline
 S(q_0,\neg q_1)&\text{false}&\text{true}\\
 U(q_0,\neg q_1)&\text{true}&\text{false}
\end{array}
\]
This supplies the two distinct nonexpressibility conclusions
\cite[ch.~IV, thm.~1, pp.~95--99]{Kamp1968}.
\end{hbtheorem}

\begin{hbexample}{def:p09-kamp-iv1-integer-boundaries}{The complete truth sets in the separating interpretation}
For the \(E,H\) in \cref{thm:p09-kamp-iv1}, the generated
field consists of exactly
\[
 \varnothing,\quad\mathbb R,\quad E,\quad H,\quad\mathbb Z,
 \quad E^c,\quad H^c,\quad\mathbb R\setminus\mathbb Z.
\]
Its three atoms are \(E\), \(H\), and
\(\mathbb R\setminus\mathbb Z\). The binary truth sets are
\[
 S(E,H^c)=\bigcup_{k\in\mathbb Z}(2k,2k+1],\qquad
 U(E,H^c)=\bigcup_{k\in\mathbb Z}[2k-1,2k).
\]
Each contains some noninteger points and excludes others. At an even
integer both formulas are false; at an odd integer both are true.
The closed endpoint at an odd integer is required because the current
point lies outside the open interval on which the second operand is
tested. The annex proves these full descriptions along with the theorem.
\end{hbexample}

\begin{hbexample}{def:p09-kamp-iv1-du-boundaries}{The empty case and the role of both unbounded directions}
Both \(\varnothing\) and \(\mathbb R\) satisfy
\(\operatorname{DU}\), and their four-choice family reduces to
\(\{\varnothing,\mathbb R\}\). Thus every permitted unary
tense sends either of these two inputs to one of these two outputs.

The singleton \(A=\{0\}\) is discrete but fails the source's
unboundedness condition. The permitted unary strict-past tense
\[
 O(A)=\{t:\exists s<t\ (s\in A)\}
\]
has value \((0,\infty)\), which is none of
\(\varnothing,\mathbb R,\{0\},\mathbb R\setminus\{0\}\).
This demonstrates why the four-choice assertion cannot be extended to
all discrete propositions. Likewise, requiring isolation only at members
would be insufficient to state the source definition: the set
\(\{1/m:m\in\mathbb Z,\ m\geq1\}\) has isolated members
but violates its right-hand gap condition at zero.
\end{hbexample}

% Kamp 1968, dense-order observation and Chapter IV, Definition 2/Theorem 2.

\subsection{Dense-order separation and a unary basis on the integers}
\label{sec:p09-kamp-iv2-integer-basis}

\subsubsection{The adjacent dense-order observation}

Kamp extends his real-order separation to every particular dense linear
order with more than one point. Such an order may have endpoints, need
not be complete, and need not resemble the reals in cardinality. The
following local argument supplies the extension at those exact hypotheses
\cite[ch.~IV, pp.~99--100]{Kamp1968}.

\begin{hbdefinition}{def:p09-kamp-dense-local-language}{Unary tenses on a fixed dense order}
Fix a strict linear order \(\mathcal J=\langle T,<\rangle\).
A permitted unary \(\mathcal J\)-tense is an operation
\(O:\mathcal P(T)\to\mathcal P(T)\) defined by one source
first-order formula in the equality-free signature \(\{<,P_0\}\)
whose only free variable is the current point
\(t_0\), which contains \(P_0\), and which contains no other
unary predicate or named point. Negation and conjunction act on sets by
complement in \(T\) and intersection. Write
\(\mathsf{TM}_{\mathcal J}\) for the finite formulas generated
from proposition letters \(q_i\) by these operations, with every
permitted unary tense available. As before, every finite
truth-functional connective can be admitted as an abbreviation.

For a convex subset \(I\subseteq T\), complements of subsets of
\(I\) in the next lemma are taken relative to \(I\).
A two-sided sequence of markers in \(I\) is a family
\((d_n)_{n\in\mathbb Z}\) with \(d_n<d_{n+1}\) such that
every \(x\in I\) lies strictly between some two markers.
In particular the markers are cofinal in both directions inside \(I\).
The interpretation of a unary predicate outside \(I\) is unrestricted.
\end{hbdefinition}

\begin{hblemma}{thm:p09-kamp-dense-local-unary}{proved=annex:p9.kamp.dense.local.unary}{Local four-choice closure without completeness}
Let \(T\) be densely ordered, and let \(I\) be a nonempty
convex subset with no first or last point. Let \(A\subseteq T\).
Suppose that \(A\cap I\) is \(\varnothing\), \(I\), a
two-sided marker set \(D=\{d_n:n\in\mathbb Z\}\), or
\(I\setminus D\). Then every permitted unary tense satisfies
\[
 O(A)\cap I\in
       \{\varnothing,I,A\cap I,I\setminus(A\cap I)\}.
\]
More precisely, two points of \(I\) with the same membership in
\(A\) satisfy exactly the same first-order formulas in the signature
\(\{<,P_0\}\) with only the current point free. The formula may
quantify over all of \(T\), including its endpoints and the
unrestricted interpretation outside \(I\).
\end{hblemma}

\begin{hbtheorem}{thm:p09-kamp-dense-unary-separation}{proved=annex:p9.kamp.dense.unary.separation}{Unary separation on every nontrivial dense order}
For every particular dense strict linear order
\(\mathcal J=\langle T,<\rangle\) with more than one element,
neither of the strict operations
\[
\begin{aligned}
 S(A,B)&=\{t\in T:\exists s<t\ (s\in A\land(s,t)\subseteq B)\},\\
 U(A,B)&=\{t\in T:\exists s>t\ (s\in A\land(t,s)\subseteq B)\}
\end{aligned}
\]
is expressible by a finite formula in
\(\mathsf{TM}_{\mathcal J}\) using only its two input letters.
Equivalence is required for every \(A,B\subseteq T\).
The unary tenses may be chosen for this particular order; no uniform
definition across a class of orders is imposed as an additional
hypothesis. Density and more than one element are the only order
hypotheses. On a one-point order the two strict operations are constantly
false, so that excluded boundary behaves differently
\cite[ch.~IV, pp.~99--100]{Kamp1968}.
\end{hbtheorem}

\subsubsection{The integer basis and its equality-free definitions}

\begin{hbdefinition}{def:p09-kamp-iv2-basis}{The four unary integer tenses}
For the remainder of this section the time order is
\(\mathsf{In}=\langle\mathbb Z,<\rangle\).
To display the source's equality-free definitions, abbreviate
\(P_0(z)\land\neg P_0(z)\) by \(\bot_{P_0}(z)\).
This is a contradictory formula, not an extra constant in the language.
The defining formulas at \(t_0\) are
\[
\begin{aligned}
 \pi_-(t_0)={}&\exists t_1\bigl(t_1<t_0\land P_0(t_1)\land
       \forall t_2(t_1<t_2<t_0\to\bot_{P_0}(t_2))\bigr),\\
 \pi_+(t_0)={}&\exists t_1\bigl(t_0<t_1\land P_0(t_1)\land
       \forall t_2(t_0<t_2<t_1\to\bot_{P_0}(t_2))\bigr).
\end{aligned}
\]
The isolation formulas are
\[
\begin{aligned}
 \iota_-(t_0)={}&\exists t_1\bigl(t_1<t_0\land P_0(t_1)\land\\
 &\quad\exists t_2\bigl(t_2<t_1\land\neg P_0(t_2)\land
       \forall t_3(t_2<t_3<t_1\to\bot_{P_0}(t_3))\bigr)\land\\
 &\quad\forall t_4(t_1<t_4<t_0\to\neg P_0(t_4))\bigr),\\
 \iota_+(t_0)={}&\exists t_1\bigl(t_0<t_1\land P_0(t_1)\land\\
 &\quad\exists t_2\bigl(t_1<t_2\land\neg P_0(t_2)\land
       \forall t_3(t_1<t_3<t_2\to\bot_{P_0}(t_3))\bigr)\land\\
 &\quad\forall t_4(t_0<t_4<t_1\to\neg P_0(t_4))\bigr).
\end{aligned}
\]
They define respectively \(\mathsf{PAST}_1\),
\(\mathsf{FUTURE}_1\), \(\mathsf{ISOLP}\), and
\(\mathsf{ISOLF}\). The basis is exactly
\[
 \mathcal S=\{\mathsf{ISOLP},\mathsf{ISOLF},
              \mathsf{PAST}_1,\mathsf{FUTURE}_1,
              \mathsf{NOT},\mathsf{AND}\}.
\]
Its sentential language \(\mathsf{TS}\) uses \(P_1,F_1\)
for immediate past and future, and \(S_1,U_1\) for past and
future isolation. Its finite formulas start with proposition letters
\(q_i\) and apply these four unary constructors and the
sentential constructors \(\neg,\land\).
Disjunction is a Boolean abbreviation. Juxtaposition of unary symbols
denotes nesting; for example \(F_1F_1\alpha=F_1(F_1\alpha)\).
The subscripted \(S_1,U_1\) are unary isolation operators;
the unsubscripted \(S,U\) remain the binary target operations
\cite[ch.~IV, def.~2 and thm.~2, p.~100]{Kamp1968}.
\end{hbdefinition}

\begin{hblemma}{thm:p09-kamp-iv2-unary-semantics}{proved=annex:p9.kamp.iv2.unary.semantics}{The exact integer clauses for the source basis}
For any \(\mathsf{TS}\)-formula \(\alpha\), any valuation,
and any \(i\in\mathbb Z\), write \(\alpha(i)\) for its
truth at \(i\). Then
\[
\begin{aligned}
 (P_1\alpha)(i)&\Longleftrightarrow\alpha(i-1),\\
 (F_1\alpha)(i)&\Longleftrightarrow\alpha(i+1),\\
 (S_1\alpha)(i)&\Longleftrightarrow
   \exists j<i\bigl(\alpha(j)\land\neg\alpha(j-1)\land
       \forall k(j<k<i\to\neg\alpha(k))\bigr),\\
 (U_1\alpha)(i)&\Longleftrightarrow
   \exists j>i\bigl(\alpha(j)\land\neg\alpha(j+1)\land
       \forall k(i<k<j\to\neg\alpha(k))\bigr).
\end{aligned}
\]
Thus an isolation operator requires the nearest strict occurrence
in its direction to have a false neighbor on its far side.
There is no condition on \(\alpha(i)\). The intervening
universal condition is vacuous when that nearest occurrence is an
immediate neighbor
\cite[ch.~IV, pp.~101--102]{Kamp1968}.
\end{hblemma}

\subsubsection{The complete two-input construction}

\begin{hbdefinition}{def:p09-kamp-iv2-marker-formulas}{Markers that distinguish a successful endpoint from a failure}
Use \(q_0\) for the witness proposition and \(q_1\) for
the open-interval proposition. Define the following formulas of
\(\mathsf{TS}\):
\[
\begin{aligned}
 C&=\neg q_0\land q_1,\\
 B_P&=\neg q_0\land\neg q_1\land P_1\neg q_0
                                      \land P_1P_1\neg q_0,\\
 B_F&=\neg q_0\land\neg q_1\land F_1\neg q_0
                                      \land F_1F_1\neg q_0,\\
 \eta_1&=(q_0\land P_1\neg q_0\land P_1q_1)
                                      \lor B_P\lor F_1B_P,\\
 \eta_2&=(q_0\land F_1\neg q_0\land F_1q_1)
                                      \lor B_F\lor P_1B_F.
\end{aligned}
\]
The marker \(\eta_1\) is used for a future search and inspects
the two preceding positions. The marker \(\eta_2\) is used
for a past search and inspects the two following positions.
The final shifted disjunct places an extra marker immediately before
a failed future search, or immediately after a failed past search.

For \(\delta\in\{+1,-1\}\), define the directional notation
\[
\begin{array}{c|ccc}
 \delta&N_\delta&J_\delta&\eta_\delta\\ \hline
 +1&F_1&U_1&\eta_1\\
 -1&P_1&S_1&\eta_2
\end{array}
\]
where \(\eta_\delta\) is only a directional abbreviation in
this table: \(\eta_{+1}=\eta_1\) and
\(\eta_{-1}=\eta_2\). The symbol \(N_\delta\) means one
step in the chosen direction; \(N_{-\delta}\) means one step
back toward the starting point.
\end{hbdefinition}

\begin{hblemma}{thm:p09-kamp-iv2-marker-lemma}{proved=annex:p9.kamp.iv2.marker.lemma}{The marker calculation in both directions}
Fix any integer \(t\), either \(\delta\in\{+1,-1\}\),
and any interpretations of \(q_0,q_1\). Suppose that \(C\)
holds at \(t+\delta\) and \(t+2\delta\). Then
\[
 (J_\delta\eta_\delta)(t)
 \quad\Longleftrightarrow\quad
 \exists m\geq3\bigl(q_0(t+\delta m)\land
       \forall r\,(1\leq r<m\to q_1(t+\delta r))\bigr).
\]
Here \(m,r\) range over integers.
The proof includes a successful first endpoint, a first failed
interval point, and the case with no endpoint or failure anywhere
in the chosen direction. It also verifies the source's alternative
case in which the witness proposition never occurs in that direction.
\end{hblemma}

\begin{hbdefinition}{def:p09-kamp-iv2-target-formulas}{Explicit expressions for strict UNTIL and SINCE}
Retain both conjuncts of the source construction. Put
\[
\begin{aligned}
 G_F&=F_1q_0\lor
       \bigl(F_1q_1\land(F_1F_1q_0\lor F_1F_1q_1)\bigr),\\
 G_P&=P_1q_0\lor
       \bigl(P_1q_1\land(P_1P_1q_0\lor P_1P_1q_1)\bigr),\\
 \zeta_1&=G_F\land\bigl[
       F_1q_0\lor(F_1q_1\land F_1F_1q_0)\\
       &\hspace{38mm}\lor
          (F_1C\land F_1F_1C\land U_1\eta_1)\bigr],\\
 \zeta_2&=G_P\land\bigl[
       P_1q_0\lor(P_1q_1\land P_1P_1q_0)\\
       &\hspace{38mm}\lor
          (P_1C\land P_1P_1C\land S_1\eta_2)\bigr].
\end{aligned}
\]
All abbreviations expand to finite formulas over exactly
\(\mathcal S\), with no proposition letters beyond \(q_0,q_1\).
These give the corrected source construction: the marker inspections
face toward the starting point, and the repeated shifts in each
outer formula face in its search direction
\cite[ch.~IV, thm.~2, pp.~100--105]{Kamp1968}.
\end{hbdefinition}

\begin{hbtheorem}{thm:p09-kamp-iv2}{proved=annex:p9.kamp.iv2}{Kamp's integer-order unary basis theorem}
For every \(A,B\subseteq\mathbb Z\), every
\(t\in\mathbb Z\), and the valuation \(q_0\mapsto A\),
\(q_1\mapsto B\),
\[
\begin{aligned}
 \zeta_1(t)&\Longleftrightarrow
       \exists s>t\bigl(s\in A\land
                       \forall u(t<u<s\to u\in B)\bigr),\\
 \zeta_2(t)&\Longleftrightarrow
       \exists s<t\bigl(s\in A\land
                       \forall u(s<u<t\to u\in B)\bigr).
\end{aligned}
\]
Consequently both strict UNTIL and strict SINCE are expressible
on \(\mathsf{In}\) from
\(\{\mathsf{ISOLP},\mathsf{ISOLF},\mathsf{PAST}_1,
\mathsf{FUTURE}_1,\mathsf{NOT},\mathsf{AND}\}\).
The first operand remains the witness proposition and the second
the open-interval proposition. An immediate witness makes the
interval requirement vacuous. The conclusion holds at every integer,
including when either predicate is empty, full, or unbounded in
one or both directions
\cite[ch.~IV, thm.~2, pp.~100--105]{Kamp1968}.
\end{hbtheorem}

\begin{hbexample}{def:p09-kamp-iv2-examples}{Successful endpoints, paired failures and immediate witnesses}
At \(t=0\), take \(A=\{3,4,5,\ldots\}\) and
\(B=\mathbb Z\). The first two future points satisfy
\(C\). The marker \(\eta_1\) is true at three and
false at every other positive integer. Thus \(U_1\eta_1\)
and \(\zeta_1\) are true, with three a witness for strict
UNTIL. Reflecting these sets through zero gives the corresponding
SINCE example with \(A=\{\ldots,-5,-4,-3\}\).

If instead \(A=\varnothing\) and \(B=\{1,2\}\), the
first failed future point is three. The marker is true at both two
and three, so its nearest future occurrence is not isolated on its
far side. Thus \(U_1\eta_1\) and \(\zeta_1\) are false.
If \(A=\varnothing\) and \(B=\mathbb Z\), there is no
future marker at all, and the same conclusion follows.

For an immediate witness take \(A=\{1\}\) and any \(B\),
even \(B=\varnothing\). Then \(F_1q_0\) makes both
conjuncts of \(\zeta_1\) true, matching the empty interval
\((0,1)\cap\mathbb Z\). For \(A=\{2\}\),
\(\zeta_1(0)\) is true exactly when \(1\in B\).
Neither case imposes a condition on membership of the current point
or the witness in \(B\). The annex verifies the examples and
their reflected past cases.
For comparison, \(A=\{0\}\) and \(B=\{-1,-2\}\)
make \(\zeta_2(0)\) false: the current point cannot be
used as a strict past witness.
\end{hbexample}


\paragraph{Why the order hypothesis cannot be deleted.}
\Cref{thm:p09-kamp-iv5} supplies the rational-order obstruction using
the periodic and accumulating interpretations of
\cref{def:p09-kamp-rational-cut-models}.
The first-order formula of
\cref{def:p09-kamp-rational-distinguishing-formula} distinguishes them
at the rational current point three, whereas
\cref{thm:p09-kamp-iv5-temporal-induction} proves agreement there
for every one-input formula built from negation, conjunction and
the source's strict SINCE and UNTIL.
Thus the complete-order basis does not extend unchanged to all
linear orders \cite[ch.~IV, pp.~109--112]{Kamp1968}.

\paragraph{The literal Chapter IV Theorem 3.}
Kamp writes
\[
 \mathsf{In}=\langle\mathbb Z,<\rangle.
\]
The source displays equality-free first-order formulas for four unary tenses.
On \(\mathsf{In}\), those formulas are equivalent to the following
integer-specialized clauses for \(P\subseteq\mathbb Z\):
\[
\begin{aligned}
 \mathsf{PAST}_1(P)(i)
 &\Longleftrightarrow P(i-1),\\
 \mathsf{FUTURE}_1(P)(i)
 &\Longleftrightarrow P(i+1),\\
 \mathsf{ISOLP}(P)(i)
 &\Longleftrightarrow
 \exists j<i\,
 \bigl(P(j)\land\neg P(j-1)\land
       \forall k\,(j<k<i\Rightarrow\neg P(k))\bigr),\\
 \mathsf{ISOLF}(P)(i)
 &\Longleftrightarrow
 \exists j>i\,
 \bigl(P(j)\land\neg P(j+1)\land
       \forall k\,(i<k<j\Rightarrow\neg P(k))\bigr).
\end{aligned}
\]
In the source formulas, immediate adjacency is expressed without equality or
a successor symbol: an existentially chosen neighbor is accompanied by a
universal clause whose consequent is \(P_0(t)\land\neg P_0(t)\), forcing the
open interval to be empty.  Thus the first two operations inspect the
immediate neighbor.  The latter two
locate an isolated edge of the next past or future block of \(P\): they stand
as unary tenses on the integer order, differ from ordinary strict
past/future by inspecting adjacency and the absence of intervening \(P\)-points,
and are used by Theorems~IV.2--IV.4.

The four remaining operations in the proposed basis are, for
\(P,Q\subseteq\mathbb Z\),
\[
\begin{aligned}
 \mathsf{PAST}(P)(i)&\Longleftrightarrow \exists j<i\,P(j),&
 \mathsf{FUTURE}(P)(i)&\Longleftrightarrow \exists j>i\,P(j),\\
 \mathsf{NOT}(P)&=\mathbb Z\setminus P,&
 \mathsf{AND}(P,Q)&=P\cap Q.
\end{aligned}
\]
The first two are strict unbounded past and future, not the immediate-neighbor
operations \(\mathsf{PAST}_1,\mathsf{FUTURE}_1\).  Put
\[
 \mathcal S'
 =
 \{\mathsf{PAST}_1,\mathsf{FUTURE}_1,
   \mathsf{PAST},\mathsf{FUTURE},
   \mathsf{NOT},\mathsf{AND}\}.
\]
This is a metalinguistic basis of operations, not another object-language
connective.  A unary operation \(F:\mathcal P(\mathbb Z)\to
\mathcal P(\mathbb Z)\) is \emph{expressible from \(\mathcal S'\)} exactly
when a finite term with one set variable, formed using the displayed unary
and binary operations, denotes \(F(P)\) for every
\(P\subseteq\mathbb Z\).  This definition is needed to state Chapter~IV
Theorem~3: neither \(\mathsf{ISOLP}\) nor \(\mathsf{ISOLF}\) is expressible on
\(\mathsf{In}\) from \(\mathcal S'\).  It is an integer-order
non-expressibility theorem, not the rational-order failure in Theorem~IV.5
\cite[ch.~IV, pp.~100--108]{Kamp1968}.

\paragraph{The sparse-square proof.}
Interpret the sole proposition letter by
\[
 P_0
 =
 \{n^2:n\in\omega\}
 \cup
 \{(2n)^2+1:n\in\omega\}.
\]
The source calculates
\[
\begin{aligned}
 \mathsf{ISOLF}(P_0)
 &=
 \{i\in\omega:
   \exists n\in\omega\,((2n)^2\le i<(2n+1)^2)\},\\
 \mathsf{ISOLP}(P_0)
 &=
 \{1\}\cup
 \{i\in\omega:
   \exists n\in\omega\,((2n+3)^2<i\le(2n+4)^2+1)\}.
\end{aligned}
\]
Write \(A(i)\) for the metalinguistic assertion that \(A\) is true at
\(i\in\mathbb Z\) under this fixed interpretation.  The source claims by
structural induction that natural numbers \(n_A,m_A\), with
\(m_A<n_A^2\), can be chosen so that two eventual uniformities hold.  Its
first clause is:

\begin{enumerate}
\item if \(n,n'>n_A\) and
      \[
       n^2+m_A<i<(n+1)^2-m_A,
       \qquad
       (n')^2+m_A<i'<(n'+1)^2-m_A,
      \]
      then \(A(i)\Longleftrightarrow A(i')\).
\end{enumerate}

The second printed clause omits a necessary hypothesis.  It says that, for
all \(n,n'>n_A\) and \(-m_A\le k\le m_A\),
\(A(n^2+k)\Longleftrightarrow A((n')^2+k)\).  This is false already in the
printed atomic case \(A=q_0\), \(n_A=m_A=2\): if \(n>2\) is even,
\(n'>2\) is odd, and \(k=1\), then \(n^2+1\in P_0\) but
\((n')^2+1\notin P_0\).  The error occurs on the source page itself, not only
in a later retelling.

The non-expressibility argument is repaired by adding exactly the missing
parity premise:
\begin{enumerate}
\setcounter{enumi}{1}
\item if \(n,n'>n_A\), \(n\equiv n'\pmod 2\), and
      \(-m_A\le k\le m_A\), then
      \(A(n^2+k)\Longleftrightarrow A((n')^2+k)\).
\end{enumerate}
This is a handbook-supplied correction, not a claim about what Kamp printed.
Here is the verification needed before using it.

For \(q_0\), take \(n_A=m_A=2\).  The central portions in clause~1 contain
no member of \(P_0\); within distance two of a sufficiently large square,
membership is determined by \(k=0\), or by \(k=1\) together with the square
index being even.  Hence the repaired clause~2 holds for indices of the same
parity.

An elementary enlargement observation handles conjunction.  Suppose clauses
1--2 hold for \(A\) with \((n_A,m_A)=(N,m)\), and let \(M\ge m\).  Choose
\(N'\ge N+1\) so large that \(2N'-1>M+m\).  Clause~1 with margin \(M\) is a
restriction of clause~1 with margin \(m\).  In clause~2, offsets of absolute
value at most \(m\) use the old clause~2; every new offset
\(m<|k|\le M\) lies, after this increase of \(N'\), in a central portion of
the gap immediately before or after the relevant square and is therefore
covered by clause~1.  Thus any finite increase of the offset bound is valid
after a corresponding increase of the threshold.  Negation preserves a
common pair of bounds.  For conjunction, take the maximum of the two offset
bounds, enlarge both thresholds as just described, and then take their
maximum.

For \(\mathsf{PAST}_1 A\) and \(\mathsf{FUTURE}_1 A\), truth at \(i\) is
truth of \(A\) at \(i-1\) and \(i+1\), respectively.  Increase the margin by
one and the threshold far enough that translating a central point by one
keeps it in the old central portion.  Matching offsets that remain in the old
bounded window use repaired clause~2; the finitely many new outer offsets lie
in an adjacent central portion and use clause~1.  This proves both repaired
clauses for the two immediate-neighbor operations while preserving the
same-parity condition.

The two unbounded operations need no defective square argument.  If the
truth set of \(A\) is empty, \(\mathsf{PAST}(A)\) is always false; otherwise
it is true at every point later than any chosen member of that truth set.
Likewise, if the truth set of \(A\) is unbounded above,
\(\mathsf{FUTURE}(A)\) is always true, whereas if it is bounded above then
\(\mathsf{FUTURE}(A)\) is false above an upper bound.  Each result is thus
eventually constant, so clauses~1--2 hold with central margin zero and a
sufficiently large threshold.  These cases complete a valid structural
induction for the repaired pair of uniformities.  Kamp instead prints the
parameter choices \((n_A+1,m_A+1)\) for immediate neighbors and
\((n_A+2,0)\) for strict past and future, together with threshold case
splits; because those lines rely on the false unrestricted second clause,
they are source evidence for the intended proof shape, not a valid proof as
printed.

The two displayed isolation sets violate the first uniformity: their truth on
the increasingly long central parts of consecutive-square gaps continues to
alternate with the parity of the gap index.  Hence no formula over
\(\mathcal S'\) defines either isolation tense.  The source's load-bearing
claim is the in-proof induction labelled (1.a) and (1.b); the thesis does not
package it as a separate lemma, and the handbook uses only the corrected
version proved above
\cite[ch.~IV, pp.~105--108]{Kamp1968}.

% Kamp 1968, Chapter IV, Theorems 4 and 5, printed pp. 108--112.

\subsection{Integer irredundancy and the rational-cut obstruction}
\label{sec:p09-kamp-iv4-iv5-separations}

\subsubsection{The converse integer irredundancy}

\begin{hbdefinition}{def:p09-kamp-iv4-basis}{Isolation and unbounded past/future without immediate neighbors}
On \(\mathsf{In}=\langle\mathbb Z,<\rangle\), retain the
unary operations \(\mathsf{ISOLP},\mathsf{ISOLF}\) with the
exact integer clauses in \cref{thm:p09-kamp-iv2-unary-semantics}.
For \(A\subseteq\mathbb Z\), define
\[
\begin{aligned}
 \mathsf{PAST}(A)&=\{i:\exists j<i\ (j\in A)\},\\
 \mathsf{FUTURE}(A)&=\{i:\exists j>i\ (j\in A)\}.
\end{aligned}
\]
These are strict existential operations without a bound on the
distance to the witness. The basis for Theorem IV.4 is exactly
\[
 \mathcal S''=
   \{\mathsf{ISOLP},\mathsf{ISOLF},\mathsf{PAST},
                       \mathsf{FUTURE},\mathsf{NOT},\mathsf{AND}\}.
\]
Its finite sentential formulas start with proposition letters and
use these four unary operations, negation, and conjunction.
The targets \(\mathsf{PAST}_1,\mathsf{FUTURE}_1\) inspect
the immediate integer neighbors, so
\(\mathsf{PAST}_1(A)=\{i:i-1\in A\}\) and
\(\mathsf{FUTURE}_1(A)=\{i:i+1\in A\}\).

For each integer \(r\), put
\(R_r=\{4k+r:k\in\mathbb Z\}\), so that
\(R_{r+4}=R_r\). Write
\(V=R_0\cup R_2\) for the even integers and
\(W=R_1\cup R_3\) for the odd integers.
Complements in this integer subsection are relative to \(\mathbb Z\)
\cite[ch.~IV, thm.~4, pp.~108--109]{Kamp1968}.
\end{hbdefinition}

\begin{hblemma}{thm:p09-kamp-iv4-field}{proved=annex:p9.kamp.iv4.field}{The complete eight-set calculation modulo four}
The family
\[
 \mathcal B_4=
 \{\mathbb Z,\varnothing,R_2,R_2^c,R_0,R_0^c,V,W\}
\]
is a field of subsets of \(\mathbb Z\). It is closed under
every operation of \(\mathcal S''\). The complete values of
the four unary operations on this family are
\[
\begin{array}{c|cc|cc}
 A&\mathsf{ISOLP}(A)&\mathsf{ISOLF}(A)
     &\mathsf{PAST}(A)&\mathsf{FUTURE}(A)\\ \hline
 \mathbb Z&\varnothing&\varnothing&\mathbb Z&\mathbb Z\\
 \varnothing&\varnothing&\varnothing&\varnothing&\varnothing\\
 R_2&\mathbb Z&\mathbb Z&\mathbb Z&\mathbb Z\\
 R_2^c&R_0&R_0&\mathbb Z&\mathbb Z\\
 R_0&\mathbb Z&\mathbb Z&\mathbb Z&\mathbb Z\\
 R_0^c&R_2&R_2&\mathbb Z&\mathbb Z\\
 V&\mathbb Z&\mathbb Z&\mathbb Z&\mathbb Z\\
 W&\mathbb Z&\mathbb Z&\mathbb Z&\mathbb Z
\end{array}
\]
The three atoms of this field are \(R_0,R_2,W\), so no
member distinguishes residue one from residue three. In particular,
if the only occurring input letter is interpreted by \(R_2\),
every finite \(\mathcal S''\)-formula has its truth set in
\(\mathcal B_4\). The same statement holds when every
occurring atomic interpretation already belongs to the field
\cite[ch.~IV, pp.~108--109]{Kamp1968}.
\end{hblemma}

\begin{hbtheorem}{thm:p09-kamp-iv4}{proved=annex:p9.kamp.iv4}{Immediate past and future are not definable from the remaining integer basis}
Neither \(\mathsf{PAST}_1\) nor \(\mathsf{FUTURE}_1\)
is expressible on \(\mathsf{In}\) by a finite one-input
formula over \(\mathcal S''\). Equivalence would have to
hold for every \(A\subseteq\mathbb Z\) at every integer.
The fixed input \(A=R_2\) already prevents either definition:
\[
              \mathsf{PAST}_1(R_2)=R_3,
       \qquad \mathsf{FUTURE}_1(R_2)=R_1,
\]
and neither target set belongs to \(\mathcal B_4\).
At the two current points one and three, every candidate formula
has equal truth values, whereas each target has unequal truth values
\cite[ch.~IV, thm.~4, pp.~108--109]{Kamp1968}.
\end{hbtheorem}

\subsubsection{Two interpretations on the rational order}

\begin{hbdefinition}{def:p09-kamp-rational-cut-models}{The periodic and accumulating rational propositions}
Now the fixed time domain is
\(\mathsf{Ra}=\langle\mathbb Q,<\rangle\).
For real numbers \(a<b\), write
\((a,b)_{\mathbb Q}=\{r\in\mathbb Q:a<r<b\}\).
Thus real interval endpoints specify subsets of the rational domain;
they are not extra time points or constants in the first-order language.
Put \(\alpha=\sqrt2\) and \(b_n=\alpha+1/n\) for
integers \(n\geq1\). All union indices below are integers.
Define
\[
\begin{aligned}
 Q&=\bigcup_{k\in\mathbb Z}
                   (\alpha+2k,\alpha+2k+1)_{\mathbb Q},\\
 Q'&=\bigcup_{\substack{k\in\mathbb Z\\k\ne0}}
                   (\alpha+2k,\alpha+2k+1)_{\mathbb Q}\\
   &\quad\cup\
       \bigcup_{j\geq0}(b_{2j+2},b_{2j+1})_{\mathbb Q}.
\end{aligned}
\]
The second interpretation replaces the original band
\((\alpha,\alpha+1)_{\mathbb Q}\) by alternating bands
whose endpoints decrease toward the irrational cut \(\alpha\).
In the displayed intervals the smaller endpoint is always first.

Let \(\mathcal M=\langle\mathbb Q,<,Q\rangle\) and
\(\mathcal M'=\langle\mathbb Q,<,Q'\rangle\), interpreting
the sole unary predicate \(P_0\) by these sets. The same sets
interpret \(q_0\) in the source's temporal language
\(\mathsf{TL}_1\), whose constructors are negation,
conjunction, strict \(S\), and strict \(U\), with the
witness operand first and the open-interval operand second.
Complements from this point onward are relative to \(\mathbb Q\)
\cite[ch.~IV, thm.~5 construction, pp.~109--110]{Kamp1968}.
\end{hbdefinition}

\begin{hblemma}{thm:p09-kamp-rational-local-table}{proved=annex:p9.kamp.rational.local.table}{All temporal operations on the four local truth-set choices}
Each of \(Q,Q'\) is locally constant on the rational domain:
every rational point has a rational open neighborhood on which its
membership is constant. Each set and its complement are unbounded
in both directions. Neither an irrational band endpoint nor the
accumulation cut is a point of the time domain.

More generally, let \(A\subseteq\mathbb Q\) have these
local-constancy and two-sided unboundedness properties. For
\(X,Y\in\{\varnothing,A,A^c,\mathbb Q\}\), both
strict temporal operations have the following complete table,
with the witness set \(X\) indexing rows and the interval
set \(Y\) indexing columns:
\[
\begin{array}{c|cccc}
 S(X,Y)=U(X,Y)&\varnothing&A&A^c&\mathbb Q\\ \hline
 \varnothing&\varnothing&\varnothing&\varnothing&\varnothing\\
 A&\varnothing&A&\varnothing&\mathbb Q\\
 A^c&\varnothing&\varnothing&A^c&\mathbb Q\\
 \mathbb Q&\varnothing&A&A^c&\mathbb Q
\end{array}
\]
The same table therefore governs the two specified interpretations
\cite[ch.~IV, pp.~110--112]{Kamp1968}.
\end{hblemma}

\begin{hblemma}{thm:p09-kamp-iv5-temporal-induction}{proved=annex:p9.kamp.iv5.temporal.induction}{The full one-input temporal indistinguishability induction}
For every finite \(\mathsf{TL}_1\)-formula \(\eta\) with
no proposition letters other than \(q_0\), the pair of truth sets
\[
  \bigl(\llbracket\eta\rrbracket_{\mathcal M},
        \llbracket\eta\rrbracket_{\mathcal M'}\bigr)
\]
is one of
\[
 (\varnothing,\varnothing),\quad
 (\mathbb Q,\mathbb Q),\quad (Q,Q'),\quad (Q^c,(Q')^c).
\]
Consequently all three source comparison assertions hold for any
\(r,s\in\mathbb Q\):

\begin{enumerate}
\item If \(r\in Q\) exactly when \(s\in Q\), then
      \(\eta\) has equal truth values at \(r,s\) in
      \(\mathcal M\).
\item If \(r\in Q'\) exactly when \(s\in Q'\), then
      \(\eta\) has equal truth values at \(r,s\) in
      \(\mathcal M'\).
\item If \(r\in Q\) exactly when \(s\in Q'\), then
      \(\eta\)'s truth at \(r\) in \(\mathcal M\)
      equals its truth at \(s\) in \(\mathcal M'\).
\end{enumerate}

In particular every such formula agrees at the rational current point
three in the two interpretations
\cite[ch.~IV, thm.~5 proof, pp.~110--112]{Kamp1968}.
\end{hblemma}

\begin{hbdefinition}{def:p09-kamp-rational-distinguishing-formula}{A first-order condition for the accumulating obstruction}
The following equality-free formula uses only \(<,P_0\), and
has exactly the current point \(t_0\) free:
\[
\begin{aligned}
 \beta(t_0)=\exists t_1\bigl[&t_1<t_0\land\neg P_0(t_1)\land
           \exists t_2(t_1<t_2<t_0\land P_0(t_2))\\
 &\land\forall t_2\bigl((t_1<t_2<t_0\land P_0(t_2))\to\\
 &\hspace{18mm}\exists t_3\exists t_4
      (t_1<t_3<t_4<t_2\land P_0(t_3)\land\neg P_0(t_4))\bigr)
 \bigr].
\end{aligned}
\]
All quantified variables range over \(\mathbb Q\). Put
\(\varphi(t_0)=\neg\beta(t_0)\).
The witness \(t_1\) for \(\beta\) is a false point with
a true point between it and the current time. Every true point
in that interval must have, still earlier in the interval, a
true point followed by a false point. Thus no first true block
after \(t_1\) can account for all subsequent true points.
The existential true-point clause excludes a vacuous instance.

This gives a corrected distinguishing formula for the source's
two-interpretation construction. Its definition contains no
\(\sqrt2\), real endpoint, or named rational parameter;
those enter only the chosen interpretations used to prove separation.
\end{hbdefinition}

\begin{hbtheorem}{thm:p09-kamp-iv5}{proved=annex:p9.kamp.iv5}{A first-order unary rational tense beyond SINCE and UNTIL}
The formula \(\varphi\) defines a unary
\(\mathsf{Ra}\)-tense which is not expressible in terms of
\(\mathsf{NOT},\mathsf{AND},\mathsf{SINCE},\mathsf{UNTIL}\).
Precisely, there is no finite one-input \(\mathsf{TL}_1\)
formula equivalent to \(\varphi\) for every unary predicate
interpretation on \(\mathbb Q\) and every rational current point.
For the explicit interpretations above,
\[
                     \mathcal M\models\varphi[t_0\mapsto3],
            \qquad \mathcal M'\not\models\varphi[t_0\mapsto3],
\]
while every one-input temporal formula agrees there. The order
domain remains \(\mathbb Q\) throughout; the real cut is not
adjoined as a possible time point. Hence the complete-order
expressive-completeness theorem cannot be extended unchanged to
all linear orders
\cite[ch.~IV, thm.~5, pp.~109--112]{Kamp1968}.
\end{hbtheorem}

\begin{hbexample}{def:p09-kamp-iv5-cut-example}{Concrete witnesses around the missing point}
Since \(1<\sqrt2<3/2\), the rational point one lies in
the false band immediately below \(\alpha\), and two lies
in the true band \((b_2,b_1)_{\mathbb Q}\) in both models.
The current point three lies in
\((\alpha+1,\alpha+2)_{\mathbb Q}\), which is false in
both. In \(\mathcal M'\), take \(t_1=1\) as the witness
for \(\beta(3)\). Before every true point between one and
three, sufficiently small reciprocal bands supply a true point
followed by a false point. In \(\mathcal M\), any false
starting point with a later true point has a first subsequent
true band; a point inside that band prevents the universal clause
of \(\beta\). The annex proves these assertions for every
quantified point, including all rational positions near the cut.
\end{hbexample}

% Kamp 1968, final substantive argument, printed pp. 112--117.

\subsection{Present progressive and the limits of ordinary past and future}
\label{sec:p09-kamp-present-progressive}

\paragraph{The formal question and its linguistic motivation.}
The thesis ends its substantive text by returning to the progressive
operation announced on pp.~18--19. Kamp considers two possible
formalizations and suggests that the use of progressive constructions in
natural language might reflect a conception of time as dense. That is
an interpretive suggestion, without an empirical demonstration in this
argument. The mathematical results below concern the explicitly defined
set operations and their expressibility; they do not establish a claim
about how speakers of any language represent time
\cite[pp.~18--19, 112--113]{Kamp1968}.

\begin{hbdefinition}{def:p09-kamp-progressive-language}{The present atom and the ordinary past/future language}
Fix a nonempty strict linear order \(\mathcal J=\langle T,<\rangle\).
For \(A\subseteq T\), put
\[
 \mathsf{PAST}(A)=\{x:\exists y<x\ (y\in A)\},\qquad
 \mathsf{FUTURE}(A)=\{x:\exists y>x\ (y\in A)\}.
\]
The source's \(\mathsf{TL}_0\) consists of finite formulas
generated from proposition letters by \(\neg,\land,P,F\),
where the unary symbols \(P,F\) denote these strict existential
operations. A one-input schema contains \(q_0\) and no other
proposition letter. The atom \(q_0\) denotes the input at the
current point, so the present is available throughout. Other finite
truth-functional connectives are Boolean abbreviations. There are
no named time points, additional predicate parameters, immediate-neighbor
operators, or isolation operators in this basis.

For \(a\leq b\) in \(T\), the symbols \((a,b)\) and
\([a,b]\) denote the open and closed intervals in the actual
order. They do not assert density or the existence of points between
the endpoints. For a formula \(\eta\) and input \(A\), write
\(\llbracket\eta\rrbracket_A\subseteq T\) for its truth set
\cite[pp.~18--19, 114]{Kamp1968}.
\end{hbdefinition}

\begin{hbdefinition}{def:p09-kamp-progressive-two-definitions}{Two distinct proposed progressive operations}
The initial proposal is
\[
\begin{aligned}
 \rho_0(t_0)={}&P_0(t_0)\land
   \exists t_1\bigl(t_1<t_0\land
        \forall t_2(t_1<t_2<t_0\to P_0(t_2))\bigr)\\
 &\land\exists t_3\bigl(t_0<t_3\land
        \forall t_4(t_0<t_4<t_3\to P_0(t_4))\bigr).
\end{aligned}
\]
Its operation is denoted \(\mathsf{PR}_0\). The revised
proposal additionally requires the two strict witnesses to satisfy
the input predicate:
\[
\begin{aligned}
 \rho(t_0)={}&P_0(t_0)\land
   \exists t_1\bigl(t_1<t_0\land P_0(t_1)\land
        \forall t_2(t_1<t_2<t_0\to P_0(t_2))\bigr)\\
 &\land\exists t_3\bigl(t_0<t_3\land P_0(t_3)\land
        \forall t_4(t_0<t_4<t_3\to P_0(t_4))\bigr).
\end{aligned}
\]
Write \(\mathsf{PR}\) for the operation defined by \(\rho\).
Both are equality-free first-order unary tenses with only \(t_0\)
free and only \(P_0\) in their unary signature. Their quantifiers
range over \(T\). In particular the revised operation is
\[
 \mathsf{PR}(A)=\{x\in T:\exists a<x<b\quad[a,b]\subseteq A\}.
\]
It is this revised operation, already specified in the initial
announcement on p.~19, which the final nonexpressibility result
addresses
\cite[pp.~19, 112--113]{Kamp1968}.
\end{hbdefinition}

\begin{hblemma}{thm:p09-kamp-progressive-definitions}{proved=annex:p9.kamp.progressive.definitions}{Dense, discrete and endpoint behavior of the two proposals}
On every strict linear order,
\(\mathsf{PR}(A)\subseteq\mathsf{PR}_0(A)\subseteq A\),
and both operations are false at a first or last time point.
On a dense order, \(\mathsf{PR}=\mathsf{PR}_0\), including
when the order has endpoints.

On a discrete order in the source's conditional-neighbor sense,
let \(T_*\) be the set of points having both an earlier and a
later point. Each \(x\in T_*\) has an immediate predecessor
\(x^-\) and successor \(x^+\), and
\[
 \mathsf{PR}_0(A)=A\cap T_*,\qquad
 \mathsf{PR}(A)=\{x\in T_*:x^-,x,x^+\in A\}.
\]
Thus the weak proposal equals the present operation on an endpoint-free
discrete order, whereas the revised one also tests the adjacent points.
On a discrete order with endpoints, the weak proposal has the expression
\[
 q_0\land P\neg(q_0\land\neg q_0)
       \land F\neg(q_0\land\neg q_0),
\]
which reduces to \(q_0\) when there are no endpoints.
For example, on \(\mathbb Z\) the singleton input \(\{0\}\)
has weak progressive value \(\{0\}\) and revised value
\(\varnothing\).

Using the source's strict, witness-first binary operators, on all
linear orders the revised operation has the expression
\[
                         q_0\land S(q_0,q_0)\land U(q_0,q_0).
\]
This expression belongs to \(\mathsf{TL}_1\), which has the
binary constructors \(S,U\); the nonexpressibility question below
concerns the smaller \(\mathsf{TL}_0\) basis
\cite[pp.~112--113]{Kamp1968}.
\end{hblemma}

\begin{hbdefinition}{def:p09-kamp-progressive-depths}{The two source directional depths}
Define even natural numbers \(d_-(\eta),d_+(\eta)\) by
\[
\begin{array}{c|cc}
 \eta&d_-(\eta)&d_+(\eta)\\ \hline
 q_0&0&0\\
 \neg\xi&d_-(\xi)&d_+(\xi)\\
 \xi\land\theta&\max(d_-(\xi),d_-(\theta))
                    &\max(d_+(\xi),d_+(\theta))\\
 P\xi&d_-(\xi)+2&d_+(\xi)\\
 F\xi&d_-(\xi)&d_+(\xi)+2
\end{array}
\]
The quantity \(d_-\) is the source's depth for its increasing
sequence argument, and \(d_+\) is its depth for the reversed
argument. In particular future adds no cost to \(d_-\), and
past adds no cost to \(d_+\). These are defined numerical
measures of finite syntax, not distances between time points
\cite[pp.~114--117]{Kamp1968}.
Unless a two-sided index set \(\mathbb Z\) is displayed,
all indices \(i,j,m,n\) below are nonnegative integers.
\end{hbdefinition}

\subsubsection{The source's sequence constructions with explicit local scope}

\begin{hblemma}{thm:p09-kamp-progressive-two-sided}{proved=annex:p9.kamp.progressive.two.sided}{The two-sided sequence construction}
Suppose a strictly increasing sequence \((a_n)_{n\in\mathbb Z}\)
of points of \(T\) is given. Put
\[
\begin{aligned}
 I&=\{x:\exists m,n\in\mathbb Z\ (a_m<x<a_n)\},\\
 A&=[a_{-2},a_2]\cup\{a_{2n}:n\in\mathbb Z\}.
\end{aligned}
\]
For every one-input \(\mathsf{TL}_0\)-formula \(\eta\),
\[
 \llbracket\eta\rrbracket_A\cap I
          \in\{\varnothing,I,A\cap I,I\setminus A\}.
\]
But \(\mathsf{PR}(A)=(a_{-2},a_2)\). It contains
\(a_0\) and excludes \(a_4\), although both belong to
\(A\cap I\). Thus \(\mathsf{PR}\) is not expressible
from the stated basis on any order with this sequence.
The sequence is not assumed cofinal in the whole order: the invariant
is on its convex hull \(I\), and temporal quantifiers still range
over all of \(T\)
\cite[pp.~113--114]{Kamp1968}.
\end{hblemma}

\begin{hblemma}{thm:p09-kamp-progressive-one-sided}{proved=annex:p9.kamp.progressive.one.sided}{Both one-sided sequence and depth constructions}
For a strictly increasing sequence \((a_n)_{n\geq0}\), put
\[
\begin{aligned}
 I_+&=\{x:a_0\leq x<a_n\text{ for some }n\geq0\},\\
 H_m^+&=\{x\in I_+:a_m<x\},\\
 A_i^+&=\{a_{2n}:n\geq0\}\cup[a_{i+2},a_{i+4}],
                   \qquad i\text{ even}.
\end{aligned}
\]
If \(i\geq d_-(\eta)\), put \(H=H_{d_-(\eta)}^+\).
Then
\[
 \llbracket\eta\rrbracket_{A_i^+}\cap H
             \in\{\varnothing,H,H\cap A_i^+,H\setminus A_i^+\}.
\]
The exact progressive image is
\(\mathsf{PR}(A_i^+)=(a_{i+2},a_{i+4})\).

For a strictly decreasing sequence \((b_n)_{n\geq0}\), instead put
\[
\begin{aligned}
 I_-&=\{x:b_n<x\leq b_0\text{ for some }n\geq0\},\\
 H_m^-&=\{x\in I_-:x<b_m\},\\
 A_i^-&=\{b_{2n}:n\geq0\}\cup[b_{i+4},b_{i+2}],
                   \qquad i\text{ even}.
\end{aligned}
\]
If \(i\geq d_+(\eta)\), put \(K=H_{d_+(\eta)}^-\).
Then
\[
 \llbracket\eta\rrbracket_{A_i^-}\cap K
             \in\{\varnothing,K,K\cap A_i^-,K\setminus A_i^-\},
\]
and
\(\mathsf{PR}(A_i^-)=(b_{i+4},b_{i+2})\).
Either sequence therefore supplies a counterinterpretation for any
proposed one-input definition. The inequalities on the even block
indices are part of the construction. A first or last point of the
whole order need not belong to either local tail
\cite[pp.~114--116]{Kamp1968}.
\end{hblemma}

\subsubsection{A finite version of the two-ended construction}

\begin{hbdefinition}{def:p09-kamp-progressive-window}{Separated marker chains and their filled intervals}
Choose a strictly increasing chain \((a_n)\) and a strictly decreasing chain
\((b_n)\), with \(a_n<b_m\) for every pair of available
indices. The indices may run over all nonnegative integers, as in
the source's two-chain construction, or over a finite segment
\(0,\ldots,N\). Choose even nonnegative \(i,j\); in the
finite case require \(\max(i,j)+6\leq N\). Put
\[
\begin{aligned}
 A_{ij}={}&\{a_{2n}:2n\text{ is available}\}
       \cup\{b_{2n}:2n\text{ is available}\}\\
       &\cup[a_{i+2},a_{i+4}]
        \cup[b_{j+4},b_{j+2}],\\
 C_{mn}={}&\{x\in T:a_m<x<b_n\}.
\end{aligned}
\]
The intervals include all their points in \(T\), not just the
selected markers. Both chains and their interval endpoints are
data for the interpretation; they are not constants available in
candidate formulas. The finite version is sufficient for the
unrestricted infinite-order conclusion.
\end{hbdefinition}

\begin{hblemma}{thm:p09-kamp-progressive-window}{proved=annex:p9.kamp.progressive.window}{The complete two-depth window invariant}
For the construction above, if
\(d_-(\eta)\leq i\) and \(d_+(\eta)\leq j\), put
\(C=C_{d_-(\eta),d_+(\eta)}\). Then
\[
 \llbracket\eta\rrbracket_{A_{ij}}\cap C
              \in\{\varnothing,C,C\cap A_{ij},C\setminus A_{ij}\}.
\]
This holds for both finite and infinite marker chains, with temporal
quantification over the whole order. Moreover,
\[
 \mathsf{PR}(A_{ij})
       =(a_{i+2},a_{i+4})\cup(b_{j+4},b_{j+2}).
\]
Both intervals lie inside the displayed central window. The points
\(a_{i+3},a_{i+6}\) both belong to the input and the window,
but only the first belongs to the progressive image
\cite[p.~117]{Kamp1968}.
\end{hblemma}

\begin{hbtheorem}{thm:p09-kamp-present-progressive}{proved=annex:p9.kamp.present.progressive}{Present progressive is not expressible by ordinary past and future on an infinite order}
For every infinite strict linear order \(\mathcal J=\langle T,<\rangle\),
there is no finite one-input \(\mathsf{TL}_0\)-schema \(\eta\)
such that
\[
                  \forall A\subseteq T\qquad
                      \llbracket\eta\rrbracket_A=\mathsf{PR}(A).
\]
The permitted primitive operations are exactly negation, conjunction,
strict existential \(\mathsf{PAST}\), and strict existential
\(\mathsf{FUTURE}\), with the present atom \(q_0\) available.
The order may have neither endpoint, just a first endpoint, just a
last endpoint, or both. No density, discreteness, completeness,
countable cofinality, or existence of separated infinite chains is
assumed in this conclusion.

For a candidate formula, it suffices to choose two separated finite
chains of length \(N+1\), where
\(N=\max(d_-(\eta),d_+(\eta))+6\), and use
\(A_{ij}\) with \(i=d_-(\eta)\), \(j=d_+(\eta)\).
Thus the general conclusion uses a construction from finitely many
selected points inside any infinite order
\cite[pp.~19, 112--117]{Kamp1968}.
\end{hbtheorem}

\paragraph{Relation to the complete-order expressive theorem.}
The revised progressive operation is first-order definable, and its
\(q_0\land S(q_0,q_0)\land U(q_0,q_0)\) representation works
on every linear order. The result above concerns the smaller ordinary
past/future language. It is consistent with, and illustrates the need
for a richer basis in, the earlier expressive-completeness result over
\(\mathsf{Com}\).

\begin{hbexample}{def:p09-kamp-progressive-boundaries}{Parity, endpoints and the bounded infinite order}
On \(\mathbb N\), let \(a_n=n\) and choose the even block
index \(i=0\). Then
\(A_0^+=\{0,2,4,6,\ldots\}\cup[2,4]\) has
\(\mathsf{PR}(A_0^+)=\{3\}\). Both three and six satisfy
the input, while only three satisfies the progressive operation.
If the block index were instead odd, the exact-image statement would
need additional qualifications: with \(i=1\), the same construction
contains every integer from two through six, and its progressive image
is \(\{3,4,5\}\).

At the first point of \(\mathbb N\), both progressive proposals
are false even when the input contains that point. On the bounded
infinite order \(\omega+1\), a decreasing sequence indexed by all
nonnegative integers does not exist, although both a first and a last
point exist. The finite-window construction still applies, by choosing
the required markers among finitely many natural numbers below the
last point. The annex verifies these boundary examples.
\end{hbexample}


\paragraph{Adjacent separations and proof status.}
The complete real-order nonexpressibility proof is
\cref{thm:p09-kamp-iv1}; its three lemmas use the even/odd-integer
interpretation inside \(\mathbb R\).
\Cref{thm:p09-kamp-dense-unary-separation} supplies the separate
extension to every nontrivial dense order under its stated language.
Theorem~IV.2 positively expresses \(S,U\) on the integers from a basis which
includes the two isolation tenses, with corrected explicit expressions and
complete proofs in \cref{thm:p09-kamp-iv2}.
Theorem~IV.4 gives the converse
irredundancy: a modulo-four interpretation shows that immediate past/future
cannot be recovered from isolation plus ordinary strict past/future,
as proved in \cref{thm:p09-kamp-iv4}.
The rational-cut counterexample is proved in \cref{thm:p09-kamp-iv5}.
None may
be relabelled as Chapter~IV Theorem~3
\cite[ch.~IV, pp.~95--112]{Kamp1968}.

The thesis supplies the Chapter~III construction, while abbreviating symmetric
or routine subcases.  It also prints the structural induction intended to
prove Theorem~IV.3, but the missing parity premise makes that induction
invalid as written; the corrected argument above is the handbook's repair,
not an attribution to Kamp
\cite[ch.~III, pp.~39--94; ch.~IV, pp.~105--108]{Kamp1968}.

\subsubsection{Pnueli's 1977 transition system and proof systems}

\paragraph{Source and program model.}
Pnueli's NYU memorial publication list supplies an author-hosted copy of the
twelve-page FOCS paper, printed pp.~46--57.  Its May 1981 archival wrapper is
not a second edition.  Citations use the printed page numbers.  Neither the
wrapper nor the printed body grants an open-content or republication licence,
and the presentation below is independently worded \cite{Pnueli1977}.

The paper begins with a dynamic discrete system
\[
 \Sigma=\langle S,R,s_0\rangle,
 \qquad R\subseteq S\times S,
\]
where \(S\) may be infinite and \(s_0\) is the initial state.  An execution is
an infinite sequence
\[
 \sigma=s_0,s_1,\ldots
 \quad\text{such that}\quad
 R(s_i,s_{i+1})\ \text{for every }i\ge 0.
\]
The transition relation may be nondeterministic.  To type the paper's
sequential specialization, let \(L\) be a finite set of control locations and
\(D\) a data domain.  It takes
\[
 S=L\times D,
 \qquad
 N:L\times D\to L,
 \qquad
 T:L\times D\to D,
\]
and defines
\[
 R(\langle\pi,u\rangle,\langle\pi',u'\rangle)
 \quad\Longleftrightarrow\quad
 \pi'=N(\pi,u)\ \land\ u'=T(\pi,u).
\]
The source assumes deterministic \(N,T\) only to simplify notation and says
that the restriction is not essential.

For \(n\) concurrent processors, take finite location sets \(L_i\),
\(S=L_1\times\cdots\times L_n\times D\), and maps
\[
 N_i:L_i\times D\to L_i,
 \qquad
 T_i:L_i\times D\to D.
\]
The exact atomic interleaving is
\[
\begin{split}
&R(\langle\pi_1,\ldots,\pi_n;u\rangle,
   \langle\pi'_1,\ldots,\pi'_n;u'\rangle)\\
&\quad\Longleftrightarrow\quad
 \exists i\in\{1,\ldots,n\}\,
 \left(
   \pi'_i=N_i(\pi_i,u)\land
   u'=T_i(\pi_i,u)\land
   \bigwedge_{h\ne i}\pi'_h=\pi_h
 \right).
\end{split}
\]
Thus one processor is selected nondeterministically and completes one
statement while all other control coordinates remain fixed.  The paper calls
this atomic granularity a modelling choice: a statement may instead be split
into smaller atomic steps when interference inside it matters.  Its early
progress discussion informally requires that a processor which remains enabled
not be postponed indefinitely.  The distinct formal action-occurrence condition
printed later is stated separately below
\cite[pp.~46--47]{Pnueli1977}.

\paragraph{Invariance and the source's temporal implication.}
Define the finite powers and closure of \(R\) by
\[
\begin{aligned}
 R^0(s,t)&\Longleftrightarrow s=t,\\
 R^{m+1}(s,t)&\Longleftrightarrow
   \exists v\in S\,(R^m(s,v)\land R(v,t)),\\
 R^*&=\bigcup_{m\ge0}R^m,
 \qquad
 X=\{s\in S:R^*(s_0,s)\}.
\end{aligned}
\]
Thus \(R^*\) is reflexive--transitive and \(s_0\in X\); this unpacks what the
paper informally calls the ``transitive closure.''  A \emph{state assertion}
is a map
\(p:S\to\{\mathsf{false},\mathsf{true}\}\).  It describes one state and is
not yet a tense formula.  Such an assertion is invariant exactly when
\(\forall s\in X\,p(s)\).

Partial correctness constrains the output if an exit is reached.  The source's
separate \emph{clean execution} condition is an invariant prohibition on
specified illegal intermediate operations or states, such as division by
zero, numerical overflow, or an out-of-range subscript.  The distinction is
needed because a satisfactory exit value does not show that the intervening
execution was legal.  Mutual exclusion and absence of deadlock are further
invariant examples; the producer--consumer calculation below uses this same
invariant standing.

Within this subsubsection, \(\varphi\leadsto\psi\) names only Pnueli's printed
temporal-implication statement family: it is a relation between state
assertions, not a new primitive in the handbook's LTL grammar.  The separate
notation is needed because the paper prints a strict-future clause:
\[
 \varphi\leadsto\psi
 \quad\Longleftrightarrow\quad
 \forall\sigma=s_0s_1\cdots\ \forall i\ge0\
 \bigl(\varphi(s_i)\Rightarrow
       \exists j>i\ \psi(s_j)\bigr).
\]
Total correctness is the instance taking entry-with-precondition to
exit-with-postcondition.  For cyclic programs, accessibility says that a
process which announces an intention to enter a critical region eventually
does so; responsiveness says that every request is eventually answered.
The source expressly leaves specifications involving three or more ordered
events and quantitative response times outside its two basic cases
\cite[pp.~48--49]{Pnueli1977}.

\paragraph{Invariant induction, ranks, and ER.}
The source's first general principle, \(({\rm P1})\), is ordinary
reachable-state induction:
\[
 \frac{
   p(s_0)
   \qquad
   \forall s,s'\bigl(p(s)\land R(s,s')\Rightarrow p(s')\bigr)}
  {\forall s\in X\ p(s)}.
\]
For the second principle, the source says that \(A(s,n)\) uses a positive
natural rank, but its displayed decrement \(A(s',n-1)\) requires a value at
zero when \(n=1\).  To type the printed expression, let
\[
 A:S\times\{0,1,2,\ldots\}
   \longrightarrow\{\mathsf{false},\mathsf{true}\}
\]
be a metalinguistic rank predicate while rule instances quantify \(n>0\).
Allowing the zero-rank argument and supplying the unprinted universal closures
are handbook typing repairs; the current-state target and rule conclusion below
remain literal.  The displayed natural-number rule is
\[
 \frac{
  \begin{gathered}
   \forall s\bigl(\varphi(s)\Rightarrow
                    \exists n>0\ A(s,n)\bigr)\\
   \forall s,s'\ \forall n>0\,
   \bigl(A(s,n)\land R(s,s')\Rightarrow
         A(s',n-1)\lor\psi(s)\bigr)
  \end{gathered}}
  {\varphi\leadsto\psi}\ ({\rm P2}).
\]
The occurrence \(\psi(s)\) is the source's current-state occurrence, not a
handbook replacement by \(\psi(s')\).  The paper presents only the
natural-number version and refers to the literature for an extension to other
well-founded sets \cite[p.~49]{Pnueli1977}.

Totalizing \(A\) at zero makes the display well typed but does not make it
sound.  Even if eventuality included the present, take
\(S=\{a,b\}\), \(s_0=a\), \(R=\{(a,b),(b,b)\}\), let \(\varphi\) hold only at
\(a\), let \(\psi\) be false everywhere, and let \(A\) hold exactly at
\((a,1)\) and \((b,0)\).  The first premise holds, and the only true
positive-rank transition premise is discharged by \(A(b,0)\), yet no execution
state satisfies \(\psi\).  This is a handbook counterexample to the zero-rank
totalization, not a defect caused by strict futurity.  One sufficient handbook
repair for inclusive eventuality would add
\[
 \forall s\in S\ \neg A(s,0).
\]
If \(\psi\) never held, the transition premise would then lower a positive rank
at every step until it demanded the forbidden rank zero.  The paper prints no
such boundary premise, and the handbook does not add it to P2.

Literal P2 independently conflicts with the printed strict-future meaning of
\(\leadsto\).  Let \(S=\{a,b\}\), \(s_0=a\),
\(R=\{(a,b),(b,b)\}\), let \(\varphi\) and \(\psi\) hold exactly at \(a\),
and let \(A\) hold exactly at \((a,1)\).  Both displayed premises hold: at the
only ranked transition the disjunct \(\psi(a)\) is true.  But the unique
execution has no \(\psi\)-state strictly after position \(0\), so
\(\varphi\leadsto\psi\) fails.  This is a handbook counterexample to the
literal combination of P2 with \(j>i\); the handbook does not silently repair
the source rule.

Its third principle is the system ER.  A1--R3 use state assertions
\(p,q,r:S\to\{\mathsf{false},\mathsf{true}\}\).  R4 additionally needs the
parameter suppressed by the source display.  Fix a nonempty set \(U\), take
\[
 \widehat p:S\times U\to\{\mathsf{false},\mathsf{true}\},
 \qquad q:S\to\{\mathsf{false},\mathsf{true}\},
\]
where the hat keeps this parameterized family distinct from the ordinary state
assertion \(p\).  For \(u\in U\), define the state assertion
\(\widehat p_u(s)=\widehat p(s,u)\) and the pointwise existential state assertion
\[
 (\exists_U\widehat p)(s)\quad\Longleftrightarrow\quad
 \exists u\in U\ \widehat p(s,u).
\]
These parameterized maps are needed only to type R4; they do not add a tense
constructor.  With that convention, the displayed ER core is
\[
\begin{array}{c@{\qquad}c}
 \displaystyle
 \frac{\forall s,s'\,
       (p(s)\land R(s,s')\Rightarrow q(s'))}
      {p\leadsto q}\ ({\rm A1})
 &
 \displaystyle
 \frac{\forall s\,(p(s)\Rightarrow q(s))}
      {p\leadsto q}\ ({\rm A2})
 \\[2.2ex]
 \displaystyle
 \frac{p\leadsto q\quad
       \forall s,s'\,
       (r(s)\land R(s,s')\Rightarrow r(s'))}
      {(p\land r)\leadsto(q\land r)}\ ({\rm R1})
 &
 \displaystyle
 \frac{p\leadsto q\quad q\leadsto r}
      {p\leadsto r}\ ({\rm R2})
 \\[2.2ex]
 \displaystyle
 \frac{p_1\leadsto q\quad p_2\leadsto q}
      {(p_1\lor p_2)\leadsto q}\ ({\rm R3})
 &
 \displaystyle
 \frac{\widehat p_u\leadsto q\quad(u\in U\text{ arbitrary})}
      {(\exists_U\widehat p)\leadsto q}\ ({\rm R4}).
\end{array}
\]
The R4 premise means a derivation uniform in the arbitrary parameter \(u\):
\(q\) and every open assumption are independent of \(u\).  The source prints
neither these parameter types nor the latter variable-hygiene condition; the
handbook supplies them to state its first-order reading.  ER also takes the
theorems of first-order predicate calculus as axioms.  R1 frames an eventuality
with an invariant, R2 composes inevitable stages, R3 joins cases, and R4 hides
a witness.

The source next invokes ordinary integer induction at the metalevel.  To type
that use, specialize the parameter domain to
\(\mathbb N_0=\{0,1,2,\ldots\}\) and put
\(\widehat p_n(s)=\widehat p(s,n)\).  Every \(\widehat p_n\) is a state assertion,
and \(q\) is independent of \(n\).  Its derived induction principle is
\[
 \frac{
   \vdash \widehat p_0\leadsto q
   \qquad
   \widehat p_n\leadsto q\ \vdash\ \widehat p_{n+1}\leadsto q}
  {\vdash \widehat p_n\leadsto q\quad\text{for every }n\in\mathbb N_0}
  \ ({\rm I}).
\]
Here the second premise is a derivation under the displayed temporary
hypothesis for arbitrary \(n\); \((\mathrm I)\) is a proof metarule, not an
additional temporal connective.  The already defined pointwise assertion is
now \((\exists_{\mathbb N_0}\widehat p)(s)\), equivalent to
\(\exists n\in\mathbb N_0\,\widehat p(s,n)\).  After I yields
\(\widehat p_n\leadsto q\) for every \(n\), the typed R4 instance gives
\((\exists_{\mathbb N_0}\widehat p)\leadsto q\).

Theorem~1 calls ER sound and complete for statements
\(p\leadsto q\).  The paper labels its completeness argument a proof sketch.
It assumes an assertion language expressive enough to define
\(P:S\times\mathbb N_0\to\{\mathsf{false},\mathsf{true}\}\), where the state
assertion \(P_n(s)=P(s,n)\) means that every execution beginning at \(s\)
reaches a \(q\)-state within at most \(n\) steps.  Its ``bounded
nondeterminism'' hypothesis is the
pointwise finite-branching condition
\[
 \forall s\in S\quad
 \{s'\in S:R(s,s')\}\ \text{is finite};
\]
no single uniform numerical bound over all states is asserted.  If no finite
reachability bound existed, K\"onig's infinity lemma would give an infinite
execution avoiding \(q\).  Principle~I derives
\(P_n\leadsto q\) for every \(n\), R4 forms the pointwise assertion
\(s\mapsto\exists n\,P(s,n)\), and R2 with A2
concludes the requested implication
\cite[pp.~49--50]{Pnueli1977}.  These expressiveness and branching hypotheses,
and the fact that the proof is a sketch, are part of the source theorem.

The sketch has two independent unspoken gaps.  First, its step from
\(p\leadsto q\) to the pointwise assertion
\(p(s)\Rightarrow\exists n\,P(s,n)\) ranges over arbitrary \(s\), although the
paper defined system executions only from \(s_0\).  Let
\[
 S=\{s_0,u\},\qquad
 R=\{(s_0,s_0),(u,u)\},
\]
let \(p\) hold only at \(u\), and let \(q\) be false everywhere.  The only
system execution from \(s_0\) never satisfies \(p\), so \(p\leadsto q\) holds
vacuously.  But the infinite path \(u,u,\ldots\) avoids \(q\), so \(P(u,n)\)
is false for every \(n\), and the pointwise assertion fails at \(u\).  This
counterexample is serial and finitely branching.  The printed step therefore
needs either \(S=X\), a consistent restriction of assertion and rule semantics
to \(X\), or a theorem semantics quantifying over executions from every start
state; the paper states none.

Second, executions were defined to be infinite, but seriality has not been
assumed at this point.  Let \(S=\{d\}\), \(R=\varnothing\), and let \(q(d)\)
be false.  There is no infinite execution beginning at \(d\), so the English
definition of \(P(d,0)\) as a universal assertion about all such executions is
vacuously true, while the proof's base claim
\(P(d,0)\Rightarrow q(d)\) is false.  Seriality appears only in the later
finite-state specialization.  Even after the first gap is repaired, this base
step therefore needs either existence of an infinite execution from every
state or a nonvacuous finite-prefix or recursive definition of \(P\); the paper
states neither.  Both defects are source-proof gaps, not hypotheses silently
added to Theorem~1 by the handbook.

The zero-rank defect is independent of several strictness tensions in the
author copy which must not be repaired silently.  The definition prints
\(j>i\), while P2 can discharge
its rank obligation at a current \(\psi(s)\)-state and A2 treats pointwise
implication as a temporal implication.  Later, \(F\) explicitly includes the
present and Theorem~3 represents \(p\leadsto q\) by \(p\Rightarrow Fq\).
The paper gives no reconciliation of these strict and inclusive clauses.  The
handbook records each at its printed standing instead of replacing them all by
the modern non-strict LTL convention.

\paragraph{Sequential and concurrent applications.}
For a sequential program, fix
\[
 L=\{\ell_0,\ldots,\ell_m\},
 \qquad s_0=\langle\ell_0,u_0\rangle,
\]
and, for every \(\ell\in L\), a location-indexed assertion
\[
 q_\ell:D\longrightarrow
   \{\mathsf{false},\mathsf{true}\}.
\]
The location index is necessary because \(N(\ell,u)\) returns an element of
\(L\), not an integer; \(q_i\) may abbreviate \(q_{\ell_i}\) only when the
enumeration is in view.  The assertions are collected into
\[
 q(\pi,u)=
 \bigwedge_{i=0}^{m}
  \bigl((\pi=\ell_i)\Rightarrow q_{\ell_i}(u)\bigr).
\]
Substitution into P1 reduces the proof to obligations, rather than deriving
those obligations from induction:
\[
 q_{\ell_0}(u_0),
 \qquad
 \forall\ell\in L\ \forall u\in D\,
 \bigl(q_\ell(u)\Rightarrow
       q_{N(\ell,u)}(T(\ell,u))\bigr).
\]
Once the initial and per-instruction obligations are proved, P1 yields
invariance of the attachment and hence the exit assertion.  This is the
Floyd-style full annotation.  ER's finite-path form is likewise a convenient
summary: formally the path is decomposed into finitely many A1 instances and
then composed by R2
\cite[p.~50]{Pnueli1977}.

For two concurrent processes, choose finite index sets \(I,J\), enumerate
\[
 L_1=\{\ell_i:i\in I\},
 \qquad L_2=\{m_j:j\in J\},
\]
 and take a predicate
\[
 h_{ij}:D\to\{\mathsf{false},\mathsf{true}\}
 \qquad(i\in I,\ j\in J)
\]
for every control-location pair.  The source's \emph{full attachment} is the
metalinguistic representation
\[
 q_{\rm full}(\pi_1,\pi_2;u)
 =\bigwedge_{i\in I}\bigwedge_{j\in J}
 \left(
   ((\pi_1=\ell_i)\land(\pi_2=m_j))\Rightarrow h_{ij}(u)
 \right).
\]
It represents one global state assertion, not a change to program semantics.
Its name is needed to distinguish its \(|I||J|\) pair-indexed components from
the two one-indexed families below, and that distinction is used immediately
in the size comparison.

The source's \emph{partial attachment} instead takes assertion families
\[
 p_i:L_2\times D\to\{\mathsf{false},\mathsf{true}\},
 \qquad
 q_j:L_1\times D\to\{\mathsf{false},\mathsf{true}\}.
\]
Its metalinguistic representation is
\[
\begin{split}
q(\pi_1,\pi_2;u)
={}&\bigwedge_{i\in I}
  \bigl((\pi_1=\ell_i)\Rightarrow p_i(\pi_2,u)\bigr)\\
 &{}\land
 \bigwedge_{j\in J}
  \bigl((\pi_2=m_j)\Rightarrow q_j(\pi_1,u)\bigr).
\end{split}
\]
It is not a different program semantics: it differs from full attachment by
organizing a global assertion into two location-indexed families while each
family may still mention the other process's location.  The distinction is
used here to compare proof-obligation size.  With loose interaction the top
level has \(|I|+|J|\) written assertions.  With strong interaction, unfolding
all opposite-location dependencies may recover \(|I||J|\) cases, the same
product size as the two-process full attachment.  That is product, or at most
quadratic, in the two displayed location-set sizes; an exponential
cross-product arises when the number of processes also varies, since a full
attachment then has \(\prod_h|L_h|\) control tuples.  This parameterization
makes precise the paper's informal claim that strong interaction can regain
exponential complexity.  Global arithmetic invariants can avoid materializing
that full cross-product attachment, but their initial truth and preservation
across each relevant processor transition must still be proved
\cite[p.~51]{Pnueli1977}.

The producer--consumer example makes the global-invariant method concrete.
To prevent collision with the next-location map \(N\), write
\(C\in\mathbb N\) with \(C>0\) for the source's capacity \(N\), and write
\(b:S\to\mathbb Z\) for its load variable \(n\).  Let
\(m_i,r_i,s_i:S\to\{0,1\}\) be the source's three location indicators for
process \(i\), and let
\(\mu,e,f:S\to\mathbb N\) give the values of the source semaphores
\({\it MUTEX},{\it IS\_EMPTY},{\it IS\_FULL}\), respectively.  These are
numerical state components and indicators, not tense formulas; the renaming is
needed only to keep their types distinct from the transition maps.  The printed
calculation says that each indicator is one on some control locations and zero
elsewhere, but it does not enumerate those support sets.  The deductions below
use only the indicator consequences that the source itself states:
simultaneous critical-region membership gives \(m_1=m_2=1\); conversely,
\(m_i=1\) means that process \(i\) is at its critical-section location and
therefore is not waiting on a semaphore operation; if process \(i\) is blocked
on \({\it MUTEX}\), then \(m_i=0\) and \(\mu=0\); and simultaneous producer
and consumer waits on \({\it IS\_EMPTY}\) and
\({\it IS\_FULL}\), respectively, give \(r_1=r_2=e=f=0\).  Together with
\(0\le s_i\le r_i\le1\), these are all the location-indicator facts used; no
unprinted support set is invented.  Pointwise on reachable states, the three
printed invariants are
\[
\begin{aligned}
 m_1+m_2+\mu&=1,\\
 r_1+r_2+e+f&=C,\\
 s_1+s_2+e&=b,
\end{aligned}
\qquad 0\le s_i\le r_i\le1.
\]
If both processes were in their critical regions, \(m_1=m_2=1\) would force
\(\mu=-1\), contradicting nonnegativity.  The third invariant gives
\(b\ge0\), and the remaining two give
\[
 b=s_1+s_2+e
 \le r_1+r_2+e=C-f\le C.
\]
For deadlock, a processor blocked on the mutex has its corresponding
\(m_i=0\) and \(\mu=0\); the first invariant gives \(m_{3-i}=1\), so the
stated indicator meaning puts the other process in its critical region and
shows that it is not blocked.  In the only remaining
simultaneous-wait case, the producer waits on \({\it IS\_EMPTY}\) and the
consumer on \({\it IS\_FULL}\), whence
\(r_1=r_2=e=f=0\).  The second invariant then forces \(C=0\), contradicting
the positive-capacity hypothesis.  Thus the source calculation proves mutual
exclusion, \(0\le b\le C\), and absence of deadlock for \(C>0\); it is not a
later automata-theoretic model-checking example
\cite[pp.~51--52]{Pnueli1977}.

\paragraph{The future-only tense presentation.}
For an execution \(\sigma=s_0s_1\cdots\), the source's
\(H(t,p)\) is the metalinguistic evaluation abbreviation \(p(s_t)\); it is not
an additional object-language operator and applies only to an atemporal state
assertion \(p\).  To type nested tense formulas, let the propositional core be
generated from state assertions by Boolean connectives and unary \(F,G\), and
define recursively
\[
\begin{aligned}
 (\sigma,n)\models p
 &\Longleftrightarrow p(s_n),\\
 (\sigma,n)\models FA
 &\Longleftrightarrow
   \exists t\ge n\ (\sigma,t)\models A,\\
 (\sigma,n)\models GA
 &\Longleftrightarrow
   \forall t\ge n\ (\sigma,t)\models A,
\end{aligned}
\]
with the ordinary recursive Boolean clauses.  Both operators include the
present.  The source also allows first-order quantifiers and prints their
interaction schemata below, but it does not supply a separate domain-semantics
discussion at this point; none is imported here.  No next or until operator
occurs in this presentation.  The source takes
\(FA:=\neg G\neg A\).  For a system formula, it identifies the absolute
present with the initial state \(s_0\); the recursive clauses also make sense
at each later position.

Its pure propositional basis consists of the following axiom schemata,
\[
\begin{aligned}
 G(A\Rightarrow B)&\Rightarrow(GA\Rightarrow GB) &&({\rm G1}),\\
 GA&\Rightarrow A &&({\rm G2}),\\
 GA&\Rightarrow GGA &&({\rm G3}),
\end{aligned}
\]
and the proof rules
\[
 \frac{A\text{ is a classical tautology}}{\vdash A}\ ({\rm RT}),
 \qquad
 \frac{\vdash A}{\vdash GA}\ ({\rm RG}),
 \qquad
 \frac{\vdash A\quad\vdash A\Rightarrow B}{\vdash B}\ ({\rm MP}).
\]
The dual definition gives the printed derived consequences
\[
 A\Rightarrow FA\quad({\rm F2}),
 \qquad
 FFA\Rightarrow FA\quad({\rm F3}).
\]
For the source's formula metavariable \(p\), with object variable \(x\) free,
the paper then prints four quantified axiom schemata at the following exact
strengths:
\[
\begin{aligned}
 G(\forall x\,p)&\Longleftrightarrow\forall x\,Gp &&({\rm Q1}),\\
 F(\exists x\,p)&\Longleftrightarrow\exists x\,Fp &&({\rm Q2}),\\
 F(\forall x\,p)&\Rightarrow\forall x\,Fp &&({\rm Q3}),\\
 \exists x\,Gp&\Rightarrow G(\exists x\,p) &&({\rm Q4}).
\end{aligned}
\]
Thus Q1--Q2 are equivalences, while Q3--Q4 are only the displayed
implications.  The paper does not say whether \(p\) here must be atemporal or
may itself contain \(F,G\), and it supplies no separate domain semantics.
Accordingly, the displayed schema range remains source-unresolved.  Restricting
later handbook uses to atemporal instances is a conservative handbook policy,
not a restriction attributed to Pnueli.

RG belongs without qualification to the pure propositional basis just listed.
The paper warns that once other means can introduce facts true only at the
present, RG must be restricted; the finite-state rule MRG below makes that
restriction formal.  In this presentation, \(K_b\) denotes exactly the pure
propositional future calculus with derived \(F:=\neg G\neg\), axiom schemata
G1--G3, and rules RT, RG, and MP displayed above.  It names a theorem system,
not a formula class or model class, and it is distinct from the later GPSS
systems \(DX\) and \(DUX\).  Theorem~2 calls this future-restricted calculus
complete in the absolute sense and decidable.  Its proof is not expanded: the
paper points to a modification of a cited completeness proof for a broader
external system, whose axioms it does not reproduce, and invokes that system's
S4 correspondence for decidability.

The completeness assertion is model-class-sensitive.  The recursive clauses
above evaluate formulas on one linearly ordered execution, whereas S4 permits
arbitrary reflexive--transitive branching frames.  For example,
\[
 (Fp\land Fq)\Rightarrow
 \bigl(F(p\land Fq)\lor F(q\land Fp)\bigr)
\]
is valid on every linear future: compare a future \(p\)-position with a future
\(q\)-position and use whichever occurs first.  Under the S4 accessibility
reading of \(F\), it fails in the Kripke model
\(M=(W,\rho,V)\) with
\[
 W=\{r,u,v\},\qquad
 \rho=\{(r,r),(u,u),(v,v),(r,u),(r,v)\},\qquad
 V(p)=\{u\},\quad V(q)=\{v\}.
\]
The relation is reflexive and transitive; \(u\) and \(v\) are incomparable and
have no successors other than themselves.  Thus the root \(r\) satisfies
neither atom but satisfies \(Fp\land Fq\).  At \(u\) no \(q\)-world is
reachable, at \(v\) no \(p\)-world is reachable, and at \(r\) neither atom is
true, so neither disjunct in the consequent holds at \(r\).  This handbook
countermodel shows that the S4 referral does not establish completeness for the
preceding linear-execution semantics.  The paper does not reconcile those two
readings, so the handbook records Theorem~2 as a source assertion with an
unresolved model class, not as an established completeness theorem for the
displayed semantics.  The paper also says only that reference~[29] contains
complexity results for the decision procedure; it states no bound or complexity
class
\cite[pp.~52--53]{Pnueli1977}.

The system-dependent axioms are kept separate from that pure basis.  The
source partitions \(R\) into finitely many action relations.  To avoid
collision with tense-formula metavariables, write this partition as
\[
 \mathcal A=\{a_1,\ldots,a_k\},
 \qquad R=\bigcup_{h=1}^{k}a_h,
 \qquad a_h\cap a_\ell=\varnothing\quad(h\ne\ell).
\]
At printed p.~53 the paper restricts the execution class by
\[
 \forall a\in\mathcal A\ \forall i\ge0\
 \exists j\ge i\ a(s_j,s_{j+1}),
\]
so every action occurs infinitely often.  This literal action condition is
the formal premise used for N.  It is recorded separately from the earlier
informal enabled-processor wording; the paper does not prove those two
formulations equivalent.

For a designated \(a\in\mathcal A\), the system-dependent displayed forms are
\[
\begin{array}{c@{\qquad}c}
\displaystyle
\frac{p(s_0)\quad
       \forall s,s'\,
       (p(s)\land R(s,s')\Rightarrow p(s'))}
      {\vdash Gp}\ ({\rm I1})
&
\displaystyle
\frac{\forall s\,(p(s)\Rightarrow q(s))\quad
       \forall s,s'\,
       (q(s)\land R(s,s')\Rightarrow q(s'))}
      {\vdash p\Rightarrow Gq}\ ({\rm I2})
\\[2.4ex]
\displaystyle
\frac{\forall s,s'\,
       (p(s)\land R(s,s')\Rightarrow q(s'))}
      {\vdash p\Rightarrow Fq}\ ({\rm E})
&
\displaystyle
\frac{\begin{gathered}
       \forall s,s'\,
       (p(s)\land R(s,s')\land\neg a(s,s')\Rightarrow p(s'))\\
       \forall s,s'\,
       (p(s)\land a(s,s')\Rightarrow q(s'))
       \end{gathered}}
      {\vdash p\Rightarrow Fq}\ ({\rm N}).
\end{array}
\]
I1 internalizes invariant induction from \(s_0\); I2 proves an invariant
starting when a condition first holds; E internalizes a one-step eventuality.
N's conclusion is interpreted only over the executions satisfying the
displayed action-occurrence restriction.  Without it N is unsound even when
\(a\) is a nonempty member of the partition.  Let
\(S=\{s,t\}\), \(s_0=s\),
\(a=\{(t,t)\}\), \(c=\{(s,s)\}\), and
\(\mathcal A=\{a,c\}\).  If \(p\) holds only at \(s\) and \(q\) nowhere,
both local premises hold, but the raw execution
\(s,s,\ldots\) never takes \(a\) and falsifies \(p\Rightarrow Fq\).
Proofs therefore have a
system-specific phase establishing these transition and fairness premises and
a pure tense-logical phase.

The paper next prints two stronger rules used in its second example.  For state
assertions \(p_1,p_2,q\), a designated \(a\in\mathcal A\), and universally
closed state variables, their typed forms are
\[
\frac{\begin{gathered}
 \forall s\,(p_1(s)\Rightarrow q(s))\\
 \forall s,s'\,
 (q(s)\land p_2(s')\land R(s,s')\Rightarrow q(s'))
 \end{gathered}}
 {\vdash(p_1\land Gp_2)\Rightarrow Gq}\ ({\rm DI}),
\]
\[
\frac{\begin{gathered}
 \forall s,s'\,
 (p_1(s)\land R(s,s')\land\neg a(s,s')\land p_2(s')
    \Rightarrow p_1(s'))\\
 \forall s,s'\,
 (p_1(s)\land a(s,s')\land p_2(s')\Rightarrow q(s'))
 \end{gathered}}
 {\vdash(p_1\land Gp_2)\Rightarrow Fq}\ ({\rm DN}).
\]
DI and DN are derived strengthenings of I2 and N, respectively, not new
primitive axioms.  Their additional \(Gp_2\) premise is needed to retain a
future side condition while propagating \(q\) or waiting for action \(a\).
Printed Example~2 uses DI to preserve such a side condition through the
mutual-exclusion accessibility argument and DN when the designated action
establishes its target.

Theorem~3 says \(K_b\) is at least as strong as ER.  Its entire local proof is
skeletal: the paper instructs the reader to express \(p\leadsto q\) as
\(p\Rightarrow Fq\), then asserts that every ER axiom can be shown to be a
theorem of \(K_b\).  Under the exact pure-calculus definition above, that
assertion is under-specified.  In particular, translated A1 needs the
system-dependent rule E, which is not a rule of pure \(K_b\); the paper does
not say whether Theorem~3 tacitly adds E or any of I1, I2, and N.  It prints
none of the individual ER derivations.  The handbook therefore records the
comparison and following relative-completeness corollary only as unresolved
source assertions, not as an embedding theorem for bare \(K_b\).  The paper's
mutual-exclusion example then uses invariant and inevitability lemmas to prove
accessibility of the critical region
\cite[pp.~53--55]{Pnueli1977}.

\paragraph{Finite-state decision arguments.}
For the final specialization, \(S\) is finite and the paper assumes for
simplicity that there are no halting states or deadlocks:
\[
 \forall s\in S\ \exists t\in S\ R(s,t).
\]
Thus no state is terminal, and every finite \(R\)-path can be extended to an
infinite execution.  The finite action family \(\mathcal A\) labels the
directed graph.  For any \(s\in S\), an infinite \(R\)-path from \(s\) is a
sequence \(\rho=t_0t_1\cdots\) with \(t_0=s\) and
\(R(t_i,t_{i+1})\) for every \(i\ge0\).  It satisfies the action condition when
\[
 \forall a\in\mathcal A\ \forall i\ge0\
 \exists j\ge i\ a(t_j,t_{j+1}).
\]
This arbitrary-start path property is needed because the semantic procedure
below tests continuations from individual states.  A \emph{proper execution}
of \(\Sigma\) is exactly such a path with \(t_0=s_0\).  The source name thus
denotes initial-state paths, not every arbitrary-start continuation, and is
used in both decision theorems and Appendix~B.

Because validity here is restricted to proper executions rather than all
\(R\)-paths, define
\[
 \Sigma\models_{\rm prop}A
 \quad\Longleftrightarrow\quad
 \forall\sigma\in S^\omega\,
 \bigl(\sigma\text{ is a proper execution of }\Sigma
       \Rightarrow(\sigma,0)\models A\bigr).
\]
This is a semantic relation between a fixed system and a tense formula, not a
new formula connective; the separate notation is necessary to state the two
finite-system theorems at their restricted path standing.

In Theorem~4 the source reuses \(A,B\) for propositional expressions.  To avoid
collision with its action relations, write these state assertions as \(p,q\).
The theorem decides
\[
 \Sigma\models_{\rm prop}G(p\Rightarrow Fq)
\]
on the fixed finite system.  The paper gives the sufficient procedure of
checking every \(p\)-state.  If the graph is first restricted to \(X\), the
same procedure is exact for initial-state validity; that reachable-state
pruning is a handbook observation, not wording attributed to the paper.

For a relevant \(p\land\neg q\)-state \(s\), delete all \(q\)-states and
decompose the remaining graph into strongly connected components.  The source
says to label a component by actions occurring on edges among its constituent
states, but prints neither a set formula nor the two-direction proof.
Formalizing that sentence, the handbook criterion is: an infinite
\(q\)-avoiding \(R\)-path from \(s\) satisfying the action condition exists
exactly when some cyclic component \(C\), reachable from \(s\), satisfies
\[
 a\cap(C\times C)\ne\varnothing
 \qquad\text{for every }a\in\mathcal A.
\]
The equivalence proof is handbook-supplied.  Choose one internal edge for each
action.  Strong connectivity supplies finite paths joining the head of each
chosen edge to the tail of the next and the last back to the first; repeating
the resulting closed walk traverses every action infinitely often.  Conversely,
an infinite path through the finite acyclic condensation graph eventually
remains in one component, and its action-satisfying tail supplies an internal
edge of every action there.  This formalizes the source's semantic decision
route
\cite[p.~55]{Pnueli1977}.

To obtain a proof system for an arbitrary finite system, the paper adds the
initial-state axiom and restricts generalization:
\[
 \frac{p(s_0)}{\vdash p}\ ({\rm P}),
 \qquad
 \frac{\vdash A}{\vdash GA}\ ({\rm MRG})
 \quad
 \begin{gathered}
  \text{only when the derivation of }\vdash A\\
  \text{used no instance of }({\rm P}).
 \end{gathered}
\]
P is needed to admit system-specific initial truth.  MRG is nonredundant:
on the path \(s_0\to s_1\to s_1\), if \(p\) holds only at \(s_0\), P yields
\(p\), while unrestricted RG would yield the false formula \(Gp\).

Theorem~5 has two clauses.  For every propositional tense formula \(A\), whether
\(\Sigma\models_{\rm prop}A\) holds is decidable; and whenever it holds, \(A\)
is derivable in the \(K_b\) calculus extended by the system-dependent axioms,
P, and MRG.  The paper asserts the adequacy clause but prints no general
construction of a derivation from semantic validity.  Appendix~B addresses
only the decidability clause.

\paragraph{Appendix A proof ownership.}
Appendix~A attributes the following source-labelled theorems T2 and T3 to
reference~[23] and prints their adjacent corollaries:
\[
\begin{aligned}
 Gp\land Fq&\Rightarrow F(p\land q) &&({\rm T2}),\\
 Gp\land Fq&\Rightarrow F(Gp\land q) &&({\rm corollary}),\\
 F(p\lor q)&\Rightarrow Fp\lor Fq &&({\rm T3}),\\
 Gp\land Gq&\Rightarrow G(p\land q) &&({\rm corollary}).
\end{aligned}
\]
Those results are referred to, not reproved, in this paper.  Pnueli separately
prints a local proof of the inference
\[
 \frac{\vdash(p\land G\neg q)\Rightarrow Fq}
      {\vdash p\Rightarrow Fq}.
\]
Writing its six steps at their proof-theoretic standing gives
\[
\begin{array}{rll}
1.&\vdash(p\land G\neg q)\Rightarrow Fq&\text{assumption},\\
2.&\vdash Fq\lor\neg Fq&\text{classical tautology},\\
3.&\vdash Fq\lor G\neg q&\text{definition of }F,\\
4.&\vdash(p\land Fq)\Rightarrow Fq&\text{classical tautology},\\
5.&\vdash p\land(Fq\lor G\neg q)\Rightarrow Fq&1,4,\\
6.&\vdash p\Rightarrow Fq&3,5.
\end{array}
\]
Thus the local lemma proof is not transferred to the externally owned T2 and
T3 \cite[p.~57]{Pnueli1977}.

\paragraph{Appendix B decision boundary.}
Appendix~B discusses two semantic decision routes for Theorem~5's decidability
clause: reintroduce explicit time variables and reduce to the monadic
second-order theory of successor, or use languages of infinite state words.
To state the second route without
supplying an automaton tuple that the paper does not print, define
\[
 L_{\rm prop}(\Sigma)
 =\{\sigma\in S^\omega:\sigma\text{ is a proper execution of }\Sigma\}.
\]
This language name is needed only to separate the already defined execution
class from the source's automaton notation \(A_\Sigma\); it stands as a
subset of \(S^\omega\) and is used in the inclusion below.  Pnueli locally
asserts that a finite-state infinite-word acceptor for exactly this language
can be constructed, but prints neither its tuple, acceptance condition,
construction, nor proof.  That assertion is not assigned to reference~[32].

For an atemporal assertion \(p\), identify each state \(s\in S\) with the
one-letter word \((s)\in S^1\subseteq S^*\), and put
\(\llbracket p\rrbracket=\{(s)\in S^1:p(s)\}\).  If
\(U\subseteq S^*\) and \(V\subseteq S^\omega\), define prefix concatenation by
\(UV=\{uv:u\in U,\ v\in V\}\).  Appendix~B prints the recursion
\[
\begin{aligned}
 L(p)&=\llbracket p\rrbracket S^\omega,\\
 L(\neg W)&=S^\omega\setminus L(W),\\
 L(W_1\land W_2)&=L(W_1)\cap L(W_2),\\
 L(W_1\lor W_2)&=L(W_1)\cup L(W_2),\\
 L(FW)&=S^*L(W),\\
 L(GW)&=L(\neg F\neg W).
\end{aligned}
\]
The source calls these languages \(\omega\)-regular.  Here that imported label
is used only as the eliminable metalinguistic abbreviation ``accepted by some
nondeterministic B\"uchi word automaton of
Definition~\ref{def:p09-omega-automata}.''  It names a class of languages over
infinite words, not a new automaton, and is needed only to state the external
closure and inclusion results.  Reference~[32] accompanies the imported
\(\omega\)-regular language machinery used for the recursion and effective
closure step.  Reference~[33] supports the decision of
\[
 L_{\rm prop}(\Sigma)\subseteq L(W).
\]
The local acceptor assertion is unproved here; the imported language/closure
machinery and inclusion decision are likewise not proved in this paper.
Pnueli prints the recursion and route, not a later LTL-to-B\"uchi translation
or a complexity classification
\cite[pp.~55--57]{Pnueli1977,McNaughton1966,Landweber1969}.

\subsubsection{The GPSS presentation of fairness and expressive completeness}

\paragraph{Published source and rights boundary.}
Shelah's author archive identifies the linked file as the published
eleven-page POPL version, pp.~163--173, and records the DOI, publisher, series,
and ISBN.  The printed title page names all four authors; the file's embedded
\texttt{Author} field names Pnueli, Shelah, and Stavi but omits Dov Gabbay, so
the printed page controls the authorship record.
Printed p.~163 carries the 1980 ACM copyright and conditional personal or
classroom copying notice, not an open-content or republication licence.  The
account below is an independent reconstruction rather than a reproduction
\cite[pp.~163--173]{GabbayPnueliShelahStavi1980}.

\paragraph{The source model and the strict-operator convention.}
The paper takes
\[
 \mathcal M=(T,<,Q_0,Q_1,\ldots),
 \qquad Q_i\subseteq T,
\]
where \((T,<)\) is a linear order and \(Q_i\) is the truth set of atom
\(q_i\).  The source's initial language \(L(F,X,U)\) is generated by
\[
 A,B::=q_i\mid\neg A\mid(A\land B)\mid(A\lor B)\mid(A\to B)
       \mid FA\mid XA\mid(A\,U\,B),
 \qquad i\in\mathbb N.
\]
Thus \(q_i\) is a formula symbol and \(Q_i\) its interpretation; \(G\) is
initially the derived constructor \(GA\defeq\neg F\neg A\), not an additional
grammar production.  No past constructor occurs in this initial grammar.
The informal clauses on p.~163 make \(F\) and \(G\) quantify only
over points strictly later than the current point.  Formal clause~(3) on
p.~164 gives the positive case for \(XA\) when an immediate successor satisfies
\(A\), but does not separately print the false case when no successor exists.

The strict-future truth condition is
\[
 \lVert FA\rVert_m=\mathsf{true}
 \quad\Longleftrightarrow\quad
 \exists m_1>m\;\lVert A\rVert_{m_1}=\mathsf{true}.
\]
Thus \(FA\equiv\topv U A\), with the source's infix operand order.
On an arbitrary order, \(XA\) is false when no immediate successor
exists. The later \(DX/DUX\) results concern \(\omega\)-models, where
every point has an immediate successor
\cite[pp.~163--164]{GabbayPnueliShelahStavi1980}.

One overload must be removed before the source formulas can be read safely.
Within this subsection only, write \(A\mathbin{\mathsf U_s}B\) for the paper's
printed infix \(A\,U\,B\), with the exact clause
\[
 \mathcal M,t\entails A\mathbin{\mathsf U_s}B
 \quad\Longleftrightarrow\quad
 \exists u>t\,
 \bigl(\mathcal M,u\entails B\ \land\
   \forall v\,(t<v<u\to\mathcal M,v\entails A)\bigr).
\]
The article later names \(S\) as the past counterpart of \(U\), but it prints
no grammar production, arity declaration, or truth clause for \(S\).  For the
convention bridge only, the handbook supplies the exact binary order dual
\[
 \mathcal M,t\entails A\mathbin{\mathsf S_s}B
 \quad\Longleftrightarrow\quad
 \exists u<t\,
 \bigl(\mathcal M,u\entails B\ \land\
   \forall v\,(u<v<t\to\mathcal M,v\entails A)\bigr).
\]
Thus \(\mathsf U_s\) transcribes the source clause, whereas \(\mathsf S_s\) is
the handbook-supplied past constructor just defined, not a source
transcription.  The subscript is an eliminable local warning, not a new
temporal logic: \(A\mathbin{\mathsf U_s}B\) is Kamp's prefix
\(U\,B\,A\) from Definition~\ref{def:p08-tense-until-since}, with the operand
order reversed, and similarly in the past.  The warning is necessary because
Definition~\ref{def:p09-ltl} instead uses inclusive strong until, and it is used
throughout the formulas below.  The source identities
\[
 FA\equiv \topv\mathbin{\mathsf U_s}A,
 \qquad
 XA\equiv \bot\mathbin{\mathsf U_s}A,
 \qquad
 GA\equiv\neg F\neg A
\]
hold over every linear order under that coherent future reading, not merely
over \(\mathbb N\).  On p.~164, \(L(U)\) first names the syntactic sublanguage
obtained by deleting \(X,F\); the later equation \(L(U)=L(X,F,U)\) identifies
the two expressive powers.  It is not a literal equality of the two grammars.
An \(\omega\)-model is the special case \((T,<)=(\mathbb N,<)\),
and \(\omega\)-validity quantifies over every such valuation and every
position \cite[pp.~163--164]{GabbayPnueliShelahStavi1980}.

On \(\omega\)-models, the exact bridge to the handbook's inclusive operators
is
\[
\begin{aligned}
 A\mathbin{\mathsf U_s}B
   &\equiv \mathsf X(A\mathbin{\mathsf U}B),
 &A\mathbin{\mathsf U}B
   &\equiv B\lor\bigl(A\land
      (A\mathbin{\mathsf U_s}B)\bigr),\\
 A\mathbin{\mathsf S_s}B
   &\equiv \mathsf Y(A\mathbin{\mathsf S}B),
 &A\mathbin{\mathsf S}B
   &\equiv B\lor\bigl(A\land
      (A\mathbin{\mathsf S_s}B)\bigr).
\end{aligned}
\]
Consequently strict \(U_s,S_s\) are interdefinable with the handbook's
inclusive \(U,S\) only when the discrete next/previous operators remain
available.  In particular, the source equation
\(XA\equiv\bot\mathbin{\mathsf U_s}A\) must not be rewritten as
\(XA\equiv\bot\mathbin{\mathsf U}A\): the latter is just false unless \(A\)
already holds at the current point.

\paragraph{Program executions, correctness, and three strengths of response.}
Programs are treated as generators of execution sequences; a nondeterministic
program may generate several sequences from one input, and the current program
location is included in each state.  The state proposition
\(\mathit{at}_{\ell}\) says that the instruction labelled \(\ell\) is about to
execute.  Let \(\bar y\) be the tuple of current program-variable values and
\(\bar\xi\) a same-arity tuple of logical variables retaining the input values;
\(\varphi(\bar\xi)\) is the precondition and
\(\psi(\bar\xi,\bar y)\) the postcondition.  These are atemporal program-state
predicates, value tuples, and formula metavariables inside the temporal
formula, not new temporal connectives.  If \(\ell_0\) and \(\ell_e\) are entry
and exit locations, the source writes partial correctness as
\[
 \bigl(\mathit{at}_{\ell_0}\land \bar y=\bar\xi
       \land\varphi(\bar\xi)\bigr)
 \to
 G\bigl(\mathit{at}_{\ell_e}\to
        \psi(\bar\xi,\bar y)\bigr).
\]
This constrains an exit if one occurs; it does not imply termination.  Total
correctness replaces the conclusion by
\[
 F\bigl(\mathit{at}_{\ell_e}\land
        \psi(\bar\xi,\bar y)\bigr).
\]
These source displays use the coherent future reading just recorded
\cite[pp.~164--165]{GabbayPnueliShelahStavi1980}.

The paper then separates invariance formulas \(p\to Gq\) from eventuality
formulas \(p\to Fq\), listing safety, partial correctness, mutual exclusion,
clean performance, and deadlock freedom under the former heading, and
liveness, total correctness, accessibility, and livelock freedom under the
latter.  The lists are source classifications, not definitions of the later
semantic safety/liveness classes.  It next orders three response conditions by
how long the trigger must persist:
\[
\begin{array}{lll}
\text{insistence}   & Gp\to Fq,
                    & G(Gp\to Fq)\text{ at every strictly later reference},\\
\text{persistence}  & GFp\to Fq,
                    & G(GFp\to Fq)\text{ at every strictly later reference},\\
\text{impulse}      & p\to Fq,
                    & G(p\to Fq)\text{ at every strictly later reference}.
\end{array}
\]
At the current reference point, the first antecedent requires \(p\) at every
strictly later point; it is not the weaker eventual-permanence condition
\(FGp\).  The second asks when \(p\) recurs arbitrarily far into the future,
and the third after a single occurrence.  A prefixed \(G\) repeats the local
obligation at every strictly later reference point.  The source treats the
impulse condition as natural for internal events but potentially too strong
for an external signal that may occur only once and escape detection.  With
the source's strict convention, the first two local
forms may also be written
\[
 \neg G(p\land\neg q),
 \qquad
 \neg G(Fp\land\neg q),
\]
respectively.

For a busy-wait instruction at location \(m\), ``if \(x=0\), go to \(m\),''
the insistence condition is
\[
 \neg G(\mathit{at}_m\land x>0).
\]
It rules out remaining at \(m\) while the enabling condition stays true.  The
paper calls this the weakest acceptable semaphore behaviour and the
corresponding scheduling assumption the minimal fairness requirement: every
process is eventually scheduled.  It explicitly warns that, for semaphores,
this still permits infinite overtaking.  For a semaphore instruction
\(m:P(x)\), the stronger persistence condition is
\[
 \neg G\bigl(F(x>0)\land\mathit{at}_m\bigr),
\]
which rules out indefinite waiting when positive semaphore values recur
arbitrarily late.  The paper stresses that this does not say the process itself
is scheduled infinitely often: its scheduling instants could always coincide
with \(x=0\).  For a guarded region
\(m:\text{with }r\text{ when }B\text{ do }S\), the paper gives
\[
 \neg G\bigl(F(r>0\land B)\land\mathit{at}_m\bigr)
\]
and contrasts it with the weaker
\[
 \neg G\bigl(F(r>0)\land\mathit{at}_m\land B\bigr).
\]
The former reacts when free-resource-and-guard states recur together; the
latter applies when the guard has become permanently true while free-resource
states recur.  These are semantic obligations on executions, not scheduler
algorithms supplied by the logic.  The paper reports that the former guarded
condition can be implemented with semaphores plus a queue or an unbounded
counter, but not with semaphores alone
\cite[pp.~165--167]{GabbayPnueliShelahStavi1980}.

\paragraph{Why strict until is used.}
Response alone says that a requested event eventually occurs; it does not say
that a response was solicited or preserve the order of two requests.  Ignoring
the current-point boundary exactly as the paper does, absence of an unsolicited
response is
\[
 Fq\to(\neg q\mathbin{\mathsf U_s}p).
\]
For repeated use, the source introduces the eliminable formula abbreviation
\[
 \operatorname{Pr}(A,B)
 \defeq
 FB\to(\neg B\mathbin{\mathsf U_s}A).
\]
Here \(\operatorname{Pr}\) is a metalinguistic abbreviation from pairs of
formulas to formulas, not a new connective or relation on executions.  It is
used in the three following displays: the general FIFO condition, its
implementation-specific semaphore instance, and the repeated-response
refinement.  In each use it retains the first-occurrence ordering condition;
unlike \(A\to FB\), it constrains a future \(B\) to have an \(A\) before it or
at the same witness point.

For requests \(p_1,p_2\) and corresponding responses \(q_1,q_2\), the printed
FIFO condition is
\[
 \operatorname{Pr}(p_1,p_2)\to\operatorname{Pr}(q_1,q_2).
\]
For a particular FIFO semaphore implementation, the paper proposes the
following as a possible basic implementation axiom:
\[
 \operatorname{Pr}(\mathit{at}_m,\mathit{at}_n)
 \to
 \operatorname{Pr}(\mathit{at}_{m'},\mathit{at}_{n'}).
\]
Its use is assumption-relative; temporal logic alone does not prove that the
implementation is FIFO.
The single precedence formula still permits one \(p\) to license many later
\(q\)'s.  Requiring a fresh intervening \(p\) between successive \(q\)'s gives
the stronger source formula
\[
 \operatorname{Pr}(p,q)\land
 G\bigl(q\to\operatorname{Pr}(p,q)\bigr).
\]
Thus the demonstrated contribution of until is an ordered interval condition,
not merely another spelling of eventuality
\cite[pp.~167--168]{GabbayPnueliShelahStavi1980}.

The paper then poses four questions: whether \(U\) is eliminable in favour of
\(F,X\); whether the resulting \(L(U)\) is expressively complete; and which
deductive systems serve \(L(F,X)\) and \(L(F,X,U)\).  It attributes a negative
answer to the first question to Kamp, answers the second with the qualified
results below, and answers the last two with \(DX\) and \(DUX\).  These are
expressiveness and theoremhood questions; no computational-complexity question
is asked there \cite[pp.~168--169]{GabbayPnueliShelahStavi1980}.

The Kamp attribution is compressed.  Literal Chapter~IV Theorem~3 of Kamp's
thesis proves that the integer-order tenses \(\mathsf{ISOLP}\) and
\(\mathsf{ISOLF}\) are not expressible from his specified unary basis; GPSS
does not print the reduction from that result to nondefinability of \(U\) from
\(F,X\).  The handbook therefore reports the attribution without treating the
cited theorem itself as the missing reduction
\cite[ch.~IV, pp.~100--108]{Kamp1968}.

\paragraph{First-order comparison and four distinct result domains.}
The comparison language \(L_1\) is first-order logic with equality, \(<\), and
the unary predicate symbols \(q_i\), interpreted by the sets \(Q_i\).  A model
is \emph{complete} in the source's order
sense when every nonempty bounded subset has both a least upper bound and a
greatest lower bound.  It is \emph{discrete} when every nonmaximal point has an
immediate successor and every nonminimal point has an immediate predecessor.
The natural numbers and integers are discrete and complete; the reals are
complete but not discrete; the rationals are not complete.

A future \(L_1\)-formula \(\varphi(x)\) is one whose quantifiers are restricted
to \([x,\infty)\), printed as \((\forall y\ge x)\) and
\((\exists y\ge x)\).  This is a syntactic restriction relative to the
distinguished free variable, not a claim that arbitrary first-order formulas
forget the past.  The paper merely calls the temporal-to-first-order
translation ``natural'' and does not print its recursion.  The following
load-bearing strict clauses are the handbook's spelling-out.  Write
\(A\mapsto A^*(x)\) for this eliminable metalinguistic translation, not for an
object-language connective:
\[
\begin{aligned}
 (q_i)^*(x)
   &\defeq q_i(x),\\
 (FA)^*(x)
   &\defeq \exists y\,(x<y\land A^*(y)),\\
 (XA)^*(x)
   &\defeq \exists y\,
      \bigl(x<y\land A^*(y)\land
        \neg\exists z\,(x<z\land z<y)\bigr),\\
 (A\mathbin{\mathsf U_s}B)^*(x)
   &\defeq \exists y\,
       \bigl(x<y\land B^*(y)\land
         \forall z\,(x<z\land z<y\to A^*(z))\bigr),\\
 (A\mathbin{\mathsf S_s}B)^*(x)
   &\defeq \exists y\,
       \bigl(y<x\land B^*(y)\land
         \forall z\,(y<z\land z<x\to A^*(z))\bigr).
\end{aligned}
\]
The Boolean clauses are homomorphic.
This direction works on every linear order.  The reverse direction is divided
as follows.

\begin{enumerate}
\item Theorem~2.1: for every \(L_1\)-formula \(\varphi(x)\), some
  \(\theta\in L(U)\) agrees with \(\varphi\) at point \(0\) of every
  \(\omega\)-model.  Point \(0\) is essential because a future-only formula at
  an arbitrary point cannot inspect its past. The full construction is
  \cref{thm:p09-gpss-discrete-future-completeness}.
\item Theorem~2.2: for every future \(L_1\)-formula \(\varphi(x)\), some
  \(\theta\in L(U)\) agrees with it at every point of every discrete complete
  model, as proved in \cref{thm:p09-gpss-discrete-future-completeness}.
\item Theorem~2.3: over discrete complete models, every \(L(S,U)\)-formula is
  equivalent to a Boolean combination of formulas from the past-only
  \(L(S)\) and future-only \(L(U)\) fragments, with the complete
  effective construction in \cref{thm:p09-gpss-past-future-separation}.
\item Theorem~2.5: for arbitrary linearly ordered models, every
  \(L_1\)-formula \(\varphi(x)\) is said to have an equivalent formula in
  \(L(S,U,S',U')\).
\end{enumerate}

The last statement says only that the additional printed symbols \(U'\)
and \(S'\) ``handle gaps.''  This paper supplies no grammar productions,
arity declarations, or satisfaction clauses for them.  They therefore have
standing here only as source-reported formula symbols in Theorem~2.5, not as
defined handbook operators from which further claims may be derived.  ``At
zero,'' ``future,'' ``discrete complete,'' and
``arbitrary with additional symbols'' are four different hypothesis packages
\cite[pp.~168--171]{GabbayPnueliShelahStavi1980}.

The later defined four-operator basis and its full arbitrary-order
completeness construction are supplied in
\cref{sec:p09-stavi-gap-temporal-logic,thm:p09-stavi-completeness},
following Rabinovich's revised 2018 presentation
\cite{Rabinovich2018Stavi}.

\paragraph{What the proof outline actually contains.}
GPSS labels the passage ``Proof of 2.2 (outline)'' and says that the proof
presented there follows its reference [SS], identified in the article's
bibliography as the forthcoming Shelah--Stavi manuscript
\emph{Expressive Completeness for Temporal Logic}.  This records only the
article's attribution: the [SS] body is not reproduced there, and the handbook
does not import or claim independent verification of that manuscript's full
proof.
For a variable \(z\), the source sets
\[
 \operatorname{At}(z)
 =\{q_i(z),\neg q_i(z):i\in\omega\}.
\]
For distinct \(x,y\), it lets \(\operatorname{Bet}_0(x,y)=\varnothing\).
The set \(\operatorname{Bet}_{n+1}(x,y)\) contains each formula
\[
 \chi=\exists z\,
 \bigl(x<z\land z<y\land
       \varphi_1\land\cdots\land\varphi_m\bigr)
\]
and its negation, where \(m\in\omega\) and every \(\varphi_i\) belongs to
\[
 \operatorname{Bet}_n(x,z)\cup
 \operatorname{At}(z)\cup
 \operatorname{Bet}_n(z,y).
\]
Then
\[
 \operatorname{Bet}(x,y)=\bigcup_{n\in\omega}\operatorname{Bet}_n(x,y),
 \quad
 \operatorname{Aft}(x)=\operatorname{Bet}(x,\infty),
 \quad
 \operatorname{Bef}(x)=\operatorname{Bet}(-\infty,x),
\]
where \(\pm\infty\) are metalinguistic interval endpoints, not members of the
model.

A source ``special'' formula on \(x_1,\ldots,x_k\) is exactly a conjunction
\[
 \pi\land\psi_1\land\cdots\land\psi_m,
\]
where \(\pi\) fixes the relative order of the free variables and every
\(\psi_i\) belongs to the corresponding sequence of before, atomic, between,
and after sets
\[
 \operatorname{Bef}(x_1)\cup\operatorname{At}(x_1)\cup
 \operatorname{Bet}(x_1,x_2)\cup\cdots\cup
 \operatorname{At}(x_k)\cup\operatorname{Aft}(x_k).
\]
Lemma~1 says that every \(L_1\)-formula is equivalent to a finite disjunction
of such formulas.  This name is nonredundant only as the stated normal-form
class used in the following induction; it does not name a new formula
language. The full induction is supplied in
\cref{thm:p09-gpss-special-normal-form}. Its explicit bridge to Kamp's
Chapter~III normal form is
\cref{thm:p09-gpss-kamp-normal-form-bridge}; the two syntactic
presentations retain their distinct constructors and empty cases
\cite[ch.~III]{Kamp1968}.

A source decomposition formula on \((x,y)\) has the exact displayed shape
\[
\begin{split}
 x<y\ \land\
 \exists z_0,\ldots,z_n\,[x=z_0<z_1<\cdots<z_n=y]
 \ \land\
 \bigwedge_{i=1}^{n-1}\theta_i(z_i)\\
 {}\land
 \bigwedge_{i=1}^{n}
 \forall u\,
 \bigl(z_{i-1}<u\land u<z_i\to\varphi_i(u)\bigr),
\end{split}
\]
where every \(\theta_i\) and \(\varphi_i\) is the \(L_1\)-translation of an
\(L(U)\)-formula.  When \(y=\infty\), the paper also permits finitely many
cofinal-occurrence conjuncts
\[
 \forall u\,\exists v>u\;\psi_j(v),
\]
again with \(\psi_j\) translated from \(L(U)\).
On discrete complete orders, with actual lower endpoint \(x\) and
either actual \(y>x\) or formal \(y=+\infty\), every regional
formula is equivalent to a finite disjunction of these decompositions.
The full induction is part of
\cref{thm:p09-gpss-discrete-future-completeness}.

The negation step, corresponding to Proposition~2.4, uses relative
complementation \(x<y\land\neg\delta(x,y)\) for finite endpoints.
\Cref{thm:p09-gpss-determining-events} constructs the finite list
of partial event expressions, including missing first and last occurrences.
\Cref{thm:p09-gpss-bounded-complement} translates every relevant
defined, undefined and weak-order case back into the exact decomposition
grammar. The future-ray construction is
\cref{thm:p09-gpss-suffix-complement}. All point and gap labels
remain translations of future-only \(L(U)\) formulas.

Combining this construction with the full special normal form and
suffix translation proves Theorem~2.2 at every point of a discrete
complete order. Relativization at the initial point of an
\(\omega\)-model gives Theorem~2.1 with its different evaluation scope.
Formal-left-endpoint complementation fails in this unchanged label
class, by \cref{thm:p09-gpss-prefix-complement-obstruction}.
The announced complete-order modification concerns Kamp's \(S,U\)
basis; it does not assert a \(U\)-only completeness theorem on every
complete order. The local boundary and positive-chain constructions
are \cref{thm:p09-gpss-infimum-case-compiler,thm:p09-gpss-complete-positive-chain}, with the bounded obstruction
in \cref{thm:p09-gpss-complete-bounded-obstruction}.
The complete discrete separation and its guarded complete-order
comparison are \cref{thm:p09-gpss-past-future-separation,thm:p09-gpss-complete-order-separation}.
\Cref{thm:p09-gpss-pure-separation-obstruction} proves why the
pure-fragment version does not extend to every complete order.
The gap operators of Theorem~2.5 retain their own construction
\cite[pp.~169--171]{GabbayPnueliShelahStavi1980}.

% GPSS 1980, Lemma 1 and positive decomposition preparation, pp. 169--171.

\subsection{The GPSS special normal form and positive decomposition construction}
\label{sec:p09-gpss-special-normal-form}

This presentation supplies the first-order normal form and the positive
construction steps in the GPSS expressive-completeness outline. The
normal form works on every nonempty strict linear order. The subsequent
negation proposition is a separate obligation in the source's proof for
discrete complete orders
\cite[pp.~169--171]{GabbayPnueliShelahStavi1980}.

\begin{hbdefinition}{def:p09-gpss-regional-stages}{The literal GPSS atoms and signed regional stages}
The first-order language \(L_1\) has equality, strict order \(<\),
unary predicates \(q_i\), and the classical connectives and quantifiers.
Its models have nonempty domain \(T\), with arbitrary unary
interpretations. The subscript \(G\) below distinguishes the paper's
classes from Kamp's different classes. Put
\[
              \operatorname{At}_G(z)=\{q_i(z),\neg q_i(z):i\geq0\}.
\]
These are individual signed literals, not arbitrary Boolean formulas
or arbitrary one-variable first-order conditions.

An interval endpoint may be a variable, the formal lower endpoint
\(-\infty\), or the formal upper endpoint \(+\infty\), in the
appropriate position. The formal endpoints are never quantified or
interpreted as elements of \(T\). Let \(g_{ab}(z)\) be
\(a<z\land z<b\), omitting a comparison when its endpoint is
formal. Thus \(g_{-\infty,+\infty}(z)\) is true. Define
\[
 \operatorname{Bet}_{G,0}(a,b)=\varnothing,
\]
and let \(\operatorname{Bet}_{G,n+1}(a,b)\) consist of
\[
 e=\exists z\left(g_{ab}(z)\land\bigwedge_{j=1}^{m}\gamma_j\right)
                  \quad\text{and}\quad\neg e,
\]
where \(m\geq0\), the bound variable is fresh, and every
\(\gamma_j\) belongs to
\[
 \operatorname{Bet}_{G,n}(a,z)\cup\operatorname{At}_G(z)
                                  \cup\operatorname{Bet}_{G,n}(z,b).
\]
An empty conjunction is true. Each index is a nonnegative integer.
For actual endpoint variables use distinct names; their assigned values
need not be distinct unless an order prefix requires it. Write
\[
\begin{aligned}
 \operatorname{Bet}_G(a,b)&=\bigcup_{n\geq0}\operatorname{Bet}_{G,n}(a,b),\\
 \operatorname{Bef}_{G,n}(x)&=\operatorname{Bet}_{G,n}(-\infty,x),\qquad
 \operatorname{Aft}_{G,n}(x)=\operatorname{Bet}_{G,n}(x,+\infty).
\end{aligned}
\]
Unindexed ray classes are their stage unions. Using both formal endpoints
also gives whole-order components, needed below for sentences. This makes
the zero-free-variable extension of the source convention explicit
\cite[p.~170]{GabbayPnueliShelahStavi1980}.
\end{hbdefinition}

\begin{hbdefinition}{def:p09-gpss-special-form}{Special clauses with complete weak order types}
For a finite nonempty list \(X\) of distinct parameter variables,
a complete weak order type is an ordered partition
\((B_1,\ldots,B_s)\) of \(X\). Choose a representative
\(y_j\in B_j\). Its prefix \(\pi\) conjoins
\(y_1<\cdots<y_s\) and the equalities of every variable
in \(B_j\) with \(y_j\). For one representative use its
self-equality when the prefix otherwise has no conjunct.

A GPSS special clause on this type is
\(\pi\land\bigwedge_j\gamma_j\), where its local conjuncts
come from
\[
 \operatorname{Bef}_G(y_1)\cup
 \bigcup_{j=1}^{s}\operatorname{At}_G(y_j)\cup
 \bigcup_{j=1}^{s-1}\operatorname{Bet}_G(y_j,y_{j+1})\cup
 \operatorname{Aft}_G(y_s).
\]
The list of local conjuncts may be empty. A special normal form is a
finite disjunction of such clauses; the empty disjunction is false.
Coincident parameters form one block and do not create a spurious
strict interval. For an empty parameter list, a special sentence
clause is a finite conjunction of whole-order components, with no
prefix; finite disjunctions and their empty cases have the same meanings.
All substitutions concern free occurrences and avoid variable capture
\cite[Lemma~1, p.~170]{GabbayPnueliShelahStavi1980}.
\end{hbdefinition}

\begin{hblemma}{thm:p09-gpss-kamp-normal-form-bridge}{proved=annex:p9.gpss.kamp.normal.form.bridge}{The exact bridge to Kamp's earlier regional grammar}
The GPSS regional stages are increasing, and signed atoms and regional
components have a complement operation given by reversing their sign.
Finite Boolean combinations of them can therefore be put effectively
into finite disjunctive normal form using the same permitted components.

After renaming \(q_i\) as \(P_i\), replacing equality by
\(\neg(x<y)\land\neg(y<x)\), and using the starred
Kamp classes of \cref{def:p09-kamp-iii1-components}, each GPSS
before, after or between component with at least one actual endpoint
variable translates to a corresponding
Kamp component. Conversely each starred Kamp component translates
to a finite disjunction of conjunctions of GPSS components in the
same region. A Kamp pointwise literal conjunction becomes a list
of GPSS point literals; its distinguished true and false slots
become respectively an empty conjunction and empty disjunction.

Both translations preserve meaning on every nonempty strict linear
order and introduce no new unary predicate. They provide an explicit
bridge to \cref{thm:p09-kamp-iii1}; they do not identify the two
syntactic classes. In particular Kamp's later arbitrary elementary
one-point formulas are not GPSS \(\operatorname{At}_G\) literals.
Whole-order components use the separate sentence convention.
\end{hblemma}

\begin{hbtheorem}{thm:p09-gpss-special-normal-form}{proved=annex:p9.gpss.special.normal.form}{The complete GPSS special-form normal-form lemma}
Every \(L_1\)-formula is effectively equivalent, on every
nonempty strict linear order and every unary interpretation, to a
finite disjunction of GPSS special clauses on its declared parameter
list. The construction retains all weak order types, including
equal-point blocks, and handles every existential and universal
quantifier. Sentences have the explicit whole-order version above.
No unary predicate absent from the input is needed.

The normal form still contains nested first-order regional
quantification. This result is GPSS's Lemma~1, with its induction,
equality cases, sentence convention and empty-list cases supplied;
it is not yet the later conversion to future temporal syntax
\cite[Lemma~1, p.~170]{GabbayPnueliShelahStavi1980}.
\end{hbtheorem}

\begin{hbcorollary}{thm:p09-gpss-future-special-form}{proved=annex:p9.gpss.future.special.form}{The exact present-and-future specialization}
Suppose \(\varphi(x)\) is a source future formula: every
quantifier is restricted to \(y\geq x\), relative to the
same distinguished free variable. Then it is equivalent on every
linear-order model to a finite disjunction of conjunctions from
\(\operatorname{At}_G(x)\cup\operatorname{Aft}_G(x)\).
The inclusive lower bound retains the present atomic conditions.
No before component, past-dependent one-point label, or completeness
hypothesis is needed for this first-order specialization
\cite[pp.~169--170]{GabbayPnueliShelahStavi1980}.
\end{hbcorollary}

\subsubsection{The exact future-labelled decomposition class}

\begin{hbdefinition}{def:p09-gpss-future-labels}{Permitted labels and their strict future semantics}
A permitted label is the first-order translation \(B^*(u)\)
of a finite \(L(U)\)-formula \(B\), with no free individual
parameter except its current point \(u\). Its strict until clause is
\[
 (B\mathbin{\mathsf U_s}C)^*(u)=
 \exists v>u\left(C^*(v)\land
                  \forall w(u<w<v\to B^*(w))\right).
\]
Atomic and Boolean clauses are homomorphic. Put
\(FB=\top\mathbin{\mathsf U_s}B\) and \(GB=\neg F\neg B\).
True and false are ordinary Boolean abbreviations, for example
\(q_0\lor\neg q_0\) and \(q_0\land\neg q_0\), with
no semantic dependence on that inessential letter.
Before copying a label to a chain point, rename its bound variables
apart from every endpoint and chain variable, and replace only its
free current-point variable.
All labels are evaluated in the original whole model. A label at an
interior point may inspect points beyond the displayed interval's
right endpoint; it is not implicitly relativized to that interval
\cite[pp.~164, 169--170]{GabbayPnueliShelahStavi1980}.
\end{hbdefinition}

\begin{hblemma}{thm:p09-gpss-future-label-closure}{proved=annex:p9.gpss.future.label.closure}{Boolean closure and the exact future-only restriction}
Permitted labels are closed under all finite Boolean operations.
Their truth at a point depends only on the ordered labelled tail
at and after that point. Thus a general past-dependent formula
cannot be inserted as a decomposition label merely because it has
one free variable. In particular \(\exists v<u\,q_0(v)\)
has no equivalent future label over all linear-order models.
\end{hblemma}

\begin{hbdefinition}{def:p09-gpss-decomposition}{Decompositions with finite or formal endpoints}
For admissibly ordered endpoints \(a<b\), choose \(r\geq0\)
interior points \(z_1<\cdots<z_r\) of \(T\), all between the
endpoints. There are \(r\) point labels \(\theta_i\), and
\(r+1\) gap labels \(\phi_i\), all permitted future labels.
Put \(z_0=a,z_{r+1}=b\) only as endpoint abbreviations. The
decomposition formula is
\[
\begin{aligned}
 \delta(a,b)={}&a<b\ \land\
 \exists z_1\cdots z_r\bigl[
       a<z_1<\cdots<z_r<b\\
 &\qquad\land\bigwedge_{i=1}^{r}\theta_i(z_i)
 \land\bigwedge_{i=0}^{r}
       \forall u(z_i<u<z_{i+1}\to\phi_i(u))\bigr].
\end{aligned}
\]
Formal endpoint comparisons and quantification are omitted as in the
regional grammar. For \(r=0\), the internal chain and point-label
conjunction are empty and there is just the one gap universal.
In the paper's finite-endpoint notation this is \(n=r+1\),
with \(z_0=x,z_n=y\). Endpoints have no point labels. All
open gaps may be empty even when their endpoints are distinct.
The finite-endpoint spelling with \(n=0\) would require
\(x=y\) as well as \(x<y\), and contributes only false.

If the right endpoint is \(+\infty\), also permit a finite
list of global cofinal-occurrence conditions
\[
          \operatorname{Inf}(\psi)
             :=\forall u\in T\ \exists v>u\ \psi(v),
\]
where \(\psi\) is a permitted future label. The quantifier
over \(u\) is over the whole domain. Empty lists impose no
condition. Finite disjunctions of decompositions, including the
empty disjunction, will be used. Statements about complements on
a finite interval are relative to its guard \(a<b\)
\cite[pp.~170--171]{GabbayPnueliShelahStavi1980}.
\end{hbdefinition}

\begin{hbtheorem}{thm:p09-gpss-positive-decomposition}{proved=annex:p9.gpss.positive.decomposition}{Finite positive closure and existential assembly}
Finite disjunctions of the stated decompositions are closed under
finite conjunction on a fixed admissible interval. The conjunction
construction enumerates all order-preserving merges of the two
interior witness chains, including coincident points from different
chains. It conjoins point labels and transfers each covering
gap label onto newly inserted points and subintervals. Every resulting
label remains future-only. Cofinal-condition lists are combined by
union when the upper endpoint is formal.

If \(L(a,z)\) and \(R(z,b)\) have such finite disjunctions,
and \(\theta(z)\) is a permitted future label, then so does
\[
 \exists z\bigl(g_{ab}(z)\land L(a,z)\land\theta(z)\land R(z,b)\bigr).
\]
The construction concatenates the two witness chains and puts
\(\theta\) at their common interior point. It includes both rays
and the whole-order endpoint convention.

The empty conjunction is represented, on the admissible interval,
by the decomposition with no interior point and gap label true;
the empty disjunction is false. On admissible endpoints, every member of
\(\operatorname{Bet}_{G,1}\), of either sign, already has a
decomposition presentation. More generally the positive regional
constructor is handled whenever the presentations of all its
regional children are supplied. No arbitrary negation closure is
assumed in these claims
\cite[pp.~170--171]{GabbayPnueliShelahStavi1980}.
\end{hbtheorem}

\begin{hblemma}{thm:p09-gpss-suffix-translation}{proved=annex:p9.gpss.suffix.translation}{The complete translation of a future decomposition}
For a decomposition \(\delta(x,+\infty)\) with actual lower
endpoint \(x\), let \(\Theta_i,\Phi_i,\Psi_j\) be
the \(L(U)\)-formulas supplying its labels. Define
\[
 H_r=G\Phi_r,\qquad
 H_i=\Phi_i\mathbin{\mathsf U_s}(\Theta_{i+1}\land H_{i+1})
                    \quad(0\leq i<r).
\]
Then the exact future expression is
\[
                 H_0\land\bigwedge_j(F\Psi_j\land GF\Psi_j).
\]
It is equivalent to the decomposition at every point of every
nonempty linear-order model, including a last point. For \(r=0\)
the first factor is simply \(G\Phi_0\). This proves the
source's positive return from future decompositions to \(L(U)\)
\cite[p.~171]{GabbayPnueliShelahStavi1980}.
\end{hblemma}

\begin{hbexample}{def:p09-gpss-positive-merge-example}{Equal witnesses and the labels at inserted points}
Let \(p,q,r,s\) be proposition letters. A decomposition with
one interior point labelled \(p\) and both gaps labelled true
expresses \(D_p(x,y)=\exists z(x<z<y\land p(z))\) on
\(x<y\). The conjunction \(D_p\land D_q\) has three
decomposition cases: one point labelled \(p\land q\); two
points labelled \(p,q\) in that order; or two points labelled
\(q,p\). All their gap labels are true. On the three-point
order \(\{0,1,2\}\), with endpoints zero and two and both
letters true at one, only the coincident-witness case can hold.

If instead the first decomposition has point label \(p\)
and gap labels \(r,s\), and the other has no interior point
and gap label \(q\), their intersection has point label
\(p\land q\) and gap labels \(r\land q,s\land q\).
The \(q\)-condition at the merged point is necessary: in
the same three-point order the two smaller open gaps are empty,
but the original single gap still contains the middle point.
\end{hbexample}

\begin{proofboundary}{def:p09-gpss-negation-boundary}{The exact endpoint scope of complementation}
On discrete complete orders, an actual lower endpoint permits the
complete relative-complement constructions of
\cref{thm:p09-gpss-bounded-complement,thm:p09-gpss-suffix-complement}.
For finite endpoints the target is
\(x<y\land\neg\delta(x,y)\); for a future ray it is
\(\neg\delta(x,+\infty)\).
The next presentation supplies the event grammar, determining recursion
and complete translation of its defined, undefined and weak-order cases.
These results close the regional induction and yield
\cref{thm:p09-gpss-discrete-future-completeness}.

The unchanged future-label class is not closed under complementation
at a formal left endpoint. The prefix and whole-order counterexamples
are proved in \cref{thm:p09-gpss-prefix-complement-obstruction}.
Arbitrary complete orders require their separate nonattainment and
cut constructions
\cite[p.~171]{GabbayPnueliShelahStavi1980}.
\end{proofboundary}

% Constructive discrete-complete GPSS complementation and exact endpoint scope.

\subsection{Constructive complementation of GPSS decompositions}
\label{sec:p09-gpss-decomposition-complement}

Use the exact future-labelled decomposition class of
\cref{def:p09-gpss-decomposition}. Point and gap labels are supplied
with their finite \(L(U)\) formulas; they are evaluated in the
whole model and have no individual parameter other than their
current point. This section proves relative complementation for
bounded intervals and for future rays with an actual lower endpoint.
It also settles why the same future-labelled class cannot support
the broader prefix and whole-order complement claim
\cite[Proposition~2.4 and its outline, p.~171]{GabbayPnueliShelahStavi1980}.

\begin{hblemma}{thm:p09-gpss-finite-bounded-interval}{proved=annex:p9.gpss.finite.bounded.interval}{Bounded intervals in discrete complete orders are finite}
In a nonempty discrete complete strict linear order, every closed
interval \([a,b]\), with actual endpoints \(a\leq b\),
is finite. Consequently every nonempty subset of an open bounded
interval has a first and last member. An adjacent pair has an empty
open interval. No finiteness of the entire order is asserted.
\end{hblemma}

\begin{hbdefinition}{def:p09-gpss-bounded-events}{A finite partial event language on a bounded interval}
Fix a discrete complete model and actual endpoints \(x<y\).
The root terms are \(x,y\).
For a future temporal label \(A\), form
\[
 \operatorname{First}_A(s,t),\qquad
 \operatorname{Last}_A(s,t),\qquad
 \operatorname{Test}_A(s).
\]
First and last return the first or last point satisfying \(A\)
strictly between the values of their two bounds. They are undefined
if a bound is undefined, the bounds are not strictly ordered, or the
occurrence set is empty. Tests are applied only to nonroot terms:
\(\operatorname{Test}_A(s)\) returns the value of \(s\) if
that value exists and satisfies \(A\), and is undefined otherwise.
A test is a first-occurrence lookup on the singleton consisting of
the already selected point.

Every term is a finite expression in earlier terms. Write
\(\operatorname{Def}(t)\) for its Boolean definedness status.
The undefined symbol is not a domain point and is never compared
with points by \(<\). Every defined nonroot term lies in
\((x,y)\). These are metalinguistic event terms, not additional
temporal connectives or predicates in the output formula. For example,
the source's last \(q\) before the first \(p\) in \((x,y)\)
is \(\operatorname{Last}_q(x,\operatorname{First}_p(x,y))\).
\end{hbdefinition}

\begin{hbdefinition}{def:p09-gpss-canonical-prefix}{The canonical positive-chain event test}
For a chain requiring point labels \(\Theta_1,\ldots,\Theta_m\)
and the preceding gap labels \(\Phi_0,\ldots,\Phi_{m-1}\),
with no final gap requirement, define future formulas
\[
 \Lambda_m=\top,\qquad
 \Lambda_i=\Phi_i\mathbin{\mathsf U_s}
                   (\Theta_{i+1}\land\Lambda_{i+1})
                       \quad(0\leq i<m).
\]
For bounds \(a<b\) within the root interval, define
\[
 \mu_1=\operatorname{First}_{\Theta_1\land\Lambda_1}(a,b),\qquad
 \mu_i=\operatorname{First}_{\Theta_i\land\Lambda_i}(\mu_{i-1},b)
                       \quad(2\leq i\leq m).
\]
The Boolean event test is
\[
 V_m(a,b)=
   \bigwedge_{i=1}^{m}\operatorname{Def}(\mu_i)
   \land\neg\operatorname{Def}
          (\operatorname{Last}_{\neg\Phi_0}(a,\mu_1)).
\]
For \(m=0\), set \(V_0=\mathsf{true}\), using no event
terms. Each use of \(V_m\) retains its own displayed point and
gap lists. The \(\Lambda_i\) formulas do not contain the
bounds \(a,b\) as parameters.
\end{hbdefinition}

\begin{hblemma}{thm:p09-gpss-canonical-prefix}{proved=annex:p9.gpss.canonical.prefix}{Correctness and minimality of the canonical chain}
On a bounded interval in a discrete complete order, \(V_m(a,b)\)
holds exactly when there is a strict chain of the required \(m\)
points and preceding gaps inside \((a,b)\). When it holds,
the selected points \(\mu_i\) satisfy those conditions and are
coordinatewise no later than the points of any other satisfying
chain. Future viability in \(\Lambda_i\) may inspect beyond
\(b\); the bounded event selectors still enforce every actual
witness to be strictly before \(b\).
\end{hblemma}

\begin{hbdefinition}{def:p09-gpss-determining-events}{The finite determining-list recursion}
Write a bounded decomposition as
\(D_W(a,b)\), where
\(W=(\Theta_1,\ldots,\Theta_r;\Phi_0,\ldots,\Phi_r)\).
Its rank is the number \(r\) of internal points. For \(0\leq j<r\),
let \(W_j\) be its prefix with \(j\) points and gap labels
\(\Phi_0,\ldots,\Phi_j\). Let \(V^{j}(a,b)\) be the
canonical positive-chain test for points \(\Theta_{j+1},\ldots,
\Theta_r\) with preceding gaps \(\Phi_j,\ldots,\Phi_{r-1}\);
thus \(V^{r}=\mathsf{true}\).

For rank zero define the decision expression
\[
 T_W(a,b)=\neg\operatorname{Def}
                    (\operatorname{Last}_{\neg\Phi_0}(a,b)).
\]
For positive rank put \(\ell=\operatorname{Last}_{\neg\Phi_r}(a,b)\)
and define
\[
\begin{aligned}
 T_W(a,b)={}&
  \bigl(\neg\operatorname{Def}(\ell)\land V^0(a,b)\bigr)\\
 &\lor\Bigl(\operatorname{Def}(\ell)\land
        \bigl[\ \bigvee_{j=0}^{r-1}
          \bigl(T_{W_j}(a,\ell)\land
             \operatorname{Def}(\operatorname{Test}_{\Phi_j}(\ell))
             \land V^{j}(\ell,b)\bigr)\\
 &\hspace{31mm}\lor\bigvee_{j=1}^{r}
          \bigl(T_{W_{j-1}}(a,\ell)\land
             \operatorname{Def}(\operatorname{Test}_{\Theta_j}(\ell))
             \land V^{j}(\ell,b)\bigr)\ \bigr]\Bigr).
\end{aligned}
\]
The first inner disjunction puts \(\ell\) in a gap; the second
puts it at an internal point. Form the event list \(E_W\) from
all terms in this expression and every subterm, including \(x,y\).
Each recursive decomposition has smaller rank. Identical endpoint
subterms can be shared as a finite acyclic expression graph.
\end{hbdefinition}

\begin{hblemma}{thm:p09-gpss-determining-events}{proved=annex:p9.gpss.determining.events}{The determining list is finite and decides the decomposition}
The recursion terminates and gives a finite event list whose definedness
statuses and values determine the truth of \(D_W(x,y)\). More
precisely, on actual endpoints \(x<y\),
\[
                             D_W(x,y)\longleftrightarrow T_W(x,y).
\]
All labels in the list are finite future-only \(L(U)\) formulas.
The construction includes the last-gap failure being absent, lying
in any gap, or coinciding with any internal witness. No label is
tested at either external endpoint.
\end{hblemma}

\begin{hbtheorem}{thm:p09-gpss-bounded-complement}{proved=annex:p9.gpss.bounded.complement}{Constructive relative complement on bounded intervals}
For every GPSS decomposition \(D_W(x,y)\), a finite procedure
constructs decompositions \(C_1(x,y),\ldots,C_N(x,y)\), with
the same permitted future-only label class, such that over every
discrete complete model
\[
             x<y\land\neg D_W(x,y)
                        \longleftrightarrow\bigvee_{i=1}^{N}C_i(x,y).
\]
The procedure enumerates all defined/undefined assignments to the
finite event list and all weak equality/order types of the defined
values, with \(x<y\). It compiles every consistent case into one
decomposition by the following exact constraints:

\begin{enumerate}
\item A defined first occurrence has its label true and has no earlier
      occurrence between its lower bound and its value; a defined last
      occurrence has no later occurrence before its upper bound.
\item An undefined first or last occurrence with defined ordered bounds
      has no occurrence anywhere between those bounds. Undefined or
      unordered bounds force an undefined result without this extra test.
\item A defined test equals its argument and has its label true; an
      undefined test with a defined argument has the label false there.
\item Universal interval constraints are transferred to every merged
      gap and every inserted point they cover. All point constraints
      are conjoined at their value's equality class.
\end{enumerate}

Retain exactly the cases where \(T_W\) is false. Impossible order
or status cases contribute no disjunct; contradictory label cases
may remain as unsatisfiable decompositions. Empty output and empty
open intervals are included. This supplies the finite determining
and back-translation steps of the event-expression outline for
bounded intervals
\cite[Proposition~2.4, p.~171]{GabbayPnueliShelahStavi1980}.
\end{hbtheorem}

\begin{hbtheorem}{thm:p09-gpss-suffix-complement}{proved=annex:p9.gpss.suffix.complement}{Complementation of future rays, including last points}
For an actual lower endpoint \(x\), the complement of every
\(D(x,+\infty)\) has a finite disjunction of exact GPSS
decompositions on \((x,+\infty)\). This includes all the
allowed global cofinal-occurrence conjuncts.

Take its already proved future expression \(B(x)\) from
\cref{thm:p09-gpss-suffix-translation} and put \(E=\neg B\).
Its outer Boolean construction has only constants and strict
future modalities as leaves. At a last point evaluate those leaves
by \(F A=\mathsf{false}\), \(G A=\mathsf{true}\), and
\(A\mathbin{\mathsf U_s}C=\mathsf{false}\), obtaining a
Boolean value \(e_{\mathrm{last}}\). For a point having a next point,
form the future label \(L_E\) by Boolean recursion using
\[
\begin{aligned}
 L_{F A}&=A\lor F A,\\
 L_{G A}&=A\land G A,\\
 L_{A\mathbin{\mathsf U_s}C}
       &=C\lor\bigl(A\land(A\mathbin{\mathsf U_s}C)\bigr).
\end{aligned}
\]
One output decomposition has one internal point labelled \(L_E\),
first gap labelled false and final gap labelled true. If
\(e_{\mathrm{last}}\) is true, add the zero-point decomposition
whose only gap is labelled false. These are all the output cases.
\end{hbtheorem}

\begin{hbtheorem}{thm:p09-gpss-prefix-complement-obstruction}{proved=annex:p9.gpss.prefix.complement.obstruction}{Why formal left endpoints require a different claim}
The same future-labelled prefix class is not closed under
complement, even over \(\langle\mathbb Z,<\rangle\).
The legal decomposition
\[
 D(-\infty,y)=\exists z<y\ \forall u<z\ q_0(u)
\]
has one internal point labelled true, initial gap labelled \(q_0\),
and final gap labelled true. Its complement has no finite disjunction
of prefix decompositions whose labels remain in \(L(U)\).
The analogous whole-order decomposition
\(D(-\infty,+\infty)=\exists z\,\forall u<z\,q_0(u)\)
has the same failure, even with the permitted future-cofinal conjuncts.

Thus the future proof uses bounded intervals and future rays with an
actual lower endpoint. A claim about all formal-endpoint configurations
in the unchanged future-label class would be false; changing to a
past-label or enlarged coinitial grammar would be a separate construction.
\end{hbtheorem}

\begin{hbtheorem}{thm:p09-gpss-discrete-future-completeness}{proved=annex:p9.gpss.discrete.future.completeness}{Completion of the discrete-complete future construction}
On discrete complete orders, every member of
\(\operatorname{Bet}_G(x,y)\) on \(x<y\), and every
member of \(\operatorname{Aft}_G(x)\), has a finite
disjunction of the exact GPSS decompositions. Here \(x\)
is an actual point variable. Consequently every source future
\(L_1\)-formula at \(x\), with quantifiers restricted to
\([x,+\infty)\), has an equivalent \(L(U)\)-formula at
every point of every discrete complete model. This is Theorem~2.2.
Every \(L_1\)-formula with only \(x\) free therefore also has
an equivalent \(L(U)\)-formula at point zero of every
\(\omega\)-model, the exact conclusion of Theorem~2.1.

The bounded failure on arbitrary complete orders is proved in
\cref{thm:p09-gpss-complete-bounded-obstruction}; the enlarged local
boundary tools and positive-chain construction are
\cref{thm:p09-gpss-infimum-case-compiler,thm:p09-gpss-complete-positive-chain}.
The discrete \(U\)-only conclusion above is not asserted beyond
its stated order class
\cite[Theorems~2.1--2.2, pp.~170--171]{GabbayPnueliShelahStavi1980}.
\end{hbtheorem}

\begin{hbexample}{def:p09-gpss-last-failure-example}{A last failure may be the witness itself}
On \(\{0,1,2,3\}\), let the endpoints be zero and three.
Take a rank-one word with first gap always true, final gap false
only at two, and its point label true only at two. Its last final-gap
failure is \(\ell=2\), and the decomposition is true with
its witness equal to \(\ell\). The final gap excludes that
witness, so the failed gap label there is harmless. If the point
label is instead true only at one, the decomposition is false.
Both outcomes are represented by the equality and gap cases of the
finite determining expression. On an adjacent endpoint interval,
every positive-rank decomposition is false and its relative complement
is true.
\end{hbexample}

% GPSS complete-order boundary constructions and the necessary endpoint extension.

\subsection{Future limits and complete-order boundary constructions}
\label{sec:p09-gpss-complete-order-boundaries}

The extension from discrete complete orders needs formulas for nonattained
bounds. It also needs a change to the intermediate decomposition grammar:
the endpoint-free class in \cref{def:p09-gpss-decomposition} is not closed
under bounded relative complement on all complete orders. The counterexample
below has two actual endpoints. This section supplies future-only limit
formulas, an infimum-event compiler with a labelled actual lower endpoint,
and the bounded positive-chain result at its stated scope.
The modification announced after Proposition~2.4 concerns Kamp's
\(S,U\) result. It does not assert an arbitrary-complete-order
\(U\)-only future theorem. The local constructions below retain
their explicit language and endpoint restrictions
\cite[pp.~169--171]{GabbayPnueliShelahStavi1980}.

\begin{hblemma}{thm:p09-gpss-complete-order-bridge}{proved=annex:p9.gpss.complete.order.bridge}{The exact complete-order and temporal bridges}
The source condition that every nonempty bounded subset has an infimum and
supremum is equivalent to Kamp's \(\mathsf{Com}\) condition: every
nonempty lower-bounded subset has an infimum and every nonempty
upper-bounded subset has a supremum. The orders are nonempty and strict;
endpoints and adjacent points are permitted.

For formulas made from unary letters, Boolean operations and Kamp's
future constructor \(U(C,B)\), renaming its letters to the corresponding
GPSS letters and replacing that constructor by
\(B\mathbin{\mathsf U_s}C\) preserves truth on every linear order.
In particular the nested future word construction in
\cref{thm:p09-kamp-iii2-lemma1} supplies an \(L(U)\) formula when its
labels already have \(L(U)\) presentations. This bridge does not convert
an arbitrary elementary Kamp label, or a formula containing \(S\), into
a future-only label.
\end{hblemma}

\begin{hbdefinition}{def:p09-gpss-right-boundary-labels}{Nonempty right limits and right neighborhoods}
For a permitted future label \(A\), define the following actual \(L(U)\)
formulas:
\[
\begin{aligned}
 N^+A&=A\mathbin{\mathsf U_s}\top,\\
 \partial^+A&=\neg\bigl((\neg A)\mathbin{\mathsf U_s}\top\bigr),\\
 K^+A&=F\top\land\partial^+A.
\end{aligned}
\]
Here \(N^+A\) says that some strict right neighborhood has \(A\)
throughout its open interval. The symbol \(\partial^+\) agrees with
the unary future operation in \cref{def:p09-kamp-limit-operators};
\(K^+\) additionally requires a future point. These are abbreviations
for the displayed finite formulas, not new primitive connectives.
\end{hbdefinition}

\begin{hblemma}{thm:p09-gpss-right-limit}{proved=annex:p9.gpss.right.limit}{The precise right-limit semantics}
On every nonempty linear-order model,
\[
\begin{split}
 K^+A(x)\quad\Longleftrightarrow\quad&
 \exists b\,(x<b)\ \land\\
 &\forall b\bigl(x<b\longrightarrow
            \exists z\,(x<z<b\land A(z))\bigr).
\end{split}
\]
The predicate is tested at \(z\), with no condition on \(A(x)\).
Both \(\partial^+A\) and \(K^+A\) are false at a point with an
immediate successor. At a last point \(\partial^+A\) is true and
\(K^+A\) is false. For every actual \(x<y\), \(K^+A(x)\)
is equivalent to nonemptiness and coinitiality at \(x\) of the
\(A\)-points in \((x,y)\). This equivalence uses the whole-model
label \(A\), without restricting the quantifiers inside \(A\).
\end{hblemma}

\begin{hbdefinition}{def:p09-gpss-future-word-limits}{A future-only construction for finite-word accumulation}
A word is
\[
 W=(A_0,B_1,A_1,\ldots,B_n,A_n),\qquad n\geq0,
\]
with all point labels \(A_i\) and gap labels \(B_i\) in \(L(U)\).
Its realization is \(z_0<\cdots<z_n\), with \(A_i(z_i)\)
and \(B_i\) throughout \((z_{i-1},z_i)\). At \(n=0\)
there is just one labelled point. Define its starting formula by
\[
 E^+_{(A_0)}=A_0,\qquad
 E^+_W=A_0\land
             \bigl(B_1\mathbin{\mathsf U_s}E^+_{W^-}\bigr)
                    \quad(n\geq1),
\]
where \(W^-=(A_1,B_2,A_2,\ldots,B_n,A_n)\).
The formula \(E^+_W\) says that a realization starts at the current
point and places no upper bound on its end.

Define a second finite recursion:
\[
 \Gamma_{(A)}=K^+A,\qquad
 \Gamma_W=K^+(E^+_W)\land\Gamma_{W^-}\quad(n\geq1).
\]
Every recursive call removes the first point and first gap. Each resulting
formula uses only future until and Boolean operations on the supplied
labels; no past operator occurs.
\end{hbdefinition}

\begin{hblemma}{thm:p09-gpss-future-word-limit}{proved=annex:p9.gpss.future.word.limit}{Finite patterns arbitrarily close to the right}
For every word \(W\), \(\Gamma_W(x)\) is equivalent, over every
linear-order model, to
\[
 \exists d\,(x<d)\ \land\
 \forall b\left(x<b\to
   \exists z_0<\cdots<z_n
       \bigl(x<z_0\land z_n<b\land
                    (z_0,\ldots,z_n)\text{ realizes }W\bigr)\right).
\]
All realization conditions expand to the finite point and gap formula
in the definition.

Consequently Kamp's vacuous future limiting condition for the same
labelled word has the future-only presentation
\[
                       \neg F\top\ \lor\ \Gamma_W.
\]
For the original predicate-letter words this is exactly the semantic
condition of \(\operatorname{Lf}_n\) in
\cref{def:p09-kamp-limiting-classes}, including \(n=0\).
The proof supplies a different finite splice from the displayed
\(\Lambda_W^+\) recursion: that earlier recursion may contain \(S\)
inside a formula for a realization ending at the current point.
Here the formula is proved directly to remain in \(L(U)\).
\end{hblemma}

\begin{hblemma}{thm:p09-gpss-fixed-start-limit}{proved=annex:p9.gpss.fixed.start.limit}{Accumulation while retaining the initial point}
Let \(n\geq1\) and keep the word \(W\) above. The condition that
there is a later point and, for every \(b>x\), a realization of
\(W\) starts at \(x\) and ends strictly before \(b\), has the
future-only expression
\[
                      A_0\land N^+B_1\land\Gamma_{W^-}.
\]
The corresponding condition without the nonemptiness requirement has
expression
\[
        \neg F\top\ \lor\
                      (A_0\land N^+B_1\land\Gamma_{W^-}).
\]
Thus the fixed-start right-limit support used in Kamp's boundary
construction has an \(L(U)\) presentation whenever its labels do.
This statement retains the original initial point and the initial gap.
\end{hblemma}

\begin{hbdefinition}{def:p09-gpss-left-labelled-decomposition}{Decompositions with a labelled actual lower endpoint}
For actual \(x<y\), or actual \(x\) and formal \(+\infty\),
a \emph{left-labelled decomposition} is
\[
                           \alpha(x)\land D(x,y),
\]
where \(\alpha\) is a permitted future label and \(D\) is exactly a
decomposition of \cref{def:p09-gpss-decomposition}. The neutral
choice is \(\alpha=\top\). The same definition applies to a future
ray, with its permitted cofinal conditions. There is no label at a
formal point and no new individual parameter inside a label.

This is an explicitly enlarged intermediate grammar. The original
decomposition class has no point label at \(x\). All uses below
identify which of these two classes is the output class.
\end{hbdefinition}

\begin{hblemma}{thm:p09-gpss-left-labelled-closure}{proved=annex:p9.gpss.left.labelled.closure}{Positive assembly and future-ray complementation}
Finite disjunctions of left-labelled decompositions are closed under
finite conjunction and under the guarded existential concatenation of
\cref{thm:p09-gpss-positive-decomposition}, when the lower endpoint
is actual. At a concatenated midpoint \(z\), the right operand's
left label is conjoined with the required point label at \(z\).
All empty Boolean cases are retained.

On every linear order, every finite disjunction of left-labelled
future-ray decompositions has an \(L(U)\) expression, and its
complement has a left-labelled future-ray decomposition. Explicitly,
translate the disjunction to \(B(x)\) by
\cref{thm:p09-gpss-suffix-translation}, retaining every left label;
return
\[
              (\neg B)(x)\land D_{\top}(x,+\infty),
\]
where \(D_{\top}\) has no interior point, the single gap label
true, and no cofinal conjunct. Last points are included. The output
class here is the enlarged class just defined.
\end{hblemma}

\begin{hbdefinition}{def:p09-gpss-infimum-events}{Bounded infimum events and their exact statuses}
In a complete model let \(s<t\) be defined actual values, and put
\[
                     A_{s,t}=\{u:s<u<t,\ A(u)\}.
\]
The event \(\operatorname{Inf}_A(s,t)\) is undefined if a bound is
undefined, the bounds are not strictly ordered, or this set is empty.
Otherwise it returns \(c=\inf A_{s,t}\). It records one of three
statuses:
\[
\begin{array}{c|l}
 \mathrm{initial}&c=s,\\
 \mathrm{attained}&s<c<t,\quad A(c),\\
 \mathrm{unattained}&s<c<t,\quad\neg A(c).
\end{array}
\]
The initial status does not test \(A(s)\): \(s\) is excluded from
the occurrence set regardless of that truth value. Defined bounds and
values are domain points. Undefinedness is a separate symbol with no
place in their strict order.
\end{hbdefinition}

\begin{hblemma}{thm:p09-gpss-infimum-descriptors}{proved=annex:p9.gpss.infimum.descriptors}{All infimum statuses with future labels}
For actual \(s<t\), the following are exact descriptions of the
event and its status:
\[
\begin{array}{c|l}
 \mathrm{undefined}&\forall u\,(s<u<t\to\neg A(u)),\\
 \mathrm{initial}&c=s,\quad K^+A(s),\\
 \mathrm{attained}&s<c<t,\quad A(c),\
                    \forall u\,(s<u<c\to\neg A(u)),\\
 \mathrm{unattained}&s<c<t,\quad\neg A(c)\land K^+A(c),\
                    \forall u\,(s<u<c\to\neg A(u)).
\end{array}
\]
On a complete order exactly one row holds, with a unique value in
each defined row. Each row's equivalence to its specified infimum
property holds on any linear order on which that bound exists.
The point labels in these descriptions are actual \(L(U)\) formulas.
The initial row is precisely the new lower-endpoint label when \(s=x\).
An empty interval contributes only the undefined row.

For the unbounded future occurrence set \(\{u:s<u,\ A(u)\}\),
the same four statuses have the following formulas at \(s\),
in the same order:
\[
\begin{array}{c|l}
 \mathrm{undefined}&G\neg A,\\
 \mathrm{initial}&K^+A,\\
 \mathrm{attained}&(\neg A)\mathbin{\mathsf U_s}A,\\
 \mathrm{unattained}&(\neg A)\mathbin{\mathsf U_s}
                                      (\neg A\land K^+A).
\end{array}
\]
Over complete orders these are mutually exclusive and exhaustive,
including a last point. The interior boundary is the strict until
witness in each of the last two rows.
\end{hblemma}

\begin{hbtheorem}{thm:p09-gpss-infimum-case-compiler}{proved=annex:p9.gpss.infimum.case.compiler}{Finite infimum-event cases have explicit left-labelled decompositions}
Fix actual roots \(x<y\) and a finite acyclic list of events.
Each event is \(\operatorname{Inf}_A(s,t)\), where both arguments
are roots or earlier events, and \(A\) is a supplied future label.
One may also include tests \(\operatorname{Test}_A(s)\), returning
\(s\) if defined and satisfying \(A\), and undefined otherwise,
provided \(s\) is the root \(x\) or an earlier event, never \(y\).

Every Boolean condition on the definedness, the three infimum statuses,
and the equality/order of the defined values has an effective finite
disjunction of left-labelled decompositions equivalent to it on
complete models under \(x<y\).

The construction enumerates all definedness and status assignments and
all complete weak orders of the defined terms and roots. It retains
only the requested cases and compiles the displayed infimum descriptors
and test clauses. Equal values share one point. All labels at \(x\)
are conjoined into its left label. Every other nonroot value is
strictly between \(x,y\). Universal interval constraints are placed
on every merged gap and interior point they cover. No event symbol,
extra proposition letter, or endpoint parameter survives inside a label.
\end{hbtheorem}

\begin{hbdefinition}{def:p09-gpss-complete-positive-chain}{The complete-order positive-chain event recursion}
Fix a list of \(m\geq1\) required points
\(\Theta_1,\ldots,\Theta_m\) and preceding gaps
\(\Phi_0,\ldots,\Phi_{m-1}\), with no final-gap condition.
Let \(V_{\Theta;\Phi}(a,b)\) be the following Boolean event
expression. For the empty point list it is true. For a nonempty
list and undefined or non-strictly ordered bounds it is false.
Otherwise form the point word
\[
                     W=(\Theta_1,\Phi_1,\Theta_2,\ldots,
                                      \Phi_{m-1},\Theta_m)
\]
and the events
\[
\begin{aligned}
 c&=\operatorname{Inf}_{E^+_W}(a,b),&
 h&=\operatorname{Inf}_{\neg\Phi_0}(a,c),\\
 j&=\operatorname{Test}_{N^+\Phi_0}(c),&
 k&=\operatorname{Test}_{\Phi_0\land N^+\Phi_0}(c).
\end{aligned}
\]
For \(m=1\), \(W=(\Theta_1)\). For \(m>1\) let
\[
 R=V_{(\Theta_2,\ldots,\Theta_m);
                  (\top,\Phi_2,\ldots,\Phi_{m-1})}(c,b);
\]
for \(m=1\) let \(R=\mathsf{true}\). Thus the shortened
list has a true initial gap and retains every gap between
its remaining points. Put
\[
\begin{split}
 V_{\Theta;\Phi}(a,b)={}&
 \operatorname{Def}(c)\land\neg\operatorname{Def}(h)\land R\\
 &{}\land\bigl[
       \operatorname{Att}(c)
       \lor(\operatorname{Init}(c)\land\operatorname{Def}(j))
       \lor(\operatorname{Unatt}(c)\land\operatorname{Def}(k))
                    \bigr].
\end{split}
\]
Here \(\operatorname{Att},\operatorname{Init},\operatorname{Unatt}\)
are the recorded statuses of this particular infimum event, false
when it is undefined. The recurrence removes one point per call,
even if an initial event returns the same lower endpoint.
All symbols in this display are metalinguistic event expressions
to be compiled by \cref{thm:p09-gpss-infimum-case-compiler}.
\end{hbdefinition}

\begin{hbtheorem}{thm:p09-gpss-complete-positive-chain}{proved=annex:p9.gpss.complete.positive.chain}{A determining construction for bounded positive chains on complete orders}
On actual \(a<b\) in a complete model, the event expression just
defined holds exactly when there are points
\[
                           a<z_1<\cdots<z_m<b
\]
with \(\Theta_i(z_i)\), with \(\Phi_0\) throughout
\((a,z_1)\), and with \(\Phi_i\) throughout
\((z_i,z_{i+1})\) for \(1\leq i<m\).
The empty point list imposes no condition.

For \(m\) points the event list contains at most \(4m\)
nonroot nodes in addition to its two roots. It uses only
future labels, infima and tests. Its Boolean decision
therefore has an explicitly constructible finite
disjunction of left-labelled decompositions, as does
its relative complement under \(a<b\).

The attained infimum supplies a first point. An initial
or interior unattained infimum instead supplies points
arbitrarily close on its right; the neighborhood and
interior point tests in the displayed recurrence preserve
the whole initial gap. The suffix can be realized before
the upper bound independently of the provisional global
future realization, and the proof supplies the exact
splice joining them. This replaces the finite-interval
first-occurrence assumption of
\cref{thm:p09-gpss-canonical-prefix} for this positive-chain
component. A final-gap requirement remains separate.
\end{hbtheorem}

\begin{hbtheorem}{thm:p09-gpss-complete-bounded-obstruction}{proved=annex:p9.gpss.complete.bounded.obstruction}{Why the original bounded output class is insufficient}
The original endpoint-free decomposition class is not closed under
relative complement on arbitrary complete orders, even with actual
bounded endpoints. It already fails for
\[
 D_A(x,y)=x<y\land
       \exists z\bigl(x<z<y\land
                \forall u(x<u<z\to A(u))\bigr).
\]
This is an original decomposition with one true-labelled interior
point, initial gap \(A\), and final gap true.

Take the complete order
\[
                  T=\{0,2\}\cup\{1/n:n\geq1\},
                  \qquad x=0,\quad y=2.
\]
Vary \(q_0(1/n)\) independently, take \(A=q_0\), fix \(q_0\)
at the endpoints, and fix all other predicates throughout \(T\).
No finite disjunction of original decompositions agrees with
\(x<y\land\neg D_A(x,y)\) for every such interpretation.
This is a failure at an actual lower endpoint with no immediate
successor, distinct from the formal-left-endpoint obstruction in
\cref{thm:p09-gpss-prefix-complement-obstruction}.
\end{hbtheorem}

\begin{hblemma}{thm:p09-gpss-initial-gap-complement}{proved=annex:p9.gpss.initial.gap.complement}{The complete initial-neighborhood complement with its endpoint label}
On every linear-order model, for actual \(x<y\),
\[
 \neg D_A(x,y)\quad\Longleftrightarrow\quad
       \operatorname{Adj}(x,y)\ \lor\ K^+(\neg A)(x),
\]
where \(\operatorname{Adj}(x,y)\) means
\(x<y\land\neg\exists z(x<z<y)\).
Thus its guarded complement is the disjunction of two explicitly
left-labelled decompositions: the zero-point word with false gap and
true left label; and the zero-point word with true gap and left label
\(K^+(\neg A)\). These cases include adjacency, a nonattained
initial limit, and every truth value of \(A(x)\) and \(A(y)\).
\end{hblemma}

\begin{hbexample}{def:p09-gpss-boundary-examples}{Attained and unattained beginnings}
On \(\mathbb R\), with bounds \(0<2\), take successively
\(A=\varnothing\), \(A=(0,1)\), \(A=\{1\}\), and
\(A=(1,3/2)\). The infimum-event statuses are respectively
undefined, initial, attained at 1, and unattained at 1. Replacing
\(A(0)\) alone changes none of these statuses. On \(0<1\) with
no intermediate point the event is undefined for every interpretation.

For \(W=(A,B,C)\), the word-limit formula is
\[
       K^+\bigl(A\land(B\mathbin{\mathsf U_s}C)\bigr)
                                      \land K^+C.
\]
On the real line with \(A=(0,1)\), \(B=\mathbb R\), and
\(C=\{2\}\), its first conjunct holds at zero but its second
does not; the required two-point patterns cannot occur arbitrarily
close to zero. With \(A=C=(0,1)\) both conjuncts hold and every
right neighborhood contains a realizing pair.

For an event-list example take
\(i=\operatorname{Inf}_P(x,y)\) and
\(j=\operatorname{Inf}_Q(i,y)\).
The case that \(i\) is initial and \(j\) is attained compiles to
\[
 K^+P(x)\land x<y\land
 \exists c\bigl(x<c<y\land Q(c)\land
                   \forall u(x<u<c\to\neg Q(u))\bigr).
\]
Its left label is \(K^+P\), its one point label is \(Q\),
and its two gaps are \(\neg Q,\top\). The infimum descriptors
identify \(i=x\) and \(j=c\) in both directions.

In the positive-chain recurrence with one point, bounds \(0<2\),
and \(\Theta_1=(1,3/2)\) on \(\mathbb R\), the candidate
infimum is the unattained interior point \(c=1\). With
\(\Phi_0=\mathbb R\) the chain exists. With
\(\Phi_0=\mathbb R\setminus\{1\}\) it does not: every
possible first point lies above 1, so its initial gap contains
the failed point 1. The interior test in \(k\) distinguishes
these interpretations even though both open side intervals
can satisfy their gap conditions.
\end{hbexample}

\begin{proofboundary}{def:p09-gpss-complete-boundary-continuation}{What the boundary constructions establish}
The infimum compiler translates a supplied finite list; it does not
by itself construct a list deciding every decomposition.
\Cref{thm:p09-gpss-complete-positive-chain} supplies such a list
for the whole positive chain without its final gap.
A general final gap involves upper-boundary information that is
not covered by those finite-list premises. Consequently the proved
local statements do not license a full regional negation induction
in this future-only grammar.

For a nonempty occurrence set in \((a,b)\), an unattained supremum
\(c\) is either the outer endpoint \(b\) or an interior point.
Its left approach requires
\[
 \forall r<c\ \exists u\,
          (a<u<c\land r<u\land A(u)).
\]
This includes a real witness by taking \(r=a\). An interior
supremum also requires \(\neg A(c)\) and absence of occurrences
in \((c,b)\); an attained interior supremum instead requires
\(A(c)\) and that same absence. These are first-order boundary
conditions, not admitted future labels at \(c\).
In particular left accumulation at \(c\) cannot itself be a
future-only label there: on \(\mathbb R\), interpretations
\(A=(-1,0)\) and \(A=\varnothing\) agree at and after zero
and disagree on left accumulation at zero.

Kamp's complete-order construction permits past-dependent elementary
labels at such boundaries, and proves completeness for the
\(S,U\) basis. That valid result is not a promise that the past
operators can be removed from every complete-order future presentation.
The countertheorem above also prevents an inference that the original
endpoint-free GPSS Main Lemma extends unchanged. The source's announced
complete-order modification and any separately stated future-language
comparison must retain their actual operator bases
\cite[p.~171]{GabbayPnueliShelahStavi1980}.
\end{proofboundary}

% Corrected complete-order target and the Hirshfeld--Rabinovich obstruction.

\subsection{Why future until completeness stops at discrete orders}
\label{sec:p09-gpss-complete-order-complement}

GPSS Theorem~2.2 concerns discrete complete models. The modification
announced on its following page yields another route to Kamp's
\(S,U\) completeness result; it is not a theorem that \(U\)
alone expresses every future first-order formula on arbitrary complete
orders \cite[pp.~170--171]{GabbayPnueliShelahStavi1980}.
The latter statement is false even on the real line. We give the
Hirshfeld--Rabinovich nonexpressibility argument, including the ordinal
game used by its real-order construction
\cite[secs.~2--4 and appendix~A]{HirshfeldRabinovich2003Future}.
The source used here is the authors' version dated 6 March 2003;
its page locators refer to that sixteen-page version.
The source used here is the authors' version dated 6 March 2003;
its page locators refer to that sixteen-page version.

\begin{hbdefinition}{def:p09-hr-future-basis}{First-order future modalities and their depth}
A first-order truth table for an \(r\)-ary modality is a monadic
order formula \(\mu(t;X_1,\ldots,X_r)\) with only \(t\)
free as an individual variable. Its arguments are interpreted as
arbitrary unary sets. A truth table is \emph{future} when its value
at \(t\) agrees with its value in the induced labelled tail
\([t,+\infty)\), for every nonempty linear order and interpretation.
Equivalently it has a presentation with every quantifier restricted
to that tail. The forward implication follows by restricting all
quantifiers in the original formula; the reverse follows from the
semantics of that restriction.

The quantifier depth of a fixed truth table is the maximum number
of nested quantifiers. Write \(\mathsf{FTL}_d\), \(d\geq1\),
for the temporal language allowing every finite-arity future modality
with a supplied truth table of depth at most \(d\), together with
constants, proposition letters and Boolean connectives. A formula
has finite syntax but arbitrary modal nesting; the basis itself
may be infinite. The depth bound concerns individual truth tables,
not the first-order expansion of a whole temporal formula.
The strict until of \cref{def:p09-gpss-future-labels} belongs
to \(\mathsf{FTL}_2\).
\end{hbdefinition}

\begin{hbdefinition}{def:p09-hr-ordinal-models}{Finite powers of the natural order}
For \(k\geq1\), let \(O_k\) be the lexicographically ordered
set of \(k\)-tuples of natural numbers, with the leftmost
coordinate most significant. Its order type is \(\omega^k\).
Use finite Cantor polynomials
\[
                 \omega^{k-1}a_{k-1}+\cdots+\omega a_1+a_0
\]
for its points. The zero tuple is its least point.
Order sums and products in the following statements are ordinal
operations, not commutative numerical operations.

Two structures are \(n\)-equivalent, written \(\equiv_n\),
if they satisfy the same first-order sentences of quantifier
depth at most \(n\). An \(n\)-round comparison game chooses
one point in either structure at each round, answered by a
point in the other; the answering player wins when the
correspondence preserves equality, order and all unary
predicates on the selected points.
\end{hbdefinition}

\begin{hblemma}{thm:p09-hr-ordinal-game}{proved=annex:p9.hr.ordinal.game}{The complete ordinal comparison strategy}
Let \(n\geq0\), let \(\alpha,\beta\) be nonzero ordinals
below \(\omega^\omega\), and let \(\rho<\omega^n\).
The answering player has a winning \(n\)-round strategy on
\[
                   \omega^n\alpha+\rho,\qquad
                   \omega^n\beta+\rho.
\]
In particular \(O_k\equiv_n O_m\) whenever \(k,m\geq n\)
and \(k,m\geq1\). The annex gives the interval strategy
and proves its quantifier-depth consequence directly.
\end{hblemma}

\begin{hblemma}{thm:p09-hr-ordinal-tail}{proved=annex:p9.hr.ordinal.tail}{Future homogeneity of the ordinal models}
Every tail of \(O_k\) with an actual first point is isomorphic
to \(O_k\). Every future formula without unary predicates,
including a temporal formula made from constants and future
modalities, consequently has a constant truth value on \(O_k\).
For every \(\mathsf{FTL}_d\) formula without proposition
letters, this value is the same in \(O_k,O_m\) whenever
\(k,m>d+1\).
\end{hblemma}

\begin{hbdefinition}{def:p09-hr-future-limit-hierarchy}{A future hierarchy with its reference point retained}
For distinct current and tested variables \(x,t\), put
\[
\begin{aligned}
 L_0(x,t)&=(t=t),\\
 L_{j+1}(x,t)&=
 x<t\land
 \forall v\left(x\leq v<t\longrightarrow
       \exists u\,(v<u<t\land L_j(x,u))\right).
\end{aligned}
\]
All bound names are fresh. The fixed parameter \(x\) remains
the same throughout the recursion. Define
\[
                         \varphi_n(x)=\exists t\,(x<t\land L_n(x,t)),
                         \qquad n\geq1.
\]
These are future first-order formulas: every variable quantified
in a relevant branch is at or above \(x\). The explicit
condition \(x<t\) prevents a vacuous initial-point limit.

For a unary predicate \(P\), let \(\varphi_n^P(x)\) be
the formula obtained by restricting every quantified variable
in \(\varphi_n\) to \(P\), and conjoining \(P(x)\).
Existentials become \(\exists u(P(u)\land\cdots)\)
and universals become \(\forall u(P(u)\to\cdots)\);
all original order bounds remain. Thus
\(\varphi_n^P\) is still a future formula.
\end{hbdefinition}

\begin{hblemma}{thm:p09-hr-limit-separator}{proved=annex:p9.hr.limit.separator}{The exact threshold of the ordinal separators}
For \(k,n\geq1\),
\[
                         O_k,0\models\varphi_n
                                  \quad\Longleftrightarrow\quad k>n.
\]
More precisely, at a nonzero point of \(O_k\), \(L_n(0,t)\)
holds exactly when the last \(n\) Cantor coefficients of
that point are zero, with the convention that no nonzero
point below \(\omega^k\) has that property when \(n\geq k\).
For an ordered unary subset \(P\) and \(x\in P\),
\(\varphi_n^P(x)\) has exactly the value of \(\varphi_n\)
in the induced order on \(P\) at \(x\).
\end{hblemma}

\begin{hbdefinition}{def:p09-hr-continuous-models}{Real models with an ordinal set of marked points}
For \(k\geq1\), form the lexicographic order
\[
                C_k=\omega^k\times[0,1),\qquad
                P_k=\{(\gamma,0):\gamma<\omega^k\}.
\]
The only variable predicate is \(P_k\); every other proposition
letter, if used, is fixed false. The interval \([0,1)\)
here is real. These are ordered labelled structures, with
least point \((0,0)\).
\end{hbdefinition}

\begin{hblemma}{thm:p09-hr-continuous-models}{proved=annex:p9.hr.continuous.models}{Completeness, two tail types, and the lifted comparison game}
The order \(C_k\) is isomorphic to the nonnegative real line.
Its marked subset has type \(\omega^k\). Tails beginning at
marked points are mutually isomorphic, and tails beginning
at unmarked points are mutually isomorphic. Thus any
one-predicate future formula on \(C_k\) is equivalent
there to one of
\[
                              \top,\quad\bot,\quad P,\quad\neg P.
\]

If \(k,m\geq n\), \(k,m\geq1\), the labelled structures
\((C_k,P_k)\) and \((C_m,P_m)\) are \(n\)-equivalent.
The strategy plays the ordinal game in the first coordinate
and copies the chosen second coordinate exactly.
\end{hblemma}

\begin{hblemma}{thm:p09-hr-modal-invariance}{proved=annex:p9.hr.modal.invariance}{Arbitrary temporal nesting preserves the same two truth values}
For \(d\geq1\), every \(\mathsf{FTL}_d\) formula in the
single proposition letter \(P\) determines the same member
of \(\{\top,\bot,P,\neg P\}\) in every \(C_k\) with \(k>d\).
In particular no such formula distinguishes the least points
of \(C_k,C_m\) when \(k,m>d\).

The induction includes modalities of every finite arity and
arbitrarily deep nesting. In its modality step the operand
truth sets are replaced by their four possible unary
descriptions. The resulting truth table still has depth
at most \(d\); its uniform value on each of the two point
types is transferred by sentences of depth at most \(d+1\).
\end{hblemma}

\begin{hbtheorem}{thm:p09-hr-no-finite-future-basis}{proved=annex:p9.hr.no.finite.future.basis}{The finite-basis obstruction over real time}
For every \(d\geq1\), the future first-order formula
\(\varphi_{d+1}^{P}\) has no equivalent
\(\mathsf{FTL}_d\) formula over the real line.
It distinguishes \(C_{d+1}\) from \(C_{d+2}\) at their
least marked points, whereas the preceding lemma prevents
every such temporal distinction.

Consequently no finite set of first-order future modalities
is expressively complete for all future monadic first-order
formulas over the real line. The same conclusion holds on
the nonnegative reals, on the rational line and its
nonnegative part, and over the class of ordinals.
The finite-basis quantifier requires finite arities and
first-order future truth tables, as in the source
\cite[Theorems~3.6, 4.1 and 5.1]{HirshfeldRabinovich2003Future}.
\end{hbtheorem}

\begin{hbcorollary}{thm:p09-gpss-until-complete-order-obstruction}{proved=annex:p9.gpss.until.complete.order.obstruction}{An explicit counterinstance to until-only complete-order expressiveness}
The strict future-until language \(L(U)\) does not express
every future \(L_1\)-formula on complete orders.
The finite first-order formula \(\varphi_3^P\) is one
counterinstance: it is false at the least point of
\(C_3\) and true at the least point of \(C_4\),
while every \(L(U)\) formula in \(P\) agrees there.
The models may equivalently be placed on the full real
line and compared at zero.

Furthermore, even the enlarged left-labelled bounded
decomposition class of
\cref{def:p09-gpss-left-labelled-decomposition} cannot
be closed under all relative complements over complete
orders. If it were, its proved positive operations and
future-ray translation would give the false \(L(U)\)
completeness assertion. The obstruction therefore cannot
be removed by finishing a general final-gap recursion
while requiring all labels to remain in that class.
\end{hbcorollary}

\begin{hbexample}{def:p09-hr-until-table}{A three-quantifier future condition already separates until}
On every \(C_k\), the strict-until operation on the four truth sets
has this same table. Rows give the interval operand \(A\), columns
the witness operand \(B\), and entries give \(A\mathbin{\mathsf U_s}B\):
\[
\begin{array}{c|cccc}
 &\bot&\top&P&\neg P\\ \hline
 \bot&\bot&\bot&\bot&\bot\\
 P&\bot&\bot&\bot&\bot\\
 \top&\bot&\top&\top&\top\\
 \neg P&\bot&\top&\top&\top
\end{array}
\]
Every open nonempty interval has an unmarked point, and every point
has a next marked point with an entirely unmarked open gap before it.
These facts prove the table and are independent of \(k\).
Boolean induction therefore makes every until-only formula the
same one of the four truth sets on every \(C_k\).

Nevertheless the future formula
\[
\begin{split}
 P(x)\land\exists t\bigl(x<t\land P(t)\land
  \forall v\bigl((P(v)\land x\leq v<t)\to
        \exists u\,(P(u)\land v<u<t)\bigr)\bigr)
\end{split}
\]
is false at the least point of \(C_1\) and true at that of
\(C_2\). It is \(\varphi_1^P\): in the latter model the
marked point with index \(\omega\) is approached by the
marked points with finite indices. In \(C_1\), each
future marked point has a preceding marked point.
The appendix gives the table and this shorter counterinstance
within the proof of
\cref{thm:p09-gpss-until-complete-order-obstruction}.
\end{hbexample}

\begin{hbdefinition}{def:p09-hr-guarded-past-limit}{The valid guarded past-limit syntax}
The source's unary past-limit operation has truth table
\[
 K^-A(t)=
          \forall r\left(r<t\longrightarrow
                     \exists u\,(r<u<t\land A(u))\right).
\]
It is equivalent to \(\neg((\neg A)\mathbin{\mathsf S_s}\top)\),
where infix strict since puts its interval operand first.
This convention is vacuously true at a first point; it
is distinct from adding a nonemptiness conjunct.
\(K^-\) is a past-dependent operation.

Let \(\mathcal A\) contain all finite formulas built
using proposition letters, Boolean operations, strict
until and \(K^-\). Its \emph{guarded future fragment}
is the Boolean closure of proposition letters and
formulas \(B\mathbin{\mathsf U_s}C\), where
\(B,C\in\mathcal A\). Equivalently, every \(K^-\)
occurrence must lie under some strict until occurrence.
Thus its operands may contain a past-limit operation,
although every complete formula in the guarded
fragment is future. This is not the grammar \(L(U)\).
\end{hbdefinition}

\begin{hblemma}{thm:p09-hr-guarded-future}{proved=annex:p9.hr.guarded.future}{Why the guarded past-limit fragment is future}
Every formula of \(\mathcal A\) has the same truth
value at every \(t>x\) in a linear model and its
induced tail \([x,+\infty)\). Hence every formula
of the guarded future fragment is future at \(x\).
The strict guard is essential: \(K^-P\) itself
can distinguish two real-order interpretations
which agree at and after the current point.
\end{hblemma}

\begin{proofboundary}{def:p09-hr-positive-source-boundary}{The separate positive comparison}
The original paper's concluding discussion attributes
expressive completeness of this guarded \(U,K^-\)
fragment to Gabbay--Hodkinson--Reynolds, Theorem~10.3.20.
Its proof is not contained in the acquired paper.
The guarded-future lemma above proves its semantic scope.
The complete constructive equivalence over Dedekind-complete
orders is supplied in \cref{thm:p09-guarded-future-completeness},
using the distinct later primary presentation
\cite{Rabinovich2014Kamp}.
The original historical referral retains its own source identity
\cite[sec.~5, item~3]{HirshfeldRabinovich2003Future}.

GPSS's complete-order announcement remains the
\(S,U\) construction. The discrete \(U\)-only
theorems remain valid. The later almost-future
completeness result remains a distinct presentation.
Neither it nor the proved guarded comparison asserts
\(U\)-only completeness on arbitrary complete orders.
\end{proofboundary}

% Rabinovich 2014, complete guarded future translation and interval complement.

\subsection{Complete-order expressiveness with a guarded past limit}
\label{sec:p09-guarded-future-completeness}

Every future monadic first-order formula has a guarded
\(U,K^-\) expression on Dedekind-complete orders. The interval
construction below supplies the negation step and the final syntactic
translation. It uses the actual past-limit operation, including its
value at an actual endpoint; it does not return to the insufficient
\(U\)-only basis. The primary is Rabinovich's Theorem~7.4 and its
supporting interval proof
\cite[secs.~5 and 7, pp.~7--15]{Rabinovich2014Kamp}.

\begin{hbdefinition}{def:p09-gfc-canonical-words}{The canonical label language and its interval words}
Let \(\mathcal A\) be the full finite \(U,K^-\) language of
\cref{def:p09-hr-guarded-past-limit}, with Boolean constants,
arbitrary Boolean operations and the specified strict until. Thus
\[
 K^-A(t)=\forall r\left(r<t\to
                 \exists u\,(r<u<t\land A(u))\right).
\]
In particular \(K^-A\) is true at a first point. A canonical
label is a supplied finite formula of \(\mathcal A\), evaluated
in the whole original model at its sole current-point argument.
Naming such a formula as a predicate is only an abbreviation;
all names are inlined by capture-free substitution in the output.

For \(n\geq1\) and a word
\(W=(A_0,B_1,A_1,\ldots,B_n,A_n)\), define
\[
\begin{split}
 [W](a,b)={}&a<b\land
 \exists z_1\cdots z_{n-1}\bigl[
      a=z_0<z_1<\cdots<z_n=b\\
 &{}\land\bigwedge_{i=0}^n A_i(z_i)
       \land\bigwedge_{i=1}^n
          \forall u\,(z_{i-1}<u<z_i\to B_i(u))\bigr].
\end{split}
\]
The displayed \(z_0,z_n\) are endpoint abbreviations, not
extra quantified variables. A one-interval word has no interior
witness. Both actual endpoints have labels. Finite word lists
denote disjunctions, with the empty list false. This is the
ordered case of the source's canonical
\((a,b)\)-decomposition grammar. Its additional cases \(a=b\)
and \(b<a\) are kept separately when complements are unguarded.
No formal infinity is ever a predicate argument.
\end{hbdefinition}

\begin{hblemma}{thm:p09-gfc-label-transfer}{proved=annex:p9.gfc.label.transfer}{Transferring the completed local construction to canonical labels}
The infimum descriptors, finite event-case compiler and bounded
positive-chain construction of
\cref{thm:p09-gpss-infimum-descriptors,thm:p09-gpss-infimum-case-compiler,thm:p09-gpss-complete-positive-chain}
remain valid when their supplied labels belong to \(\mathcal A\),
and their output is interpreted as canonical words.

The finite construction treats each distinct input label as a
temporary unary letter, applies the proved uniform construction,
then substitutes its original \(\mathcal A\) formula. Its added
Boolean and strict-future operations preserve \(\mathcal A\).
The former left label becomes \(A_0\), and its unlabelled upper
endpoint receives true. This is a substitution theorem for
these supplied constructions, not an identification of the
\(U\)-only and \(U,K^-\) languages.
\end{hblemma}

\begin{hblemma}{thm:p09-gfc-positive-assembly}{proved=annex:p9.gfc.positive.assembly}{Positive canonical assembly with all endpoint labels}
Finite canonical word lists on the same ordered endpoints
are closed under conjunction. Enumerate every order-preserving
merge of their interior chains, including equal points from
different chains. Conjoin labels at each endpoint and each
merged point, and transfer every covering gap label onto
inserted points and subgaps.

Guarded concatenation at \(a<c<b\) joins a word on \((a,c)\)
and one on \((c,b)\), conjoining their two endpoint labels
and every additional required label at \(c\). Distribute over
all finite choices. These constructions also apply when the
right operand is a future ray. An empty conjunction is the
one-interval true word, or the zero-witness true ray.
\end{hblemma}

\begin{hbdefinition}{def:p09-gfc-prefix-handler}{The supplied complement handler when the last gap is uniform}
Let \(W\) have \(n\geq2\) intervals. Define \(Q_W(a,b)\)
to say that a realization of
\[
                 (A_0,B_1,A_1,\ldots,B_{n-1},A_{n-1})
\]
starts at \(a\) and ends at a point strictly before \(b\).
Its explicit event decision is
\[
 A_0(a)\land
 V_{(A_1,\ldots,A_{n-1});(B_1,\ldots,B_{n-1})}(a,b),
\]
using \cref{def:p09-gpss-complete-positive-chain}.
Use the transferred event compiler to obtain a finite canonical
list \(\overline Q_W\) for its relative complement.

Write \(\mathsf{All}_B(a,b)\) for
\(\forall u(a<u<b\to B(u))\).
Define \(\mathsf{UC}_W(a,b)\) by the canonical list for
\[
 \begin{cases}
   \neg A_0(a)\lor\neg A_1(b),&n=1,\\
   \neg A_n(b)\lor\overline Q_W(a,b),&n\geq2.
 \end{cases}
\]
Point-only alternatives have a one-interval true gap and
true at the other endpoint.
\end{hbdefinition}

\begin{hblemma}{thm:p09-gfc-uniform-final-gap}{proved=annex:p9.gfc.uniform.final.gap}{The complete uniform-final-gap reduction}
For every canonical word \(W\), on actual \(a<b\),
\[
 \mathsf{All}_{B_n}(a,b)\ \longrightarrow\
             \bigl(\neg[W](a,b)\longleftrightarrow
                         \mathsf{UC}_W(a,b)\bigr).
\]
The handler is fully constructed from the positive-chain
event recursion and requires no interval-complement
induction hypothesis of the same length.
For \(n\geq2\) it also covers adjacent outer endpoints,
where \(Q_W\) has no possible interior witness.
\end{hblemma}

\begin{hbdefinition}{def:p09-gfc-upper-boundary}{The final-gap failure set and its actual supremum}
For \(a<b\), set
\[
                 S_{a,b}=\{u:a<u<b,\ \neg B_n(u)\}.
\]
For a proposed interior boundary \(c\), define
\[
\begin{split}
 \Sigma_W(a,c,b)={}&a<c<b\land
               \mathsf{All}_{B_n}(c,b)\\
       &{}\land\bigl(\neg B_n(c)\lor K^-(\neg B_n)(c)\bigr).
\end{split}
\]
The existential \(\exists c\,\Sigma_W(a,c,b)\) is a
canonical two-interval word: the first gap is true, the
second \(B_n\), the middle label is the displayed
disjunction, and both outer labels are true. When \(c\)
remains named, \(\Sigma_W\) is that word's matrix at
the specified middle point; the same \(c\) is retained
in every additional condition.
\end{hbdefinition}

\begin{hblemma}{thm:p09-gfc-upper-boundary}{proved=annex:p9.gfc.upper.boundary}{Every final-gap boundary case with its exact formula}
In a complete order, exactly one of the following holds:
the set \(S_{a,b}\) is empty; it is nonempty with supremum
\(b\); or it has a unique supremum \(c\) with \(a<c<b\).
The first two cases are expressed respectively by
\[
              \mathsf{All}_{B_n}(a,b),\qquad
              K^-(\neg B_n)(b).
\]
The third is expressed by \(\exists c\,\Sigma_W(a,c,b)\).
At an interior supremum the attained subcase has
\(\neg B_n(c)\); the unattained subcase has
\(B_n(c)\land K^-(\neg B_n)(c)\).
The final-gap universal in \(\Sigma_W\) is retained
in both subcases. Every past-limit label is applied
to an actual point having the earlier point \(a\),
so its witness condition is nonvacuous here.
\end{hblemma}

\begin{hbdefinition}{def:p09-gfc-word-cuts}{Every point and gap cut of a word}
For \(n\geq2\), set
\[
\begin{aligned}
 P_i&=(A_0,B_1,\ldots,B_i,A_i),&
 S_i&=(A_i,B_{i+1},\ldots,B_n,A_n)
                          &&(1\leq i<n),\\
 L_i&=(A_0,B_1,\ldots,B_{i-1},A_{i-1},B_i,B_i),&
 R_i&=(B_i,B_i,A_i,B_{i+1},\ldots,B_n,A_n)
                          &&(1\leq i\leq n).
\end{aligned}
\]
An omitted prefix or suffix has no entries; thus
\(L_1=(A_0,B_1,B_1)\) and
\(R_n=(B_n,B_n,A_n)\).
These are words, not temporal operators. Their interval
counts are \(i,n-i,i,n-i+1\), respectively.
\end{hbdefinition}

\begin{hblemma}{thm:p09-gfc-word-cuts}{proved=annex:p9.gfc.word.cuts}{The exact cut identity, including both extreme gaps}
For every actual \(a<c<b\),
\[
 [W](a,b)\longleftrightarrow
 \bigvee_{i=1}^{n-1}\bigl([P_i](a,c)\land[S_i](c,b)\bigr)
 \ \lor\
 \bigvee_{i=1}^{n}\bigl([L_i](a,c)\land[R_i](c,b)\bigr).
\]
Point cuts make \(c\) a witness and retain its point label.
Gap cuts put \(c\) inside one original open gap and therefore
require \(B_i(c)\), as well as both subgap universals.
The formula holds for each chosen interior \(c\), including
one defined by a final-gap supremum.
\end{hblemma}

\begin{hbmethod}{def:p09-gfc-complement-algorithm}{The terminating canonical complement recursion}
For a word \(W\) with \(n\geq1\) intervals, construct a list
\(\mathsf N_W(a,b)\) for its complement under \(a<b\).
At \(n=1\), return the lists for
\[
 \neg A_0(a)\ \lor\ \neg A_1(b)\
              \lor\exists u\,(a<u<b\land\neg B_1(u)).
\]
The last alternative has one interior point labelled
\(\neg B_1\), true gaps, and true endpoint labels.

For \(n\geq2\), define the three branch conditions
\[
\begin{aligned}
 E_1&=\neg A_n(b)\lor K^-(\neg B_n)(b),\\
 E_2&=A_n(b)\land\mathsf{All}_{B_n}(a,b),\\
 E_3(a,c,b)&=A_n(b)\land\neg K^-(\neg B_n)(b)
                                      \land\Sigma_W(a,c,b).
\end{aligned}
\]
Each is an explicit canonical condition. Using the words
in the cut definition, form
\[
\begin{aligned}
 D_i(a,c,b)&=\mathsf N_{P_i}(a,c)\lor\mathsf N_{S_i}(c,b)
                                      &&(1\leq i<n),\\
 H_i(a,c,b)&=\mathsf N_{L_i}(a,c)\lor\mathsf N_{R_i}(c,b)
                                      &&(2\leq i<n),\\
 H_1(a,c,b)&=\mathsf N_{L_1}(a,c)\lor\mathsf{UC}_{R_1}(c,b).
\end{aligned}
\]
Here every \(\mathsf N\) call has fewer than \(n\) intervals.
The handler \(\mathsf{UC}_{R_1}\) is the already completed
uniform-final-gap construction, not a recursive call of
length \(n\). Put
\[
 M_W(a,c,b)=
       \bigwedge_{i=1}^{n-1}D_i(a,c,b)
       \land H_1(a,c,b)
       \land\bigwedge_{i=2}^{n-1}H_i(a,c,b),
\]
where an empty conjunction is true. Return the finite
canonical expansion of
\[
 E_1\ \lor\
       (E_2\land\mathsf{UC}_W(a,b))
 \ \lor\
       \exists c\bigl(a<c<b\land E_3(a,c,b)\land M_W(a,c,b)\bigr).
\]
To expand a branch, distribute every finite list and split
any condition on \((a,b)\) at the named \(c\), using all
point and gap cuts. Merge the resulting lists separately
on \((a,c)\) and \((c,b)\), and concatenate at \(c\).
All point conditions at the three actual endpoints are
conjoined there. The proof justifies the omitted extreme
gap cut \(i=n\).
\end{hbmethod}

\begin{hbtheorem}{thm:p09-gfc-interval-complement}{proved=annex:p9.gfc.interval.complement}{Full bounded relative complement in the permitted canonical basis}
The finite algorithm above terminates and satisfies
\[
       a<b\land\neg[W](a,b)
                            \longleftrightarrow\mathsf N_W(a,b)
\]
on every Dedekind-complete model with arbitrary unary
interpretations. It retains all attained and unattained
supremum cases, witness coincidences, empty intervals,
endpoint labels and shorter-prefix conditions.

For an unguarded complement, add \(a=b\) and \(b<a\)
as separate source-order cases. Finite canonical lists
are closed under relative complement by complementing
each member and applying finite conjunction. The label
class is \(\mathcal A=TL(U,K^-)\), as required by
Rabinovich's Lemma~7.8(2), with its permitted actual
endpoint labels \cite[pp.~9--11, 14]{Rabinovich2014Kamp}.
\end{hbtheorem}

\begin{hblemma}{thm:p09-gfc-ray-translation}{proved=annex:p9.gfc.ray.translation}{The complete return from canonical future rays}
A canonical ray with actual points
\(x=z_0<z_1<\cdots<z_m\), point labels \(A_i\),
gap labels \(B_1,\ldots,B_m\), and final gap label
\(B_{m+1}\) on \((z_m,+\infty)\), has the
\(\mathcal A\) expression
\[
 T_m=A_m\land G B_{m+1},\qquad
 T_i=A_i\land(B_{i+1}\mathbin{\mathsf U_s}T_{i+1})
                              \quad(i=m-1,\ldots,0).
\]
For \(m=0\) this is \(A_0\land GB_1\).
Each permitted cofinal conjunct is translated by
\(F\Psi\land GF\Psi\). Finite ray lists translate
by disjunction; their complements are zero-witness
rays with left label equal to the negated resulting
\(\mathcal A\) formula and final gap true.
The equivalences hold on every linear order and at
last points as well.
\end{hblemma}

\begin{hblemma}{thm:p09-gfc-regional-induction}{proved=annex:p9.gfc.regional.induction}{All future regional constructors in the canonical language}
Every \(\operatorname{Bet}_G(a,b)\) component on
actual \(a<b\), and every \(\operatorname{Aft}_G(a)\)
component, has an effectively constructed finite
canonical word or ray list over complete orders.
The induction is on the exact source stage
\(\operatorname{Bet}_{G,n}\). It includes positive
and negative signed existentials, finite conjunctions
with no children, empty disjunctions, and every
permitted point literal.

Together with the complete GPSS future special normal
form, this supplies a formula \(\theta\in\mathcal A\)
for every future first-order formula. Its label names
are inlined. Only actual lower endpoints and bounded
intervals or future rays occur; no formal left ray or
unbounded-from-below sentence is assumed expressible
in this grammar.
\end{hblemma}

\begin{hbdefinition}{def:p09-gfc-initial-normalizer}{Removing unguarded past limits at the initial point}
For \(\theta\in\mathcal A\), define \(g(\theta)\)
by commuting with atoms and Boolean operations,
retaining a whole strict-until subformula unchanged,
and replacing an encountered \(K^-B\) by true.
Thus the recursion stops at each until or past-limit
node. Every surviving past-limit occurrence lies
under strict until. The true replacement is fixed
by the displayed vacuous \(K^-\) convention at a
first point.

To place an arbitrary future first-order input in
the source's strict-quantifier form, first restrict
all quantifiers to the tail \([x,+\infty)\), then
use the capture-free identities
\[
\begin{aligned}
 \exists y\geq x\,\chi&\longleftrightarrow
                \chi[x/y]\lor\exists y>x\,\chi,\\
 \forall y\geq x\,\chi&\longleftrightarrow
                \chi[x/y]\land\forall y>x\,\chi.
\end{aligned}
\]
All bound variables are first made fresh from \(x\).
This explicitly connects GPSS's inclusive future
quantifiers with the strict convention of the
positive source's Definition~7.11.
\end{hbdefinition}

\begin{hbtheorem}{thm:p09-guarded-future-completeness}{proved=annex:p9.guarded.future.completeness}{The full guarded future completeness theorem}
For every future monadic first-order formula
\(\varphi(x)\), the finite normal-form, interval,
ray and normalization constructions above produce
a guarded \(\mathcal A\) formula \(g(\theta)\)
such that
\[
     \mathcal M,x\models\varphi
                 \quad\Longleftrightarrow\quad
     \mathcal M,x\models g(\theta)
\]
for every nonempty Dedekind-complete linear-order
model and every unary interpretation.
This includes real, finite, integer and natural
orders, with no assumption of density or absence
of endpoints.

Conversely every guarded \(\mathcal A\) formula
has an explicit future first-order translation
on every linear order: expand the strict until
and past-limit clauses with fresh variables,
then restrict all quantifiers to the current
tail. Thus the guarded fragment and future
first-order logic are expressively equivalent
over complete orders.

On discrete complete orders, each outermost past-limit
subformula has only Boolean and future-until ancestors,
so it is evaluated after the current point, where an
immediate predecessor makes it false. Replacing those
entire subformulas by false, without descending into
their operands, recovers the \(U\)-only future conclusion
without an exception at an initial point.

This is the positive result with the actual
\(U,K^-\) basis. It asserts no corresponding
completeness theorem over the rationals, no
\(U\)-only completeness over the reals, and no
unrestricted almost-future completeness theorem
\cite[Theorem~7.4, pp.~13--15]{Rabinovich2014Kamp}.
\end{hbtheorem}

\begin{hbexample}{def:p09-gfc-attainment-example}{A supremum may coincide with a required witness}
On \(\mathbb R\), take \(a=0,b=2\), a two-interval
word with true endpoint labels, \(B_1=\top\),
and its sole interior point label true only at 1.
If \(B_2\) fails only at 1, the failure supremum
is attained there, yet the word is true with its
interior witness equal to that supremum: the final
gap excludes the witness.
If instead \(B_2\) fails throughout \((1/2,1)\)
and holds at 1, the supremum is unattained,
\(K^-(\neg B_2)(1)\) is true, and the same
witness still realizes the word.
If failures are cofinal below 2, no last witness
strictly below 2 can satisfy its final gap.

The guarded formula
\[
          F\bigl(P\land K^-P\bigr)
\]
expresses the existence of a future marked point
approached from below by marked points. The
past-limit occurrence is evaluated strictly after
the current point. Its first-order future form
keeps only the interval after that current point,
and conjoining \(P\) at the current point gives
the counterexample in
\cref{def:p09-hr-until-table}.
\end{hbexample}

% GPSS Theorem 2.3, complete constructive separation and its exact comparison.

\subsection{Separating past and future}
\label{sec:p09-gpss-past-future-separation}

GPSS Theorem~2.3 states that every \(L(S,U)\) formula is
equivalent, on discrete complete orders, to a Boolean
combination of \(L(S)\) and \(L(U)\) formulas. The
construction here passes through the complete first-order
normal form and the already proved future translation.
The paper separately refers to Gabbay's September 1979
manuscript for a direct syntactic proof. This supplied
construction proves the stated theorem without identifying
that unacquired manuscript with the present proof
\cite[Theorem~2.3 and its referral, p.~170; reference G2, p.~173]{GabbayPnueliShelahStavi1980}.

\begin{hbdefinition}{def:p09-gpss-separation-languages}{The two pure operator fragments}
Use the source's strict infix convention:
\(A\mathbin{\mathsf U_s}B\) requires \(B\) at a
strict future witness and \(A\) on the open gap;
\(A\mathbin{\mathsf S_s}B\) uses the corresponding
strict past witness. The language \(L(S,U)\)
contains proposition letters, Boolean constants,
Boolean operations, and those two constructors.
\(L(U)\) permits only the until constructor,
and \(L(S)\) only the since constructor.

Here \emph{pure} means that the other temporal
operator is absent. It does not exclude current-point
atoms: a Boolean combination of proposition letters
belongs to both fragments. A future formula may
depend on the present and future; a past formula
may depend on the present and past.

Derived abbreviations may first be expanded:
\[
\begin{aligned}
 FA&=\top\mathbin{\mathsf U_s}A,&
 PA&=\top\mathbin{\mathsf S_s}A,\\
 GA&=\neg F\neg A,&
 HA&=\neg P\neg A,\\
 XA&=\bot\mathbin{\mathsf U_s}A,&
 YA&=\bot\mathbin{\mathsf S_s}A.
\end{aligned}
\]
The last two use an immediate successor or predecessor,
respectively, if it exists. Empty open gaps retain
their ordinary vacuous universal meaning.
\end{hbdefinition}

\begin{hbdefinition}{def:p09-gpss-separation-translation}{The full first-order translation at a named point}
For a current variable \(x\), define \(A^*(x)\)
by translating \(q_i\) to \(q_i(x)\), commuting
with Boolean operations, and using fresh bound
variables in
\[
\begin{aligned}
 (A\mathbin{\mathsf U_s}B)^*(x)
 &=\exists y\left(x<y\land B^*(y)\land
                  \forall z(x<z<y\to A^*(z))\right),\\
 (A\mathbin{\mathsf S_s}B)^*(x)
 &=\exists y\left(y<x\land B^*(y)\land
                  \forall z(y<z<x\to A^*(z))\right).
\end{aligned}
\]
Each copied constituent receives fresh bound names;
only its free current-point variable is replaced.
Logical constants may be translated as \(x=x\)
and \(\neg(x=x)\).
\end{hbdefinition}

\begin{hblemma}{thm:p09-gpss-separation-translation}{proved=annex:p9.gpss.separation.translation}{The exact temporal-to-first-order bridge}
For every linear-order model, every point \(t\),
and every \(A\in L(S,U)\),
\[
                  \mathcal M,t\models A
                    \quad\Longleftrightarrow\quad
                  \mathcal M\models A^*(x)[x\mapsto t].
\]
The translation is finite and introduces no unary
predicate absent from the input. Its source argument
order agrees with Kamp's prefix notation by
\[
       U(B,A)=A\mathbin{\mathsf U_s}B,\qquad
       S(B,A)=A\mathbin{\mathsf S_s}B.
\]
Pure \(L(U)\) formulas depend only on the labelled
tail at and after the current point; pure \(L(S)\)
formulas depend only on the labelled initial segment
at and before it.
\end{hblemma}

\begin{hbdefinition}{def:p09-gpss-separation-dual}{Typed order reversal}
For a model \(\mathcal M\), let
\(\mathcal M^{\mathrm{op}}\) have the same domain
and unary interpretations, with the order reversed.
For a first-order formula \(\rho\), define
\(\rho^{\mathrm d}\) by reversing every order
atom, retaining equality and unary atoms, and
commuting with connectives and quantifiers.
For a temporal formula define
\[
\begin{aligned}
 (A\mathbin{\mathsf U_s}B)^{\mathrm d}
       &=A^{\mathrm d}\mathbin{\mathsf S_s}B^{\mathrm d},\\
 (A\mathbin{\mathsf S_s}B)^{\mathrm d}
       &=A^{\mathrm d}\mathbin{\mathsf U_s}B^{\mathrm d},
\end{aligned}
\]
with atoms and constants fixed and Boolean operations
retained. Reversal changes the time direction,
not the interval/witness operand positions.
Both operations are involutions, up to fresh
bound-variable names and the standard spelling
of order guards.
\end{hbdefinition}

\begin{hblemma}{thm:p09-gpss-separation-dual}{proved=annex:p9.gpss.separation.dual}{Reversal preserves the exact model and regional classes}
For a first-order formula \(\rho\) and assignment \(\nu\),
\[
           \mathcal M,\nu\models\rho^{\mathrm d}
                \quad\Longleftrightarrow\quad
           \mathcal M^{\mathrm{op}},\nu\models\rho.
\]
For a temporal formula \(A\) and actual point \(t\),
\[
           \mathcal M,t\models A^{\mathrm d}
                \quad\Longleftrightarrow\quad
           \mathcal M^{\mathrm{op}},t\models A.
\]
The first-order translation commutes with
these two reversals up to bound renaming.
Linearity, discreteness and Dedekind completeness
are preserved; first and last points, and
successors and predecessors, exchange roles.

On the exact GPSS regional stages with actual
endpoint variables \(a,b\), reversing
\(\operatorname{Bet}_{G,n}(a,b)\) yields
\(\operatorname{Bet}_{G,n}(b,a)\), after
standardizing the reversed guard. Consequently
\(\operatorname{Bef}_{G,n}(x)\) and
\(\operatorname{Aft}_{G,n}(x)\) exchange roles.
Formal ray endpoints exchange only as notation;
they never become domain points or predicate
arguments. Every before component is past,
and every after component is future.
\end{hblemma}

\begin{hbdefinition}{def:p09-gpss-separation-algorithm}{The effective discrete-complete separation algorithm}
Let \(T_F^{\mathrm{dc}}\) be the actual
finite future translation supplied by
\cref{thm:p09-gpss-discrete-future-completeness}.
For a before component define
\[
        T_P^{\mathrm{dc}}(\beta)
                   =\bigl(T_F^{\mathrm{dc}}
                                      (\beta^{\mathrm d})\bigr)^{\mathrm d}.
\]
The inner reversal is first-order; the outer
one is temporal. Its input is an after
component and its output belongs to \(L(S)\).

For \(A\in L(S,U)\), perform these finite steps:
\begin{enumerate}
\item Form \(A^*(x)\).
\item Apply \cref{thm:p09-gpss-special-normal-form}
      with declared parameter list \((x)\).
      Group each resulting clause as
      \[
       \bigwedge_j\beta_{ij}(x)\land
                     C_i(x)\land
       \bigwedge_k\gamma_{ik}(x),
      \]
      where the \(\beta_{ij}\) are before
      components, \(C_i\) is a conjunction
      of current signed literals, and the
      \(\gamma_{ik}\) are after components.
      The unique one-point order prefix \(x=x\)
      is a tautology.
\item Return
      \[
       \operatorname{Sep}_{\mathrm{dc}}(A)=
       \bigvee_i\left[
        \bigwedge_j T_P^{\mathrm{dc}}(\beta_{ij})
        \land C_i
        \land\bigwedge_k T_F^{\mathrm{dc}}(\gamma_{ik})
                    \right].
      \]
\end{enumerate}
Each empty conjunction is true and the empty
outer disjunction is false. The current literal
formula uses the same proposition letters as
its first-order clause.
\end{hbdefinition}

\begin{hbtheorem}{thm:p09-gpss-past-future-separation}{proved=annex:p9.gpss.past.future.separation}{GPSS Theorem 2.3 with its complete construction}
For every \(A\in L(S,U)\), the displayed
algorithm produces a Boolean combination
of pure \(L(S)\) and pure \(L(U)\) formulas
equivalent to \(A\) at every point of every
nonempty discrete complete model, for every
unary interpretation.

This includes finite, natural and integer
orders, endpoint points and empty past or
future regions. It is not restricted to
the initial point of an \(\omega\)-model.
All calls are to supplied finite constructions;
no search for a semantically equivalent formula
is part of the algorithm
\cite[Theorem~2.3, p.~170]{GabbayPnueliShelahStavi1980}.
\end{hbtheorem}

\begin{hbdefinition}{def:p09-gpss-guarded-separation-bases}{The two guarded bases for complete-order separation}
Let \(\mathcal A_F=TL(U,K^-)\) with the
vacuous past-limit convention already fixed.
Define its dual limit by
\[
 K^+_{\mathrm v}B(t)=
       \forall r\left(t<r\to
                    \exists u(t<u<r\land B(u))\right)
       =\neg\bigl((\neg B)\mathbin{\mathsf U_s}\top\bigr)(t).
\]
The subscript distinguishes this vacuous
right limit from a nonempty right limit.
It is true at a last point.
Put \(\mathcal A_P=TL(S,K^+_{\mathrm v})\).

Let \(\mathcal F\) be the Boolean closure
of atoms, constants and \(B\mathbin{\mathsf U_s}C\)
with \(B,C\in\mathcal A_F\); let \(\mathcal P\)
be the corresponding closure using
\(B\mathbin{\mathsf S_s}C\) with
\(B,C\in\mathcal A_P\).
Equivalently every \(K^-\) in \(\mathcal F\)
lies under until, and every \(K^+_{\mathrm v}\)
in \(\mathcal P\) lies under since.

Extend temporal reversal by
\[
               (K^-B)^{\mathrm d}
                    =K^+_{\mathrm v}(B^{\mathrm d}),\qquad
               (K^+_{\mathrm v}B)^{\mathrm d}
                    =K^-(B^{\mathrm d}).
\]
The guarded future translation from
\cref{thm:p09-guarded-future-completeness}
is denoted \(T_F^{\mathrm{com}}\).
For a before component use
\[
        T_P^{\mathrm{com}}(\beta)=
           \bigl(T_F^{\mathrm{com}}
                            (\beta^{\mathrm d})\bigr)^{\mathrm d}.
\]
\end{hbdefinition}

\begin{hblemma}{thm:p09-gpss-guarded-separation-dual}{proved=annex:p9.gpss.guarded.separation.dual}{The guarded past translation and its semantic scope}
The extended reversal retains its satisfaction
identity and maps \(\mathcal F\) onto
\(\mathcal P\). Every formula in
\(\mathcal F\) is future and every formula
in \(\mathcal P\) is past, on all linear
orders. On complete orders,
\(T_P^{\mathrm{com}}(\beta)\) expresses its
before component at every point.

Each guarded formula expands finitely into
\(L(S,U)\), using the displayed right-limit
clause and
\(K^-B=\neg((\neg B)\mathbin{\mathsf S_s}\top)\).
Thus the guarded comparison introduces no
undefined temporal operation.
\end{hblemma}

\begin{hbtheorem}{thm:p09-gpss-complete-order-separation}{proved=annex:p9.gpss.complete.order.separation}{The correct complete-order comparison}
Replacing the two discrete translators in
\(\operatorname{Sep}_{\mathrm{dc}}\) by
\(T_F^{\mathrm{com}},T_P^{\mathrm{com}}\)
gives an effective separation of every
\(L(S,U)\) formula into a Boolean combination
of \(\mathcal F\) and \(\mathcal P\) formulas
over all nonempty Dedekind-complete orders.
The components are semantically future and
past, with current atoms permitted.

Their internal limit operators need not
belong to the pure \(L(U)\) or \(L(S)\)
grammar. The result uses the exact guarded
completeness theorem and its order dual;
it does not extend GPSS's pure-operator
conclusion unchanged to all complete orders
\cite[Theorem~7.4]{Rabinovich2014Kamp}.
\end{hbtheorem}

\begin{hbtheorem}{thm:p09-gpss-pure-separation-obstruction}{proved=annex:p9.gpss.pure.separation.obstruction}{Pure since/until separation fails over the real line}
The mixed formula
\[
 \Xi=
 \top\mathbin{\mathsf U_s}
       \left(P\land
          \neg\bigl((\neg P)\mathbin{\mathsf S_s}\top\bigr)\right)
\]
has no equivalent Boolean combination of
pure \(L(S)\) and pure \(L(U)\) formulas
over all complete orders, and already
has no such equivalent over the real line.

Use the models \(C_1,C_2\) from
\cref{def:p09-hr-continuous-models}, with
their block beginnings marked. At their
least points every pure-since formula
agrees, and every pure-until formula
agrees by the complete four-truth-set
calculation of \cref{def:p09-hr-until-table}.
The formula \(\Xi\) is false in \(C_1\)
and true in \(C_2\). Extending both
models by the same unmarked negative
real half-line proves the full-real-line
claim while preserving the pure-past
and pure-future agreements.
\end{hbtheorem}

\begin{hbexample}{def:p09-gpss-separation-worked}{A future point at which something is already past}
On every linear order,
\[
 F(Pq)\longleftrightarrow
       \bigl((Pq\lor q)\land F\top\bigr)
                    \lor F(q\land F\top).
\]
Here \(P\) is the strict-past modality
in the derived abbreviation \(Pq\);
\(q\) is a proposition letter.
The right side is a Boolean combination
of a pure-past formula, current atoms
and pure-future formulas.

On the two-point order \(0<1\), with
\(q\) true only at 1, both sides are
false at zero. The future occurrence
of \(q\) needs a still later point
to count as past there. At a last
point the entire expression is false,
even if \(q\) has occurred in its past.
The annex proves the identity by the
three exact positions of the past witness.
\end{hbexample}

% Rabinovich 2018 Stavi presentation: exact operators and partition foundation.

\subsection{Temporal operators for proper gaps}
\label{sec:p09-stavi-gap-temporal-logic}

GPSS announces an arbitrary-order completeness theorem using
\(U',S'\), without giving their satisfaction clauses.
The distinct later presentation used here is Rabinovich's
\emph{A Proof of Stavi's Theorem}, in its revised published
2018 version. It defines two binary gap operators and
develops partition formulas for the full past/future
basis \cite[secs.~2--5]{Rabinovich2018Stavi}.
We write these operators as \(\mathsf U_{\mathrm{gap}}\)
and \(\mathsf S_{\mathrm{gap}}\). The notation
\(\mathsf U_s,\mathsf S_s\) continues to mean ordinary
strict until and since. No literal identification with
the undefined 1980 symbols is assumed.

\begin{hbdefinition}{def:p09-stavi-proper-gaps}{Proper gaps and the point/gap distinction}
In a nonempty strict linear order \(T\), a proper gap
is a partition \(T=L\mathbin{\dot\cup}R\) into
nonempty sets with every point of \(L\) below every
point of \(R\), where \(L\) has no greatest point
and \(R\) has no least point. Identify the gap
\(g\) with its downward-closed left side \(L\).
Write \(t<g\) when \(t\in L\), and \(g<t\)
when \(t\in R\). Intervals such as \((t,g)\)
and \((g,t)\) contain only actual points of \(T\).
For gaps \(g,h\), \(g<h\) means proper inclusion
of their left sides.

The source permits downward-closed proper cuts
whose supremum, if present, belongs to the cut.
Its proper gaps are exactly the two-sided gaps
just defined. A gap is not a temporal evaluation
point, and neither operand of a modality is
evaluated at a gap.
\end{hbdefinition}

\begin{hblemma}{thm:p09-stavi-gap-cuts}{proved=annex:p9.stavi.gap.cuts}{Gap boundaries and order classes}
The two definitions of proper gap agree.
A linear order has no proper gaps exactly
when it is Dedekind complete. If \(t<g\),
there is an actual point strictly between
\(t\) and \(g\); if \(g<t\), there is one
strictly between \(g\) and \(t\).
Distinct gaps have an actual point between
them. Order reversal exchanges their left
and right sides and preserves properness.
First and last points are not proper gaps.
\end{hblemma}

\begin{hbdefinition}{def:p09-stavi-definable-gap}{The gap determined by the first failure of a label}
A \(P\)-gap succeeding \(t\) is a proper gap
\(g>t\) such that \(P\) holds throughout
\((t,g)\) and
\[
        \forall r>g\ \exists u\,
                     (g<u<r\land\neg P(u)).
\]
Define the actual temporal formula
\[
\begin{split}
 \gamma^+(P)={}&
 (P\mathbin{\mathsf U_s}P)
 \land\neg(P\mathbin{\mathsf U_s}\neg P)
 \land F\neg P\\
 &{}\land\neg\bigl(P\mathbin{\mathsf U_s}
                        (P\land\partial^+\neg P)\bigr),
\end{split}
\]
where
\[
 \partial^+B(t)=
 \forall r\left(t<r\to
                    \exists u(t<u<r\land B(u))\right)
       =\neg((\neg B)\mathbin{\mathsf U_s}\top)(t).
\]
This is the vacuous right-limit convention.
The mirror formula \(\gamma^-(P)\) uses
strict since, strict past existence and
the corresponding left limit.
\end{hbdefinition}

\begin{hblemma}{thm:p09-stavi-definable-gap}{proved=annex:p9.stavi.definable.gap}{An explicit formula for the unique succeeding gap}
\(\gamma^+(P)\) holds at \(t\) exactly
when a \(P\)-gap succeeds \(t\), and that
gap is unique. The formula uses only
ordinary until and Boolean operations.
It tests \(P\) strictly after \(t\),
not at \(t\) itself. The mirror statement
holds for \(\gamma^-\).

For a first-order spelling put
\[
 \operatorname{Run}_P(t,u)=
          t<u\land\forall v(t<v<u\to P(v)).
\]
Then \(\Gamma_P^+(t)\), the first-order
translation of \(\gamma^+(P)\), is the
conjunction of
\[
\begin{gathered}
 \exists u\,(\operatorname{Run}_P(t,u)\land P(u)),\qquad
 \neg\exists u\,(\operatorname{Run}_P(t,u)\land\neg P(u)),\\
 \exists u\,(t<u\land\neg P(u)),\qquad
 \neg\exists u\,(\operatorname{Run}_P(t,u)\land
                              P(u)\land\partial^+\neg P(u)).
\end{gathered}
\]
All bound variables in its expanded constituents
are made fresh \cite[pp.~4--5]{Rabinovich2018Stavi}.
\end{hblemma}

\begin{hbdefinition}{def:p09-stavi-modalities}{The exact binary Stavi clauses}
The formula \(P\mathbin{\mathsf U_{\mathrm{gap}}}Q\)
holds at an actual point \(t\) exactly when
a \(P\)-gap \(g\) succeeds \(t\) and there
is an actual \(d>g\) such that \(Q\) holds
throughout \((g,d)\).
No condition is imposed on \(P(t)\), \(Q(t)\)
or \(Q(d)\).

The binary past dual
\(P\mathbin{\mathsf S_{\mathrm{gap}}}Q\)
holds at \(t\) exactly when a proper gap
\(g<t\) has \(P\) throughout \((g,t)\),
\(\neg P\) occurs arbitrarily close to
the left of \(g\), and some actual
\(d<g\) has \(Q\) throughout \((d,g)\).

The full language \(\mathcal T_{\mathrm{gap}}\)
is generated by proposition letters, Boolean
constants and operations, and the four
binary constructors
\(\mathsf U_s,\mathsf S_s,
\mathsf U_{\mathrm{gap}},\mathsf S_{\mathrm{gap}}\).
These definitions are the source's
\(\mathrm{Until}^{s},\mathrm{Since}^{s}\)
presentation, with different typography
to avoid confusing Stavi's superscript
with the handbook's strict-operator subscript.
\end{hbdefinition}

\begin{hbtheorem}{thm:p09-stavi-first-order-clauses}{proved=annex:p9.stavi.first.order.clauses}{First-order truth tables and all temporal constructors}
The future gap clause is exactly the
following first-order formula:
\[
\begin{split}
 \Gamma_P^+(t)\ \land\
 \exists x_1\bigl(t<x_1\land\neg P(x_1)\land
  \forall y\bigl(t<y<x_1\land\neg P(y)\to
                    \forall z(y<z<x_1\to Q(z))\bigr)\bigr).
\end{split}
\]
Writing \(\Gamma_P^-\) for the order
reverse of \(\Gamma_P^+\), the past
gap clause is explicitly
\[
\begin{split}
 \Gamma_P^-(t)\ \land\
 \exists x_1\bigl(x_1<t\land\neg P(x_1)\land
  \forall y\bigl(x_1<y<t\land\neg P(y)\to
                    \forall z(x_1<z<y\to Q(z))\bigr)\bigr).
\end{split}
\]
This reverses every order relation
without exchanging the two operands.
Together with the ordinary strict
since/until clauses, these give an
effective translation of every
\(\mathcal T_{\mathrm{gap}}\) formula
to monadic first-order logic, valid
on every nonempty linear order.

Both gap modalities are false on
Dedekind-complete orders. The future
one is false at a last point and
the past one at a first point.
Discreteness alone does not make
them false: an incomplete discrete
order can have a proper gap.
\end{hbtheorem}

\begin{hbdefinition}{def:p09-stavi-partitions}{Finite partitions and singleton flags}
For \(k\geq1\), labels \(D_1,\ldots,D_k\)
and \(O\subseteq\{1,\ldots,k\}\), write
\[
                  \operatorname{Part}(\langle D_1,\ldots,D_k\rangle,O)(I)
\]
when \(I\) is the union of \(k\) nonempty
convex intervals \(I_1<\cdots<I_k\), every
point of \(I_i\) satisfies \(D_i\), and
\(I_i\) is a singleton for each \(i\in O\).
An index outside \(O\) may also have a
singleton cell. Empty \(I\) cannot satisfy
a positive-length partition.

Initially the labels may be quantifier-free
one-variable formulas. In the canonical
construction they are supplied finite
\(\mathcal T_{\mathrm{gap}}\) formulas,
treated as names for their whole-model
truth sets and eventually inlined.
No interval endpoint is added as a free
parameter inside a label.

Basic partition formulas use \(x=y\),
\(x<y\), and these partitions on
\([x,y]\), \([x,+\infty)\),
\((-\infty,x]\), or the whole order.
Simple formulas are their finite
positive Boolean combinations.
Open or half-open bounded intervals
are expressed by padding the excluded
endpoints with true-labelled forced
singleton cells, with every original
flag shifted to its new position.
\end{hbdefinition}

\begin{hbdefinition}{def:p09-stavi-two-block}{A two-cell transition allowing a gap boundary}
Let \(P\mathbin{\mathsf U_*}Q\) abbreviate
\[
\begin{split}
 P\land\bigl[
    &(P\mathbin{\mathsf U_s}Q)
     \lor(Q\mathbin{\mathsf U_s}Q)\\
    &{}\lor
       P\mathbin{\mathsf U_s}
                   (P\land(Q\mathbin{\mathsf U_s}Q))
     \lor(P\mathbin{\mathsf U_{\mathrm{gap}}}Q)
                                                        \bigr].
\end{split}
\]
This is a finite formula in the declared
basis, not an additional primitive.
\end{hbdefinition}

\begin{hblemma}{thm:p09-stavi-two-block}{proved=annex:p9.stavi.two.block}{Every attained or missing two-cell boundary}
At \(t\), \(P\mathbin{\mathsf U_*}Q\)
holds exactly when some actual \(b>t\)
has a partition of \([t,b]\) into
two nonempty consecutive cells, labelled
\(P,Q\), with neither cell required
to be a singleton.
The four displayed alternatives cover
a first point of the second cell,
a last point of the first cell equal
to \(t\), a later last point of the
first cell, and a proper-gap boundary.
\end{hblemma}

\begin{hbdefinition}{def:p09-stavi-prefix-recursion}{The four-case finite partition recurrence}
For a labelled tuple and flags \(O\), put
\(F_k=D_k\), and for \(i=k,\ldots,2\) put
\[
 F_{i-1}=D_{i-1}\land
 \begin{cases}
    \bot\mathbin{\mathsf U_s}F_i,&i-1\in O,\ i\in O,\\
    D_i\mathbin{\mathsf U_s}F_i,&i-1\in O,\ i\notin O,\\
    D_{i-1}\mathbin{\mathsf U_s}F_i,&i-1\notin O,\ i\in O,\\
    D_{i-1}\mathbin{\mathsf U_*}F_i,&i-1\notin O,\ i\notin O.
 \end{cases}
\]
Every occurrence of \(\mathsf U_*\) is
expanded by its formula above.
For a future-ray translation replace
only the terminal seed by
\[
 F_k=
 \begin{cases}
   D_k\land G\bot,&k\in O,\\
   D_k\land GD_k,&k\notin O,
 \end{cases}
\]
retaining the original \(D_i\) in the
recurrence's interval operands.
\end{hbdefinition}

\begin{hbtheorem}{thm:p09-stavi-partition-prefix-rays}{proved=annex:p9.stavi.partition.prefix.rays}{Complete finite-prefix and ray translations}
With the finite seed, \(F_1(t)\) is
equivalent to the existence of \(b\geq t\)
such that the given partition holds
on \([t,b]\). More generally its suffix
formula \(F_i\) holds at every point
of the \(i\)-th cell of any realizing
partition, and expresses the corresponding
finite suffix starting at that point.

With the ray seed, \(F_1(t)\) is
equivalent to the partition on
\([t,+\infty)\), including a last
point and every singleton-flag case.
Reversing the order, tuple and flags,
and dualizing its labels,
gives the full past-ray translation.
These equivalences hold on every
linear order, with all labels evaluated
in the original whole model.
\end{hbtheorem}

\begin{hbtheorem}{thm:p09-stavi-positive-partitions}{proved=annex:p9.stavi.positive.partitions}{Finite refinement, anchoring and existential assembly}
Simple partition formulas are effectively
closed under finite conjunction,
disjunction and existential quantification.
The construction refines overlapping
partitions by their nonempty cell
intersections, enumerates weak order
types of named points, and splits
cells at each named point.

Each actual endpoint is represented
by its own forced singleton when
two distinct endpoints must be
retained. Existential elimination
then glues the two partitions at
the quantified singleton and
conjoins its labels. This preserves
the strict inequalities to surviving
endpoints. Every original singleton
flag, empty intersection case and
coincident named-point case is retained.
This proves the positive closure
needed before the source's negation
proposition.
\end{hbtheorem}

\begin{hblemma}{thm:p09-stavi-basic-partition-translations}{proved=annex:p9.stavi.basic.partition.translations}{Basic first-order definitions and the one-variable return}
Every basic partition formula has an
effective first-order presentation.
For a bounded closed interval, translate
its future-ray formula to first-order
logic and restrict its quantifiers
to that interval, while keeping the
input labels as unary names until
after restriction.

Every simple partition formula with
at most one free individual variable
has an effective
\(\mathcal T_{\mathrm{gap}}\) expression.
For a whole-order partition, choose
an actual \(z\) in its first cell,
combine its one-cell past-prefix
translation with its full future-ray
translation, and assert that this
combined formula holds somewhere:
\[
                               A\lor FA\lor PA.
\]
The past one-cell flag set is
\(O\cap\{1\}\). This restriction
retains a forced singleton first
cell without using an out-of-range
flag.
\end{hblemma}

\begin{hbdefinition}{def:p09-stavi-common-gap}{The same definable gap for two labels}
Define
\[
 \operatorname{UG}(D,E)=
     \gamma^+(D)\land\gamma^+(D\land E)
              \land\neg((D\land E)\mathbin{\mathsf U_{\mathrm{gap}}}D).
\]
\end{hbdefinition}

\begin{hblemma}{thm:p09-stavi-common-gap}{proved=annex:p9.stavi.common.gap}{The exact gap-capped interval test}
\(\operatorname{UG}(D,E)\) holds at \(t\)
exactly when a \(D\)-gap \(g\) succeeds
\(t\) and \(E\) holds throughout
\((t,g)\). The formula imposes no
condition on \(E(t)\).
The proof distinguishes equality of
the \(D\)-gap and \((D\land E)\)-gap
from a strictly earlier latter gap.
\end{hblemma}

\begin{hbdefinition}{def:p09-stavi-gap-partition}{The corrected open and closed gap-partition translations}
For \(k\geq1\), define
\(\operatorname{OG}_D(\langle E_1,\ldots,E_k\rangle,O)\).
If \(k\in O\), its value is false.
Otherwise form the tuple
\[
 \left\langle
   \top,\ D\land E_1,\ldots,D\land E_{k-1},
                D\land E_k\land\operatorname{UG}(D,E_k)
 \right\rangle
\]
with flags
\(\{1\}\cup\{i+1:i\in O\}\).
Apply the finite-prefix recurrence
to this tuple, and conjoin
\(\gamma^+(D)\) at its initial
point. The first cell is the
explicit initial singleton and
does not require \(D\) there.

The closed version
\(\operatorname{CG}_D(\langle E_1,\ldots,E_k\rangle,O)\)
is \(E_1\) conjoined with the
disjunction of these permitted
alternatives:
\[
\begin{array}{c|l}
 1\notin O&
  \operatorname{OG}_D(\langle E_1,\ldots,E_k\rangle,O),\\
 k\geq2&
  \operatorname{OG}_D(\langle E_2,\ldots,E_k\rangle,
                         \{i-1:i\in O,\ i\geq2\}).
\end{array}
\]
An absent alternative contributes
false. The second alternative
makes the original first cell
the current point alone.
\end{hbdefinition}

\begin{hbtheorem}{thm:p09-stavi-gap-partition}{proved=annex:p9.stavi.gap.partition}{Full partitions ending at a proper definable gap}
\(\operatorname{OG}_D\) holds at \(t\)
exactly when a \(D\)-gap \(g\)
succeeds \(t\) and the specified
partition holds on \((t,g)\).
\(\operatorname{CG}_D\) expresses
the corresponding condition on
\([t,g)\).

A forced singleton last cell
is impossible at a proper gap.
The current point in the closed
version must satisfy its own
label, but \(D(t)\) is not
required by a succeeding gap.
Order reversal gives both
right-open and right-closed
past versions with the same
explicit flag transformations.
\end{hbtheorem}

\begin{hbexample}{def:p09-stavi-gap-examples}{Dense and discrete incomplete examples}
On \(\mathbb Q\), let \(g\) be
the cut at \(\sqrt2\), take
\(P=\{q:q<\sqrt2\}\) and
\(Q=\{q:q>\sqrt2\}\).
At \(t=0\),
\(P\mathbin{\mathsf U_{\mathrm{gap}}}Q\)
holds. Every operand is evaluated
at a rational point; the missing
irrational is only a description
of the cut. Changing \(P(0)\)
alone does not change this value.

On the discrete order
\(\mathbb Z+\mathbb Z\), let
\(P\) be the first copy and
\(Q\) the second. The cut
between the copies is proper,
and the same gap formula holds
at every point of the first
copy. Thus local predecessor
and successor existence does
not eliminate proper gaps.
On a complete order both gap
modalities are false under
every interpretation.
\end{hbexample}

\begin{proofboundary}{def:p09-stavi-completeness-continuation}{The complete arbitrary-order construction}
For every monadic first-order formula \(\varphi(x)\),
\cref{thm:p09-stavi-completeness} constructs an equivalent formula
of \(\mathcal T_{\mathrm{gap}}\) over every nonempty linear order,
following the full target of the source
\cite[Theorems~2.3 and 4.4]{Rabinovich2018Stavi}.
The defined modalities, positive partitions and temporal returns here
provide the inputs to that construction.

The complete canonical negation of open partitions is
\cref{thm:p09-stavi-open-partition-complement}. It includes the
plain labelled-point-chain complement and the bounded prefix reduction,
resolves forced endpoint singleton cells before uniform-cell extension,
and retains all actual-boundary, ordered-gap and equal-gap cases.
\Cref{thm:p09-stavi-simple-negation} supplies Boolean negation closure;
the final theorem supplies every first-order induction step.
All gap formulas use the defined 2018 operators and retain the same
named cut whenever their conjuncts refer to a common missing point.
\end{proofboundary}

% Full arbitrary-order partition negation and Stavi completeness.

\subsection{The full arbitrary-order Stavi completeness construction}
\label{sec:p09-stavi-completeness}

We complete the partition-negation argument for the explicitly
defined basis \(\mathcal T_{\mathrm{gap}}\). The labels below
are supplied finite formulas of that language, interpreted
in the original whole model. The construction works on
every nonempty strict linear order; no order-completeness,
density or discreteness assumption is introduced
\cite[Propositions~5.3--5.4 and Theorem~4.4]{Rabinovich2018Stavi}.

\begin{hbdefinition}{def:p09-stavi-completion-primitives}{The elementary interval conditions}
Throughout a bounded construction assume actual \(a<b\).
Write \(E(a,b)\) for emptiness of \((a,b)\),
\(\operatorname{Occ}_B(a,b)\) for
\(\exists z(a<z<b\land B(z))\), and
\(\operatorname{All}_B(a,b)\) for
\(\forall z(a<z<b\to B(z))\).
For a tuple with flags \(W=(\langle D_1,\ldots,D_N\rangle,O)\),
write \(\Pi_W(a,b)\) for its partition of \((a,b)\).
Other bracket choices on \(\Pi_W\) indicate
the corresponding closed or half-open interval.
Its subtuple \(\Pi_{i:j}\) uses flags
\[
                     O_{i:j}=\{h:1\leq h\leq j-i+1,\ h+i-1\in O\}.
\]
For \(i>j\), \(\Pi_{i:j}(a,b)\) is the bookkeeping
condition \(E(a,b)\), not a positive-length partition.
Its complement under \(a<b\) is
\(\operatorname{Occ}_{\top}(a,b)\).
\end{hbdefinition}

\begin{hblemma}{thm:p09-stavi-completion-primitives}{proved=annex:p9.stavi.completion.primitives}{Simple conditions and exhaustive point-or-gap boundaries}
The elementary conditions and their negations have
effective simple partition presentations:
\[
\begin{aligned}
 E(a,b)&=\operatorname{Part}(\langle\top,\top\rangle,\{1,2\})[a,b],\\
 \operatorname{All}_B(a,b)&=
       E(a,b)\lor\operatorname{Part}(\langle B\rangle,\varnothing)(a,b),\\
 \neg\operatorname{All}_B&=\operatorname{Occ}_{\neg B},\qquad
 \neg\operatorname{Occ}_B=\operatorname{All}_{\neg B}.
\end{aligned}
\]
Exact finite cardinalities and their complements are
also simple, using forced singleton cells and increasing
tuples of true-labelled points.

A nonempty occurrence set in \((a,b)\) has its lower
boundary either at \(a\), at an actual point \(c\in(a,b)\),
or at a proper gap \(g\) with \(a<g<b\).
The first case is expressed by \(\partial^+B(a)\);
the actual-point descriptor is
\[
 \operatorname{Inf}_B(a,c,b)=
     a<c<b\land\operatorname{All}_{\neg B}(a,c)
                          \land(B(c)\lor\partial^+B(c));
\]
the proper-gap case is
\(\gamma^+(\neg B)(a)\land\operatorname{Occ}_B(a,b)\).
There are exact order-dual upper-boundary formulas,
using \(\partial^-\), \(\gamma^-\) and
\[
 \operatorname{Sup}_B(a,c,b)=
     a<c<b\land\operatorname{All}_{\neg B}(c,b)
                          \land(B(c)\lor\partial^-B(c)).
\]
The point and gap possibilities are exhaustive
even when an infimum or supremum is missing from
the domain.
\end{hblemma}

\begin{hbdefinition}{def:p09-stavi-point-chain-recursion}{The plain labelled-chain complement}
Put
\[
 A_m(\mathbf P;a,b)=
       \exists x_1\cdots x_m
          (a<x_1<\cdots<x_m<b\land\bigwedge_iP_i(x_i)).
\]
Construct \(O_m\) for its negation under \(a<b\).
At \(m=1\), let \(O_1=\operatorname{All}_{\neg P_1}\).
For \(m\geq2\) retain the following alternatives:
\[
\begin{gathered}
 \operatorname{All}_{\neg P_1}(a,b)\lor
                  \operatorname{All}_{\neg P_m}(a,b),\\
 \partial^+P_1(a)\land O_{m-1}(P_2,\ldots,P_m;a,b),\\
 \partial^-P_m(b)\land O_{m-1}(P_1,\ldots,P_{m-1};a,b),\\
 \exists c\,(\operatorname{Inf}_{P_1}(a,c,b)\land
                              O_{m-1}(P_2,\ldots,P_m;c,b)),\\
 \exists c\,(\operatorname{Sup}_{P_m}(a,c,b)\land
                              O_{m-1}(P_1,\ldots,P_{m-1};a,c)).
\end{gathered}
\]
For the remaining two-gap cases put
\[
\begin{split}
 G_{\mathbf P}(a,b)={}&
 \gamma^+(\neg P_1)(a)\land\operatorname{Occ}_{P_1}(a,b)\\
 &{}\land\gamma^-(\neg P_m)(b)\land\operatorname{Occ}_{P_m}(a,b),\\
 B_{\mathbf P}(a,z,b)={}&
     a<z<b\land\operatorname{Occ}_{P_1}(a,z)
                         \land\operatorname{Occ}_{P_m}(z,b).
\end{split}
\]
Add
\[
 G_{\mathbf P}(a,b)\land
        \operatorname{Part}(\langle\neg P_1,\neg P_m\rangle,\varnothing)(a,b)
\]
and \(G_{\mathbf P}(a,b)\land\exists z(B_{\mathbf P}(a,z,b)\land C_{\mathbf P})\),
where
\[
\begin{split}
 C_{\mathbf P}={}&
 \bigwedge_{i=1}^{m-1}
   \bigl(O_i(P_1,\ldots,P_i;a,z)\lor
                O_{m-i}(P_{i+1},\ldots,P_m;z,b)\bigr)\\
 &{}\land\bigwedge_{i=2}^{m-1}
   \bigl(O_{i-1}(P_1,\ldots,P_{i-1};a,z)
          \lor\neg P_i(z)
          \lor O_{m-i}(P_{i+1},\ldots,P_m;z,b)\bigr).
\end{split}
\]
Empty conjunctions are true. Every recursive chain
has fewer labels. Compile the finite disjunction
using the positive closure of
\cref{thm:p09-stavi-positive-partitions}.
\end{hbdefinition}

\begin{hbtheorem}{thm:p09-stavi-point-chain-complement}{proved=annex:p9.stavi.point.chain.complement}{The complete plain-chain negation construction}
The recursion terminates and constructs a simple
formula \(O_m\) equivalent to \(\neg A_m\)
under \(a<b\), on every linear order.
For the two proper gaps, the first added
partition expresses that the lower gap is
at or after the upper gap; the last
alternative handles their strict opposite
order. All labels, endpoint equalities
at actual bounds, and internal witness
coincidences in the cut are retained.
This proves the source's Proposition~5.4
with the full remaining label lists.
\end{hbtheorem}

\begin{hblemma}{thm:p09-stavi-bounded-sampling}{proved=annex:p9.stavi.bounded.sampling}{The full bounded sampling lemma}
Let \(F_1,\ldots,F_k\) be the finite-seed
formulas of \cref{def:p09-stavi-prefix-recursion}.
For actual \(a_1\leq d\), the following
are equivalent:
\[
\begin{gathered}
 \exists a_2\cdots a_k\,
       (a_1<a_2<\cdots<a_k\leq d
                         \land\bigwedge_{i=1}^kF_i(a_i)),\\
 \exists b\,(a_1\leq b\leq d\land
                 \operatorname{Part}(\langle D_1,\ldots,D_k\rangle,O)[a_1,b]).
\end{gathered}
\]
The proof preserves \(D_i\land F_i\)
on every constructed cell and supplies
all four singleton-flag cases, including
each possible location of the truncation
point. It holds without an attainment
or density assumption.
\end{hblemma}

\begin{hbdefinition}{def:p09-stavi-uniform-handlers}{Independent bounded-prefix and uniform-cell handlers}
For \(W=(\mathbf D,O)\), let \(L_W(a,b)\)
be the simple complement of
\(\exists z(a<z<b\land\Pi_W(a,z])\).
Construct it by padding a true initial
singleton: use \((\top,D_1,\ldots,D_N)\)
with flags \(\{1\}\cup\{i+1:i\in O\}\)
in the finite-prefix recurrence, obtaining
\(F_1,\ldots,F_{N+1}\), and return
\[
              L_W(a,b)=\neg F_1(a)\lor O_N(F_2,\ldots,F_{N+1};a,b).
\]
Let \(R_W(a,b)\) be the exact order dual,
the complement of
\(\exists z(a<z<b\land\Pi_W[z,b))\).
The tuple, flags and all label formulas
are reversed together, and the output
is reversed back.
\end{hbdefinition}

\begin{hblemma}{thm:p09-stavi-uniform-handlers}{proved=annex:p9.stavi.uniform.handlers}{The bounded interfaces and their precise flag hypotheses}
The displayed handlers have their stated
meanings. Moreover,
\[
\begin{aligned}
 1\notin O,\quad \operatorname{All}_{D_1}(a,b)
 &\ \Longrightarrow\
             (\neg\Pi_W(a,b)\leftrightarrow R_W(a,b)),\\
 N\notin O,\quad \operatorname{All}_{D_N}(a,b)
 &\ \Longrightarrow\
             (\neg\Pi_W(a,b)\leftrightarrow L_W(a,b)).
\end{aligned}
\]
They are constructed from plain-chain negation,
not from a same-rank partition-complement
hypothesis. The unflagged endpoint-cell
conditions are essential.
\end{hblemma}

\begin{hbdefinition}{def:p09-stavi-partition-cut}{The full point cut of a partition}
For \(N\geq2\) and \(a<z<b\), define
\[
 A_i(a,z,b)=
 \begin{cases}
  \Pi_{1:i-1}(a,z)\land D_i(z)\land\Pi_{i+1:N}(z,b),&i\in O,\\
  (\Pi_{1:i-1}\lor\Pi_{1:i})(a,z)\land D_i(z)
       \land(\Pi_{i+1:N}\lor\Pi_{i:N})(z,b),&i\notin O.
 \end{cases}
\]
The empty subtuples have the \(E\) meaning
declared above. For an interior index
\(2\leq i<N\), let \(\overline A_i\)
be its explicit De Morgan complement:
replace each nonempty subtuple by its
already constructed smaller-rank complement,
replace \(E\) by \(\operatorname{Occ}_{\top}\),
and replace \(D_i\) by \(\neg D_i\).
In the unflagged case the two alternatives
on each side become a conjunction of
their complements. Put
\[
                       M_W(a,z,b)=\bigwedge_{i=2}^{N-1}\overline A_i(a,z,b).
\]
\end{hbdefinition}

\begin{hblemma}{thm:p09-stavi-partition-cut}{proved=annex:p9.stavi.partition.cut}{Cut equivalence at the same actual point}
For every fixed actual \(a<z<b\),
\[
                      \Pi_W(a,b)\longleftrightarrow\bigvee_{i=1}^N A_i(a,z,b).
\]
A flagged cut cell is exactly \(\{z\}\);
an unflagged one may have a nonempty
portion on either side, or neither.
Consequently any retained condition
\(H(a,z,b)\) can be conjoined with
the complement by replacing it at
that same \(z\) with
\(\bigwedge_i\neg A_i\).
Each interior \(\overline A_i\)
uses only ranks below \(N\).
\end{hblemma}

\begin{hbdefinition}{def:p09-stavi-partition-negation}{The full open-partition negation algorithm}
Construct \(N_W(a,b)\) under the guard
\(a<b\), by induction on \(N\).
Every returned expression is conjoined with
that guard, including each recursive call.
For \(N=1\), if \(1\notin O\), return
\(E(a,b)\lor\operatorname{Occ}_{\neg D_1}(a,b)\).
If \(1\in O\), return emptiness, at
least two interior points, or an
interior failure of \(D_1\), explicitly
\[
 E(a,b)\lor\exists u\,\exists v\,(a<u<v<b)
                      \lor\operatorname{Occ}_{\neg D_1}(a,b).
\]

For \(N\geq2\), first remove forced
endpoint cells. If \(1\in O\), return
\[
 \neg X\top(a)\lor E(a,b)\ \lor\
 \exists z\bigl(a<z<b\land E(a,z)\land
                  (\neg D_1(z)\lor N_{2:N}(z,b))\bigr).
\]
If \(1\notin O\) and \(N\in O\),
return its explicit past counterpart
\[
 \neg Y\top(b)\lor E(a,b)\ \lor\
 \exists z\bigl(a<z<b\land E(z,b)\land
                  (\neg D_N(z)\lor N_{1:N-1}(a,z))\bigr).
\]
Here \(X\top=\bot\mathbin{\mathsf U_s}\top\)
and \(Y\top=\bot\mathbin{\mathsf S_s}\top\).
All recursive partitions have fewer cells.

It remains that \(1,N\notin O\).
Include \(E(a,b)\). On a nonempty
open interval include the following
uniform and endpoint-limit alternatives:
\[
\begin{gathered}
 \operatorname{All}_{D_1}(a,b)\land R_W(a,b),\qquad
 \operatorname{All}_{D_N}(a,b)\land L_W(a,b),\\
 \partial^+\neg D_1(a)\lor\partial^-\neg D_N(b).
\end{gathered}
\]
For actual boundary points put
\[
\begin{aligned}
 I_W(a,z,b)&=\operatorname{Inf}_{\neg D_1}(a,z,b),\\
 J_W(a,z,b)&=\operatorname{Sup}_{\neg D_N}(a,z,b).
\end{aligned}
\]
Include
\[
\begin{split}
 \exists z\bigl(I_W\land M_W
  &{}\land(\neg D_1(z)\lor N_{2:N}(z,b))\\
  &{}\land[
    \neg D_N(z)\lor\operatorname{Occ}_{\neg D_N}(z,b)
         \lor(N_{1:N-1}(a,z)\land R_W(a,z))]\bigr)
\end{split}
\]
and
\[
\begin{split}
 \exists z\bigl(J_W\land M_W
  &{}\land(\neg D_N(z)\lor N_{1:N-1}(a,z))\\
  &{}\land[
    \neg D_1(z)\lor\operatorname{Occ}_{\neg D_1}(a,z)
         \lor(N_{2:N}(z,b)\land L_W(z,b))]\bigr).
\end{split}
\]
All occurrences of \(I_W,J_W,M_W\)
in these displays have the same
arguments \(a,z,b\).

For two proper failure gaps, define
\[
\begin{split}
 G_W(a,b)={}&
 \gamma^+(D_1)(a)\land\operatorname{Occ}_{\neg D_1}(a,b)\\
 &{}\land\gamma^-(D_N)(b)\land\operatorname{Occ}_{\neg D_N}(a,b),\\
 V_W(a,z,b)={}&G_W\land a<z<b
       \land\operatorname{All}_{D_1}(a,z)\land D_1(z)
       \land D_N(z)\land\operatorname{All}_{D_N}(z,b),\\
 B_W(a,z,b)={}&G_W\land a<z<b
       \land\operatorname{Occ}_{\neg D_1}(a,z)
       \land\operatorname{Occ}_{\neg D_N}(z,b).
\end{split}
\]
Include the two alternatives
\[
\begin{split}
 \exists z\bigl(V_W\land M_W
       &{}\land N_{2:N}(z,b)\land L_W(z,b)\\
       &{}\land N_{1:N-1}(a,z)\land R_W(a,z)\bigr),\\
 &\exists z\,(B_W\land M_W).
\end{split}
\]
Their meanings are respectively \(d<z<c\)
and \(c<z<d\), where \(c\) is the first
failure gap of \(D_1\) and \(d\) the last
failure gap of \(D_N\).

The final equal-gap alternative is defined
in the next display. Every output uses
finite positive assembly; no unresolved
negation of a partition remains.
\end{hbdefinition}

\begin{hbdefinition}{def:p09-stavi-equal-gap-branch}{The named common gap and the exact final-label test}
Keep the same \(W\) with unflagged
first and last cells. Define
\[
\begin{aligned}
 T_W&=D_1\mathbin{\mathsf U_s}
                   (D_1\land\operatorname{UG}(D_1,D_N)),\\
 Q_W(a,b)&=
 G_W(a,b)\land
 \operatorname{Part}(\langle D_1,D_N\rangle,\varnothing)(a,b)
                                                  \land\neg T_W(a).
\end{aligned}
\]
The labels in the common-gap test are the
actual first and final labels \(D_1,D_N\).
For \(1\leq i\leq N\) let
\[
 \lambda_i=\operatorname{OG}_{D_1}
                        (\langle D_1,\ldots,D_i\rangle,O_{1:i}),
\]
and let \(\rho_j\), \(1\leq j\leq N\),
be the past dual of the open-gap translator
for \((D_j,\ldots,D_N)\), using the
gap defined from \(b\) by \(D_N\).
Explicitly reverse the tuple, dualize
its labels, replace every flag \(\ell\geq j\)
by \(N+1-\ell\), use
\(\operatorname{OG}_{D_N^{\mathrm d}}\),
then dualize the resulting formula.

Set
\[
 Z_W(a,b)=
 \bigwedge_{i=1}^{N-1}
              (\neg\lambda_i(a)\lor\neg\rho_{i+1}(b))
 \land
 \bigwedge_{i\notin O}
              (\neg\lambda_i(a)\lor\neg\rho_i(b)),
\]
where the second conjunction ranges over
\(1\leq i\leq N\). Include
\(Q_W(a,b)\land Z_W(a,b)\) as the
last branch of \(N_W\).
\end{hbdefinition}

\begin{hblemma}{thm:p09-stavi-equal-gap-branch}{proved=annex:p9.stavi.equal.gap.branch}{The common-gap case without a point at the gap}
\(Q_W(a,b)\) holds exactly when the two
failure gaps in \(G_W\) coincide. Under
this condition the \(\lambda_i,\rho_j\)
refer to the same gap, and
\[
                              \neg\Pi_W(a,b)\longleftrightarrow Z_W(a,b).
\]
The first conjunction complements all
between-cell cuts; the second complements
all cuts through an unflagged cell.
A singleton cell cannot cross a proper
gap. No predicate is evaluated at that
missing point.
\end{hblemma}

\begin{hbtheorem}{thm:p09-stavi-open-partition-complement}{proved=annex:p9.stavi.open.partition.complement}{Full canonical negation of open partitions}
The algorithm constructs a finite simple
formula \(N_W\) satisfying
\[
                         a<b\land\neg\Pi_W(a,b)\longleftrightarrow N_W(a,b)
\]
on every nonempty linear order.
The positive-length partition is false
when \(a\geq b\), so an unguarded
complement adds those order cases.

Every recursive partition call has
strictly fewer cells. Same-rank uniform
subintervals use the independently proved
\(L_W,R_W\) handlers. The empty,
forced-singleton, actual-boundary,
ordered-gap and equal-gap cases are
exhaustive. This completes
Proposition~5.3 in the declared basis.
\end{hbtheorem}

\begin{hbtheorem}{thm:p09-stavi-simple-negation}{proved=annex:p9.stavi.simple.negation}{All simple partition formulas are closed under negation}
The open-partition result yields effective
negation of every basic and every simple
partition formula in the canonical
\(\mathcal T_{\mathrm{gap}}\) expansion.
For a closed interval, enumerate coincident
or strictly ordered endpoints, retain
their labels, and enumerate whether the
first and last cells have nonempty
remainders after removing the endpoints.
Use \(N_W\) on every resulting open
subtuple, with \(E\) for an empty tuple.

Unary ray formulas and whole-order
partition sentences use their completed
temporal translations and Boolean
negation in the canonical label language.
Finite De Morgan distribution and positive
assembly finish the construction.
\end{hbtheorem}

\begin{hbtheorem}{thm:p09-stavi-completeness}{proved=annex:p9.stavi.completeness}{The full past/future Stavi theorem on arbitrary linear orders}
Every monadic first-order formula with
one free variable has an effectively
constructed equivalent formula of
\[
              TL(\mathsf U_s,\mathsf S_s,
                         \mathsf U_{\mathrm{gap}},\mathsf S_{\mathrm{gap}})
\]
over all nonempty strict linear orders
and all unary interpretations.
Conversely, every formula of this
explicitly defined temporal language
has an equivalent monadic first-order
formula over the same model class.

The forward construction is the structural
first-order induction using simple
partition atoms, the complete negation
operation, positive existential elimination,
and the full one-variable temporal return.
All canonical label names are inlined.
The converse uses the four proved
first-order truth tables. Thus both
directions are uniform and constructive
\cite[Theorems~2.3 and 4.4]{Rabinovich2018Stavi}.
This is the later defined Stavi basis;
it does not assign an unprinted literal
definition to GPSS's \(U',S'\).
\end{hbtheorem}

\begin{hbexample}{def:p09-stavi-completion-examples}{A forced point and a missing cut}
On \(\mathbb Q\), take \(a=0,b=2\), let \(g\)
be the cut at \(\sqrt2\), and put
\(D_1(x)\equiv x<g\), \(D_2=\top\),
\(D_3(x)\equiv g<x\). These are arbitrary
unary interpretations; \(g\) is not a
parameter of the first-order language.
The two failure gaps coincide. With
no flags, three cells
\((0,1),[1,3/2],(3/2,2)\) realize the
partition. Its middle cell crosses \(g\),
so \(\lambda_2(0)\land\rho_2(2)\) holds.
If the middle cell is forced to be
a singleton, no partition exists:
placing that point below \(g\) leaves
points below \(g\) in the final cell,
and placing it above \(g\) leaves
points above \(g\) in the initial cell.
The cut point itself is absent.
The equal-gap formula handles these
two outcomes by retaining or excluding
the within-cell index \(2\).

Here the test \(T_W\) uses \(D_3\)
and is false. Substituting \(D_2=\top\)
would make it true and discard the
required equal-gap case. Independently,
on \(\mathbb R\) the two-cell word
\((\top,\top)\) with the first cell
forced to be a singleton cannot
partition \((0,1)\). The formula
\(\neg X\top(0)\) detects the obstruction;
the uniform-first handler requires
its stated unflagged hypothesis.
\end{hbexample}


\paragraph{The printed systems \(DX\) and \(DUX\).}
Page~169 says that deductive systems for \(L(X,F)\) and \(L(X,F,U)\) had been
suggested, respectively, in its references [P2] and [MY]: Amir Pnueli,
\emph{The Temporal Semantics of Concurrent Programs}, International Symposium
on Semantics of Concurrent Computation, Evian, July 1979; and T.~J. Myers and
Amir Pnueli, \emph{The Temporal Situation Until Now}, University of
Pennsylvania technical note, 1979.  It calls the respective systems \(DX\) and
\(DUX\) and says completeness had been claimed there without proofs.  These are
GPSS's reported proposal and completeness attributions; the handbook does not
import or claim independent verification of either external body here.  GPSS
supplies the outline below, so this article is not the first source-reported
statement of those systems.
For the remainder of the paper, validity means validity at every point of every
\(\omega\)-model.  The source takes \(G,X,\mathsf U_s\) as primitive and defines
\(FA\defeq\neg G\neg A\).  Its axiom schemes are
\[
\begin{array}{rcl}
(A1)&G(A\to B)&\to(GA\to GB),\\
(A2)&X\neg A&\leftrightarrow\neg XA,\\
(A3)&X(A\to B)&\to(XA\to XB),\\
(A4)&GA&\to(XA\land XGA),\\
(A5)&G(A\to XA)&\to(XA\to GA),\\[2pt]
(U1)&A\mathbin{\mathsf U_s}B&\to FB,\\
(U2)&A\mathbin{\mathsf U_s}B
 &\leftrightarrow
 XB\lor\bigl(XA\land X(A\mathbin{\mathsf U_s}B)\bigr).
\end{array}
\]
The rules admit every propositional tautology, modus ponens, and theorem-only
\(G\)-generalization.  \(DX\) consists of \((A1)\)--\((A5)\) and those three
rules for \(L(F,G,X)\); \(DUX\) adds \((U1)\)--\((U2)\) for
\(L(F,G,X,\mathsf U_s)\).  The names stand for these two presented theorem systems,
not for semantic model classes.  Theorem~3.1 states their soundness, but no
soundness proof follows.  Theorem~3.2 states weak theoremhood completeness:
\[
\begin{aligned}
 A\in L(F,G,X)\text{ is \(\omega\)-valid}
   &\Longrightarrow \vdash_{DX}A,\\
 A\in L(F,G,X,\mathsf U_s)\text{ is \(\omega\)-valid}
   &\Longrightarrow \vdash_{DUX}A.
\end{aligned}
\]

The completeness construction starts from a consistent formula \(A_0\).  A
formula set is \emph{maximally consistent for the relevant system}---\(DX\) in
the language without \(\mathsf U_s\), and \(DUX\) in the language with it---if
it is consistent in that system and no proper superset is consistent in that
system.  Let \(S\) be the
set of all such formula sets.  GPSS asserts that a standard maximal-extension
argument supplies some \(\Delta_0\in S\) with \(A_0\in\Delta_0\), but it does
not print that argument.  For \(\Delta\in S\), it defines
\[
 \Delta^+=\{A:XA\in\Delta\},
 \qquad
 \Delta<\Theta
 \Longleftrightarrow
 \forall A\,(GA\in\Delta\to A\in\Theta).
\]
Without printing the supporting derivations, the source asserts
\[
\begin{gathered}
 \forall\Delta\in S\;(\Delta^+\in S),\\
 \forall x,y,z\in S\;
 \bigl(x<y\land x<z\land y\ne z
       \to y<z\lor z<y\bigr),\\
 \forall\Delta\in S\;\forall A\;
 \bigl(FA\in\Delta\to
       \exists\Theta\in S\,
       (\Delta<\Theta\land A\in\Theta)\bigr).
\end{gathered}
\]
It concludes only that \(<\) is a linear order consistent with \(F\) when
restricted to the exact domain
\[
 \{\Delta\in S:\Delta=\Delta_0\ \text{or}\ \Delta_0<\Delta\}.
\]
It then extends \(\{A_0\}\) to a formula set \(\Gamma\), closed under
subformulas and Boolean operations, and with the additional obligations
\[
\begin{array}{lll}
 FA\in\Gamma&\Longrightarrow&XA,XFA\in\Gamma,\\
 XA\in\Gamma&\Longrightarrow&X\neg A\in\Gamma,\\
 A\mathbin{\mathsf U_s}B\in\Gamma&\Longrightarrow&
 FB,XA,XB,X(A\mathbin{\mathsf U_s}B)\in\Gamma.
\end{array}
\]
It says \(\Gamma\) can be chosen with only finitely many formulas up to logical
equivalence and then forms what it asserts is the finite quotient
\[
 \bar S=\{\Delta\cap\Gamma:\Delta\in S\}.
\]
The source does not print the argument from finite equivalence classes to the
finiteness of \(\bar S\), so that step remains an assertion rather than a
handbook-supplied representative construction.
On \(\bar S\), let \(t\mathrel{\rho_+}s\) when some \(\Delta\) has
\(t=\Delta\cap\Gamma\) and \(s=\Delta^+\cap\Gamma\); let
\(t\mathrel{\rho_<}s\) when some \(\Delta<\Theta\) represent \(t,s\).
The next lemma says \(\rho_<\) is the transitive closure of \(\rho_+\), using
\((A5)\).

Starting from \(s_0=\Delta_0\cap\Gamma\), which contains \(A_0\), the paper
threads an
infinite path \((s_n)_{n\in\mathbb N}\) such that, if a state occurs infinitely
often, each of its \(\rho_+\)-successors also occurs infinitely often.  A path
can therefore be chosen ultimately periodic inside the recurrent finite
component.  The earlier prose calls this sequence a path, which requires
\(s_n\mathrel{\rho_+}s_{n+1}\) for every \(n\); the later formal selection
condition states only the recurrence requirement and omits that adjacency
clause.  The handbook retains the intended path condition while recording the
omission.  The source then sets
\[
 Q_i=\{n:q_i\in s_n\}
\]
and asserts the truth lemma
\[
 \mathcal M,n\entails A
 \quad\Longleftrightarrow\quad
 A\in s_n
 \qquad(A\in\Gamma).
\]
This yields an \(\omega\)-model satisfying \(A_0\).  The source does not print
the maximal-extension proof, the relation lemmas, the fair-thread existence
argument, or the connective cases of the truth lemma.  Its completeness proof
is therefore the announced outline.  The decidability corollary appeals to a
bound on the finite representation but gives neither the bound nor a
complexity class \cite[pp.~171--172]{GabbayPnueliShelahStavi1980}.

\paragraph{The final convention warning and proof ownership.}
The paper closes by comparing its strict operators with the earlier inclusive
future convention
\[
 F^0A\equiv A\lor FA,
 \qquad
 G^0A\equiv A\land GA,
\]
and prints for an earlier symbol \(\mathsf U^0\)
\[
 A\mathbin{\mathsf U^0}B
 \equiv
 GA\lor B\lor\bigl(A\land(A\mathbin{\mathsf U_s}B)\bigr).
\]
The prose calls this symbol present-inclusive and universal, and the \(GA\)
disjunct permits \(B\) never to occur.  But the literal equation is not
compatible with the recursion scheme printed below.  At time \(0\) of an
\(\omega\)-model, let \(A\) be false at \(0\) and true at every later point,
and let \(B\) be false everywhere.  Strict \(GA\) makes the displayed
definition true, whereas
\[
 B\lor\bigl(A\land X(A\mathbin{\mathsf U^0}B)\bigr)
\]
is false at \(0\).  Replacing the first disjunct \(GA\) by
\(G^0A\equiv A\land GA\) would give the usual present-inclusive weak fixed
point, but GPSS does not print that replacement.  The handbook therefore
retains the two source displays as conflicting artifacts and does not adopt
\(\mathsf U^0\) as one coherent defined operator.

The source's transformed schemes are
\[
\begin{array}{rcl}
(A^0 1)&G^0(A\to B)&\to(G^0A\to G^0B),\\
(A^0 2)&X\neg A&\leftrightarrow\neg XA,\\
(A^0 3)&X(A\to B)&\to(XA\to XB),\\
(A^0 4)&G^0A&\to(A\land XG^0A),\\
(A^0 5)&G^0(A\to XA)&\to(A\to G^0A),\\[2pt]
(U^0 1)&G^0A&\to A\mathbin{\mathsf U^0}B,\\
(U^0 2)&A\mathbin{\mathsf U^0}B
 &\leftrightarrow
 B\lor\bigl(A\land X(A\mathbin{\mathsf U^0}B)\bigr).
\end{array}
\]
With propositional tautologies, modus ponens, and \(G^0\)-generalization,
\((A^0 1)\)--\((A^0 5)\) form \(D^0X\), which the source claims complete for
\(L(G^0,X)\); adding \((U^0 1)\)--\((U^0 2)\) forms \(D^0UX\), which it claims
complete for \(L(G^0,X,\mathsf U^0)\).  Page~173 prints the transformation
claim and schemes, not a separate completeness proof.  The \(D^0UX\) claim
must in addition be read through the unresolved semantic/recursion conflict
above.  These systems are not aliases for \(DX,DUX\); the paper says an
equivalent \(D^0UX\) appears in the Myers--Pnueli University of Pennsylvania
technical note just identified as its reference [MY].  It does not print the
strong-until induction rule of the Part~VIII sheet.  No exact-calculus attribution may
therefore move between these presentations without the displayed translations.

The source does not use the later name ``LTL,'' proves no PSPACE bound, and
gives no temporal-formula-to-B\"uchi translation.  Those belong to later
presentations.  The public discrete-word theorem in
Theorem~\ref{thm:p09-ltl-fo} uses the strict/inclusive bridge stated there; it
does not convert either source outline into a fully expanded proof
\cite[p.~173]{GabbayPnueliShelahStavi1980}.

\subsubsection{Sistla--Clarke's operator-sensitive complexity boundary}

\paragraph{Journal source and repaired version boundary.}
A.~Prasad Sistla and Edmund~M. Clarke,
\emph{The Complexity of Propositional Linear Temporal Logics},
\emph{Journal of the ACM} 32(3) (July 1985), pp.~733--749, DOI
\(10.1145/3828.3837\), is the controlling work.  The printed title is plural
and prints Sistla before Clarke; Clarke's Carnegie Mellon University index
uses a singular catalogue title and displays the authors in the opposite
order, so the facsimile title page controls.  That index links the complete
seventeen-page image scan of the journal article.
The first printed page carries the 1985 ACM copyright notice and conditional
copy terms, not a Creative Commons or public-domain dedication.  The
handbook independently reconstructs the definitions, proofs, and tables and
does not reproduce the scanned pages or figures.
The paper records receipt in April 1982, revision in October 1984, and
acceptance in January 1985 \cite[p.~749]{SistlaClarke1985}.  Its distinct 1982
STOC predecessor occupies pp.~159--168 and has noninterchangeable pagination
\cite[pp.~159--168]{SistlaClarke1982}.  The journal conclusion says that a
theorem used by the conference version for the regular-operator extension
contained an error and that the journal submission corrected it
\cite[p.~748]{SistlaClarke1985}.  The journal body therefore controls the
presentation below.  In particular, conference Theorem~4.6
\cite[p.~165]{SistlaClarke1982} is strengthened and renumbered as journal
Theorem~4.7 \cite[pp.~743--744]{SistlaClarke1985}, and the repaired
regular-operator argument is reported only in its journal form.

\paragraph{Source languages and inclusive temporal clauses.}
Fix the source's finite set \(\mathcal P\) of atomic propositions.  Its full
linear temporal formulas are generated by
\[
 f::=P\mid\neg f\mid(f\land f)\mid Xf\mid(f\,U\,f)\mid(f\,S\,f),
 \qquad P\in\mathcal P,
\]
with the ordinary Boolean abbreviations and
\[
 Ff:=\top\,U\,f,
 \qquad
 Gf:=\neg F\neg f.
\]
For distinct operators \(O_1,\ldots,O_k\) chosen from
\(\{X,F,G,U,S\}\), the source notation \(L(O_1,\ldots,O_k)\) names the
Boolean temporal language whose available temporal operators are restricted
to that list; it does not require every listed operator to occur in every
formula.
This is a metalinguistic language class, not an additional connective, and it
is needed because every complexity theorem below depends on the available
operator basis.

The source defines \(\operatorname{length}(f)\) as the number of symbols in
\(f\).  Its finite closure \(\operatorname{SF}(f)\) is obtained from the
syntactic subformulas of \(f\) by adjoining each subformula's
single-negation mate: for a subformula \(h\) not headed by \(\neg\), adjoin
\(\neg h\); for \(h=\neg h_0\), adjoin \(h_0\), so double negations are
eliminated.  Hence
\[
 \operatorname{card}(\operatorname{SF}(f))
 \le 2\operatorname{length}(f).
\]
This is an eliminable metalinguistic formula closure, not a temporal
operator or a set of formulas assumed true at a state.  Its negation mates
distinguish it from the ordinary subformula set, and it is needed below by
the finite certificate and the periodic truth sets.  The cardinality of this
closure is also distinct from the later count of realizable truth sets
\cite[p.~735]{SistlaClarke1985}.

A source structure is \(S=(s,\xi)\), where
\(s=(s_0,s_1,\ldots)\) is an \(\omega\)-sequence of distinct states and
\(\xi(s_i)\subseteq\mathcal P\) is the valuation at \(s_i\).  At an
interpretation \((S,s_i)\),
\[
\begin{aligned}
 S,s_i\models Xf
 &\Longleftrightarrow S,s_{i+1}\models f,\\
 S,s_i\models f_1\,U\,f_2
 &\Longleftrightarrow
   \exists k\ge i\,
   \bigl(S,s_k\models f_2\land
         \forall j\,(i\le j<k\Rightarrow S,s_j\models f_1)\bigr),\\
 S,s_i\models f_1\,S\,f_2
 &\Longleftrightarrow
   \exists k\le i\,
   \bigl(S,s_k\models f_2\land
         \forall j\,(k<j\le i\Rightarrow S,s_j\models f_1)\bigr).
\end{aligned}
\]
Thus \(U,S\), and the derived \(F\) permit the witness to be the current
point.  These are not the strict GPSS clauses immediately above.

The separate notation \(\widetilde L(F,X)\) permits only
\(\land,\lor,F,X\), with negation applied to atomic propositions alone.  Its
standing is again a restricted formula class.  It differs from \(L(F,X)\)
because arbitrary negation in the latter recovers \(G\) from \(F\); that
nonredundancy is used by the NP/PSPACE separation below
\cite[pp.~734--735, 748]{SistlaClarke1985}.

\paragraph{The two decision problems and the path quantifier.}
Satisfiability asks whether some interpretation \((S,s_i)\) satisfies the
input formula.  An \(R\)-structure is
\[
 T=(N,R,\eta),
 \qquad
 |N|<\infty,\quad
 R\subseteq N\times N,\quad
 \eta:N\longrightarrow 2^{\mathcal P},
\]
where \(R\) is total.  A path from \(\delta\in N\) is an infinite sequence
\(p_0p_1\ldots\) with \(p_0=\delta\) and \(p_iRp_{i+1}\).  It induces the
linear structure \(S_p\) by the valuation sequence
\(\eta(p_0),\eta(p_1),\ldots\).

The source's \emph{determination of truth in an \(R\)-structure} is the
existential-path problem
\[
 \text{given }T,\delta,f,\qquad
 \exists p\text{ from }\delta\quad S_p,s_0\models f\;?
\]
It is therefore the counterexample-search half of modern universal-path model
checking, not the assertion that every path satisfies \(f\).  Universal
correctness for \(f\) is obtained by asking whether the existential problem
for \(\neg f\) has no witness.  This quantifier distinction is essential in
the NP rows below.  Table~I also has a validity column.  The 1982 conference
predecessor formally defines validity at p.~160
\cite[p.~160]{SistlaClarke1982}.  The controlling journal article does not
restate that clause on pp.~734--735 and only retains the validity column on
p.~748 \cite[pp.~735, 748]{SistlaClarke1985}.  Relative to the journal, the
identical universal clause made explicit here is therefore a reconstruction:
\[
 f\text{ is valid}
 \quad\Longleftrightarrow\quad
 \forall S=(s,\xi)\ \forall i\in\mathbb N\quad S,s_i\models f.
\]
The conference wording is predecessor evidence; it does not make this a
definition literally printed in the journal.  The handbook-explicit clause
types the table's universal structure-and-point problem.

\paragraph{Why \(L(F)\) has a linear small structure.}
For \(S=(s,\xi)\), the source defines \(\operatorname{final}(S)\) to be the
unique maximal suffix \(s''\) of \(s\) in which every valuation that appears
appears infinitely often.  Writing \(s=s'\mathbin{\cdot}s''\), it sets
\[
 \operatorname{init}(S)=s',
 \qquad
 \operatorname{final}(S)=s'',
 \qquad
 \operatorname{range}(S)
 =\{\xi(s_k):s_k\text{ occurs in }s''\},
\]
and measures
\[
 \operatorname{size}(S)
 =
 \operatorname{length}(\operatorname{init}(S))
 +
 \operatorname{card}(\operatorname{range}(S)).
\]
This is a measure of a finite prefix plus recurring valuation types; it does
not make the underlying temporal sequence finite.

Lemmas~3.1--3.3 establish the pruning argument.  Lemma~3.1 says that positions
in the recurring suffix with the same valuation satisfy the same
\(L(F)\)-formulas.  Lemma~3.2 has the following exact hypotheses.  For
structures \(S=(s,\xi)\) and \(T=(t,\pi)\), put
\(m=\operatorname{length}(\operatorname{init}(S))\).  If
\[
\begin{gathered}
 m=\operatorname{length}(\operatorname{init}(T)),\qquad
 \xi(s_j)=\pi(t_j)\quad(0\le j<m),\\
 m=0\Longrightarrow\xi(s_0)=\pi(t_0),\qquad
 \operatorname{range}(S)=\operatorname{range}(T),
\end{gathered}
\]
then, for every \(f\in L(F)\),
\[
 S,s_0\models f\quad\Longleftrightarrow\quad T,t_0\models f.
\]
The separate first-valuation clause supplies the boundary case in which the
common initial length is zero.  Thus the lemma preserves the entire ordered
finite prefix but only the set, not the order, of recurring valuations.

Here a subsequence \(t\) of the source sequence is \emph{acceptable} exactly
when retaining any state token from \(\operatorname{final}(S)\) also retains
every token in \(\operatorname{final}(S)\) having the same valuation.  This is
a source-defined condition on a subsequence, not a new kind of temporal model;
ordinary subsequences need not satisfy it.

The source's measures for such a subsequence depend on the ambient structure
\(S\).  To expose that otherwise hidden parameter, the handbook writes
\(\operatorname{init}_{S}(t)\), \(\operatorname{final}_{S}(t)\),
\(\operatorname{range}_{S}(t)\), and \(\operatorname{size}_{S}(t)\) for what
the source writes without subscripts.  This is an eliminable metalinguistic
disambiguation, not a new temporal object or relation: the subscript is
necessary only because the same subsequence notation may be read against
different ambient sequences, and it differs from the source notation only by
making that dependence visible.  Exactly,
\(\operatorname{init}_{S}(t)\) is the longest prefix of \(t\) that is a
subsequence of \(\operatorname{init}(S)\); writing
\[
 t=\operatorname{init}_{S}(t)\mathbin{\cdot}
   \operatorname{final}_{S}(t)
\]
defines the remaining suffix \(\operatorname{final}_{S}(t)\), and
\[
\begin{aligned}
 \operatorname{range}_{S}(t)
 &=\{\xi(q):q\text{ occurs in }\operatorname{final}_{S}(t)\},\\
 \operatorname{size}_{S}(t)
 &=\operatorname{length}(\operatorname{init}_{S}(t))
   +\operatorname{card}(\operatorname{range}_{S}(t)).
\end{aligned}
\]
These definitions are used in the extraction bound below.  The preparatory
source paragraph also says that, if \(t_1,t_2\) are acceptable and \(v\)
contains, in source order, exactly the state tokens that occur in either one,
then \(v\) is acceptable.  It literally prints
\[
 \operatorname{size}(v)
 <
 \operatorname{size}(t_1)+\operatorname{size}(t_2).
\]
That strict form already fails when both inputs are the empty subsequence.
The following conjunction induction instead uses the proof-supported
subadditivity fact
\[
 \operatorname{size}_{S}(v)
 \le
 \operatorname{size}_{S}(t_1)+\operatorname{size}_{S}(t_2).
\]

The source measure \(\mu(f)\in\mathbb N\) is the number of occurrences of the
symbol \(F\) in \(f\).  It is a metalinguistic formula measure, not a temporal
operator or the total length of \(f\); its distinct count is required by the
extraction bound.  The exact preservation content of Lemma~3.3 is stronger
than ordinary thinning.  If \(t\) is an acceptable subsequence of \(s\),
\(f\in L(F,G)\) is in negation-normal form, and \(f\) holds at every
represented position of \(t\), then there is an acceptable subsequence \(u\)
of \(s\) satisfying the lemma's literally printed, unsubscripted bound
\[
 \operatorname{size}(u)<\mu(f),
\]
which in the handbook's explicit notation is
\(\operatorname{size}_{S}(u)<\mu(f)\),
and the following universal preservation clause.  For every intermediate
sequence \(w\) such that \(u\) is a subsequence of \(w\) and \(w\) is a
subsequence of \(s\), let \(W=(w,\varphi)\), where \(\varphi\) is the
restriction of \(\xi\) to the states in \(w\).  Then
\[
 w_r\text{ occurs in }t\quad\Longrightarrow\quad W,w_r\models f
 \qquad\text{for every represented position }w_r\text{ of }w.
\]
Thus selected states may be reinserted into any such intermediate
supersequence without losing truth at the surviving positions inherited from
\(t\).  This quantified stability is what the later small-structure argument
uses.

The preparatory union sentence has the strict-sign defect just disclosed.
Lemma~3.3 has the same kind of conflict: its displayed statement uses the
strict inequality \(\operatorname{size}(u)<\mu(f)\), meaning
\(\operatorname{size}_{S}(u)<\mu(f)\) with the ambient parameter exposed.
In the atomic basis, however, the proof takes the empty
subsequence, so both sides are zero; every induction case then derives
\(\operatorname{size}_{S}(u)\le\mu(f)\), and Theorem~3.4 uses that non-strict
bound.  The handbook therefore records the printed strict inequality as a
strict/non-strict defect and uses the proof-supported non-strict statement
rather than changing it silently.  The evidence identifies the mathematical
conflict but does not establish whether its cause was typesetting or authorial
revision.

Theorem~3.4, whose source proof is expressly labelled
a sketch, concludes
\[
 f\in L(F)\text{ satisfiable}
 \quad\Longrightarrow\quad
 \exists S\,
 \bigl(\operatorname{size}(S)\le\operatorname{length}(f)
       \land S,s_0\models f\bigr).
\]

Theorem~3.5 then proves that both satisfiability and existential truth in an
\(R\)-structure are NP-complete for \(L(F)\).  For \(R\)-structure hardness,
the paper reduces a 3-CNF formula to a chain of binary diamonds.  Choosing the
upper or lower branch at the \(j\)-th diamond chooses \(x_j\) or
\(\neg x_j\); that branch is labelled by every clause it satisfies, and
\[
 FC_1\land\cdots\land FC_m
\]
holds on some path exactly when the original clauses have a satisfying
assignment.  For the \(R\)-structure upper bound, a nondeterministic machine
uses the source's letter \(q_0\) for the supplied start state by setting
\(q_0:=\delta\).  This is an eliminable metalinguistic renaming, not a second
graph node: it is needed only to align the certificate with the article's
notation and is used in both the prefix and acceptance clauses.  The machine
guesses a finite path prefix
\[
 p'=(x_0,\ldots,x_r),\qquad x_0=q_0,
\]
with
\[
 \operatorname{length}(p')
 \le \operatorname{length}(f)\operatorname{card}(N)
\]
and a subset \(C\subseteq N\).  It checks that \(p'\) is a path from \(q_0\),
that the graph on \(C\) with edges \(R\cap(C\times C)\) is strongly
connected, and that an edge leads from the last node of \(p'\) into \(C\).

The article next assigns one \(\operatorname{label}(x)\) to each graph node in
the prefix or in \(C\), but its \(F\)-clause refers to ``that occurrence'' of
\(x\) in the prefix.  A repeated node makes those instructions incompatible.
For example, let
\[
 \begin{gathered}
  N=\{x,u,c\},\qquad
  R=\{(x,u),(x,c),(u,x),(c,c)\},\\
  \eta(u)=\{a\},\qquad \eta(x)=\{b\},\qquad \eta(c)=\varnothing .
 \end{gathered}
\]
The path \(x,u,x,c,c,\ldots\) satisfies
\[
 Fa\land F(b\land\neg Fa).
\]
At its first occurrence of \(x\), \(Fa\) is true; at its second occurrence,
\(\neg Fa\) is true.  One node-indexed label cannot record both facts.

The handbook therefore makes the certificate occurrence-sensitive.  For every
prefix position \(0\le i\le r\), initialize a finite metalinguistic truth set
\(\mathcal L_i=\varnothing\); for every recurrent node \(x\in C\), initialize
a separate state summary \(\mathcal L_C(x)=\varnothing\).  Process
\(g\in\operatorname{SF}(f)\) in increasing order of
\(\operatorname{length}(g)\).  If \(\mathcal L\) denotes any one of these
sets, use the Boolean rules below within that set; use the two distinct
eventuality rules for prefix occurrences and recurrent nodes:
\[
\begin{array}{rcl}
 P\in\mathcal L_i
 &\Longleftrightarrow& P\in\eta(x_i),\\
 P\in\mathcal L_C(x)
 &\Longleftrightarrow& P\in\eta(x),\\
 \neg g_1\in\mathcal L
 &\Longleftrightarrow& g_1\notin\mathcal L,\\
 g_1\land g_2\in\mathcal L
 &\Longleftrightarrow& g_1,g_2\in\mathcal L,\\
 Fg_1\in\mathcal L_i
 &\Longleftrightarrow&
 \bigl(\exists j\;(i\le j\le r\land g_1\in\mathcal L_j)\bigr)
 \lor
 \bigl(\exists y\in C\;g_1\in\mathcal L_C(y)\bigr),\\
 Fg_1\in\mathcal L_C(x)
 &\Longleftrightarrow&
 \exists y\in C\;g_1\in\mathcal L_C(y).
\end{array}
\]
The paper's rule says that \(y\) is ``after'' \(x\).  Here that word is read
non-strictly.  This is forced by the inclusive semantics of \(F\) on
pp.~734--735 and by the proof's same use of ``after'' on p.~737, where the
current point may be the witness.  A strict reading would reject
\[
 N=\{q_0,d\},\qquad R=\{(q_0,d),(d,d)\},\qquad
 \eta(q_0)=\{A\},\quad\eta(d)=\varnothing
\]
for the true formula \(A\land FA\).  The displayed at-or-after clause is the
handbook's occurrence-indexed source-context reconstruction, not a silent
replacement of inclusive \(F\); it is separate from the directed-cycle
correction below.
The machine accepts exactly when
\(f\in\mathcal L_0\).  The sets \(\mathcal L_i\) and \(\mathcal L_C(x)\) are
eliminable finite metalinguistic data used only to check the guessed witness;
their separate standings are necessary because a prefix occurrence and a
recurrent graph node have different future domains.  With the directed-cycle
condition below, one can choose an infinite suffix that visits every node of
\(C\) infinitely often.  Structural induction on \(g\) then proves the two
displayed eventuality clauses and the acceptance condition.  These sets are
not the semantic source truth sets
\([s_i]_{S,f}\) introduced in Section~4.  There are
\(r+1+\operatorname{card}(C)\) such sets, each contained in
\(\operatorname{SF}(f)\); the occurrence-sensitive repair therefore preserves
the polynomial certificate bound.

There is a source-level gap in the graph check.  Under the ordinary reflexive
convention for reachability, a singleton induced subgraph is strongly
connected even when it has no self-loop.  Let
\[
 N=\{q_0,c,d\},\qquad
 R=\{(q_0,c),(c,d),(d,d)\},
\]
and let \(P\) hold at \(q_0,c\) but not at \(d\).  The only infinite path
from \(q_0\) falsifies \(GP=\neg F\neg P\), yet the literal checks accept the
guess \(p'=(q_0)\), \(C=\{c\}\): the singleton passes reflexive strong
connectivity, the entry edge exists, and the rules place
\(\neg F\neg P\) in \(\mathcal L_0\).  The paper does not add a convention
that repairs this case.  The handbook therefore adds, as its own necessary
correctness condition, that \(R\cap(C\times C)\) contain a directed cycle.
For singleton \(C\) this requires a self-loop; for a larger finite strongly
connected \(C\) it is automatic.  With that addition, the witness follows
\(p'\), takes the connecting edge, and thereafter can visit every node of
\(C\) infinitely often.
Satisfiability uses the linear-size
\(\operatorname{init}/\operatorname{range}\) witness.  For
\(f\in\widetilde L(F,X)\) with \(S,s_i\models f\), Lemma~3.6 produces a finite
subsequence \(u\) of \(s\) which contains \(s_i\) and satisfies
\(\operatorname{length}(u)\le\operatorname{length}(f)\).  Its guarantee again
ranges over every intermediate supersequence: whenever \(u\) is a subsequence
of \(v\), \(v\) is a subsequence of \(s\), and \(V=(v,\varphi)\) restricts
\(\xi\) to the states in \(v\), one has \(V,s_i\models f\).  Theorem~3.7 uses
this lemma to prove the same two NP-completeness results for
\(\widetilde L(F,X)\)
\cite[pp.~735--739]{SistlaClarke1985}.

\paragraph{The PSPACE proof chain for the richer bases.}
Journal Theorem~4.1 states that satisfiability and existential truth in an
\(R\)-structure are PSPACE-complete for each of
\[
 L(F,X),\qquad L(U),\qquad L(U,S,X).
\]
The proof separates the two directions.  Lemma~4.3 gives the reduction
explicitly.  Let \(T=(N,R,\eta)\), let \(f\in L(U,S,X)\), let
\(\mathcal P_2\) be the atoms occurring in \(f\), and introduce the disjoint
set \(\mathcal P_1=\{P_x:x\in N\}\).  For each \(x\in N\), put
\[
\begin{aligned}
 g_1^x&:=\bigwedge\{Q:Q\in\eta(x)\},&
 g_2^x&:=\bigvee\{Q:Q\in\mathcal P_2\setminus\eta(x)\},\\
 g_3^x&:=\bigvee\{P_y:(x,y)\in R\},&
 f_x&:=P_x\to(g_1^x\land\neg g_2^x\land Xg_3^x),\\
 h_1&:=\bigvee_{x\in N}P_x,&
 h_2&:=\bigwedge_{x\in N}f_x,
\end{aligned}
\]
with empty conjunction \(\top\) and empty disjunction \(\bot\).  Let \(h_3\)
be the formula the paper describes as asserting that exactly one member of
\(\mathcal P_1\) is true at each point.  On that description \(h_1\) repeats
the at-least-one part of \(h_3\), but both are retained because the printed
reduction explicitly includes both.  Finally set
\[
 f':=G(h_1\land h_2\land h_3),
 \qquad
 f'':=f'\land f\land P_q
\]
for the supplied start node \(q\).  Then \(f''\) is satisfiable exactly when
some path from \(q\) in \(T\) satisfies \(f\).  The displayed
\(g_i^x,f_x,h_i,f',f''\) are eliminable metalinguistic names for formulas in
this one reduction, not new temporal operators: the names are needed to keep
the original labelling, unique graph-state marker, and successor constraint
separate, and they are used only to assemble \(f''\).

If \(f\in L(F,X)\), the same construction remains in \(L(F,X)\).  For the
pure-\(U\) basis, the paper forms the same outer constraint but replaces each
\(Xg_3^x\) by
\[
 \begin{cases}
  P_x\,U\,g_3^x,
    &(x,x)\notin R,\\
  G(P_x)\lor\bigl[P_x\,U\,(g_3^x\land\neg P_x)\bigr],
    &(x,x)\in R.
 \end{cases}
\]
The first case permits finitely many repetitions of \(P_x\) before a legal
successor; Lemma~4.2 proves that such finite stuttering preserves every
\(L(U)\) formula.  The second case also permits the self-loop behaviour of
remaining at \(x\) forever.  The paper expressly does not claim Lemma~4.2 for
a language containing \(X\).

Lemma~4.4 proves PSPACE-hardness from acceptance by a polynomial-space
deterministic Turing machine.  The journal assumes that after the machine
enters its accepting or rejecting state it stays there forever without changing
the tape; this turns a halting computation into the infinite history required
by temporal semantics.  Its finite graph is a cyclic chain of \(S(n)\)
diamonds.  Each diamond represents one tape-cell position, and a complete
subpath from \(BI\) through the entire diamond chain to \(EI\) selects one
symbol at every cell and therefore represents one instantaneous description;
repeated traversals represent a sequence of descriptions.  In
\(L(F,X)\), the journal explicitly uses \(2S(n)\) nested next operators to
compare corresponding tape cells in successive descriptions
\cite[p.~741]{SistlaClarke1985}.  The conference predecessor printed \(S(n)\)
\cite[p.~164]{SistlaClarke1982}, so the journal count controls here.  A
separate \(U\)-only encoding marks cell positions and the boundaries of
descriptions.
This is a computation-history reduction, not the small \(L(F)\) diamond
construction above
\cite[pp.~741--742]{SistlaClarke1985}.

For the upper bound, use the source closure \(\operatorname{SF}(f)\) defined
above; its double-negation convention and
\(2\operatorname{length}(f)\) cardinality bound remain in force.  For a
position \(s_i\), the source set
\[
 [s_i]_{S,f}:=\{g\in\operatorname{SF}(f):S,s_i\models g\}
\]
stands only for the formulas in that closure true at the position.  Printed
p.~742 says that the number of such source truth sets is at most
\(2^{\operatorname{length}(f)}\), not
\(2^{2\operatorname{length}(f)}\): each is determined by the truth values of
the actual subformulas, because the single-negation mates add no independent
choices.  The proof of Theorem~4.7 uses the same count for the actual
subformulas of \(g\), obtaining \(2^{\operatorname{length}(g)}\).  The closure
cardinality and the realizable-truth-set count are therefore different
statements even though both are controlled by formula length.  If
\(g_1\,U\,g_2\in[s_i]_{S,f}\) and \(j\ge i\), the source says that this
\(U\)-formula is \emph{fulfilled before} \(s_j\) exactly when
\[
 \exists \ell\,(i\le \ell\le j\land S,s_\ell\models g_2).
\]
This is a source-defined finite witness obligation, not another temporal
connective; unlike membership of the until formula in the truth set, it
records that the right-hand witness appears by the proposed loop boundary.
Lemma~4.5 gives the deletion map exactly.  If \(i<j\) and
\([s_i]_{S,f}=[s_j]_{S,f}\), delete \(s_i,\ldots,s_{j-1}\) but retain
\(s_j\):
\[
 s'=(s_0,\ldots,s_{i-1},s_j,s_{j+1},\ldots),
 \qquad S'=(s',\xi'),
\]
where \(\xi'\) is \(\xi\) restricted to the state tokens retained in \(s'\).
Then every state token \(s_k\) occurring in both sequences has the same
source truth set:
\[
 [s_k]_{S,f}=[s_k]_{S',f}.
\]
The article gives only the statement that this follows by induction on
formula length.

Lemma~4.6 supplies both claims needed to repeat a reduced segment.  Suppose
\(p>0\),
\[
 [s_i]_{S,f}=[s_{i+p}]_{S,f},
\]
and every \(U\)-formula in \([s_i]_{S,f}\) is fulfilled before
\(s_{i+p}\).  Let \(S'=(s',\xi')\) be ultimately periodic with starting
index \(i\) and period \(p\), and suppose
\[
 \xi(s_k)=\xi'(s'_k)\qquad(k<i+p).
\]
Then, for every \(g\in\operatorname{SF}(f)\),
\[
\begin{array}{ll}
 \text{(a)}&
 \forall k<i+p\quad
 S,s_k\models g\ \Longleftrightarrow\ S',s'_k\models g,\\[2pt]
 \text{(b)}&
 \forall k\ge i\quad
 S',s'_k\models g\ \Longleftrightarrow\ S',s'_{k+p}\models g.
\end{array}
\]
Thus (a) transfers every subformula through the retained prefix and loop
block, while (b) proves subformula truth periodicity throughout the repeated
tail; equality of valuations alone was not substituted for either conclusion.

For \(i\in\mathbb N\) and \(p\in\mathbb N_{>0}\), the source defines a
structure \(S=(s,\xi)\) to be \emph{ultimately periodic with starting index
\(i\) and period \(p\)} exactly when
\[
 \forall k\ge i\quad \xi(s_k)=\xi(s_{k+p}).
\]
This is a condition on the structure's complete propositional valuations; the
states themselves remain distinct.  It is not equality of the restricted
source truth sets \([s_k]_{S,f}\), and neither condition by itself discharges a
pending until obligation.  The distinction is required because
Theorem~4.7 uses all three conditions separately.  The finite witness consists
of the valuation prefix before \(i\) and the nonempty block of length \(p\),
whose infinite repetition supplies the structure's valuation sequence.

Journal Theorem~4.7 is the resulting ultimately periodic model theorem:
\[
\begin{gathered}
 f\in L(U,S,X)\text{ is satisfiable}\\
 \Longleftrightarrow\\
 f\text{ is satisfiable in an ultimately periodic structure with}\\
 i\le 2^{1+\operatorname{length}(f)}
 \quad\text{and}\quad
 p\le 4^{1+\operatorname{length}(f)},
\end{gathered}
\]
where \(i\in\mathbb N\) is the starting index and
\(p\in\mathbb N_{>0}\) the period.  In addition to valuation periodicity, the
theorem requires
\[
 [s_k]_{S,f}=[s_{k+p}]_{S,f}\qquad(k\ge i),
\]
and every \(U\)-formula in \([s_k]_{S,f}\) to be fulfilled
before \(s_{k+p}\), for every \(k\ge i\).  Periodic equality alone is not the
claim.  Because a formula using \(S\) may be true only away from the beginning
of a word, the proof applies the construction to \(Ff\).

The PSPACE upper bound guesses the truth sets along the exponentially bounded
prefix and period without storing the whole word.  Write \(A,B\subseteq
\operatorname{SF}(f)\) for two consecutively guessed sets, intended for the
present and next positions.  The letters are eliminable metalinguistic
bookkeeping, needed only to distinguish the two sides of the one-step checks;
unlike \([s_i]_{S,f}\), they are guesses whose realizability is being tested.
For every relevant subformula the journal requires
\[
\begin{array}{rcl}
 Xf_1\in A
 &\Longleftrightarrow& f_1\in B,\\
 f_1\,U\,f_2\in A
 &\Longleftrightarrow&
 f_2\in A\ \lor\
 (f_1\in A\land f_1\,U\,f_2\in B),\\
 f_1\,S\,f_2\in B
 &\Longleftrightarrow&
 f_2\in B\ \lor\
 (f_1\in B\land f_1\,S\,f_2\in A),\\
 f_1\land f_2\in A
 &\Longleftrightarrow& f_1,f_2\in A,\\
 \neg f_1\in A
 &\Longleftrightarrow& f_1\notin A.
\end{array}
\]
At the initial position it separately requires
\[
 f_1\,S\,f_2\in A_0\quad\Longleftrightarrow\quad f_2\in A_0.
\]
When the guessed prefix ends, the machine saves the loop-entry set; when the
period ends, it uses that saved set as the next set, closing the loop.  For
every \(f_1\,U\,f_2\) in the saved loop-entry set, it also requires \(f_2\)
to occur in some guessed set within that period.  At every step it stores only
the present, next, and loop-entry sets plus the two bounded counters.  The
nondeterministic procedure therefore uses linear space in the formula length;
Savitch's theorem supplies the deterministic polynomial-space decision procedure
\cite[pp.~742--745]{SistlaClarke1985}.

The source-proof boundary is uneven: Lemmas~4.2 and~4.5 are accompanied only
by an appeal to induction; Lemma~4.4 prints the computation reduction but
calls its \(X\)-elimination into \(L(U)\) a brief sketch; Lemma~4.6 and
Theorem~4.7 print the main inductions while abbreviating routine cases; and
the type-guessing proof of Theorem~4.1 is printed after Theorem~4.7.  These
are not an automata-theoretic proof and are not upgraded to more detail than
the journal supplies.

\paragraph{The repaired regular-operator extension.}
Section~5 begins with regular right-linear grammars.  Such a grammar has
terminals \(a_1,\ldots,a_n\), nonterminals
\(N_1,\ldots,N_m\), and productions \(N\to aM\) or \(N\to a\), where
\(a\) is a finite word of terminals.  The paper's later consistency rule
handles only productions that emit a terminal and gives no empty-word case.
Accordingly, this presentation requires
\(a\in\{a_1,\ldots,a_n\}^{+}\), the set of nonempty finite terminal words.
This is a handbook restriction needed to make the printed proof rules coherent,
not a restriction attributed to the source; the epsilon-admitting extension is
not established here.  The source assumes that different grammars use disjoint
nonterminal names, so a displayed \(N_j\) determines its grammar.  A production
emitting \(r\ge1\) terminals can be replaced by a chain of \(r\)
single-terminal productions using fresh nonterminals.  This preserves the
generated terminal strings and increases the grammar only linearly in its
encoded production length, giving the one-terminal form used below.
It calls \(\mathsf{ETL}\) the least formula class containing the atomic
propositions and closed under Boolean negation and conjunction, the inclusive
operators \(X,U,S\) defined above, and the following rule: for every such
grammar, every nonterminal \(N_j\), and formulas
\(f_1,\ldots,f_n\in\mathsf{ETL}\), the expression
\(N_j(f_1,\ldots,f_n)\) is in \(\mathsf{ETL}\).  Disjunction, implication,
\(F\), and \(G\) remain the earlier abbreviations.  Thus \(\mathsf{ETL}\)
names a formula class extending \(L(U,S,X)\), not a grammar, automaton, or
proof system; the grammar-indexed formation rule is its distinction from the
ordinary temporal fragments and is used in Theorems~5.3--5.4 and Table~I.
Throughout this paragraph, \(n\) is both the number of grammar terminals and
the arity of its formula constructor.  The paper later uses a fresh local
letter \(q\) for that same generic arity; the handbook alpha-renames it to
\(n\) throughout.  This is an eliminable metalinguistic renaming, not a new
restriction: it is needed to avoid collision with the earlier graph start
node and is used in the constructor, \(\Gamma\), and Lemmas~5.1--5.2 and
Theorem~5.4.
For formulas \(f_1,\ldots,f_n\), the expression
\[
 N_j(f_1,\ldots,f_n)
\]
holds at \(s_k\) exactly when the grammar generates from \(N_j\) a finite or
infinite terminal sequence \(a_{i_0}a_{i_1}\ldots\) and the corresponding
argument formula holds at every represented position:
\[
 S,s_{k+\ell}\models f_{i_\ell}.
\]
For a finite generated word of length \(r\), this condition ranges over
\(0\le\ell<r\); for an infinite generated word, it ranges over every
\(\ell\ge0\).  The article literally says finite or infinite and then writes
a universal \(\ell\ge0\); its later Lemma~5.1 treats finite witnesses
separately, so the finite-domain restriction above is the handbook's minimal
typing of the two source cases, not an additional source theorem.  The
notation stands for a grammar-indexed temporal formula constructor and differs
from ordinary \(U\): its production graph can impose a regular right-linear
pattern.

The printed example has a second conflict.  For
\(N_0\to a_1a_2N_0\), the generated sequence is
\(a_1a_2a_1a_2\ldots\), and the article says that
\(N_0(\top,P)\) requires \(P\) at all even states.  Under its own zero-based
clause above, however, \(\top\) is checked at offsets \(0,2,\ldots\) and
\(P\) at offsets \(1,3,\ldots\).  Thus the displayed constructor enforces
\(P\) at odd offsets.  Replacing the word \emph{even} by \emph{odd}, swapping
the two arguments, or reindexing positions from one would each remove the
mismatch, but the article specifies none of them; no later result here depends
on the parity sentence.

The journal also enlarges \(\operatorname{SF}\) for these constructors.  If
\(N_j(f_1,\ldots,f_n)\) occurs, its closure contains the argument closures and
\(N_r(f_1,\ldots,f_n)\) for every nonterminal \(N_r\) of the same grammar,
and is then closed under single negation.  This is the closure used by the
periodic bound and by the encoded-input estimate below.

For \(S,s_j\) and \(N_i(f_1,\ldots,f_n)\), the source's labelled tree
\(\Gamma(S,j,N_i)\) has a root labelled by that headed tuple.  A node at depth
\(\ell\) labelled by \(N_k(f_1,\ldots,f_n)\) has a child labelled by
\(N_r(f_1,\ldots,f_n)\) for each continuation production
\(N_k\to a_pN_r\) whose selected argument satisfies
\(S,s_{j+\ell}\models f_p\).  This finitely branching tree records possible
continuations; it is not a temporal model, an automaton, or an object of a
proof calculus, and no additional technical standing is imposed.
Unlike a finite terminal production, it records possible continuing
derivations.  Lemma~5.1 says that
\(\neg N_i(f_1,\ldots,f_n)\) holds exactly when no finite generated word
satisfies the operator and \(\Gamma(S,j,N_i)\) is finite.

If that negated formula holds and the finite tree has maximum depth \(d\), the
article calls the obligation \emph{fulfilled before} \(s_m\) for every
\(m\ge j+d\).  This is the finite-tree termination obligation, distinct from
the until-witness obligation above.  Lemma~5.2 takes \(S,i,p,f,S'\) with all
the hypotheses of Lemma~4.6 and adds
\[
 \neg N_u(f_1,\ldots,f_n)\in[s_i]_{S,f}
 \quad\Longrightarrow\quad
 \neg N_u(f_1,\ldots,f_n)\text{ is fulfilled before }s_{i+p}
\]
for every such subformula.  It then inherits both conclusions, for every
\(g\in\operatorname{SF}(f)\):
\[
\begin{array}{ll}
 \forall k<i+p&
 S,s_k\models g\ \Longleftrightarrow\ S',s'_k\models g,\\
 \forall k\ge i&
 S',s'_k\models g\ \Longleftrightarrow\ S',s'_{k+p}\models g.
\end{array}
\]
The paper prints the new positive and negated grammar-operator induction
cases; the other cases are those of Lemma~4.6.  Theorem~5.3 then gives
a periodic model with
\[
 l\le 2^{1+\operatorname{card}(\operatorname{SF}(f))},
 \qquad
 p\le c^{\operatorname{card}(\operatorname{SF}(f))}
\]
for some fixed constant \(c\).  For every \(k\ge l\), the truth set at \(k\)
equals the one at \(k+p\), and every \(U\)-formula and negated grammar formula
in that truth set is fulfilled before \(s_{k+p}\).
Theorem~5.4 proves satisfiability and existential \(R\)-structure truth
PSPACE-complete when the grammars are encoded as part of the input.  In this
reconstruction the claim has the nonempty-emission grammar domain fixed above;
the source's unhandled epsilon extension is not included.  At this point the
paper reuses \(n\) for total encoded input length.  Because the handbook has
reserved \(n\) for terminal count and constructor arity, write
\(\ell_{\mathrm{in}}\) for that independent source variable:
\(\ell_{\mathrm{in}}\) is exactly the length of the complete encoded
formula-and-grammar input.  This is an eliminable metalinguistic size
parameter, neither the constructor arity nor the fragment notation
\(L(O_1,\ldots,O_k)\); the quantities differ, for example, when a grammar has
\(n\) terminals but the full encoding has another length.  The distinction is
needed only for the space estimate, where the paper's bound becomes
\(\operatorname{card}(\operatorname{SF}(f))\le2\ell_{\mathrm{in}}\).
With \(A\) and \(B\)
again the current and next guessed truth sets, its positive consistency rule
is exactly
\[
\begin{split}
 N_j(f_1,\ldots,f_n)\in A
 \quad\Longleftrightarrow\quad
 &\exists k\,\exists\ell\,
   \bigl((N_j\to a_kN_\ell)\text{ is a production}\ \land\
         f_k\in A\ \land\ N_\ell(f_1,\ldots,f_n)\in B\bigr)\\
 &{}\lor
 \exists k\,
   \bigl((N_j\to a_k)\text{ is a production}\ \land\ f_k\in A\bigr),
\end{split}
\]
For each \(\neg N_j(f_1,\ldots,f_n)\) in the saved loop-entry set, the paper
also stores
\(\varphi(N_j(f_1,\ldots,f_n))\), initialized to \(\{N_j\}\), and updates it
by
\[
\begin{split}
 \varphi_{\mathrm{next}}(N_j(f_1,\ldots,f_n))
 =\{N_\ell:\;&\exists N_r\,\exists k\,
 [(N_r\to a_kN_\ell)\text{ is a production}\ \land\\
 & N_r\in\varphi_{\mathrm{present}}(N_j(f_1,\ldots,f_n))
   \ \land\ f_k\in A]\}.
\end{split}
\]
Here \(\varphi\) is an eliminable metalinguistic finite set of grammar
nonterminals, not the structure valuation and not the tree \(\Gamma\).  It is
needed to distinguish successive possible continuation layers, differs from
\(\Gamma\) by recording only the current layer, and is used only for the
finite-tree check.  The machine requires this set to become empty at some
point within the guessed period.  Together with the positive consistency
rule and Lemma~5.1, that condition rules out both finite and infinite
generated witnesses for the negated formula.  These are the journal's
replacement proof steps after its stated correction
\cite[pp.~745--748]{SistlaClarke1985}.  The article does
not isolate the particular erroneous inference in the conference proof, so the
handbook does not identify any one of the new steps as \emph{the} repair.  The
uncorrected conference Section~5 argument is not used as if it were sound
\cite[pp.~166--167]{SistlaClarke1982}.

\paragraph{The source's complete classification table.}
The final table distinguishes validity from both decision problems above:
\begin{center}
\small
\begin{tabular}{@{}lccc@{}}
\toprule
source language
& satisfiability
& validity
& existential \(R\)-truth\\
\midrule
\(L(F),\ \widetilde L(F,X)\)
& \(\mathsf{NP}\)-complete
& \(\mathsf{coNP}\)-complete
& \(\mathsf{NP}\)-complete\\
\(L(F,X),\ L(U),\ L(U,X),\ L(U,S,X),\ \mathsf{ETL}\)
& \(\mathsf{PSPACE}\)-complete
& \(\mathsf{PSPACE}\)-complete
& \(\mathsf{PSPACE}\)-complete\\
\bottomrule
\end{tabular}
\end{center}
The \(L(U,X)\) row is included in the journal's Table~I even though the
displayed statement of Theorem~4.1 lists the three bases
\(L(F,X),L(U),L(U,S,X)\).  The paper separately notes a different prior
route for \(L(U,X)\); its inclusion between the proved PSPACE upper language
and a PSPACE-hard sublanguage also fixes the class.  These rows support
Theorems~\ref{thm:p09-ltl-pspace} and
\ref{thm:p09-ltl-model-checking}, while preserving the source's existential
path problem, the operator restrictions, the nonempty-emission convention for
\(\mathsf{ETL}\), and the repaired-version boundary
\cite[p.~748]{SistlaClarke1985}.

\subsubsection{Alpern--Schneider's semantic decomposition}

The source classifies semantic properties of executions, not temporal formulas
by their typography.  Let \(S\neq\varnothing\) be the set of program states,
and put
\[
 S^*=\bigcup_{n<\omega}S^n,
 \qquad
 S^\omega=\{x:\omega\to S\}.
\]
The paper calls members of \(S^\omega\) \emph{executions} and members of
\(S^*\) \emph{partial executions}.  For \(u\in S^*\) and
\(y\in S^\omega\), let
\(u\mathbin{\smallfrown}y\in S^\omega\) be their concatenation; for
\(x\in S^\omega\), let \(x{\upharpoonright}n\) contain its first \(n\)
states.  For \(s\in S\), write \(\langle s\rangle\in S^1\subseteq S^*\)
for the one-state partial execution, and write \(s^\omega\in S^\omega\) for
the constant execution defined by \(s^\omega(n)=s\) for every
\(n\in\omega\).  These notations distinguish a one-state finite prefix from
a constant infinite execution and type the example witnesses below.  Thus
\[
 u\preceq x
 \quad\Longleftrightarrow\quad
 \exists y\in S^\omega\;
 x=u\mathbin{\smallfrown}y.
\]
A property is exactly a set \(P\subseteq S^\omega\).  The source notation
\(x\models P\) means membership \(x\in P\), not satisfaction of a separately
given formula.  A program denotes its own set \(B\subseteq S^\omega\) of
executions and satisfies \(P\) when \(B\subseteq P\).  A terminating execution
is made infinite by repeating its final state
\cite[p.~181]{AlpernSchneider1985}.
The article motivates the distinction proof-theoretically: safety arguments
normally establish an invariant, whereas liveness arguments normally use a
well-founded decrease.  This is its explanation of why the semantic classes
matter, not a claim that one proof rule decides either class.

The nonemptiness of \(S\) is operationally implicit in the article and is made
explicit here.  It is needed for the identification of liveness with density:
if \(S=\varnothing\), then \(S^\omega=\varnothing\) is dense in itself, while
the continuation clause below fails for the empty partial execution.  This is
a hypothesis repair for a degenerate case, not a restriction on any ordinary
program-state example in the source.

\paragraph{Safety and liveness.}
The property \(P\) is \emph{safety} exactly when
\[
 \forall x\in S^\omega\;
 \left(
  x\notin P
  \Longrightarrow
  \exists n\in\omega\;
  \forall y\in S^\omega\;
  (x{\upharpoonright}n)\mathbin{\smallfrown}y\notin P
 \right).
\]
Thus every violating execution has a finite bad prefix which no continuation
can repair.  The source first says that a partial execution \(u\) is
\emph{live for \(P\)} exactly when some \(y\in S^\omega\) satisfies
\(u\mathbin{\smallfrown}y\in P\).  This is a relation between a finite sequence
and a property, not a further property class; it differs from liveness of
\(P\) by lacking the universal quantifier over all prefixes.  The property
\(P\) is \emph{liveness} exactly when every partial execution is live for it:
\[
 \forall u\in S^*\;
 \exists y\in S^\omega\;
 u\mathbin{\smallfrown}y\in P.
\]
Thus no partial execution is irremediable.  The suffix may depend on the finite
prefix, and the required future behavior need not reduce to one eventual
atomic event.  Because the empty sequence belongs to \(S^*\), the clause also
implies \(P\neq\varnothing\)
\cite[pp.~182--183]{AlpernSchneider1985}.

The article also gives its reason for treating this as the broadest admissible
liveness definition.  If the continuation clause fails, there is a finite
prefix after which every continuation violates \(P\); that prefix is already
an irremediable bad witness and hence safety content.  This is a formal
consequence of the failed quantifier clause.  The further judgment that every
genuine liveness property must exclude such content is the authors'
methodological argument, not a theorem proving the informal notion uniquely
correct.

The article lists mutual exclusion, deadlock freedom, partial correctness, and
first-come-first-serve as safety examples, and starvation freedom,
termination, and guaranteed service as liveness examples.  Its footnote on
stuttering gives a property on which the chosen safety definition matters: if
the value of \(x\) must differ at every pair of successive states, one repeated
adjacent value is a finite bad prefix, so the property is safety here even
though it is not invariant under inserting finite state repetitions.  The
paper reports equivalence with Lamport's formulation only under a
stuttering-invariance assumption; no such assumption is built into the two
displayed clauses \cite[p.~182, n.~1]{AlpernSchneider1985}.

These quantifiers range over the full space \(S^\omega\).  Restricting them to
the paths generated by one transition system changes both the available
continuations and the topology below and must be stated as a separate
relativization.

\paragraph{Why two stricter liveness notions are not synonyms.}
Ordinary liveness permits the witnessing suffix to depend on the partial
execution.  To ask separately whether one suffix works for every prefix, and
whether every satisfying suffix works after every prefix, the article defines
two semantic subclasses of properties.  The distinction cannot be stated by
the unqualified word \emph{liveness} alone.

A property \(P\subseteq S^\omega\) has \emph{uniform liveness} exactly when
\[
 \exists y\in S^\omega\;
 \forall u\in S^*\;
 u\mathbin{\smallfrown}y\in P.
\]
The term names a semantic class of properties: its one fixed suffix differs
from the prefix-dependent witnesses allowed by ordinary liveness.  A property
\(P\subseteq S^\omega\) has \emph{absolute liveness} exactly when
\[
 \left(\exists z\in S^\omega\;z\in P\right)
 \land
 \forall y\in S^\omega\;
 \left(
  y\in P
  \Longrightarrow
  \forall u\in S^*\;
  u\mathbin{\smallfrown}y\in P
 \right).
\]
This term also names a semantic class of properties.  It differs from uniform
liveness by requiring \emph{every} satisfying execution to survive arbitrary
finite prefixing; the existence conjunct prevents the empty property from
passing vacuously.  The definitions are used immediately to obtain
\[
 \text{absolute liveness}
 \Longrightarrow
 \text{uniform liveness}
 \Longrightarrow
 \text{liveness}.
\]
For the first implication, choose one member of the nonempty absolute
property; it is a common suffix.  The second implication reuses the uniform
suffix after each prefix \cite[p.~183]{AlpernSchneider1985}.

The source's leads-to example makes the first distinction concrete.  Let
\(E_1,E_2\subseteq S\) be state sets, and require every occurrence of a state
in \(E_1\) to be followed later by an occurrence of a state in \(E_2\).  In
its event-language, the source says that such properties are live when
\(E_2\) is satisfiable, but says without qualification that a leads-to
property is not absolutely live.  Its non-absolute witness presupposes both an
\(E_1\)-state and a satisfying execution containing neither event type; those
objects need not exist.  Here is the exact repaired state-set separation.
Assume
\[
 E_1\ne\varnothing,
 \qquad E_2\ne\varnothing,
 \qquad S\setminus(E_1\cup E_2)\ne\varnothing.
\]
Choose \(e_2\in E_2\).  The fixed suffix \(e_2^\omega\) works after every
finite prefix, even if \(e_2\in E_1\), because each new request is answered by
a later suffix occurrence.  Thus the property is uniformly live.  Choose
\(e_1\in E_1\) and \(c\in S\setminus(E_1\cup E_2)\).  Then \(c^\omega\)
satisfies the property vacuously, whereas
\(\langle e_1\rangle\mathbin{\smallfrown}c^\omega\) does not; hence the
property is not
absolutely live.  These explicit hypotheses repair the printed example rather
than silently attributing a stronger valid claim to the source.  The three
hypotheses are jointly necessary as well as sufficient for this
uniformly-live-but-not-absolutely-live separation.  If \(E_1=\varnothing\),
the property is all of \(S^\omega\) and is absolutely live.  If
\(E_2=\varnothing\) while \(E_1\neq\varnothing\), a one-state \(E_1\)-prefix
has no satisfying continuation, so the property is not live.  If
\(S\setminus(E_1\cup E_2)=\varnothing\), every satisfying execution contains
an \(E_2\)-state; hence every satisfying suffix answers all requests in any
finite prefix as well as its own, and the property is absolutely live.  In
particular, this last case includes \(S=\{e_1,e_2\}\),
\(E_1=\{e_1\}\), and \(E_2=\{e_2\}\).

The second distinction is the article's predictive example, which the source
states in prose.  The handbook renders its quantifiers as follows.  Assume that
\(A,B\subseteq S\) and that
\(A,S\setminus A,B,S\setminus B\) are all nonempty, and put
\[
\begin{split}
 P_{A,B}=\{x\in S^\omega:
 &(x_0\in A\Longrightarrow
   \exists N\;\forall n\ge N\;x_n\in B)\\
 {}\land{}&
 (x_0\notin A\Longrightarrow
   \exists N\;\forall n\ge N\;x_n\notin B)\}.
\end{split}
\]
The symbol \(P_{A,B}\) is only an eliminable name for this displayed property.
After a nonempty prefix, its first state determines whether an eventually
always-\(B\) or eventually always-not-\(B\) suffix should be chosen; the empty
prefix can first choose either compatible branch.  Hence the property is live.
No one suffix is uniform: if it is eventually always \(B\), prepend a
non-\(A\) initial state; otherwise prepend an \(A\) initial state.  The
required branch then fails.  This is the precise quantifier dependency for
which the uniform class is needed \cite[p.~183]{AlpernSchneider1985}.

\paragraph{Neither is a dichotomy.}
The source's until example requires an eventual \(E_2\)-event and only
\(E_1\)-events before it.  It is the intersection of the safety property that
forbids a non-\(E_1\) event before \(E_2\) and the liveness property that
requires eventual \(E_2\).  The combined property is generally neither class:
some finite prefixes are already bad, while an all-\(E_1\) prefix still needs
a future witness.  Total correctness likewise intersects partial correctness
with termination.  These examples refute the claim that every property must
itself be either safety or liveness \cite[pp.~183--184]{AlpernSchneider1985}.

\paragraph{Prefix topology.}
For \(u\in S^*\), let
\[
 [u]=\{x\in S^\omega:u\preceq x\}.
\]
The cylinders \([u]\) form a basis for the prefix topology on \(S^\omega\).
The two semantic definitions become exact topological characterizations:
\[
 \begin{array}{c@{\quad\Longleftrightarrow\quad}c}
 P\text{ is safety}&P\text{ is closed},\\[2pt]
 P\text{ is liveness}&P\text{ is dense}.
 \end{array}
\]
For the first equivalence, a bad prefix gives a cylinder around every
\(x\notin P\) lying wholly in the complement, so the complement is open;
conversely, a basic cylinder contained in the open complement supplies the bad
prefix.  For the second, density is equivalent to meeting every nonempty basic
cylinder, and \([u]\cap P\neq\varnothing\) is exactly the continuation clause.
The article states these equivalences; these direction proofs make its
topological step explicit \cite[p.~184]{AlpernSchneider1985}.

\paragraph{Theorem 1: every property has a safety--liveness decomposition.}
Let \(\overline P\) be the topological closure of an arbitrary property
\(P\subseteq S^\omega\), equivalently the least safety property containing
\(P\), and define
\[
 L_P
 =S^\omega\setminus(\overline P\setminus P)
 =P\cup(S^\omega\setminus\overline P).
\]
The source writes this factor simply as \(L\).  The subscript in the
handbook's \(L_P\) only exposes its dependence on \(P\) and prevents collision
with other property variables; it is an eliminable metalinguistic renaming,
not a further kind of liveness.  The closure \(\overline P\) is
safety and
\[
 \overline P\cap L_P=P.
\]
It remains to prove that \(L_P\) is live.  If it were not dense, some nonempty
open \(O\) would satisfy
\[
 O\subseteq S^\omega\setminus L_P=\overline P\setminus P.
\]
Then \(P\subseteq\overline P\setminus O\).  Because \(\overline P\) and the
complement of \(O\) are closed, \(\overline P\setminus O\) is a safety
property containing \(P\); because \(O\neq\varnothing\), it is strictly
smaller than \(\overline P\).  This contradicts the minimality of the closure.
Thus \(L_P\) is dense and therefore live
\cite[Theorem~1, p.~184]{AlpernSchneider1985}.

This is an existence theorem.  It does not say that \(P\) itself belongs to
one class, that all decompositions are unique, that either factor is
algorithmically cheap, or that a chosen temporal syntax can express the
factors.  The closure is the least safety overapproximation of \(P\), but that
minimality does not turn the entire factorization into a uniqueness theorem.

The source's Corollary~1.1 isolates the needed expressibility hypothesis.  If
a notation \(\Sigma\) is closed under complement, intersection, and
topological closure, every \(\Sigma\)-expressible property has
\(\Sigma\)-expressible safety and liveness factors: the first is
\(\overline P\), and the second is the complement of
\(\overline P\cap(S^\omega\setminus P)\).  The corollary is conditional; it
does not prove closure expressible in any particular temporal language
\cite[Corollary~1.1, p.~184]{AlpernSchneider1985}.

The article then reports, on the authority of a 1983 private communication
from Plotkin, that every temporal-logic-expressible property can be written as
the conjunction of two temporal-logic-expressible liveness properties.  This
is a source-reported private-communication claim, not a result proved in the
article or independently verified here.

\paragraph{Theorem 2: two liveness factors.}
If \(|S|>1\), choose distinct \(a,b\in S\), and put
\[
 \begin{aligned}
 E_a&=\{x\in S^\omega:\exists N\;\forall n\ge N\;x_n=a\},\\
 E_b&=\{x\in S^\omega:\exists N\;\forall n\ge N\;x_n=b\}.
\end{aligned}
\]
The article writes these properties as \(L_a,L_b\).  The handbook's
alpha-renaming to \(E_a,E_b\) avoids collision with the preceding constructed
\(L_P\); it is eliminable and used only in this proof and the following
counterexample.  The two symbols name the properties of having an eventually
constant all-\(a\) or all-\(b\) tail.  Each is live because any
finite prefix can be continued by the corresponding constant suffix, and
\(E_a\cap E_b=\varnothing\).  Hence for every
\(P\subseteq S^\omega\),
\[
 (P\cup E_a)\cap(P\cup E_b)
 =P\cup(E_a\cap E_b)
 =P,
\]
and both factors are live \cite[Theorem~2, p.~184]{AlpernSchneider1985}.
The hypothesis is material: if \(|S|=1\), the only dense subset of the
singleton space \(S^\omega\) is the whole space, so the empty property cannot
be such an intersection.  The proof also shows directly that liveness is not
closed under intersection when \(|S|>1\).

\paragraph{The printed Corollary 2.1 requires repair.}
Both the journal article and the revised Cornell report say that if a notation
is closed under \emph{intersection} and contains two expressible liveness
properties with empty intersection, then every expressible property is the
intersection of two expressible liveness properties.  That hypothesis does
not support the displayed proof construction: its factors are
\(P\cup E_a\) and \(P\cup E_b\), so their expressibility requires union, not
intersection, closure.  The minimal corrected statement replaces the printed
word \emph{intersection} by \emph{union}; closure under complement and
intersection would also suffice by De Morgan's law.  Under either corrected
hypothesis the preceding calculation proves the conclusion.  The handbook
does not silently attribute this repaired hypothesis to the authors.

The printed hypothesis is not merely too weak for the displayed proof; it is
insufficient in general.  In the same \(|S|>1\) space let
\[
 Q=\{a^\omega\},
 \qquad
\mathcal F=\{S^\omega,E_a,E_b,Q,\varnothing\}.
\]
Here \(a^\omega\) is the constant-\(a\) execution, \(Q\) is only its singleton
property used in this counterexample, and
\(\mathcal F\) is an eliminable name for a possible denotation family; a
notation with precisely these denotations and conjunction interpreted as set
intersection is closed under intersection.
Because \(Q\subseteq E_a\), \(Q\cap E_b=\varnothing\), and
\(E_a\cap E_b=\varnothing\), the family is closed under intersection.  It
contains the disjoint liveness properties \(E_a,E_b\).  Its only liveness
members are \(S^\omega,E_a,E_b\): the cylinder \([\langle b\rangle]\) misses
\(Q\), and the
empty property is not dense.  No intersection of two of those three live
members is \(Q\).  Thus its expressible property \(Q\) has no factorization of
the kind claimed by the printed corollary
\cite[Corollary~2.1, p.~184]{AlpernSchneider1985}.

\paragraph{Boolean closure.}
The resulting binary-operation table is
\begin{center}
\small
\begin{tabular}{@{}lccc@{}}
\toprule
semantic class & union & intersection & complement\\
\midrule
safety / closed & closed & closed & not closed in general\\
liveness / dense & closed & not closed in general & not closed in general\\
\bottomrule
\end{tabular}
\end{center}
Closed sets are more generally closed under arbitrary intersections and finite
unions.  A union containing a dense set is dense.  The only property that is
both safety and liveness is \(S^\omega\) itself.  These are semantic closure
facts, not claims that a particular formula language is closed under the
operations needed to denote the resulting sets
\cite[p.~184]{AlpernSchneider1985}.

\paragraph{Source and proof boundary.}
The five-page journal article and the eleven-page Cornell technical report
agree on the definitions, examples, theorems, proofs, both corollary texts, and
conclusion.  In particular, the defective Corollary~2.1 hypothesis occurs in
both versions and therefore predates journal typesetting
\cite{AlpernSchneider1985,AlpernSchneider1985TR}.  The source prints the proofs
of Theorems~1--2 but only states the safety/closed and liveness/dense
equivalences, whose elementary proofs were supplied above.  Its concluding
defense of the intuition behind liveness is an argued methodological position,
not a formal correctness or uniqueness theorem.

The copyright notice printed on p.~181 remains controlling; this section is an
independent reconstruction, not a republication of the article.

\paragraph{Boundary among the five presentations.}
Kamp's result concerns first-order expressiveness over classes of linear orders
and uses strict witness and open-interval operators.  Pnueli begins from
transition systems and proof rules: his printed leads-to relation between state
assertions uses a strict later witness, whereas later uses of \(F\) and \(G\)
are inclusive and do not settle the paper's internal tensions.  GPSS presents
a strict future language.  Over discrete time, the handbook relates its strict
future operators, together with the handbook-supplied strict past dual, to the
inclusive word language; \(\mathsf X\) and \(\mathsf Y\) remain necessary for
the exact two-way bridge.  Sistla and Clarke classify satisfiability, validity,
and existential path truth for explicitly restricted inclusive operator
bases.  Alpern and Schneider classify arbitrary subsets of \(S^\omega\)
semantically, not formulas by syntax.  None of these presentations
retroactively supplies another's definitions, proofs, or source terminology.
